\PassOptionsToPackage{table}{xcolor}
%DIF LATEXDIFF DIFFERENCE FILE
%DIF DEL main_flat.tex                    Sun Sep 27 03:21:39 2026
%DIF ADD tmp\main_revised_annotated.tex   Mon Sep 28 17:59:31 2026
\documentclass{article}
\usepackage{preprint}
\usepackage{newtxtext,newtxmath}
\usepackage{microtype}
\usepackage{graphicx}
\usepackage{booktabs}
\usepackage{multirow}   % header labels that span both rows of a two-row header, so they centre vertically
\usepackage{float}
\usepackage{placeins}
\usepackage{longtable}
\usepackage{enumitem}
\usepackage{amsmath}
\let\Bbbk\relax
\usepackage{amssymb}
\usepackage{xcolor}
\usepackage{url}
\usepackage{hyperref}
\hypersetup{hidelinks}
\input{math_commands}
\input{tables/cf_colors}
\input{tables/cf_numbers}
%DIF 23a23-27
\usepackage[scheme=plain,fontset=fandol]{ctex} %DIF > 
\newcommand{\DIFdeletedtable}[1]{\par\noindent\fcolorbox{red!70!black}{red!3}{\parbox{\dimexpr\linewidth-2\fboxsep-2\fboxrule\relax}{\color{red!55!black}{\small\sffamily\textbf{已删除的表格}}\par\smallskip\centering #1}}\par} %DIF > 
\newcommand{\DIFoldtable}[1]{\par\noindent\fcolorbox{orange!75!black}{orange!4}{\parbox{\dimexpr\linewidth-2\fboxsep-2\fboxrule\relax}{\color{orange!45!black}{\small\sffamily\textbf{修改前的表格}}\par\smallskip\centering #1}}\par} %DIF > 
\newcommand{\cfnote}[2]{\par\smallskip\noindent\fcolorbox{orange!70!black}{orange!8}{\parbox{\dimexpr\linewidth-2\fboxsep-2\fboxrule\relax}{\small\sffamily\textbf{修改说明·#1}\quad #2}}\par\smallskip} %DIF > 
\newenvironment{DIFnomarkup}{}{} %DIF > 
%DIF -------
\setlength{\textfloatsep}{6pt plus 2pt minus 2pt}
\setlength{\abovecaptionskip}{5pt}
\setlength{\belowcaptionskip}{0pt}

\title{Readability Does Not Confer Control over Clinical Concepts}

\authorline{%
Hanqi Jiang\aff{1,2,*} \quad
Weihang You\aff{1,*} \quad
Qingchan Zhu\aff{1} \quad
Junhao Chen\aff{1} \quad
Yi Pan\aff{1}\\
Zhengliang Liu\aff{1} \quad
Wei Liu\aff{3} \quad
Lei Xing\aff{2,\dag}}

\affiliations{%
\aff{1}School of Computing, University of Georgia\\
\aff{2}School of Medicine, Stanford University\\
\aff{3}Department of Radiation Oncology, Mayo Clinic, Phoenix, AZ, USA}

\titlenotes{%
\aff{*}Equal contribution.\\
\aff{\dag}Corresponding author: \texttt{lei@stanford.edu}.}
%DIF PREAMBLE EXTENSION ADDED BY LATEXDIFF
%DIF UNDERLINE PREAMBLE %DIF PREAMBLE
\RequirePackage[normalem]{ulem} %DIF PREAMBLE
\RequirePackage{color}\definecolor{RED}{rgb}{1,0,0}\definecolor{BLUE}{rgb}{0,0,1} %DIF PREAMBLE
\providecommand{\DIFaddtex}[1]{{\protect\color{blue}\uwave{#1}}} %DIF PREAMBLE
\providecommand{\DIFdeltex}[1]{{\protect\color{red}\sout{#1}}} %DIF PREAMBLE
%DIF SAFE PREAMBLE %DIF PREAMBLE
\providecommand{\DIFaddbegin}{} %DIF PREAMBLE
\providecommand{\DIFaddend}{} %DIF PREAMBLE
\providecommand{\DIFdelbegin}{} %DIF PREAMBLE
\providecommand{\DIFdelend}{} %DIF PREAMBLE
\providecommand{\DIFmodbegin}{} %DIF PREAMBLE
\providecommand{\DIFmodend}{} %DIF PREAMBLE
%DIF FLOATSAFE PREAMBLE %DIF PREAMBLE
\providecommand{\DIFaddFL}[1]{\DIFadd{#1}} %DIF PREAMBLE
\providecommand{\DIFdelFL}[1]{\DIFdel{#1}} %DIF PREAMBLE
\providecommand{\DIFaddbeginFL}{} %DIF PREAMBLE
\providecommand{\DIFaddendFL}{} %DIF PREAMBLE
\providecommand{\DIFdelbeginFL}{} %DIF PREAMBLE
\providecommand{\DIFdelendFL}{} %DIF PREAMBLE
%DIF HYPERREF PREAMBLE %DIF PREAMBLE
\providecommand{\DIFadd}[1]{\texorpdfstring{\DIFaddtex{#1}}{#1}} %DIF PREAMBLE
\providecommand{\DIFdel}[1]{\texorpdfstring{\DIFdeltex{#1}}{}} %DIF PREAMBLE
\providecommand{\DIFscaledelfig}{0.5}
%DIF HIGHLIGHTGRAPHICS PREAMBLE %DIF PREAMBLE
\RequirePackage{settobox} %DIF PREAMBLE
\RequirePackage{letltxmacro} %DIF PREAMBLE
\newsavebox{\DIFdelgraphicsbox} %DIF PREAMBLE
\newlength{\DIFdelgraphicswidth} %DIF PREAMBLE
\newlength{\DIFdelgraphicsheight} %DIF PREAMBLE
% store original definition of \includegraphics %DIF PREAMBLE
\LetLtxMacro{\DIFOincludegraphics}{\includegraphics} %DIF PREAMBLE
\providecommand{\DIFaddincludegraphics}[2][]{{\color{blue}\fbox{\DIFOincludegraphics[#1]{#2}}}} %DIF PREAMBLE
\providecommand{\DIFdelincludegraphics}[2][]{% %DIF PREAMBLE
\sbox{\DIFdelgraphicsbox}{\DIFOincludegraphics[#1]{#2}}% %DIF PREAMBLE
\settoboxwidth{\DIFdelgraphicswidth}{\DIFdelgraphicsbox} %DIF PREAMBLE
\settoboxtotalheight{\DIFdelgraphicsheight}{\DIFdelgraphicsbox} %DIF PREAMBLE
\scalebox{\DIFscaledelfig}{% %DIF PREAMBLE
\parbox[b]{\DIFdelgraphicswidth}{\usebox{\DIFdelgraphicsbox}\\[-\baselineskip] \rule{\DIFdelgraphicswidth}{0em}}\llap{\resizebox{\DIFdelgraphicswidth}{\DIFdelgraphicsheight}{% %DIF PREAMBLE
\setlength{\unitlength}{\DIFdelgraphicswidth}% %DIF PREAMBLE
\begin{picture}(1,1)% %DIF PREAMBLE
\thicklines\linethickness{2pt} %DIF PREAMBLE
{\color[rgb]{1,0,0}\put(0,0){\framebox(1,1){}}}% %DIF PREAMBLE
{\color[rgb]{1,0,0}\put(0,0){\line( 1,1){1}}}% %DIF PREAMBLE
{\color[rgb]{1,0,0}\put(0,1){\line(1,-1){1}}}% %DIF PREAMBLE
\end{picture}% %DIF PREAMBLE
}\hspace*{3pt}}} %DIF PREAMBLE
} %DIF PREAMBLE
\LetLtxMacro{\DIFOaddbegin}{\DIFaddbegin} %DIF PREAMBLE
\LetLtxMacro{\DIFOaddend}{\DIFaddend} %DIF PREAMBLE
\LetLtxMacro{\DIFOdelbegin}{\DIFdelbegin} %DIF PREAMBLE
\LetLtxMacro{\DIFOdelend}{\DIFdelend} %DIF PREAMBLE
\DeclareRobustCommand{\DIFaddbegin}{\DIFOaddbegin \let\includegraphics\DIFaddincludegraphics} %DIF PREAMBLE
\DeclareRobustCommand{\DIFaddend}{\DIFOaddend \let\includegraphics\DIFOincludegraphics} %DIF PREAMBLE
\DeclareRobustCommand{\DIFdelbegin}{\DIFOdelbegin \let\includegraphics\DIFdelincludegraphics} %DIF PREAMBLE
\DeclareRobustCommand{\DIFdelend}{\DIFOaddend \let\includegraphics\DIFOincludegraphics} %DIF PREAMBLE
\LetLtxMacro{\DIFOaddbeginFL}{\DIFaddbeginFL} %DIF PREAMBLE
\LetLtxMacro{\DIFOaddendFL}{\DIFaddendFL} %DIF PREAMBLE
\LetLtxMacro{\DIFOdelbeginFL}{\DIFdelbeginFL} %DIF PREAMBLE
\LetLtxMacro{\DIFOdelendFL}{\DIFdelendFL} %DIF PREAMBLE
\DeclareRobustCommand{\DIFaddbeginFL}{\DIFOaddbeginFL \let\includegraphics\DIFaddincludegraphics} %DIF PREAMBLE
\DeclareRobustCommand{\DIFaddendFL}{\DIFOaddendFL \let\includegraphics\DIFOincludegraphics} %DIF PREAMBLE
\DeclareRobustCommand{\DIFdelbeginFL}{\DIFOdelbeginFL \let\includegraphics\DIFdelincludegraphics} %DIF PREAMBLE
\DeclareRobustCommand{\DIFdelendFL}{\DIFOaddendFL \let\includegraphics\DIFOincludegraphics} %DIF PREAMBLE
%DIF AMSMATHULEM PREAMBLE %DIF PREAMBLE
\makeatletter %DIF PREAMBLE
\let\sout@orig\sout %DIF PREAMBLE
\renewcommand{\sout}[1]{\ifmmode\text{\sout@orig{\ensuremath{#1}}}\else\sout@orig{#1}\fi} %DIF PREAMBLE
\makeatother %DIF PREAMBLE
%DIF COLORLISTINGS PREAMBLE %DIF PREAMBLE
\RequirePackage{listings} %DIF PREAMBLE
\RequirePackage{color} %DIF PREAMBLE
\lstdefinelanguage{DIFcode}{ %DIF PREAMBLE
%DIF DIFCODE_UNDERLINE %DIF PREAMBLE
  moredelim=[il][\color{red}\sout]{\%DIF\ <\ }, %DIF PREAMBLE
  moredelim=[il][\color{blue}\uwave]{\%DIF\ >\ } %DIF PREAMBLE
} %DIF PREAMBLE
\lstdefinestyle{DIFverbatimstyle}{ %DIF PREAMBLE
	language=DIFcode, %DIF PREAMBLE
	basicstyle=\ttfamily, %DIF PREAMBLE
	columns=fullflexible, %DIF PREAMBLE
	keepspaces=true %DIF PREAMBLE
} %DIF PREAMBLE
\lstnewenvironment{DIFverbatim}[1][]{\lstset{style=DIFverbatimstyle}}{} %DIF PREAMBLE
\lstnewenvironment{DIFverbatim*}[1][]{\lstset{style=DIFverbatimstyle,showspaces=true}}{} %DIF PREAMBLE
\lstset{extendedchars=\true,inputencoding=utf8}

%DIF END PREAMBLE EXTENSION ADDED BY LATEXDIFF

\begin{document}
\maketitle

\begin{abstract}
A linear probe reads a clinical finding from the visual representation that a vision-language model's language model consumes, and activation steering suggests the converse: writing the probe direction back should move the model's answer about the finding, making the direction a handle on the concept. Evaluations test whether models encode findings and answer questions about them, not whether such a handle exists. We test it with \emph{ownership}: at a fixed dose on the same patients, does a concept's own direction move its own answer more than the directions fitted for every competing finding? Across \cfCheckpoints{} checkpoints on NIH ChestX-ray14, CheXpert Plus and COCO, clinical concepts are readable in \cfChestReadablePct{}\% of chest checkpoint--concept cells and answerable in \cfChestAnswerablePct{}\%, but owned in \cfChestOwnedPct{}\%: most clinical writes do not move their own answer beyond random directions of the same dose, and the same directions move the other findings' answers. Under the same procedure, \cfCocoOwnedPct{}\% of natural-image object cells are owned, and the gap persists among cells the model answers well (\cfStratChestTopPct{}\% against \cfStratNatTopPct{}\%). Non-clinical attributes of the radiographs share the deficit, and checkpoints that share a vision tower, and therefore write identical vectors, own different concepts: the deficit follows the chest-radiograph representation and the model that reads it. A direction fitted to the model's own answers provides the handle that the label direction lacks. A probe's readability does not confer control through its direction.
\end{abstract}

\section{Introduction}
\label{sec:introduction}

Linear probes are the standard tool for testing what a vision-language model (VLM) represents. They read out a concept by fitting a classifier to activations at a chosen layer \citep{alain2017understanding,belinkov2022probing}. In medical imaging, probes often succeed at linearly decoding clinical findings such as pleural effusion or cardiomegaly from the visual representation of general-purpose and medical VLMs \citep{zhu2026lost,pedersen2026probes}. Activation steering promises the converse: writing the probe direction back into the representation to change the model's behaviour \citep{turner2023steering,panickssery2024caa}. Together, these methods cast the probe direction as a handle on the concept: read it to detect, write it to control. We ask whether that handle exists.

A steering write that changes the answer does not settle the question. In Qwen2.5-VL-7B on NIH ChestX-ray14, a probe reads effusion from the visual tokens the language model consumes, and writing its direction back raises the model's probability of answering yes to the Effusion question by $\cfSeedWEff$, more than the 95th percentile of 119 random directions and a coordinate-permutation sham written at the same dose. By the usual steering criteria, the write works. Writing the Nodule direction into the same patients at the same dose raises the Effusion answer further, by $\cfSeedWNodule$ (Figure~\ref{fig:framework}c), although Nodule is the competing direction least similar to Effusion (Appendix~\ref{app:cf-robustness}). Random and sham directions carry no label, so they cannot expose a direction whose effect is shared with other findings \citep{pahde2025navigating}. Exposing it requires writing every competing clinical direction at the same locus and dose in the same patients and asking whether the concept's own direction moves the concept's own answer more than any other. We call the outcome \emph{ownership}. When a concept is readable but not owned, its probe direction reads the concept without writing it: reading is not writing.

We turn this comparison into an evaluation (Section~\ref{sec:framework}, Figure~\ref{fig:framework}). At the \emph{consumed block}, the final visual block whose output the language model reads, each concept receives three grades. It is \emph{readable} if its probe beats random-label probes of equal capacity, \emph{answerable} if the model's yes/no answer without any write beats chance, and \emph{owned} if a fixed-dose write along its direction clears the random and sham references and moves its own answer more than every competing clinical direction under simultaneous bounds. We apply the evaluation to \cfCheckpoints{} checkpoints from \cfFamiliesWord{} families (3B--72B; general-purpose and medical) on two chest-radiograph datasets with report-derived labels, NIH ChestX-ray14 and CheXpert Plus, and on COCO objects as a natural-image control (Section~\ref{sec:setup}).

On chest radiographs most readable and answerable concepts are not owned, while most COCO objects are, and the gap persists among concepts the model answers well (Section~\ref{sec:results}). We call this gap the \emph{attribution deficit}. Since the same write succeeds on natural images, the remaining results locate the deficit. Checkpoints that share a vision tower write identical vectors into identical representations yet own different concepts, so the model that reads the representation decides ownership. Non-clinical attributes of the same radiographs (projection, sex, and age) are owned no more often than findings, so the deficit does not depend on the concept being clinical. The reader responds to a different direction: a direction fitted to the model's own answers is owned in \cfAnsChestOwnedA{} of \cfAnsChestN{} chest checkpoint--concept cells where the label direction is owned in \cfAnsChestOwnedLabel{}, and on chest radiographs it is nearly orthogonal to the probe direction.

\begin{figure}[t]
    \centering
    \resizebox{\textwidth}{!}{\includegraphics{figures/fig1_method.pdf}}
    \caption{\textbf{Reading and writing a concept at the consumed block.} (a) The probe reads the concept from the consumed visual tokens, and its lifted normal is written back to them at dose $\alpha=+0.25$, a quarter of each token's norm. (b) The write matrix $W_{q,d}$, the mean change in $P(\mathrm{yes})$ of question $q$ under a write along direction $d$, and the three comparisons that decide ownership. (c) The Effusion case (Qwen2.5-VL-7B, NIH ChestX-ray14): a stronger competitor.}
    \label{fig:framework}
\end{figure}

\paragraph{Contributions.} (i) We introduce ownership, which separates reading a concept from writing it: a probe direction is scored against directions fitted the same way for the competing concepts, written at the same locus and dose in the same patients (Section~\ref{sec:framework}). (ii) We identify the attribution deficit: across \cfCells{} checkpoint--concept cells, clinical concepts on chest radiographs are read and answered far more often than they are owned, and the gap to natural-image objects persists at matched answer ability (Sections~\ref{sec:headline}--\ref{sec:coco}). (iii) We trace the deficit to how the model reads the chest-radiograph representation rather than to the concept's clinical status, and show that a direction fitted to the model's own answers provides the handle that the label direction lacks (Sections~\ref{sec:reader}--\ref{sec:ansdir}). We release the evaluation, scorecards, and write matrices; a new model needs one adapter naming its consumed block.

\section{Background and Related Work}
\label{sec:related}

\paragraph{Probing and its controls.}
Linear probes track information across network depth \citep{alain2017understanding}, and a probe score reflects both the representation and the learner fitted to it, which motivates matched-capacity probes and control tasks \citep{hewitt2019designing}. The probing literature distinguishes information a readout can recover from information the model uses \citep{belinkov2022probing,ravichander2021probing}; causal erasure tests use by deleting a concept \citep{ravfogel2020null,elazar2021amnesic}, in closed form for every linear reader at once \citep{belrose2023leace}. Our readability grade adopts the control-task idea, and the write test asks a question complementary to erasure: whether the concept's own direction is the one that moves its answer. In medical multimodal models, depth-wise probing maps classification information across towers, connectors, and language layers \citep{zhu2026lost}, and layer-wise probes of MedCLIP expose acquisition and device shortcuts \citep{pedersen2026probes}; we intervene on the decoded clinical direction at the locus where it is read.

\paragraph{Concept directions and steering.}
A concept activation vector is the normal of a classifier separating concept examples from a negative set, and it grades a class by the directional derivative of the logit along that normal \citep{kim2018tcav}; we write the direction forward instead. Such a normal shifts with the layer it is fitted at, entangles with correlated concepts, depends on its negative set, and carries where in the image the concept was read \citep{nicolson2025explaining}. Representation engineering and activation steering write a concept back by changing intermediate representations at inference time \citep{zou2023representation,turner2023steering,panickssery2024caa}, and a probe normal is an appealing intervention because it supplies a label and a signed vector from the same representation. A separability-trained vector can, however, absorb distractor signals and diverge from its concept \citep{pahde2025navigating}: a classifier's weights are a filter rather than the pattern they detect \citep{haufe2014interpretation}, mass-mean directions steer truth judgements better than logistic normals \citep{marks2024geometry}, and the features that identify a concept are largely not those that control the output \citep{arad2025saes}. Benchmarks accordingly score concept detection and steering as separate axes \citep{wu2025axbench}, and alignment search certifies a subspace as causal when it implements an interpretable variable \citep{geiger2024alignments}. These works judge a direction mainly against its own concept. Ownership fixes the estimator and varies the label instead: it asks whether a direction fitted the same way for another concept moves the same answer more. Section~\ref{sec:robust-geometry} runs alternative constructions, a second locus, and token-concentrated writes through this test.

\paragraph{Writing into a vision-language model.}
Image features are linearly decodable in a vision-language model's hidden layers, and targeted edits to them change downstream answers \citep{rajaram2025lineofsight}. Steering one sparse-autoencoder feature of the vision tower and asking the language model what it sees recovers what that feature encodes \citep{ferrando2026explain}; we start from the meaning, fit a direction to a clinical label, and ask whether the write keeps that label's selectivity against the other findings' directions. Causal probing of multimodal models intervenes on the language model's residual stream with difference-in-means vectors and scores each concept against its own effect \citep{deng2026causal}; we write the representation the language model consumes and score each concept against its competitors.

\section{Reading, Answering, and Writing at the Consumed Block}
\label{sec:framework}

We ask three questions about a concept $c$ in a checkpoint $M$ on a dataset with binary labels: whether $c$ is readable from the visual representation the language model consumes, whether $M$ answers the question about $c$, and whether writing $c$'s own direction moves $c$'s answer more than any alternative. A registered protocol fixes an estimator and a decision rule for each question (Appendix~\ref{app:cf-labels}).

\subsection{Consumed block and matched read--write positions}
\label{sec:locus}

Every VLM in the grid passes an image through a vision tower and hands a specific tensor to a connector or merger that produces the language model's inputs. The \emph{consumed block} is the last visual block whose output that consumer reads, restricted to the token positions the consumer actually uses (CLS or padding rows are excluded where the consumer drops them); the \emph{connector locus} is the consumer's output. Appendix~\ref{app:repro-loci} identifies the block for every family. For image $i$ and consumed token $t$, let $\hvec_{it}\in\R^D$ be the output of the consumed block, where $D$ is the tower's hidden width; mean pooling over the consumed tokens gives the read vector $\hvec_i$. A write adds the unit direction $\hat\wvec_c$ of concept $c$ (Section~\ref{sec:directions}) to the same tokens,
\begin{equation}
\hvec'_{it}=\hvec_{it}+\alpha\,\lVert\hvec_{it}\rVert_2\,\hat\wvec_c,
\label{eq:intervention}
\end{equation}
so that the dose $\alpha$ is a fraction of each token's own norm, comparable across families whose activation scales differ by orders of magnitude, and its sign selects the concept's positive or negative pole. Reading and writing therefore use the same representation and the same positions. The connector and the language model, everything downstream of the consumed block, form the \emph{reader} of that representation.

\subsection{Directions, controls, and the reference family}
\label{sec:directions}

Each concept is read by a linear probe on a shared 512-dimensional view of the representation. A fixed Gaussian matrix $R\in\R^{D\times512}$ maps the read vector to features $\boldsymbol x_i=R^{\top}\hvec_i$, which are standardised with training-row means and scales $\boldsymbol s$; an $\ell_2$-regularised logistic regression with coefficient vector $\boldsymbol\beta_c$ then scores concept $c$ (inverse regularisation strength 1, 2{,}000 iterations, three fit seeds). The probe's logit is thus an affine function of $\hvec_i$ whose gradient is $R\,\mathrm{diag}(\boldsymbol s)^{-1}\boldsymbol\beta_c$, and the write direction is that gradient at unit length,
\begin{equation}
\hat\wvec_c=\frac{\boldsymbol g_c}{\lVert\boldsymbol g_c\rVert_2},\qquad \boldsymbol g_c=R\,\mathrm{diag}(\boldsymbol s)^{-1}\boldsymbol\beta_c .
\label{eq:lift}
\end{equation}
It is the direction in the consumed representation along which the probe's evidence for $c$ grows fastest (Appendix~\ref{app:repro-probe} gives the numerical details). We call it the concept's \emph{label direction}, since it is fitted to the concept's labels.

Two families of controls bracket the label direction. For reading, 20 control probes are fitted to random labels drawn within recurring image types (view, sex, and age decade on chest sets; aspect and area buckets on COCO), so that they have the capacity and the nuisance structure of the concept probe. With $y$ the concept labels and $r$ the probe's scores on the calibration rows, and $y^{\mathrm{ctrl}}_k$ and $r^{\mathrm{ctrl}}_k$ the random labels and scores of control probe $k$, the probe's selectivity is
\begin{equation}
S=\operatorname{AUROC}(y,r)-\tfrac{1}{20}\textstyle\sum_{k=1}^{20}\operatorname{AUROC}(y^{\mathrm{ctrl}}_k,r^{\mathrm{ctrl}}_k),
\label{eq:selectivity}
\end{equation}
which is positive when the probe reads its concept better than a probe of equal capacity reads an arbitrary label with the same nuisance structure. For writing, the \emph{reference family} contains 119 random unit directions, a coordinate-permutation sham of $\hat\wvec_c$ (same coordinates, scrambled arrangement), and the five label directions of the other clinical concepts, all written at the same dose to the same images.

\subsection{Ownership}
\label{sec:ownership}

Each dataset has six concepts, and we index both the yes/no question about a concept and the direction fitted for it by that concept. For question $q$ and direction $d$, $W_{q,d}$ is the mean over test images of the change in $P(\mathrm{yes})$ for question $q$ when $\hat\wvec_d$ is written, each image scored with and without the write. Here $P(\mathrm{yes})=\sigma(\ell_{\mathrm{yes}}-\ell_{\mathrm{no}})$, where $\sigma$ is the logistic function and $\ell_{\mathrm{yes}}$ and $\ell_{\mathrm{no}}$ are the log-probabilities of the two answers at the first answer position, each the maximum over case and leading-space variants that tokenise to one token. Row $q$ of the $6\times6$ write matrix $W$ records how far each direction moves question $q$, and its diagonal entry $W_{q,q}$ is the concept's own write (Figure~\ref{fig:write-matrices}). Ownership is the diagonal margin
\begin{equation}
O_q=W_{q,q}-\max_{d\neq q}W_{q,d},
\label{eq:margin}
\end{equation}
the advantage of the concept's own direction over its strongest clinical competitor on the concept's own question. $W_{q,q}$ measures how far the concept's own direction moves its answer; $O_q$ measures whether it moves that answer further than every competitor built the same way. A large $W_{q,q}$ shows that $d=q$ can write the answer; only $O_q>0$ shows that the label names the effective direction. A concept can therefore fail ownership in two ways: its own write stays within the random and sham references, or the write clears them without holding its advantage over every competitor. Section~\ref{sec:headline} measures both. The Effusion case of the introduction illustrates the latter: the Effusion direction raises the Effusion answer by $\cfSeedWEff$, above the random 95th percentile and the sham, and the Nodule direction raises the same answer by $\cfSeedWNodule$, so $O_{\mathrm{Effusion}}=\cfSeedEffO$. Uncertainty over patients comes from 2{,}000 bootstrap draws of test rows, with the maximum in Equation~\ref{eq:margin} recomputed inside every draw. The 30 contrasts $W_{q,q}-W_{q,d}$ (six questions, five competitors each) receive simultaneous 95\% bounds from a centred studentised max-$T$ bootstrap, which hold jointly across all 30 (Appendix~\ref{app:simultaneous-bounds}).

\subsection{Graded outcomes}
\label{sec:grades}

The readable and answerable grades are computed on the 400 calibration rows, which carry no write. The write matrix and every ownership quantity are computed on the 600 test rows, which share no patient with the calibration rows.
\begin{itemize}[leftmargin=*,itemsep=1pt,topsep=2pt]
\item \textbf{Readable:} the one-sided 95\% lower bound of $S$ on the calibration rows is above zero, with at least ten known positives and ten known negatives among those 400 rows.
\item \textbf{Answerable:} the \emph{clean answer}, $P(\mathrm{yes})$ with no write, separates the concept's labels: the one-sided 95\% lower bound of its AUROC on the same calibration rows is above $0.5$, under the same ten-and-ten support rule.
\item \textbf{Owned:} on the test rows, the write meets the \emph{steering reference} ($W_{q,q}>0$, above the 95th percentile of the 119 random writes, and above the sham's absolute effect) and every simultaneous lower bound of $W_{q,q}-W_{q,d}$ is positive. A cell with any simultaneous upper bound below zero has a \emph{stronger competitor}; the rest are \emph{unresolved}.
\end{itemize}
A checkpoint--concept pair on one dataset is a \emph{cell}, and a checkpoint--dataset pair is a \emph{block}. A cell graded owned is said to be owned, and the checkpoint is said to own the concept; an alternative direction is owned in a cell when its write meets the same rule for that concept. A block contributes to the answerable and owned counts only when its write matrix and its calibration pass both completed. The grades are defined for one direction construction, one primary dose with a five-point sweep, six templates, pooled consumed-token features, and two loci, the consumed block and the connector output. Before a block is scored, seven acceptance checks on 16 held-out preflight rows confirm, at both loci, that scoring is deterministic, that $\alpha=0$ changes nothing, that the write reaches the answer and changes nothing outside the consumed tokens, and that every template maps its answers correctly without an image (Appendix~\ref{app:cf-gates}).
\label{sec:gates}

\section{Models, Datasets, and Scale}
\label{sec:setup}

\paragraph{Checkpoints.}
We evaluate \cfCheckpoints{} instruction-tuned checkpoints from \cfFamiliesWord{} families, 3B to 72B (Table~\ref{tab:cf-models}). Six families are general-purpose: Qwen2.5-VL \citep{bai2025qwen25vl}, Qwen3-VL \citep{qwen2025qwen3vl}, InternVL3.5 \citep{wang2025internvl35}, Gemma~3 \citep{gemma2025gemma3}, LLaVA-1.5 \citep{liu2024improved}, and Llama~3.2~Vision \citep{meta2024llama32vision}. Six are medical fine-tunes: MedGemma \citep{sellergren2025medgemma}, Lingshu \citep{lasa2025lingshu}, LLaVA-Med~1.5 \citep{li2023llavamed}, CheXagent \citep{chen2024chexagent}, HuatuoGPT-Vision \citep{chen2024huatuogptvision}, and LLaVA-Rad \citep{zambranochaves2025llavarad}. They span \cfTowerDesignsWord{} vision-tower designs and three consumer designs: token insertion after a merger, token insertion after a projector, and cross-attention. Three towers are trained on chest radiographs, those of CheXagent-2-3B, CheXagent-8B, and LLaVA-Rad-7B (Appendix~\ref{app:cf-grid}). Every checkpoint runs in bf16 \DIFdelbegin \DIFdel{at a pinned revision }\DIFdelend with its native processor.

\paragraph{Datasets and cohorts.}
NIH ChestX-ray14 \citep{wang2017chestxray} and CheXpert Plus \citep{chambon2024chexpertplus} provide frontal chest radiographs with report-derived labels, the CheXpert labels coming from the CheXpert labeler \citep{irvin2019chexpert}. We evaluate six findings per dataset (NIH: Effusion, Atelectasis, Pneumothorax, Cardiomegaly, Mass, Nodule; CheXpert: Effusion, Atelectasis, Pneumothorax, Cardiomegaly, Consolidation, Edema). COCO \citep{lin2014coco} provides the natural-image control, with binary presence labels for six object categories (person, dog, car, chair, bottle, bicycle). Seeded hashes of patient or image identity select one fixed cohort per dataset, shared by every checkpoint: 20{,}000 training rows (\cfCohortNihTrainRows{} on NIH) fit the probes and the answer direction, 400 calibration rows carry the readable and answerable grades, 600 test rows carry every write, and 16 preflight rows carry the acceptance checks. No patient and no image appears in two roles of a dataset (Appendix~\ref{app:repro-cohorts}).

\paragraph{Questions and scale.}
Each concept is asked as a yes/no question in six templates that cross two wordings (\emph{is there} versus \emph{does the image show}) with three answer interfaces (yes/no, A-present, B-present); the \emph{is}/yes-no template gives the primary grade (Appendix~\ref{app:protocol}). The primary dose is $\alpha=+0.25$, and a sweep adds $\alpha\in\{-0.5,-0.25,-0.1,+0.1,+0.5\}$. Each block scores the write matrix, the sham, and the 119 random directions on the 600 test rows. The campaign comprises \cfBlocks{} blocks, \cfCells{} scored concept cells, about 30 million scored answers, and roughly 3{,}000 A100-hours; \cfCoverageBlocksAllSeven{} of the \cfBlocks{} blocks carry every planned analysis (Appendix~\ref{app:cf-coverage}).

\section{Results}
\label{sec:results}

\subsection{Reading is not writing}
\label{sec:headline}

On both chest datasets, reading and answering far outpace writing (Table~\ref{tab:cf-main}, Figure~\ref{fig:overview}a--b). Clinical concepts are readable in \cfChestReadable{} of \cfChestReadableN{} cells (\cfChestReadablePct\%; median selectivity $\cfSelChestMed$), and \cfChestAnswerable{} (\cfChestAnswerablePct\%) are answerable. Yet only \cfChestOwned{} cells (\cfChestOwnedPct\%) are owned, and reading and answering together barely change this: \cfReadAnsOwned{} of the \cfReadAns{} readable and answerable cells are owned (\cfReadAnsOwnedPct\%), against \cfRestOwned{} of the remaining \cfRestN{} (\cfRestOwnedPct\%).

Ownership on chest radiographs breaks at the write itself. \cfChestRefBelow{} of \cfChestCellsN{} label directions fail the random and sham reference for their own question, and these cells carry \cfChestBelowStrong{} of the \cfChestCompetitor{} stronger-competitor verdicts (\cfChestBelowStrongPct\%\DIFdelbegin \DIFdel{; Table~\ref{tab:cf-contingency}}\DIFdelend ). Once an own write clears the reference, \cfChestOwned{} of \cfChestRefMet{} cells are owned; on COCO, \cfCocoRefMet{} of \cfCocoOwnedN{} writes clear it and \cfCocoOwned{} are owned. On chest radiographs, most label directions therefore scarcely move the answers they name, while the directions of other findings often move those answers more.

\label{sec:seed}The Effusion case reproduces on 600 new patients (Table~\ref{tab:cf-seed}). The Effusion write raises its answer by $\cfSeedWEff$, second among the 120 compared effects, the Nodule write raises it by $\cfSeedWNodule$, and $O_{\mathrm{Effusion}}=\cfSeedEffO$ $[\cfSeedEffOLo,\cfSeedEffOHi]$, matching the $\cfSeedPriorO$ $[\cfSeedPriorOLo,\cfSeedPriorOHi]$ of the two-model pilot study that motivated the evaluation (Appendix~\ref{app:seed}). Across the grid, Effusion is owned in \cfEffusionOwnedBlocks{} of \cfChestBlocks{} chest blocks, and in \cfEffusionOwnedSmall{} of these $O_{\mathrm{Effusion}}\leq0.03$.

\begin{table}[t]
\centering
\setlength{\tabcolsep}{2pt}
\renewcommand{\arraystretch}{1.08}
\input{tables/cf_colors}
\caption{\textbf{Reading, answering, and writing per checkpoint.} Per dataset: mean probe selectivity $S$ (\textsc{Read}), mean clean-answer AUROC (\textsc{Ans}), and mean ownership $O_q$ (\textsc{Own}), each superscripted with the number of the six concepts that receive the grade; then the owned share of clinical and COCO cells and their ratio, by which rows are sorted. Shading darkens with magnitude (blue for negative \textsc{Own}); ``--'' marks an undefined ratio.}
\label{tab:cf-main}
\DIFdelbeginFL %DIFDELCMD < \fitwidth{%
%DIFDELCMD < \begin{tabular}{l@{\hspace{3pt}}ccc@{\hspace{4pt}}ccc@{\hspace{4pt}}ccc@{\hspace{4pt}}ccc}
%DIFDELCMD < \toprule
%DIFDELCMD < \multirow{2}{*}{Checkpoint} & \multicolumn{3}{c}{NIH ChestX-ray14} & \multicolumn{3}{c}{CheXpert Plus} & \multicolumn{3}{c}{COCO (control)} & \multicolumn{3}{c}{Owned share}\\
%DIFDELCMD < \cmidrule(lr){2-4}\cmidrule(lr){5-7}\cmidrule(lr){8-10}\cmidrule(lr){11-13}
%DIFDELCMD < & \multicolumn{1}{c}{\textsc{Read}} & \multicolumn{1}{c}{\textsc{Ans}} & \multicolumn{1}{c}{\textsc{Own}} & \multicolumn{1}{c}{\textsc{Read}} & \multicolumn{1}{c}{\textsc{Ans}} & \multicolumn{1}{c}{\textsc{Own}} & \multicolumn{1}{c}{\textsc{Read}} & \multicolumn{1}{c}{\textsc{Ans}} & \multicolumn{1}{c}{\textsc{Own}} & \multicolumn{1}{c}{clin.} & \multicolumn{1}{c}{COCO} & \multicolumn{1}{c}{ratio}\\
%DIFDELCMD < \midrule
%DIFDELCMD < Lingshu 32B & \cellcolor{cfSage2}0.08$^{3}$ & \cellcolor{cfSage3}0.78$^{6}$ & \cellcolor{cfSlate2}-0.02$^{2}$ & \cellcolor{cfSage4}0.20$^{4}$ & \cellcolor{cfSage3}0.87$^{5}$ & \cellcolor{cfTerra2}+0.04$^{5}$ & \cellcolor{cfSage1}0.03$^{1}$ & \cellcolor{cfSage4}0.95$^{5}$ & \cellcolor{cfTerra4}+0.28$^{6}$ & \cellcolor{cfTerra4}58\% & \cellcolor{cfTerra4}100\% & 58\% \\
%DIFDELCMD < CheXagent-2-3B & \cellcolor{cfSage3}0.13$^{5}$ & \cellcolor{cfSage3}0.86$^{6}$ & \cellcolor{cfTerra1}+0.00$^{2}$ & \cellcolor{cfSage4}0.21$^{5}$ & \cellcolor{cfSage4}0.95$^{5}$ & \cellcolor{cfTerra1}+0.01$^{4}$ & \cellcolor{cfSage3}0.13$^{4}$ & \cellcolor{cfSage3}0.81$^{5}$ & \cellcolor{cfTerra2}+0.03$^{2}$ & \cellcolor{cfTerra4}50\% & \cellcolor{cfTerra2}33\% & 150\% \\
%DIFDELCMD < Qwen2.5-VL-72B & \cellcolor{cfSage2}0.07$^{2}$ & \cellcolor{cfSage2}0.62$^{4}$ & \cellcolor{cfSlate2}-0.03$^{3}$ & \cellcolor{cfSage3}0.20$^{4}$ & \cellcolor{cfSage2}0.71$^{4}$ & \cellcolor{cfSlate2}-0.09$^{1}$ & \cellcolor{cfSage1}0.01$^{1}$ & \cellcolor{cfSage4}0.99$^{5}$ & \cellcolor{cfTerra4}+0.66$^{6}$ & \cellcolor{cfTerra3}33\% & \cellcolor{cfTerra4}100\% & 33\% \\
%DIFDELCMD < Lingshu 7B & \cellcolor{cfSage2}0.08$^{3}$ & \cellcolor{cfSage3}0.76$^{6}$ & \cellcolor{cfTerra2}+0.02$^{1}$ & \cellcolor{cfSage4}0.21$^{4}$ & \cellcolor{cfSage3}0.86$^{5}$ & \cellcolor{cfTerra2}+0.04$^{3}$ & \cellcolor{cfSage1}0.02$^{1}$ & \cellcolor{cfSage4}0.97$^{5}$ & \cellcolor{cfTerra4}+0.68$^{6}$ & \cellcolor{cfTerra3}33\% & \cellcolor{cfTerra4}100\% & 33\% \\
%DIFDELCMD < Qwen3-VL-4B & \cellcolor{cfSage2}0.08$^{3}$ & \cellcolor{cfSage2}0.71$^{6}$ & \cellcolor{cfSlate3}-0.11$^{0}$ & \cellcolor{cfSage4}0.20$^{4}$ & \cellcolor{cfSage2}0.74$^{4}$ & \cellcolor{cfSlate1}-0.01$^{2}$ & \cellcolor{cfSage1}0.04$^{4}$ & \cellcolor{cfSage4}0.99$^{5}$ & \cellcolor{cfTerra4}+0.34$^{6}$ & \cellcolor{cfTerra3}17\% & \cellcolor{cfTerra4}100\% & 17\% \\
%DIFDELCMD < Qwen3-VL-32B & \cellcolor{cfSage2}0.10$^{2}$ & \cellcolor{cfSage2}0.70$^{6}$ & \cellcolor{cfSlate2}-0.02$^{1}$ & \cellcolor{cfSage3}0.17$^{5}$ & \cellcolor{cfSage3}0.82$^{5}$ & \cellcolor{cfSlate2}-0.02$^{1}$ & \cellcolor{cfSage2}0.05$^{4}$ & \cellcolor{cfSage4}0.99$^{5}$ & \cellcolor{cfTerra3}+0.15$^{6}$ & \cellcolor{cfTerra3}17\% & \cellcolor{cfTerra4}100\% & 17\% \\
%DIFDELCMD < InternVL3.5-38B & \cellcolor{cfSage2}0.10$^{5}$ & \cellcolor{cfSage3}0.78$^{6}$ & \cellcolor{cfSlate1}-0.00$^{0}$ & \cellcolor{cfSage3}0.17$^{4}$ & \cellcolor{cfSage3}0.86$^{5}$ & \cellcolor{cfSlate1}-0.00$^{2}$ & \cellcolor{cfSage2}0.08$^{4}$ & \cellcolor{cfSage4}0.99$^{5}$ & \cellcolor{cfTerra1}+0.01$^{6}$ & \cellcolor{cfTerra3}17\% & \cellcolor{cfTerra4}100\% & 17\% \\
%DIFDELCMD < Gemma 3 12B & \cellcolor{cfSage2}0.07$^{2}$ & \cellcolor{cfSage2}0.62$^{5}$ & \cellcolor{cfSlate4}-0.37$^{1}$ & \cellcolor{cfSage3}0.18$^{4}$ & \cellcolor{cfSage1}0.52$^{0}$ & \cellcolor{cfSlate4}-0.38$^{1}$ & \cellcolor{cfSage2}0.07$^{5}$ & \cellcolor{cfSage4}0.95$^{5}$ & \cellcolor{cfTerra4}+0.44$^{6}$ & \cellcolor{cfTerra3}17\% & \cellcolor{cfTerra4}100\% & 17\% \\
%DIFDELCMD < MedGemma 4B & \cellcolor{cfSage3}0.12$^{3}$ & \cellcolor{cfSage3}0.80$^{6}$ & \cellcolor{cfSlate3}-0.13$^{0}$ & \cellcolor{cfSage4}0.21$^{4}$ & \cellcolor{cfSage3}0.87$^{5}$ & \cellcolor{cfSlate3}-0.16$^{2}$ & \cellcolor{cfSage2}0.06$^{4}$ & \cellcolor{cfSage4}0.94$^{5}$ & \cellcolor{cfTerra3}+0.23$^{6}$ & \cellcolor{cfTerra3}17\% & \cellcolor{cfTerra4}100\% & 17\% \\
%DIFDELCMD < InternVL3.5-14B & \cellcolor{cfSage2}0.07$^{4}$ & \cellcolor{cfSage2}0.73$^{6}$ & \cellcolor{cfSlate1}-0.00$^{1}$ & \cellcolor{cfSage3}0.16$^{4}$ & \cellcolor{cfSage3}0.87$^{5}$ & \cellcolor{cfTerra1}+0.00$^{1}$ & \cellcolor{cfSage3}0.12$^{4}$ & \cellcolor{cfSage4}0.98$^{5}$ & \cellcolor{cfTerra1}+0.02$^{5}$ & \cellcolor{cfTerra3}17\% & \cellcolor{cfTerra3}83\% & 20\% \\
%DIFDELCMD < Qwen2.5-VL-32B & \cellcolor{cfSage2}0.09$^{3}$ & \cellcolor{cfSage2}0.63$^{4}$ & \cellcolor{cfSlate2}-0.06$^{1}$ & \cellcolor{cfSage3}0.16$^{4}$ & \cellcolor{cfSage2}0.70$^{4}$ & \cellcolor{cfSlate3}-0.11$^{0}$ & \cellcolor{cfSage1}0.02$^{1}$ & \cellcolor{cfSage4}0.99$^{5}$ & \cellcolor{cfTerra4}+0.38$^{6}$ & \cellcolor{cfTerra2}8\% & \cellcolor{cfTerra4}100\% & 8\% \\
%DIFDELCMD < InternVL3.5-8B & \cellcolor{cfSage2}0.06$^{2}$ & \cellcolor{cfSage2}0.73$^{6}$ & \cellcolor{cfSlate1}-0.01$^{0}$ & \cellcolor{cfSage3}0.19$^{4}$ & \cellcolor{cfSage3}0.86$^{5}$ & \cellcolor{cfSlate1}-0.00$^{1}$ & \cellcolor{cfSage2}0.12$^{5}$ & \cellcolor{cfSage4}0.98$^{5}$ & \cellcolor{cfTerra2}+0.03$^{6}$ & \cellcolor{cfTerra2}8\% & \cellcolor{cfTerra4}100\% & 8\% \\
%DIFDELCMD < HuatuoGPT-Vision-7B & \cellcolor{cfSage2}0.08$^{3}$ & \cellcolor{cfSage2}0.71$^{6}$ & \cellcolor{cfSlate2}-0.03$^{1}$ & \cellcolor{cfSage3}0.19$^{4}$ & \cellcolor{cfSage3}0.78$^{4}$ & \cellcolor{cfSlate1}-0.01$^{0}$ & \cellcolor{cfSage2}0.07$^{4}$ & \cellcolor{cfSage4}0.99$^{5}$ & \cellcolor{cfTerra1}+0.01$^{4}$ & \cellcolor{cfTerra2}8\% & \cellcolor{cfTerra3}67\% & 12\% \\
%DIFDELCMD < LLaVA-Med 1.5 7B & \cellcolor{cfSage2}0.08$^{3}$ & \cellcolor{cfSage1}0.53$^{0}$ & \cellcolor{cfSlate1}-0.01$^{0}$ & \cellcolor{cfSage3}0.19$^{4}$ & \cellcolor{cfSage1}0.59$^{2}$ & \cellcolor{cfSlate1}-0.01$^{1}$ & \cellcolor{cfSage2}0.07$^{4}$ & \cellcolor{cfSage3}0.76$^{5}$ & \cellcolor{cfTerra1}+0.00$^{1}$ & \cellcolor{cfTerra2}8\% & \cellcolor{cfTerra2}17\% & 50\% \\
%DIFDELCMD < CheXagent-8B & \cellcolor{cfSage3}0.16$^{5}$ & \cellcolor{cfSage2}0.73$^{6}$ & \cellcolor{cfSlate1}-0.00$^{0}$ & \cellcolor{cfSage4}0.23$^{5}$ & \cellcolor{cfSage3}0.88$^{5}$ & \cellcolor{cfSlate1}-0.02$^{1}$ & \cellcolor{cfSage2}0.05$^{3}$ & \cellcolor{cfSage1}0.52$^{1}$ & \cellcolor{cfSlate2}-0.03$^{0}$ & \cellcolor{cfTerra2}8\% & \cellcolor{cfTerra1}0\% & -- \\
%DIFDELCMD < Qwen2.5-VL-3B & \cellcolor{cfSage2}0.07$^{2}$ & \cellcolor{cfSage1}0.58$^{2}$ & \cellcolor{cfSlate2}-0.08$^{0}$ & \cellcolor{cfSage3}0.17$^{4}$ & \cellcolor{cfSage1}0.57$^{2}$ & \cellcolor{cfSlate2}-0.07$^{0}$ & \cellcolor{cfSage1}0.05$^{4}$ & \cellcolor{cfSage4}0.97$^{5}$ & \cellcolor{cfTerra4}+0.52$^{6}$ & \cellcolor{cfTerra1}0\% & \cellcolor{cfTerra4}100\% & 0\% \\
%DIFDELCMD < Qwen2.5-VL-7B & \cellcolor{cfSage2}0.08$^{3}$ & \cellcolor{cfSage2}0.62$^{3}$ & \cellcolor{cfSlate3}-0.14$^{0}$ & \cellcolor{cfSage3}0.20$^{4}$ & \cellcolor{cfSage2}0.65$^{2}$ & \cellcolor{cfSlate2}-0.08$^{0}$ & \cellcolor{cfSage1}0.02$^{1}$ & \cellcolor{cfSage4}0.98$^{5}$ & \cellcolor{cfTerra4}+0.67$^{6}$ & \cellcolor{cfTerra1}0\% & \cellcolor{cfTerra4}100\% & 0\% \\
%DIFDELCMD < Qwen3-VL-8B & \cellcolor{cfSage2}0.10$^{4}$ & \cellcolor{cfSage2}0.69$^{6}$ & \cellcolor{cfSlate1}-0.01$^{0}$ & \cellcolor{cfSage3}0.17$^{4}$ & \cellcolor{cfSage2}0.75$^{4}$ & \cellcolor{cfSlate2}-0.04$^{0}$ & \cellcolor{cfSage1}0.04$^{3}$ & \cellcolor{cfSage4}0.99$^{5}$ & \cellcolor{cfTerra4}+0.58$^{6}$ & \cellcolor{cfTerra1}0\% & \cellcolor{cfTerra4}100\% & 0\% \\
%DIFDELCMD < Gemma 3 27B & \cellcolor{cfSage2}0.07$^{2}$ & \cellcolor{cfSage2}0.60$^{3}$ & \cellcolor{cfSlate2}-0.03$^{0}$ & \cellcolor{cfSage3}0.18$^{4}$ & \cellcolor{cfSage1}0.55$^{2}$ & \cellcolor{cfSlate1}-0.01$^{0}$ & \cellcolor{cfSage2}0.07$^{5}$ & \cellcolor{cfSage4}0.96$^{5}$ & \cellcolor{cfTerra4}+0.31$^{6}$ & \cellcolor{cfTerra1}0\% & \cellcolor{cfTerra4}100\% & 0\% \\
%DIFDELCMD < Gemma 3 4B & \cellcolor{cfSage2}0.07$^{2}$ & \cellcolor{cfSage1}0.56$^{2}$ & \cellcolor{cfSlate3}-0.11$^{0}$ & \cellcolor{cfSage3}0.18$^{4}$ & \cellcolor{cfSage1}0.53$^{1}$ & \cellcolor{cfSlate4}-0.27$^{0}$ & \cellcolor{cfSage2}0.07$^{5}$ & \cellcolor{cfSage4}0.94$^{5}$ & \cellcolor{cfTerra3}+0.24$^{5}$ & \cellcolor{cfTerra1}0\% & \cellcolor{cfTerra3}83\% & 0\% \\
%DIFDELCMD < LLaVA-1.5-7B & \cellcolor{cfSage2}0.08$^{2}$ & \cellcolor{cfSage1}0.51$^{0}$ & \cellcolor{cfSlate2}-0.04$^{0}$ & \cellcolor{cfSage3}0.17$^{4}$ & \cellcolor{cfSage1}0.52$^{0}$ & \cellcolor{cfSlate2}-0.03$^{0}$ & \cellcolor{cfSage4}0.28$^{5}$ & \cellcolor{cfSage4}0.97$^{5}$ & \cellcolor{cfTerra2}+0.03$^{5}$ & \cellcolor{cfTerra1}0\% & \cellcolor{cfTerra3}83\% & 0\% \\
%DIFDELCMD < MedGemma 27B & \cellcolor{cfSage3}0.12$^{3}$ & \cellcolor{cfSage3}0.80$^{6}$ & \cellcolor{cfSlate3}-0.14$^{0}$ & \cellcolor{cfSage4}0.21$^{4}$ & \cellcolor{cfSage3}0.87$^{5}$ & \cellcolor{cfSlate2}-0.06$^{0}$ & \cellcolor{cfSage2}0.06$^{4}$ & \cellcolor{cfSage4}0.97$^{5}$ & \cellcolor{cfTerra3}+0.12$^{4}$ & \cellcolor{cfTerra1}0\% & \cellcolor{cfTerra3}67\% & 0\% \\
%DIFDELCMD < LLaVA-1.5-13B & \cellcolor{cfSage2}0.08$^{2}$ & \cellcolor{cfSage1}0.53$^{0}$ & \cellcolor{cfSlate2}-0.02$^{0}$ & \cellcolor{cfSage3}0.17$^{4}$ & \cellcolor{cfSage1}0.51$^{1}$ & \cellcolor{cfSlate1}-0.01$^{0}$ & \cellcolor{cfSage4}0.28$^{5}$ & \cellcolor{cfSage4}0.96$^{5}$ & \cellcolor{cfTerra2}+0.02$^{2}$ & \cellcolor{cfTerra1}0\% & \cellcolor{cfTerra2}33\% & 0\% \\
%DIFDELCMD < Llama 3.2 Vision 11B & \cellcolor{cfSage2}0.08$^{2}$ & \cellcolor{cfSage2}0.65$^{4}$ & \cellcolor{cfSlate2}-0.03$^{0}$ & \cellcolor{cfSage3}0.16$^{4}$ & \cellcolor{cfSage2}0.73$^{4}$ & \cellcolor{cfSlate2}-0.04$^{0}$ & \cellcolor{cfSage3}0.12$^{5}$ & \cellcolor{cfSage4}0.95$^{5}$ & \cellcolor{cfSlate1}-0.01$^{1}$ & \cellcolor{cfTerra1}0\% & \cellcolor{cfTerra2}17\% & 0\% \\
%DIFDELCMD < LLaVA-Rad-7B & \cellcolor{cfSage3}0.15$^{4}$ & \cellcolor{cfSage2}0.71$^{5}$ & \cellcolor{cfSlate1}-0.01$^{0}$ & \cellcolor{cfSage4}0.27$^{5}$ & \cellcolor{cfSage3}0.87$^{5}$ & \cellcolor{cfSlate1}-0.02$^{0}$ & \cellcolor{cfSage3}0.12$^{4}$ & \cellcolor{cfSage1}0.50$^{0}$ & \cellcolor{cfSlate2}-0.03$^{0}$ & \cellcolor{cfTerra1}0\% & \cellcolor{cfTerra1}0\% & -- \\
%DIFDELCMD < \midrule
%DIFDELCMD < \textbf{All cells} & 0.09 & 0.68 & -0.06 & 0.19 & 0.74 & -0.05 & 0.08 & 0.92 & +0.23 & 13\% & 75\% & 17\% \\
%DIFDELCMD < \quad count / cells & 74/150 & 110/150 & 13/150 & 104/150 & 89/150 & 25/150 & 90/150 & 116/150 & 113/150 & & & \\
%DIFDELCMD < \bottomrule
%DIFDELCMD < \end{tabular}}
%DIFDELCMD < %%%
\DIFdelendFL \DIFaddbeginFL \resizebox{\linewidth}{!}{%
\begin{tabular}{l@{\hspace{3pt}}ccc@{\hspace{4pt}}ccc@{\hspace{4pt}}ccc@{\hspace{4pt}}ccc}
\toprule
\multirow{2}{*}{Checkpoint} & \multicolumn{3}{c}{NIH ChestX-ray14} & \multicolumn{3}{c}{CheXpert Plus} & \multicolumn{3}{c}{COCO (control)} & \multicolumn{3}{c}{Owned share}\\
\cmidrule(lr){2-4}\cmidrule(lr){5-7}\cmidrule(lr){8-10}\cmidrule(lr){11-13}
& \multicolumn{1}{c}{\textsc{Read}} & \multicolumn{1}{c}{\textsc{Ans}} & \multicolumn{1}{c}{\textsc{Own}} & \multicolumn{1}{c}{\textsc{Read}} & \multicolumn{1}{c}{\textsc{Ans}} & \multicolumn{1}{c}{\textsc{Own}} & \multicolumn{1}{c}{\textsc{Read}} & \multicolumn{1}{c}{\textsc{Ans}} & \multicolumn{1}{c}{\textsc{Own}} & \multicolumn{1}{c}{clin.} & \multicolumn{1}{c}{COCO} & \multicolumn{1}{c}{ratio}\\
\midrule
Lingshu 32B & \cellcolor{cfSage2}0.08$^{3}$ & \cellcolor{cfSage3}0.78$^{6}$ & \cellcolor{cfSlate2}-0.02$^{2}$ & \cellcolor{cfSage4}0.20$^{4}$ & \cellcolor{cfSage3}0.87$^{5}$ & \cellcolor{cfTerra2}+0.04$^{5}$ & \cellcolor{cfSage1}0.03$^{1}$ & \cellcolor{cfSage4}0.95$^{5}$ & \cellcolor{cfTerra4}+0.28$^{6}$ & \cellcolor{cfTerra4}58\% & \cellcolor{cfTerra4}100\% & 58\% \\
CheXagent-2-3B & \cellcolor{cfSage3}0.13$^{5}$ & \cellcolor{cfSage3}0.86$^{6}$ & \cellcolor{cfTerra1}+0.00$^{2}$ & \cellcolor{cfSage4}0.21$^{5}$ & \cellcolor{cfSage4}0.95$^{5}$ & \cellcolor{cfTerra1}+0.01$^{4}$ & \cellcolor{cfSage3}0.13$^{4}$ & \cellcolor{cfSage3}0.81$^{5}$ & \cellcolor{cfTerra2}+0.03$^{2}$ & \cellcolor{cfTerra4}50\% & \cellcolor{cfTerra2}33\% & 150\% \\
Qwen2.5-VL-72B & \cellcolor{cfSage2}0.07$^{2}$ & \cellcolor{cfSage2}0.62$^{4}$ & \cellcolor{cfSlate2}-0.03$^{3}$ & \cellcolor{cfSage3}0.20$^{4}$ & \cellcolor{cfSage2}0.71$^{4}$ & \cellcolor{cfSlate2}-0.09$^{1}$ & \cellcolor{cfSage1}0.01$^{1}$ & \cellcolor{cfSage4}0.99$^{5}$ & \cellcolor{cfTerra4}+0.66$^{6}$ & \cellcolor{cfTerra3}33\% & \cellcolor{cfTerra4}100\% & 33\% \\
Lingshu 7B & \cellcolor{cfSage2}0.08$^{3}$ & \cellcolor{cfSage3}0.76$^{6}$ & \cellcolor{cfTerra2}+0.02$^{1}$ & \cellcolor{cfSage4}0.21$^{4}$ & \cellcolor{cfSage3}0.86$^{5}$ & \cellcolor{cfTerra2}+0.04$^{3}$ & \cellcolor{cfSage1}0.02$^{1}$ & \cellcolor{cfSage4}0.97$^{5}$ & \cellcolor{cfTerra4}+0.68$^{6}$ & \cellcolor{cfTerra3}33\% & \cellcolor{cfTerra4}100\% & 33\% \\
Qwen3-VL-4B & \cellcolor{cfSage2}0.08$^{3}$ & \cellcolor{cfSage2}0.71$^{6}$ & \cellcolor{cfSlate3}-0.11$^{0}$ & \cellcolor{cfSage4}0.20$^{4}$ & \cellcolor{cfSage2}0.74$^{4}$ & \cellcolor{cfSlate1}-0.01$^{2}$ & \cellcolor{cfSage1}0.04$^{4}$ & \cellcolor{cfSage4}0.99$^{5}$ & \cellcolor{cfTerra4}+0.34$^{6}$ & \cellcolor{cfTerra3}17\% & \cellcolor{cfTerra4}100\% & 17\% \\
Qwen3-VL-32B & \cellcolor{cfSage2}0.10$^{2}$ & \cellcolor{cfSage2}0.70$^{6}$ & \cellcolor{cfSlate2}-0.02$^{1}$ & \cellcolor{cfSage3}0.17$^{5}$ & \cellcolor{cfSage3}0.82$^{5}$ & \cellcolor{cfSlate2}-0.02$^{1}$ & \cellcolor{cfSage2}0.05$^{4}$ & \cellcolor{cfSage4}0.99$^{5}$ & \cellcolor{cfTerra3}+0.15$^{6}$ & \cellcolor{cfTerra3}17\% & \cellcolor{cfTerra4}100\% & 17\% \\
InternVL3.5-38B & \cellcolor{cfSage2}0.10$^{5}$ & \cellcolor{cfSage3}0.78$^{6}$ & \cellcolor{cfSlate1}-0.00$^{0}$ & \cellcolor{cfSage3}0.17$^{4}$ & \cellcolor{cfSage3}0.86$^{5}$ & \cellcolor{cfSlate1}-0.00$^{2}$ & \cellcolor{cfSage2}0.08$^{4}$ & \cellcolor{cfSage4}0.99$^{5}$ & \cellcolor{cfTerra1}+0.01$^{6}$ & \cellcolor{cfTerra3}17\% & \cellcolor{cfTerra4}100\% & 17\% \\
Gemma 3 12B & \cellcolor{cfSage2}0.07$^{2}$ & \cellcolor{cfSage2}0.62$^{5}$ & \cellcolor{cfSlate4}-0.37$^{1}$ & \cellcolor{cfSage3}0.18$^{4}$ & \cellcolor{cfSage1}0.52$^{0}$ & \cellcolor{cfSlate4}-0.38$^{1}$ & \cellcolor{cfSage2}0.07$^{5}$ & \cellcolor{cfSage4}0.95$^{5}$ & \cellcolor{cfTerra4}+0.44$^{6}$ & \cellcolor{cfTerra3}17\% & \cellcolor{cfTerra4}100\% & 17\% \\
MedGemma 4B & \cellcolor{cfSage3}0.12$^{3}$ & \cellcolor{cfSage3}0.80$^{6}$ & \cellcolor{cfSlate3}-0.13$^{0}$ & \cellcolor{cfSage4}0.21$^{4}$ & \cellcolor{cfSage3}0.87$^{5}$ & \cellcolor{cfSlate3}-0.16$^{2}$ & \cellcolor{cfSage2}0.06$^{4}$ & \cellcolor{cfSage4}0.94$^{5}$ & \cellcolor{cfTerra3}+0.23$^{6}$ & \cellcolor{cfTerra3}17\% & \cellcolor{cfTerra4}100\% & 17\% \\
InternVL3.5-14B & \cellcolor{cfSage2}0.07$^{4}$ & \cellcolor{cfSage2}0.73$^{6}$ & \cellcolor{cfSlate1}-0.00$^{1}$ & \cellcolor{cfSage3}0.16$^{4}$ & \cellcolor{cfSage3}0.87$^{5}$ & \cellcolor{cfTerra1}+0.00$^{1}$ & \cellcolor{cfSage3}0.12$^{4}$ & \cellcolor{cfSage4}0.98$^{5}$ & \cellcolor{cfTerra1}+0.02$^{5}$ & \cellcolor{cfTerra3}17\% & \cellcolor{cfTerra3}83\% & 20\% \\
Qwen2.5-VL-32B & \cellcolor{cfSage2}0.09$^{3}$ & \cellcolor{cfSage2}0.63$^{4}$ & \cellcolor{cfSlate2}-0.06$^{1}$ & \cellcolor{cfSage3}0.16$^{4}$ & \cellcolor{cfSage2}0.70$^{4}$ & \cellcolor{cfSlate3}-0.11$^{0}$ & \cellcolor{cfSage1}0.02$^{1}$ & \cellcolor{cfSage4}0.99$^{5}$ & \cellcolor{cfTerra4}+0.38$^{6}$ & \cellcolor{cfTerra2}8\% & \cellcolor{cfTerra4}100\% & 8\% \\
InternVL3.5-8B & \cellcolor{cfSage2}0.06$^{2}$ & \cellcolor{cfSage2}0.73$^{6}$ & \cellcolor{cfSlate1}-0.01$^{0}$ & \cellcolor{cfSage3}0.19$^{4}$ & \cellcolor{cfSage3}0.86$^{5}$ & \cellcolor{cfSlate1}-0.00$^{1}$ & \cellcolor{cfSage2}0.12$^{5}$ & \cellcolor{cfSage4}0.98$^{5}$ & \cellcolor{cfTerra2}+0.03$^{6}$ & \cellcolor{cfTerra2}8\% & \cellcolor{cfTerra4}100\% & 8\% \\
HuatuoGPT-Vision-7B & \cellcolor{cfSage2}0.08$^{3}$ & \cellcolor{cfSage2}0.71$^{6}$ & \cellcolor{cfSlate2}-0.03$^{1}$ & \cellcolor{cfSage3}0.19$^{4}$ & \cellcolor{cfSage3}0.78$^{4}$ & \cellcolor{cfSlate1}-0.01$^{0}$ & \cellcolor{cfSage2}0.07$^{4}$ & \cellcolor{cfSage4}0.99$^{5}$ & \cellcolor{cfTerra1}+0.01$^{4}$ & \cellcolor{cfTerra2}8\% & \cellcolor{cfTerra3}67\% & 12\% \\
LLaVA-Med 1.5 7B & \cellcolor{cfSage2}0.08$^{3}$ & \cellcolor{cfSage1}0.53$^{0}$ & \cellcolor{cfSlate1}-0.01$^{0}$ & \cellcolor{cfSage3}0.19$^{4}$ & \cellcolor{cfSage1}0.59$^{2}$ & \cellcolor{cfSlate1}-0.01$^{1}$ & \cellcolor{cfSage2}0.07$^{4}$ & \cellcolor{cfSage3}0.76$^{5}$ & \cellcolor{cfTerra1}+0.00$^{1}$ & \cellcolor{cfTerra2}8\% & \cellcolor{cfTerra2}17\% & 50\% \\
CheXagent-8B & \cellcolor{cfSage3}0.16$^{5}$ & \cellcolor{cfSage2}0.73$^{6}$ & \cellcolor{cfSlate1}-0.00$^{0}$ & \cellcolor{cfSage4}0.23$^{5}$ & \cellcolor{cfSage3}0.88$^{5}$ & \cellcolor{cfSlate1}-0.02$^{1}$ & \cellcolor{cfSage2}0.05$^{3}$ & \cellcolor{cfSage1}0.52$^{1}$ & \cellcolor{cfSlate2}-0.03$^{0}$ & \cellcolor{cfTerra2}8\% & \cellcolor{cfTerra1}0\% & -- \\
Qwen2.5-VL-3B & \cellcolor{cfSage2}0.07$^{2}$ & \cellcolor{cfSage1}0.58$^{2}$ & \cellcolor{cfSlate2}-0.08$^{0}$ & \cellcolor{cfSage3}0.17$^{4}$ & \cellcolor{cfSage1}0.57$^{2}$ & \cellcolor{cfSlate2}-0.07$^{0}$ & \cellcolor{cfSage1}0.05$^{4}$ & \cellcolor{cfSage4}0.97$^{5}$ & \cellcolor{cfTerra4}+0.52$^{6}$ & \cellcolor{cfTerra1}0\% & \cellcolor{cfTerra4}100\% & 0\% \\
Qwen2.5-VL-7B & \cellcolor{cfSage2}0.08$^{3}$ & \cellcolor{cfSage2}0.62$^{3}$ & \cellcolor{cfSlate3}-0.14$^{0}$ & \cellcolor{cfSage3}0.20$^{4}$ & \cellcolor{cfSage2}0.65$^{2}$ & \cellcolor{cfSlate2}-0.08$^{0}$ & \cellcolor{cfSage1}0.02$^{1}$ & \cellcolor{cfSage4}0.98$^{5}$ & \cellcolor{cfTerra4}+0.67$^{6}$ & \cellcolor{cfTerra1}0\% & \cellcolor{cfTerra4}100\% & 0\% \\
Qwen3-VL-8B & \cellcolor{cfSage2}0.10$^{4}$ & \cellcolor{cfSage2}0.69$^{6}$ & \cellcolor{cfSlate1}-0.01$^{0}$ & \cellcolor{cfSage3}0.17$^{4}$ & \cellcolor{cfSage2}0.75$^{4}$ & \cellcolor{cfSlate2}-0.04$^{0}$ & \cellcolor{cfSage1}0.04$^{3}$ & \cellcolor{cfSage4}0.99$^{5}$ & \cellcolor{cfTerra4}+0.58$^{6}$ & \cellcolor{cfTerra1}0\% & \cellcolor{cfTerra4}100\% & 0\% \\
Gemma 3 27B & \cellcolor{cfSage2}0.07$^{2}$ & \cellcolor{cfSage2}0.60$^{3}$ & \cellcolor{cfSlate2}-0.03$^{0}$ & \cellcolor{cfSage3}0.18$^{4}$ & \cellcolor{cfSage1}0.55$^{2}$ & \cellcolor{cfSlate1}-0.01$^{0}$ & \cellcolor{cfSage2}0.07$^{5}$ & \cellcolor{cfSage4}0.96$^{5}$ & \cellcolor{cfTerra4}+0.31$^{6}$ & \cellcolor{cfTerra1}0\% & \cellcolor{cfTerra4}100\% & 0\% \\
Gemma 3 4B & \cellcolor{cfSage2}0.07$^{2}$ & \cellcolor{cfSage1}0.56$^{2}$ & \cellcolor{cfSlate3}-0.11$^{0}$ & \cellcolor{cfSage3}0.18$^{4}$ & \cellcolor{cfSage1}0.53$^{1}$ & \cellcolor{cfSlate4}-0.27$^{0}$ & \cellcolor{cfSage2}0.07$^{5}$ & \cellcolor{cfSage4}0.94$^{5}$ & \cellcolor{cfTerra3}+0.24$^{5}$ & \cellcolor{cfTerra1}0\% & \cellcolor{cfTerra3}83\% & 0\% \\
LLaVA-1.5-7B & \cellcolor{cfSage2}0.08$^{2}$ & \cellcolor{cfSage1}0.51$^{0}$ & \cellcolor{cfSlate2}-0.04$^{0}$ & \cellcolor{cfSage3}0.17$^{4}$ & \cellcolor{cfSage1}0.52$^{0}$ & \cellcolor{cfSlate2}-0.03$^{0}$ & \cellcolor{cfSage4}0.28$^{5}$ & \cellcolor{cfSage4}0.97$^{5}$ & \cellcolor{cfTerra2}+0.03$^{5}$ & \cellcolor{cfTerra1}0\% & \cellcolor{cfTerra3}83\% & 0\% \\
MedGemma 27B & \cellcolor{cfSage3}0.12$^{3}$ & \cellcolor{cfSage3}0.80$^{6}$ & \cellcolor{cfSlate3}-0.14$^{0}$ & \cellcolor{cfSage4}0.21$^{4}$ & \cellcolor{cfSage3}0.87$^{5}$ & \cellcolor{cfSlate2}-0.06$^{0}$ & \cellcolor{cfSage2}0.06$^{4}$ & \cellcolor{cfSage4}0.97$^{5}$ & \cellcolor{cfTerra3}+0.12$^{4}$ & \cellcolor{cfTerra1}0\% & \cellcolor{cfTerra3}67\% & 0\% \\
LLaVA-1.5-13B & \cellcolor{cfSage2}0.08$^{2}$ & \cellcolor{cfSage1}0.53$^{0}$ & \cellcolor{cfSlate2}-0.02$^{0}$ & \cellcolor{cfSage3}0.17$^{4}$ & \cellcolor{cfSage1}0.51$^{1}$ & \cellcolor{cfSlate1}-0.01$^{0}$ & \cellcolor{cfSage4}0.28$^{5}$ & \cellcolor{cfSage4}0.96$^{5}$ & \cellcolor{cfTerra2}+0.02$^{2}$ & \cellcolor{cfTerra1}0\% & \cellcolor{cfTerra2}33\% & 0\% \\
Llama 3.2 Vision 11B & \cellcolor{cfSage2}0.08$^{2}$ & \cellcolor{cfSage2}0.65$^{4}$ & \cellcolor{cfSlate2}-0.03$^{0}$ & \cellcolor{cfSage3}0.16$^{4}$ & \cellcolor{cfSage2}0.73$^{4}$ & \cellcolor{cfSlate2}-0.04$^{0}$ & \cellcolor{cfSage3}0.12$^{5}$ & \cellcolor{cfSage4}0.95$^{5}$ & \cellcolor{cfSlate1}-0.01$^{1}$ & \cellcolor{cfTerra1}0\% & \cellcolor{cfTerra2}17\% & 0\% \\
LLaVA-Rad-7B & \cellcolor{cfSage3}0.15$^{4}$ & \cellcolor{cfSage2}0.71$^{5}$ & \cellcolor{cfSlate1}-0.01$^{0}$ & \cellcolor{cfSage4}0.27$^{5}$ & \cellcolor{cfSage3}0.87$^{5}$ & \cellcolor{cfSlate1}-0.02$^{0}$ & \cellcolor{cfSage3}0.12$^{4}$ & \cellcolor{cfSage1}0.50$^{0}$ & \cellcolor{cfSlate2}-0.03$^{0}$ & \cellcolor{cfTerra1}0\% & \cellcolor{cfTerra1}0\% & -- \\
\midrule
\textbf{All cells} & 0.09 & 0.68 & -0.06 & 0.19 & 0.74 & -0.05 & 0.08 & 0.92 & +0.23 & 13\% & 75\% & 17\% \\
\quad count / cells & 74/150 & 110/150 & 13/150 & 104/150 & 89/150 & 25/150 & 90/150 & 116/150 & 113/150 & & & \\
\bottomrule
\end{tabular}}
\DIFaddendFL \end{table}

\begin{figure}[t]
    \centering
    \resizebox{\textwidth}{!}{\includegraphics{figures/fig2_overview.pdf}}
    \caption{\textbf{Reading, answering, and writing across the grid.} (a) Readable, answerable, and owned cells of 150 per dataset, with descriptive 95\% bootstrap intervals over cells. (b) Mean readability against mean ownership per checkpoint. (c) Median ownership and interquartile band over 25 blocks against the dose; the dotted line marks the graded dose $\alpha=+0.25$, and the COCO band is clipped at $0.2$. (d) Cumulative share of cells by the rank of the concept write among the 120 compared effects.}
    \label{fig:overview}
\end{figure}

\subsection{Where the chest write goes}
\label{sec:deficit}

Clinical writes are not selective (Figure~\ref{fig:write-flow}). The written concept's own question receives \cfOwnShareNih\% of the positive answer change on NIH and \cfOwnShareChex\% on CheXpert, against \cfOwnShareCoco\% on COCO. The write matrices show the same structure (Figure~\ref{fig:write-matrices}). The own write is small, with median $W_{q,q}$ of $\cfMedianWqqNih$ on NIH and $\cfMedianWqqChex$ on CheXpert against $\cfMedianWqqCoco$ on COCO, and the competing directions share what effect there is: the Atelectasis, Effusion, and Consolidation directions raise each other's answers by similar amounts, and on the chest sets a row of $W$ is flat or peaks off the diagonal, whereas on COCO the diagonal dominates. At every Qwen2.5-VL size a competing direction raises the Effusion answer more than the Effusion direction does (Cardiomegaly at 3B, Nodule at 7B and 32B, Atelectasis at 72B). Single images show the pattern (Figure~\ref{fig:example}): on a radiograph, the Effusion write raises the Effusion answer from $\cfExChestClean$ to $\cfExChestOwn$ and the Nodule write raises it further, to $\cfExChestComp$, while on a natural image the dog write moves the dog answer from $\cfExCocoClean$ to $\cfExCocoOwn$ and nothing else. The chest cells that are owned are also fragile: an effect floor of $0.05$ on both $W_{q,q}$ and $O_q$ keeps \cfChestFloored{} of the \cfChestOwned{} owned chest cells against \cfCocoFloored{} of the \cfCocoOwned{} on COCO, and ownership survives both probe refits in \cfRefitGradeStableNih{} of \cfRefitGradeOwnedNih{} owned NIH cells against \cfRefitGradeStableCoco{} of \cfRefitGradeOwnedCoco{} on COCO (Table~\ref{tab:cf-refit}).

\begin{figure}[t]
    \centering
    \resizebox{\textwidth}{!}{\includegraphics{figures/fig7_write_flow.pdf}}
    \caption{\textbf{Where the write goes.} Bands join each written direction $d$ to the questions $q$ it moves, with width proportional to $F_{q,d}$, the mean over checkpoints of $\max(W_{q,d},0)$. The percentage is the own-question share $\sum_d F_{d,d}/\sum_{q,d} F_{q,d}$.}
    \label{fig:write-flow}
\end{figure}

\begin{figure}[t]
    \centering
    \resizebox{\textwidth}{!}{\includegraphics{figures/fig3_example.pdf}}
    \caption{\textbf{Two images, six questions, three writes.} $P(\mathrm{yes})$ per question for an NIH radiograph (top) and a COCO image (bottom), clean and under the concept write and the strongest competing write ($\alpha=+0.25$, Qwen2.5-VL-7B). The shaded row is the written concept, absent from the image; the dashed line marks $P(\mathrm{yes})=0.5$.}
    \label{fig:example}
\end{figure}

\begin{figure}[t]
    \centering
    \resizebox{\textwidth}{!}{\includegraphics{figures/fig4_write_matrices.pdf}}
    \caption{\textbf{Write matrices.} $W_{q,d}$, the mean change in $P(\mathrm{yes})$ of question $q$ (rows) under the write of direction $d$ (columns), for three checkpoints on NIH (top) and three on COCO (bottom). A corner mark flags owned concepts.}
    \label{fig:write-matrices}
\end{figure}

\subsection{The write mechanism is intact on natural objects}
\label{sec:coco}

On COCO, the same procedure yields the opposite result (Table~\ref{tab:cf-main}, right block). Object ownership holds in \cfCocoOwned{} of \cfCocoOwnedN{} cells (\cfCocoOwnedPct\%), \cfCocoRefMet{} cells meet the steering reference, only \cfCocoCompetitor{} have a stronger competitor, and \cfCocoBlocksFourPlus{} of the \cfCocoBlocks{} checkpoints own at least four of the six objects. Ownership also tracks the dose: on COCO, median ownership rises from $\cfDoseCocoLow$ at $\alpha=-0.1$ to $\cfDoseCocoPeak$ at $\alpha=+0.25$ and stays positive at $\alpha=+0.5$ ($\cfDoseCocoHalf$), while on NIH it stays between $\cfDoseNihMin$ and $\cfDoseNihMax$ at every dose (Figure~\ref{fig:overview}c). The write, direction construction, reference family, and decision rules thus work as intended.

The contrast survives two difficulty controls. Six small or fine-grained COCO categories (handbag, book, backpack, traffic light, knife, and cell phone) remain owned in \cfFgFineOwned{} of \cfFgFineCells{} cells, against \cfFgEasySameOwned{} of \cfFgEasySameCells{} for the easy six in the same \cfFgSameBlocks{} checkpoints (sign test $p=\cfFgSameSignP$), at a median clean-answer AUROC of \cfFgFineAuroc{} against \cfFgEasySameAuroc{}\DIFdelbegin \DIFdel{(Table~\ref{tab:cf-fgobj})}\DIFdelend . On the chest sets, ownership also stays rare where the model already answers the question well: at a clean-answer AUROC of at least \cfStratTopEdge{}, \cfStratChestTopOwned{} of \cfStratChestTopN{} chest cells are owned (\cfStratChestTopPct\%) against \cfStratNatTopOwned{} of \cfStratNatTopN{} natural-image cells (\cfStratNatTopPct\%), a difference of $\cfStratCvnDiff$ $[\cfStratCvnLo,\cfStratCvnHi]$ (Appendix~\DIFdelbegin \DIFdel{\ref{app:cf-crossed} }\DIFdelend \DIFaddbegin \DIFadd{\ref{app:cf-difficulty} }\DIFaddend states the category rule, the bins, and the bootstrap).

\subsection{The reader decides ownership}
\label{sec:reader}

Checkpoints that share a vision tower differ in what they own. Gemma~3 shares one SigLIP tower across its 4B, 12B, and 27B checkpoints, MedGemma one MedSigLIP tower across 4B and 27B, and LLaVA-1.5 one CLIP tower across 7B and 13B. Within each family, pooled features, fitted probes, and lifted write vectors are identical across sizes (Table~\ref{tab:cf-pairs}), so readability is identical by construction and only the connector and language model that read the written representation differ. Steering a feature of the vision tower and letting the language model say what it sees treats that language model as a stable read-out of the tower \citep{ferrando2026explain}; here the read-out is what varies. In \cfPairsFiveOrSix{} of the \cfPairsN{} shared-tower pair--dataset combinations, the paired patient-bootstrap interval for the between-reader difference $\Delta O_q$ excludes zero for five or six of the six concepts (Figure~\ref{fig:readers}). On NIH Effusion, the identical write on the same 600 patients yields an $O_q$ difference of $\cfPairDeltaEffFourTwelve$ $[\cfPairDeltaEffFourTwelveLo,\cfPairDeltaEffFourTwelveHi]$ between Gemma~3 4B and 12B. In 12B, competing directions win the chest answers outright ($O_{\mathrm{Effusion}}=\cfOGemmaTwelveEff$ on both chest sets), and LLaVA-1.5 7B and 13B own no chest concept yet own \cfLlavaSevenCoco{} and \cfLlavaThirteenCoco{} COCO objects with the same write (Appendix~\DIFdelbegin \DIFdel{\ref{app:cf-further}}\DIFdelend \DIFaddbegin \DIFadd{\ref{app:cf-reader}}\DIFaddend ).

A better chest representation does not supply the missing reader. The three checkpoints whose towers are trained on chest radiographs read and answer the findings best in the grid (median probe AUROC \cfSpecTowerProbeAuroc{} on NIH against \cfSpecRestProbeAuroc{} for the other NIH cells; median clean-answer AUROC \cfSpecTowerAnsAuroc{} against \cfSpecMedAnsAuroc{} and \cfSpecGenAnsAuroc{} for medically tuned and general checkpoints), yet they own \cfSpecTowerOwned{} of their \cfSpecTowerN{} chest cells, against \cfSpecMedOwned{} of \cfSpecMedN{} for medically tuned checkpoints on general towers and \cfSpecGenOwned{} of \cfSpecGenN{} for general checkpoints. Scale raises chest ownership within some families, from none at Qwen2.5-VL 3B and 7B to three NIH concepts at 72B and from three to five CheXpert concepts across Lingshu 7B and 32B, but answering alone does not produce it: InternVL3.5 answers every NIH question at every size while its chest writes stay within $\cfInternvlWqqMax$ of zero (Figure~\ref{fig:ladders}). Ownership is therefore a property of the checkpoint that reads the representation, and a probe fitted to the tower alone cannot certify it.

\begin{figure}[t]
    \centering
    \resizebox{\textwidth}{!}{\includegraphics{figures/fig5_same_write.pdf}}
    \caption{\textbf{Same write, different reader.} Ownership $O_q$ per concept on NIH, CheXpert, and COCO for Gemma~3 4B/12B/27B and MedGemma 4B/27B, whose probes and write vectors are identical across sizes within each family. Corner marks flag owned cells and hatching marks ceiling cells.}
    \label{fig:readers}
\end{figure}

\subsection{Non-clinical attributes of the radiograph share the deficit}
\label{sec:attributes}

A non-clinical property of the same film is owned as rarely as a finding. We combine projection, sex, age, and the findings into one nine-direction family. The attribute probes read their labels at least as well as the clinical probes read the findings (median probe AUROC \cfAttrAurocMin{}--\cfAttrAurocMax{} against \cfAttrClinAuroc{}\DIFaddbegin \DIFadd{; Table~\ref{tab:cf-attr}}\DIFaddend ), yet attribute cells are owned in \cfAttrAllAttrOwned{} of \cfAttrAllAttrN{} cases and finding cells in \cfAttrAllClinOwned{} of \cfAttrAllClinN{}, the same rate. At a clean-answer AUROC of at least \cfStratTopEdge{}, attributes are owned in \cfStratAttrTopOwned{} of \cfStratAttrTopN{} cells and findings in \cfStratFindTopOwned{} of \cfStratFindTopN{}\DIFdelbegin \DIFdel{(Table~\ref{tab:cf-attr})}\DIFdelend . The deficit therefore follows the chest-radiograph representation and its reader rather than the concept's clinical status.

\subsection{The direction the reader responds to}
\label{sec:ansdir}

For each question we fit the model's own clean answer margin $\ell_{\mathrm{yes}}-\ell_{\mathrm{no}}$ on training rows from the same projected, train-scaled features by ridge regression, lift the fitted coefficients with Equation~\ref{eq:lift}, and write the result on the test rows exactly as we write a label direction (Table~\ref{tab:cf-ansdir}). This \emph{answer direction} $a_q$ is a handle where the label direction is not. Across the \cfAnsChestN{} chest cells of \cfAnsChestBlocks{} blocks, $a_q$ is owned in \cfAnsChestOwnedA{} and the label direction in \cfAnsChestOwnedLabel{}, and $a_q$ remains owned in \cfAnsdirtChestKept{} of \cfAnsdirtChestKeptN{} chest (cell, template) pairs under the five held-out templates. The features that identify a concept differ from the features that control the output \citep{arad2025saes}; here that separation holds against matched clinical competitors.

The answer direction reads the model, not the label. Over the \cfValAurocAnsChestN{} chest cells, its median AUROC against the report label is $\cfValAurocAnsChest$, against $\cfValAurocProbeChest$ for the probe, while against the model's own clean answer it is $\cfValAurocAnsCleanChest$, against $\cfValAurocProbeCleanChest$ (Table~\ref{tab:cf-validation}). The two directions are far apart on chest radiographs, with a median raw model-space cosine of $\cfAnsCosChestMed$, and substantially aligned on natural images ($\cfAnsCosCocoMed$). The separation persists at four endpoints that fit neither direction (Table~\ref{tab:cf-semend}): the negated question, a forced choice against the strongest competitor in each of its two presentation orders, and the finding word's log-probability in a report continuation. Across the \cfSemCells{} endpoint cells, the answer direction is owned in \cfSemAnswerOwned{} and the label direction in \cfSemLabelOwned{}. On chest radiographs, the direction the model answers along is not the direction the probe fits.

\subsection{Controls and alternatives}
\label{sec:controls}

The reference family behaves as designed: owned concept writes rank 1--6 among the 120 compared effects, and the concept write typically ranks in the top decile on the chest sets (Figure~\ref{fig:overview}d\DIFdelbegin \DIFdel{; Table~\ref{tab:cf-controls}}\DIFdelend ). The deficit survives every alternative we tested; \DIFdelbegin \DIFdel{Table~\ref{tab:cf-ledger} gives the rows, reference , rule, and sham }\DIFdelend \DIFaddbegin \DIFadd{Appendix~\ref{app:cf-ledger} states the reference and rule }\DIFaddend of each of the \cfLedgerAnalyses{} graded analyses, and \DIFdelbegin \DIFdel{Appendices~\ref{app:cf-robustness} and~\ref{app:cf-further} give }\DIFdelend \DIFaddbegin \DIFadd{Appendix~\ref{app:cf-robustness} gives }\DIFaddend the details.
\label{sec:robust-geometry}

\DIFaddbegin \begin{DIFnomarkup}\cfnote{删表}{正文 5.1–5.7 中指向已删表格的引用改为指向附录小节（Controls and robustness、The reader、The answer direction、ledger 小节），只起索引作用的括注直接删去；指向保留的表格（Refit、Pairs、Geometry、Alternative directions、Attributes、Answer direction、Endpoints）的引用仍指向表格。}\end{DIFnomarkup}

\DIFaddend \paragraph{Geometry, scale, and numerics.} The strongest competitor matches the most similar direction or the most co-occurring label only at chance rate (Table~\ref{tab:cf-geometry}). The logit-scale grade agrees with the probability-scale grade in \cfScaleAgree{} of \cfScaleAgreeN{} cells; restricting each cell to unsaturated rows preserves the owned grade in \cfScaleUnsatGradeKept{} of \cfScaleUnsatGradeN{} owned chest cells; refitting the probes on resampled training patients changes $O_q$ by a median of \cfScaleRefitMedian{}; and rescoring in fp32 at batch size one alters \cfPrecisionGradeChanges{} of \cfPrecisionGradeCells{} grades (\DIFdelbegin \DIFdel{Table~\ref{tab:cf-scale}}\DIFdelend \DIFaddbegin \DIFadd{Appendix~\ref{app:cf-robustness}}\DIFaddend ).

\paragraph{Labels.} Re-grading the fitted directions on \cfValidRows{} CheXpert films labelled by radiologist consensus \citep{irvin2019chexpert} reproduces the owned grade in \cfValidOwnedAgree{} of \cfValidCells{} cells, and directions refitted on those labels are owned in \cfValidfitOwnedExpert{} of \cfValidfitCells{} cells, against \cfValidfitOwnedReport{} for the report-label directions (\DIFdelbegin \DIFdel{Table~\ref{tab:cf-valid}}\DIFdelend \DIFaddbegin \DIFadd{Appendix~\ref{app:cf-robustness}}\DIFaddend ).

\paragraph{Directions.} Four alternative constructions, the difference of class means, the Haufe pattern (the training covariance applied to the probe weights; \citealp{haufe2014interpretation}), the normal orthogonalised against the other five normals, and the normal refitted on features residualised against the other five labels \citep{nicolson2025explaining}, are each owned in \cfAltdirChestPctMin{}--\cfAltdirChestPctMax\% of a chest dataset's cells against \cfAltdirCocoPctMin{}--\cfAltdirCocoPctMax\% of COCO cells. Extending the competitor family to every additional dataset label with enough support retains \cfExtcompRetained{} of the \cfExtcompOwned{} owned cells, and writes concentrated on the tokens the probe reads are owned in \cfTokenwSoftChestOwned{} (softmax-weighted) and \cfTokenwTopqChestOwned{} (top quarter) of \cfTokenwChestCells{} chest cells, against \cfTokenwUniformChestOwned{} for the uniform write (Table~\ref{tab:cf-altdir}). Concentrating the write where the probe reads does not create a clinical handle.

\FloatBarrier
\section{Discussion}
\label{sec:implications}

\paragraph{Readability belongs to the probe, ownership to the reader.}
Readability describes a classifier fitted to the representation; ownership describes how the model reading that representation responds to a write. Checkpoints sharing a vision tower have identical probes and write vectors, yet own different concepts, and the chest-trained towers read and answer findings best in the grid but own only \cfSpecTowerOwned{} of \cfSpecTowerN{} chest cells. Ownership therefore belongs on the scorecard of the checkpoint to be deployed rather than on that of the vision tower it contains.

\paragraph{On chest radiographs, the label direction and the answer direction come apart.}
On natural images, the label direction and the model's answer direction are substantially aligned, and writing the label direction controls the answer. On chest radiographs the two are nearly orthogonal: the label direction barely moves its own answer while moving the other findings' answers, and the answer direction is owned in about half of the chest cells. The flat rows of the chest write matrices fit a reader that responds to a signal shared across findings rather than to each finding's own direction \citep{pahde2025navigating}; which signal that is remains open. A legible label need not mark the direction along which the reader chooses its answer, so concept control is a property of a direction and a reader together.

\paragraph{What a steering evaluation should report.}
Probe accuracy and answer accuracy, the two quantities most often reported for medical VLMs, are both high in this grid, yet ownership is \cfChestOwnedPct\%. Random and sham writes show that a direction reaches the answer; only directions fitted the same way for the competing findings, written at the same locus and dose in the same patients, show whether the direction belongs to its label. Studies that use probe normals as interventions should report the write matrix $W$, compare $O_q$ across layers, doses, and models, and pair every clinical write with a natural-image control, where the same procedure produces selective control.

\section{Conclusion}
\label{sec:conclusion}

Reading a clinical concept from a vision-language model's visual representation does not make its probe direction a handle on that concept. Across \cfCheckpoints{} checkpoints and two chest-radiograph datasets, concepts are readable in \cfChestReadablePct\% and answerable in \cfChestAnswerablePct\% of cells but owned in \cfChestOwnedPct\%, because the concept's own direction barely moves its answer while the directions fitted for the other findings move it as much or more. Under the same procedure \cfCocoOwnedPct\% of natural-object cells are owned, also at matched answer ability. Checkpoints that write identical vectors own different concepts, and the model's own answer direction, nearly orthogonal to the probe direction on chest radiographs, provides the handle that the label direction lacks. Readability is a property of a classifier fitted to the representation; control is a property of a direction and the model that reads it.

\section{Limitations}
\label{sec:limitations}

The grades describe linear writes at two loci, the consumed block and the connector output, at one primary dose with a five-point sweep; deeper language-model layers and nonlinear edits may behave differently. Each dataset contributes six concepts, and the chest labels are report-derived, with radiologist labels available for a CheXpert validation subset. The grade needs a discrete answer. MAIRA-2 \citep{bannur2024maira} emits reports rather than answers and fails the image-free mapping check on every template (Appendix~\ref{app:cf-gates}), and the discrete-diffusion report generator of \citet{chen2026echo} has the same shape, so both fall outside the evaluation.

Ownership is a verdict on the write comparison rather than a predictor of effect size. On the four semantic endpoints, adding the ownership contrast to write magnitude, probe selectivity, and clean-answer AUROC adds no held-out predictive value beyond a permutation null ($p=\cfSemCvLocoP$ and $p=\cfSemCvLoboP$ under the concept and block splits; Appendix~\DIFdelbegin \DIFdel{\ref{app:cf-crossed}}\DIFdelend \DIFaddbegin \DIFadd{\ref{app:cf-ansdir}}\DIFaddend ). The answer direction is fitted on training rows to the model's own clean answers and written on test rows and held-out templates, so it is a handle on the model's answer. Its median AUROC against the report label is $\cfValAurocAnsChest$, below the probe's $\cfValAurocProbeChest$, and answer control and clinical fidelity remain distinct objectives.
\section*{Reproducibility statement}
The registered protocol file fixes the loci, projection, probe settings, dose grid, reference family, templates, modules, gates, and decision rules\DIFdelbegin \DIFdel{; its identifier and the evaluation-code version are recorded in every block's scorecard}\DIFdelend . The release contains the evaluation code and adapters, the protocol file, the cohort manifests with their seeds, the fitted directions, the per-image outcomes of every module, every block's scorecard and coverage file, the rendered prompts, and the pinned checkpoint revisions (Table~\ref{tab:cf-models}); every table and figure in the paper is regenerated from those files by the scripts in the release. Appendix~\ref{app:repro} states the consumed block of every family, the consumed-token mask, the projection and probe construction, the control sampling procedures, the cohort roles, and the protocol amendments with their dates.

\section*{Ethics statement}
The study uses three public datasets under their published terms: NIH ChestX-ray14 (NIH Clinical Center), CheXpert Plus (Stanford AIMI, accessed through its data-use agreement), and COCO (annotations CC BY 4.0). All radiographs are de-identified by their providers; no protected health information was accessed, and no image or report is redistributed. The write test evaluates model behaviour under activation edits and is not a diagnostic tool; its grades describe what a probe direction does inside a model and carry no claim of clinical benefit.
\bibliography{refs}
\bibliographystyle{plainnat}

\clearpage
\appendix
\DIFaddbegin \begin{DIFnomarkup}\cfnote{删表}{附录表格按三条标准调整：结论两三句话能说清且文字里已有；多块拼接、列结构互不相同；平铺数字、没有视觉信息。整表删去 11 张：Label counts、Coverage、Steering reference against verdict、Reference family and controls、Ledger、Robustness of the owned grade、Attributes by AUROC、Radiologist labels、Crossing、Cohort roles、Protocol amendments（红框显示）。只保留单一结构的一块：Non-clinical attributes（属性 probe 与方向）、Answer direction（主模板下的计数）、Answer direction and column selectivity（逐 block 的 AUROC 与 cosine）、Four endpoints（ownership 计数）、Fine-grained objects（改为难度匹配小表）。Geometry、Refit、Shared-tower pairs、Alternative directions 四张保留原表。原稿这些表集中排在 A.8 末尾，现在分别放进各自话题的小节，附录 A 的 A.5、A.7–A.11 各有表；复述表内数字的句子相应删去或改为引用表格。}\end{DIFnomarkup}

\DIFaddend \section{Campaign details}
\label{app:campaign}

This appendix lists the checkpoints and their consumed loci, the acceptance decisions, the coverage of the campaign, and the per-concept results of every completed block.

\subsection{Labels, registration, and artifact}
\label{app:cf-labels}

\paragraph{Labels.} We use each dataset's supplied report-derived labels: NIH ChestX-ray14 finding labels \citep{wang2017chestxray}, the CheXpert labeler output \citep{irvin2019chexpert} distributed with CheXpert Plus \citep{chambon2024chexpertplus}, and COCO instance annotations \citep{lin2014coco}. Rows are positive for label 1, negative for label 0, and \emph{unknown} when the labeler output is uncertain or blank, which is the uncertainty convention of the CheXpert labeler \citep{irvin2019chexpert}. We exclude unknown rows from probe fitting, selectivity, and answer AUROC. We write and score them like all other rows. Fixed manifests drawn with recorded seeds define the cohorts, one manifest per dataset for every block. No patient and no image appears in two roles of a dataset. Every cohort is disjoint from every pilot-study cohort. \DIFdelbegin \DIFdel{Table~\ref{tab:cf-prevalence} reports per-concept counts for the three scored cohorts }\DIFdelend \DIFaddbegin \DIFadd{Three cohorts are scored}\DIFaddend : the training cohort for fitting the scaler, probes, and answer direction; the 400-row calibration cohort for the readable and answerable grades; and the 600-row test cohort for every write. A second CheXpert cohort contains \cfValidRows{} frontal films from the CheXpert validation set, one per patient, with the radiologists' consensus labels \citep{irvin2019chexpert}. This cohort has zero patient overlap with any campaign cohort\DIFdelbegin \DIFdel{(Table~\ref{tab:cf-valid})}\DIFdelend .

\paragraph{Registration.} We froze the protocol before scoring any grid block. The loci, projection and probe settings, dose, reference family, six templates, module set, seven gates, and three decision rules match the registered protocol file\DIFdelbegin \DIFdel{and its specified code version}\DIFdelend . We registered the pilot study's crossover and confirmation analyses before running them (Appendix~\ref{app:seed}). \DIFdelbegin \DIFdel{One analysis pass over the packaged outcomes produces all reported results for each block. }\DIFdelend After completing the grid, we added the leaderboard, write-flow share, and per-image examples as descriptive summaries of these outputs.

\paragraph{Artifact.} The release includes evaluation code, adapters, the registered protocol file, cohort manifests with their seeds, and fitted directions. It provides per-image outcomes for every module (write, dose, refit, locus, and template rows), scorecards and coverage files for every block, and rendered prompts. Each new checkpoint requires one adapter that identifies its input block and passes the gates.

\subsection{Grid and consumed loci}
\label{app:cf-grid}

\begin{table}[h]
\centering
\setlength{\tabcolsep}{3pt}
\input{tables/cf_colors}
\DIFaddbeginFL \begin{DIFnomarkup}\cfnote{删列}{删去 Revision 列（各 checkpoint 在 Hugging Face 上的提交哈希）及表注中的 pinned Hugging Face revision。代码与模型版本属于发布文件中的复现记录，不放在论文表格里。其余表格中没有版本、提交或分支信息。}\end{DIFnomarkup}

\DIFaddendFL \caption{\textbf{Checkpoints and families.} Per checkpoint: vision tower, consumed tokens per image, input size, \DIFaddbeginFL \DIFaddFL{and }\DIFaddendFL eligible templates\DIFdelbeginFL \DIFdelFL{, and pinned Hugging Face revision}\DIFdelendFL . Bold family rows give mean selectivity and clean-answer AUROC over clinical cells and the owned clinical and COCO cells, shaded by owned share.}
\label{tab:cf-models}
\DIFdelbeginFL %DIFDELCMD < \fitwidth{%
%DIFDELCMD < \begin{tabular}{llllll@{\hspace{8pt}}rrrr}
%DIFDELCMD < \toprule
%DIFDELCMD < \multirow{2}{*}{Checkpoint} & \multirow{2}{*}{Vision tower} & \multirow{2}{*}{Tokens} & \multirow{2}{*}{Image} & \multirow{2}{*}{Templates} & \multirow{2}{*}{Revision} & \multirow{2}{*}{$\bar S$} & \multirow{2}{*}{$\overline{\mathrm{AUROC}}$} & \multicolumn{2}{c}{owned cells}\\
%DIFDELCMD < \cmidrule(lr){9-10}
%DIFDELCMD < &  &  &  &  &  &  &  & \multicolumn{1}{c}{clinical} & \multicolumn{1}{c}{COCO}\\
%DIFDELCMD < \midrule
%DIFDELCMD < \multicolumn{6}{l}{\textbf{Qwen2.5-VL}~\citep{bai2025qwen25vl}} & 0.132 & 0.635 & \cellcolor{cfTerra2}5/48 & \cellcolor{cfTerra4}24/24 \\
%DIFDELCMD < \quad Qwen2.5-VL-3B & native ViT 675M & 576$\to$144 & $336^2$ & IYWYIAWA & \texttt{6628554} & & & & \\
%DIFDELCMD < \quad Qwen2.5-VL-7B & native ViT 675M & 576$\to$144 & $336^2$ & all & \texttt{cc59489} & & & & \\
%DIFDELCMD < \quad Qwen2.5-VL-32B & native ViT 675M & 576$\to$144 & $336^2$ & all & \texttt{7cfb30d} & & & & \\
%DIFDELCMD < \quad Qwen2.5-VL-72B & native ViT 675M & 576$\to$144 & $336^2$ & all & \texttt{89c8620} & & & & \\
%DIFDELCMD < \midrule
%DIFDELCMD < \multicolumn{6}{l}{\textbf{Qwen3-VL}~\citep{qwen2025qwen3vl}} & 0.137 & 0.735 & \cellcolor{cfTerra2}4/36 & \cellcolor{cfTerra4}18/18 \\
%DIFDELCMD < \quad Qwen3-VL-4B & SigLIP2-based & $\leq$1024 & $\leq 512^2$ & all & \texttt{ebb281e} & & & & \\
%DIFDELCMD < \quad Qwen3-VL-8B & SigLIP2-based & $\leq$1024 & $\leq 512^2$ & all & \texttt{0c351dd} & & & & \\
%DIFDELCMD < \quad Qwen3-VL-32B & SigLIP2-based & $\leq$1024 & $\leq 512^2$ & all & \texttt{0cfaf48} & & & & \\
%DIFDELCMD < \midrule
%DIFDELCMD < \multicolumn{6}{l}{\textbf{InternVL3.5}~\citep{wang2025internvl35}} & 0.125 & 0.807 & \cellcolor{cfTerra2}5/36 & \cellcolor{cfTerra4}17/18 \\
%DIFDELCMD < \quad InternVL3.5-8B & InternViT-300M & 1024 & $448^2$ & all & \texttt{741a7d0} & & & & \\
%DIFDELCMD < \quad InternVL3.5-14B & InternViT-300M & 1024 & $448^2$ & all & \texttt{226b96d} & & & & \\
%DIFDELCMD < \quad InternVL3.5-38B & InternViT-6B & 1024 & $448^2$ & all & \texttt{7c830fc} & & & & \\
%DIFDELCMD < \midrule
%DIFDELCMD < \multicolumn{6}{l}{\textbf{Gemma 3}~\citep{gemma2025gemma3}} & 0.124 & 0.564 & \cellcolor{cfTerra2}2/36 & \cellcolor{cfTerra4}17/18 \\
%DIFDELCMD < \quad Gemma 3 4B & SigLIP-So400m & 4096 & $896^2$ & all & \texttt{093f9f3} & & & & \\
%DIFDELCMD < \quad Gemma 3 12B & SigLIP-So400m & 4096 & $896^2$ & all & \texttt{96b6f1e} & & & & \\
%DIFDELCMD < \quad Gemma 3 27B & SigLIP-So400m & 4096 & $896^2$ & all & \texttt{005ad34} & & & & \\
%DIFDELCMD < \midrule
%DIFDELCMD < \multicolumn{6}{l}{\textbf{MedGemma} (medical)~\citep{sellergren2025medgemma}} & 0.164 & 0.835 & \cellcolor{cfTerra2}2/24 & \cellcolor{cfTerra3}10/12 \\
%DIFDELCMD < \quad MedGemma 4B & MedSigLIP & 4096 & $896^2$ & all & \texttt{290cda5} & & & & \\
%DIFDELCMD < \quad MedGemma 27B & MedSigLIP & 4096 & $896^2$ & all & \texttt{2d3e00e} & & & & \\
%DIFDELCMD < \midrule
%DIFDELCMD < \multicolumn{6}{l}{\textbf{Lingshu} (medical)~\citep{lasa2025lingshu}} & 0.142 & 0.820 & \cellcolor{cfTerra4}11/24 & \cellcolor{cfTerra4}12/12 \\
%DIFDELCMD < \quad Lingshu 7B & Qwen2.5-VL ViT & 576$\to$144 & $336^2$ & all & \texttt{b98aecd} & & & & \\
%DIFDELCMD < \quad Lingshu 32B & Qwen2.5-VL ViT & 576$\to$144 & $336^2$ & all & \texttt{36b9827} & & & & \\
%DIFDELCMD < \midrule
%DIFDELCMD < \multicolumn{6}{l}{\textbf{LLaVA-1.5}~\citep{liu2024improved}} & 0.124 & 0.516 & \cellcolor{cfTerra1}0/24 & \cellcolor{cfTerra3}7/12 \\
%DIFDELCMD < \quad LLaVA-1.5-7B & CLIP ViT-L/14-336 & 576 & $336^2$ & IYWY & \texttt{b234b80} & & & & \\
%DIFDELCMD < \quad LLaVA-1.5-13B & CLIP ViT-L/14-336 & 576 & $336^2$ & IYWYIAWA & \texttt{5dda288} & & & & \\
%DIFDELCMD < \midrule
%DIFDELCMD < \multicolumn{6}{l}{\textbf{LLaVA-Med} (medical)~\citep{li2023llavamed}} & 0.136 & 0.560 & \cellcolor{cfTerra2}1/12 & \cellcolor{cfTerra2}1/6 \\
%DIFDELCMD < \quad LLaVA-Med 1.5 7B & CLIP ViT-L/14-336 & 576 & $336^2$ & IYWY & \texttt{91bb16c} & & & & \\
%DIFDELCMD < \midrule
%DIFDELCMD < \multicolumn{6}{l}{\textbf{Llama 3.2 Vision}~\citep{meta2024llama32vision}} & 0.119 & 0.688 & \cellcolor{cfTerra1}0/12 & \cellcolor{cfTerra2}1/6 \\
%DIFDELCMD < \quad Llama 3.2 Vision 11B & ViT-H/14 & 1601 & $560^2$ & IBWB & \texttt{9eb2daa} & & & & \\
%DIFDELCMD < \midrule
%DIFDELCMD < \multicolumn{6}{l}{\textbf{CheXagent} (medical)~\citep{chen2024chexagent}} & 0.180 & 0.857 & \cellcolor{cfTerra3}7/24 & \cellcolor{cfTerra2}2/12 \\
%DIFDELCMD < \quad CheXagent-2-3B & XraySigLIP ViT-L/16 & 1024 & $512^2$ & IYWYIAWA & \texttt{8f19b53} & & & & \\
%DIFDELCMD < \quad CheXagent-8B & EVA-CLIP ViT 1.4B & 1025$\to$128 & $448^2$ & IYWYIAWA & \texttt{4934e91} & & & & \\
%DIFDELCMD < \midrule
%DIFDELCMD < \multicolumn{6}{l}{\textbf{LLaVA-Rad} (medical)~\citep{zambranochaves2025llavarad}} & 0.209 & 0.793 & \cellcolor{cfTerra1}0/12 & \cellcolor{cfTerra1}0/6 \\
%DIFDELCMD < \quad LLaVA-Rad-7B & BiomedCLIP-CXR-518 & 1369 & $518^2$ & IYWY & \texttt{dcdbc6c} & & & & \\
%DIFDELCMD < \midrule
%DIFDELCMD < \multicolumn{6}{l}{\textbf{HuatuoGPT-Vision} (medical)~\citep{chen2024huatuogptvision}} & 0.136 & 0.746 & \cellcolor{cfTerra2}1/12 & \cellcolor{cfTerra3}4/6 \\
%DIFDELCMD < \quad HuatuoGPT-Vision-7B & CLIP ViT-L/14-336 & 576 & $336^2$ & all & \texttt{34dfcdb} & & & & \\
%DIFDELCMD < \bottomrule
%DIFDELCMD < \end{tabular}}
%DIFDELCMD < %%%
\DIFdelendFL \DIFaddbeginFL \resizebox{\linewidth}{!}{%
\begin{tabular}{lllll@{\hspace{8pt}}rrrr}
\toprule
\multirow{2}{*}{Checkpoint} & \multirow{2}{*}{Vision tower} & \multirow{2}{*}{Tokens} & \multirow{2}{*}{Image} & \multirow{2}{*}{Templates} & \multirow{2}{*}{$\bar S$} & \multirow{2}{*}{$\overline{\mathrm{AUROC}}$} & \multicolumn{2}{c}{owned cells}\\
\cmidrule(lr){8-9}
&  &  &  &  &  &  & \multicolumn{1}{c}{clinical} & \multicolumn{1}{c}{COCO}\\
\midrule
\multicolumn{5}{l}{\textbf{Qwen2.5-VL}~\citep{bai2025qwen25vl}} & 0.132 & 0.635 & \cellcolor{cfTerra2}5/48 & \cellcolor{cfTerra4}24/24 \\
\quad Qwen2.5-VL-3B & native ViT 675M & 576$\to$144 & $336^2$ & IYWYIAWA & & & & \\
\quad Qwen2.5-VL-7B & native ViT 675M & 576$\to$144 & $336^2$ & all & & & & \\
\quad Qwen2.5-VL-32B & native ViT 675M & 576$\to$144 & $336^2$ & all & & & & \\
\quad Qwen2.5-VL-72B & native ViT 675M & 576$\to$144 & $336^2$ & all & & & & \\
\midrule
\multicolumn{5}{l}{\textbf{Qwen3-VL}~\citep{qwen2025qwen3vl}} & 0.137 & 0.735 & \cellcolor{cfTerra2}4/36 & \cellcolor{cfTerra4}18/18 \\
\quad Qwen3-VL-4B & SigLIP2-based & $\leq$1024 & $\leq 512^2$ & all & & & & \\
\quad Qwen3-VL-8B & SigLIP2-based & $\leq$1024 & $\leq 512^2$ & all & & & & \\
\quad Qwen3-VL-32B & SigLIP2-based & $\leq$1024 & $\leq 512^2$ & all & & & & \\
\midrule
\multicolumn{5}{l}{\textbf{InternVL3.5}~\citep{wang2025internvl35}} & 0.125 & 0.807 & \cellcolor{cfTerra2}5/36 & \cellcolor{cfTerra4}17/18 \\
\quad InternVL3.5-8B & InternViT-300M & 1024 & $448^2$ & all & & & & \\
\quad InternVL3.5-14B & InternViT-300M & 1024 & $448^2$ & all & & & & \\
\quad InternVL3.5-38B & InternViT-6B & 1024 & $448^2$ & all & & & & \\
\midrule
\multicolumn{5}{l}{\textbf{Gemma 3}~\citep{gemma2025gemma3}} & 0.124 & 0.564 & \cellcolor{cfTerra2}2/36 & \cellcolor{cfTerra4}17/18 \\
\quad Gemma 3 4B & SigLIP-So400m & 4096 & $896^2$ & all & & & & \\
\quad Gemma 3 12B & SigLIP-So400m & 4096 & $896^2$ & all & & & & \\
\quad Gemma 3 27B & SigLIP-So400m & 4096 & $896^2$ & all & & & & \\
\midrule
\multicolumn{5}{l}{\textbf{MedGemma} (medical)~\citep{sellergren2025medgemma}} & 0.164 & 0.835 & \cellcolor{cfTerra2}2/24 & \cellcolor{cfTerra3}10/12 \\
\quad MedGemma 4B & MedSigLIP & 4096 & $896^2$ & all & & & & \\
\quad MedGemma 27B & MedSigLIP & 4096 & $896^2$ & all & & & & \\
\midrule
\multicolumn{5}{l}{\textbf{Lingshu} (medical)~\citep{lasa2025lingshu}} & 0.142 & 0.820 & \cellcolor{cfTerra4}11/24 & \cellcolor{cfTerra4}12/12 \\
\quad Lingshu 7B & Qwen2.5-VL ViT & 576$\to$144 & $336^2$ & all & & & & \\
\quad Lingshu 32B & Qwen2.5-VL ViT & 576$\to$144 & $336^2$ & all & & & & \\
\midrule
\multicolumn{5}{l}{\textbf{LLaVA-1.5}~\citep{liu2024improved}} & 0.124 & 0.516 & \cellcolor{cfTerra1}0/24 & \cellcolor{cfTerra3}7/12 \\
\quad LLaVA-1.5-7B & CLIP ViT-L/14-336 & 576 & $336^2$ & IYWY & & & & \\
\quad LLaVA-1.5-13B & CLIP ViT-L/14-336 & 576 & $336^2$ & IYWYIAWA & & & & \\
\midrule
\multicolumn{5}{l}{\textbf{LLaVA-Med} (medical)~\citep{li2023llavamed}} & 0.136 & 0.560 & \cellcolor{cfTerra2}1/12 & \cellcolor{cfTerra2}1/6 \\
\quad LLaVA-Med 1.5 7B & CLIP ViT-L/14-336 & 576 & $336^2$ & IYWY & & & & \\
\midrule
\multicolumn{5}{l}{\textbf{Llama 3.2 Vision}~\citep{meta2024llama32vision}} & 0.119 & 0.688 & \cellcolor{cfTerra1}0/12 & \cellcolor{cfTerra2}1/6 \\
\quad Llama 3.2 Vision 11B & ViT-H/14 & 1601 & $560^2$ & IBWB & & & & \\
\midrule
\multicolumn{5}{l}{\textbf{CheXagent} (medical)~\citep{chen2024chexagent}} & 0.180 & 0.857 & \cellcolor{cfTerra3}7/24 & \cellcolor{cfTerra2}2/12 \\
\quad CheXagent-2-3B & XraySigLIP ViT-L/16 & 1024 & $512^2$ & IYWYIAWA & & & & \\
\quad CheXagent-8B & EVA-CLIP ViT 1.4B & 1025$\to$128 & $448^2$ & IYWYIAWA & & & & \\
\midrule
\multicolumn{5}{l}{\textbf{LLaVA-Rad} (medical)~\citep{zambranochaves2025llavarad}} & 0.209 & 0.793 & \cellcolor{cfTerra1}0/12 & \cellcolor{cfTerra1}0/6 \\
\quad LLaVA-Rad-7B & BiomedCLIP-CXR-518 & 1369 & $518^2$ & IYWY & & & & \\
\midrule
\multicolumn{5}{l}{\textbf{HuatuoGPT-Vision} (medical)~\citep{chen2024huatuogptvision}} & 0.136 & 0.746 & \cellcolor{cfTerra2}1/12 & \cellcolor{cfTerra3}4/6 \\
\quad HuatuoGPT-Vision-7B & CLIP ViT-L/14-336 & 576 & $336^2$ & all & & & & \\
\bottomrule
\end{tabular}}
\DIFaddendFL \end{table}
\begin{figure}[h]
    \centering
    \resizebox{\textwidth}{!}{\includegraphics{figures/fig6_ladders.pdf}}
    \caption{\textbf{Size ladders.} Ownership $O_q$ per concept across sizes for Qwen2.5-VL and Qwen3-VL on NIH and Lingshu on CheXpert. Solid lines mark the concept with the largest ownership at the largest size (Atelectasis, Nodule, Edema), filled markers owned cells; vertical scales differ by panel.}
    \label{fig:ladders}
\end{figure}


\DIFdelbegin \begin{table}[htbp]
\DIFdeletedtable{%
{\small\textbf{Table (deleted).} \textbf{Label counts per cohort and concept.} Positive / negative / unknown rows of the training, calibration, and test cohorts of the fixed manifests. Unknown means the CheXpert labeler was uncertain or blank.\par}\medskip
\centering
\setlength{\tabcolsep}{4pt}

\fitwidth{%
\fitwidth{%
\begin{tabular}{llrrrrrrrrr}
\toprule
Dataset & Concept & \multicolumn{3}{c}{train (pos / neg / unk)} & \multicolumn{3}{c}{calibration (pos / neg / unk)} & \multicolumn{3}{c}{test (pos / neg / unk)} \\
\cmidrule(lr){3-5}\cmidrule(lr){6-8}\cmidrule(lr){9-11}
NIH ChestX-ray14 & Effusion & 2423 & 15789 & 0 & 20 & 380 & 0 & 46 & 554 & 0 \\
 & Atelectasis & 2311 & 15901 & 0 & 33 & 367 & 0 & 57 & 543 & 0 \\
 & Pneumothorax & 1254 & 16958 & 0 & 13 & 387 & 0 & 17 & 583 & 0 \\
 & Cardiomegaly & 881 & 17331 & 0 & 15 & 385 & 0 & 23 & 577 & 0 \\
 & Mass & 1263 & 16949 & 0 & 15 & 385 & 0 & 32 & 568 & 0 \\
 & Nodule & 1549 & 16663 & 0 & 27 & 373 & 0 & 47 & 553 & 0 \\
\addlinespace
CheXpert Plus & Effusion & 5897 & 4277 & 9826 & 117 & 93 & 190 & 158 & 130 & 312 \\
 & Atelectasis & 2891 & 86 & 17023 & 67 & 4 & 329 & 84 & 2 & 514 \\
 & Pneumothorax & 1086 & 6456 & 12458 & 17 & 136 & 247 & 32 & 213 & 355 \\
 & Cardiomegaly & 2312 & 1903 & 15785 & 49 & 46 & 305 & 71 & 58 & 471 \\
 & Consolidation & 966 & 3652 & 15382 & 25 & 83 & 292 & 25 & 113 & 462 \\
 & Edema & 3814 & 2231 & 13955 & 70 & 54 & 276 & 133 & 59 & 408 \\
\addlinespace
COCO & person & 10877 & 9123 & 0 & 228 & 172 & 0 & 321 & 279 & 0 \\
 & dog & 693 & 19307 & 0 & 10 & 390 & 0 & 15 & 585 & 0 \\
 & car & 2051 & 17949 & 0 & 48 & 352 & 0 & 60 & 540 & 0 \\
 & chair & 2183 & 17817 & 0 & 59 & 341 & 0 & 63 & 537 & 0 \\
 & bottle & 1455 & 18545 & 0 & 34 & 366 & 0 & 47 & 553 & 0 \\
 & bicycle & 567 & 19433 & 0 & 9 & 391 & 0 & 14 & 586 & 0 \\
\bottomrule
\end{tabular}}}
}
\end{table}
\begin{table}[htbp]
\DIFoldtable{%
{\small\textbf{Table (before revision).} \textbf{Checkpoints and families.} Per checkpoint: vision tower, consumed tokens per image, input size, eligible templates, and pinned Hugging Face revision. Bold family rows give mean selectivity and clean-answer AUROC over clinical cells and the owned clinical and COCO cells, shaded by owned share.\par}\medskip
\centering
\setlength{\tabcolsep}{3pt}
\input{tables/cf_colors}

\fitwidth{%
\fitwidth{%
\begin{tabular}{llllll@{\hspace{8pt}}rrrr}
\toprule
\multirow{2}{*}{Checkpoint} & \multirow{2}{*}{Vision tower} & \multirow{2}{*}{Tokens} & \multirow{2}{*}{Image} & \multirow{2}{*}{Templates} & \multirow{2}{*}{Revision} & \multirow{2}{*}{$\bar S$} & \multirow{2}{*}{$\overline{\mathrm{AUROC}}$} & \multicolumn{2}{c}{owned cells}\\
\cmidrule(lr){9-10}
&  &  &  &  &  &  &  & \multicolumn{1}{c}{clinical} & \multicolumn{1}{c}{COCO}\\
\midrule
\multicolumn{6}{l}{\textbf{Qwen2.5-VL}~\citep{bai2025qwen25vl}} & 0.132 & 0.635 & \cellcolor{cfTerra2}5/48 & \cellcolor{cfTerra4}24/24 \\
\quad Qwen2.5-VL-3B & native ViT 675M & 576$\to$144 & $336^2$ & IYWYIAWA & \texttt{6628554} & & & & \\
\quad Qwen2.5-VL-7B & native ViT 675M & 576$\to$144 & $336^2$ & all & \texttt{cc59489} & & & & \\
\quad Qwen2.5-VL-32B & native ViT 675M & 576$\to$144 & $336^2$ & all & \texttt{7cfb30d} & & & & \\
\quad Qwen2.5-VL-72B & native ViT 675M & 576$\to$144 & $336^2$ & all & \texttt{89c8620} & & & & \\
\midrule
\multicolumn{6}{l}{\textbf{Qwen3-VL}~\citep{qwen2025qwen3vl}} & 0.137 & 0.735 & \cellcolor{cfTerra2}4/36 & \cellcolor{cfTerra4}18/18 \\
\quad Qwen3-VL-4B & SigLIP2-based & $\leq$1024 & $\leq 512^2$ & all & \texttt{ebb281e} & & & & \\
\quad Qwen3-VL-8B & SigLIP2-based & $\leq$1024 & $\leq 512^2$ & all & \texttt{0c351dd} & & & & \\
\quad Qwen3-VL-32B & SigLIP2-based & $\leq$1024 & $\leq 512^2$ & all & \texttt{0cfaf48} & & & & \\
\midrule
\multicolumn{6}{l}{\textbf{InternVL3.5}~\citep{wang2025internvl35}} & 0.125 & 0.807 & \cellcolor{cfTerra2}5/36 & \cellcolor{cfTerra4}17/18 \\
\quad InternVL3.5-8B & InternViT-300M & 1024 & $448^2$ & all & \texttt{741a7d0} & & & & \\
\quad InternVL3.5-14B & InternViT-300M & 1024 & $448^2$ & all & \texttt{226b96d} & & & & \\
\quad InternVL3.5-38B & InternViT-6B & 1024 & $448^2$ & all & \texttt{7c830fc} & & & & \\
\midrule
\multicolumn{6}{l}{\textbf{Gemma 3}~\citep{gemma2025gemma3}} & 0.124 & 0.564 & \cellcolor{cfTerra2}2/36 & \cellcolor{cfTerra4}17/18 \\
\quad Gemma 3 4B & SigLIP-So400m & 4096 & $896^2$ & all & \texttt{093f9f3} & & & & \\
\quad Gemma 3 12B & SigLIP-So400m & 4096 & $896^2$ & all & \texttt{96b6f1e} & & & & \\
\quad Gemma 3 27B & SigLIP-So400m & 4096 & $896^2$ & all & \texttt{005ad34} & & & & \\
\midrule
\multicolumn{6}{l}{\textbf{MedGemma} (medical)~\citep{sellergren2025medgemma}} & 0.164 & 0.835 & \cellcolor{cfTerra2}2/24 & \cellcolor{cfTerra3}10/12 \\
\quad MedGemma 4B & MedSigLIP & 4096 & $896^2$ & all & \texttt{290cda5} & & & & \\
\quad MedGemma 27B & MedSigLIP & 4096 & $896^2$ & all & \texttt{2d3e00e} & & & & \\
\midrule
\multicolumn{6}{l}{\textbf{Lingshu} (medical)~\citep{lasa2025lingshu}} & 0.142 & 0.820 & \cellcolor{cfTerra4}11/24 & \cellcolor{cfTerra4}12/12 \\
\quad Lingshu 7B & Qwen2.5-VL ViT & 576$\to$144 & $336^2$ & all & \texttt{b98aecd} & & & & \\
\quad Lingshu 32B & Qwen2.5-VL ViT & 576$\to$144 & $336^2$ & all & \texttt{36b9827} & & & & \\
\midrule
\multicolumn{6}{l}{\textbf{LLaVA-1.5}~\citep{liu2024improved}} & 0.124 & 0.516 & \cellcolor{cfTerra1}0/24 & \cellcolor{cfTerra3}7/12 \\
\quad LLaVA-1.5-7B & CLIP ViT-L/14-336 & 576 & $336^2$ & IYWY & \texttt{b234b80} & & & & \\
\quad LLaVA-1.5-13B & CLIP ViT-L/14-336 & 576 & $336^2$ & IYWYIAWA & \texttt{5dda288} & & & & \\
\midrule
\multicolumn{6}{l}{\textbf{LLaVA-Med} (medical)~\citep{li2023llavamed}} & 0.136 & 0.560 & \cellcolor{cfTerra2}1/12 & \cellcolor{cfTerra2}1/6 \\
\quad LLaVA-Med 1.5 7B & CLIP ViT-L/14-336 & 576 & $336^2$ & IYWY & \texttt{91bb16c} & & & & \\
\midrule
\multicolumn{6}{l}{\textbf{Llama 3.2 Vision}~\citep{meta2024llama32vision}} & 0.119 & 0.688 & \cellcolor{cfTerra1}0/12 & \cellcolor{cfTerra2}1/6 \\
\quad Llama 3.2 Vision 11B & ViT-H/14 & 1601 & $560^2$ & IBWB & \texttt{9eb2daa} & & & & \\
\midrule
\multicolumn{6}{l}{\textbf{CheXagent} (medical)~\citep{chen2024chexagent}} & 0.180 & 0.857 & \cellcolor{cfTerra3}7/24 & \cellcolor{cfTerra2}2/12 \\
\quad CheXagent-2-3B & XraySigLIP ViT-L/16 & 1024 & $512^2$ & IYWYIAWA & \texttt{8f19b53} & & & & \\
\quad CheXagent-8B & EVA-CLIP ViT 1.4B & 1025$\to$128 & $448^2$ & IYWYIAWA & \texttt{4934e91} & & & & \\
\midrule
\multicolumn{6}{l}{\textbf{LLaVA-Rad} (medical)~\citep{zambranochaves2025llavarad}} & 0.209 & 0.793 & \cellcolor{cfTerra1}0/12 & \cellcolor{cfTerra1}0/6 \\
\quad LLaVA-Rad-7B & BiomedCLIP-CXR-518 & 1369 & $518^2$ & IYWY & \texttt{dcdbc6c} & & & & \\
\midrule
\multicolumn{6}{l}{\textbf{HuatuoGPT-Vision} (medical)~\citep{chen2024huatuogptvision}} & 0.136 & 0.746 & \cellcolor{cfTerra2}1/12 & \cellcolor{cfTerra3}4/6 \\
\quad HuatuoGPT-Vision-7B & CLIP ViT-L/14-336 & 576 & $336^2$ & all & \texttt{34dfcdb} & & & & \\
\bottomrule
\end{tabular}}}
}
\end{table}
%DIFDELCMD < 

%DIFDELCMD < %%%
\DIFdelend The towers include CLIP \citep{radford2021clip} in the LLaVA families, SigLIP \citep{zhai2023siglip} in Gemma~3, and a SigLIP~2 design \citep{tschannen2025siglip2} in Qwen3-VL; MedGemma's MedSigLIP tower is a medically tuned SigLIP released with that model \citep{sellergren2025medgemma}. The three towers trained on chest radiographs are the XraySigLIP tower of CheXagent-2-3B \citep{chen2024chexagent}, the EVA-CLIP tower \citep{sun2023evaclip} of CheXagent-8B, and the BiomedCLIP-CXR tower of LLaVA-Rad-7B \citep{zhang2025biomedclip,zambranochaves2025llavarad}; HuatuoGPT-Vision-7B tunes on medical data over the same stock CLIP tower as LLaVA-1.5. Consumed-token counts range from \cfTokensMin{} to \cfTokensMax{} per image, and each family uses a fixed image resolution. The Llama~3.2 Vision projector also reads five intermediate local-layer outputs that bypass the final global block. LLaVA computes the final CLIP block but consumes the penultimate block. At every checkpoint, Gate (C) verifies that the write reaches the consumer and the answer logits and changes nothing outside the consumed tokens.

\subsection{Acceptance decisions}
\label{app:cf-gates}

Seven checks run on the 16 preflight rows at both loci before a block is scored: (A) bitwise determinism; (B) identity at $\alpha=0$; (C) reach, under which the write must change the consumer's input and answer logits without changing anything outside the consumed tokens; (D) batched-versus-single agreement of the answer logits within $0.25$; (E) image-free semantic mapping of every template; (F) throughput; and (G) fp32 recomputation of the answer logits. Deviations in (D) are declared and neutralised by scoring every condition of a question in one fixed batch composition. Templates that fail (E) are ineligible; a checkpoint whose yes/no template fails is graded on its eligible B-positive template (IB for Llama~3.2 Vision), and its scorecards record the template used and every declared deviation.

All \cfWriteBlocks{} completed blocks pass checks (A), (B), (C), and (G): determinism error is zero, change at $\alpha=0$ is zero, and fp32-versus-model logit differences are below \cfGateFpMax{}. In check (D), the Qwen3-VL, Gemma~3, MedGemma, InternVL, and Llama families exceed the 0.25-logit tolerance (up to \cfGateBatchMax{} logits for Gemma~3 12B). We record these as declared deviations. Fixed batch composition neutralises them. Check (E) fails across all datasets for the A/B templates for LLaVA-1.5-7B and LLaVA-Med, the B-mapped templates for Qwen2.5-VL-3B and LLaVA-1.5-13B, and the yes/no and A-mapped templates for Llama~3.2 Vision 11B. We record affected cells as ineligible. Check (E) fails for every template of MAIRA-2 \citep{bannur2024maira} on NIH: the model answers ``No'' to every finding on all \cfSpecUngradedRows{} preflight images and ``A'' to every forced choice under both letter mappings, so no template carries a graded answer and the block produces no outcomes. Two Gemma~3 checkpoints answer ``yes'' to most chest questions; Gemma~3 4B does so on \cfGemmaFourYesMin--\cfGemmaFourYesMax\% of rows (clean $P(\mathrm{yes})>0.99$). The answerability grade records this behaviour. Three checkpoints (LLaVA-1.5 7B/13B, LLaVA-Med) answer at chance.

\subsection{Coverage and cost}
\label{app:cf-coverage}

\DIFaddbegin \begin{DIFnomarkup}\cfnote{附录精简}{scorecard 内容清单缩成一句，删去与 A.3 重复的 Llama IB 模板一句。}\end{DIFnomarkup}

\DIFaddend Every checkpoint has a completed core write matrix and calibration pass on NIH ChestX-ray14, CheXpert Plus, and COCO\DIFdelbegin \DIFdel{(Table~\ref{tab:cf-coverage})}\DIFdelend . The campaign targets \cfCoveragePlannedWord{} modules per block: the write matrix, the calibration pass, the five additional templates, the dose sweep, the two probe refits, the connector locus, and the connector-locus calibration. \cfCoverageBlocksAllSeven{} of \cfBlocks{} blocks carry all seven modules. A probe-graded block whose write matrix is ineligible or incomplete contributes readability alone. \DIFdelbegin \DIFdel{Every }\DIFdelend \DIFaddbegin \DIFadd{Each }\DIFaddend block's scorecard also reports the \DIFdelbegin \DIFdel{signed }\DIFdelend dose range of $O_q$ over \DIFdelbegin \DIFdel{$\alpha\in\{-0.5,\ldots,+0.5\}$, the standard deviation of $O_q$ }\DIFdelend \DIFaddbegin \DIFadd{$\alpha\in[-0.5,0.5]$, its standard deviation }\DIFaddend over the three fit seeds\DIFdelbegin \DIFdel{(seeds 1 and 2 resample the training units with replacement and refit the scaler and the six probes), ownership and the absolute write effects }\DIFdelend \DIFaddbegin \DIFadd{, ownership }\DIFaddend at the connector locus, and the wording \DIFdelbegin \DIFdel{(}\emph{\DIFdel{is}} %DIFAUXCMD
\DIFdel{versus }\emph{\DIFdel{show}}%DIFAUXCMD
\DIFdel{) and mapping (A-present versus B-present) }\DIFdelend \DIFaddbegin \DIFadd{and mapping }\DIFaddend contrasts of $O_q$ across the six templates. \DIFaddbegin \DIFadd{The partial block is }\DIFaddend Qwen2.5-VL-72B on CheXpert Plus\DIFdelbegin \DIFdel{is the partial block: its core write matrix, calibration, dose sweep, refits, and connector-locus calibration are complete; its prompt-template and connector-locus modules are not. Llama~3.2 Vision 11B is graded on its eligible B-positive template IB. }\DIFdelend \DIFaddbegin \DIFadd{, which lacks the additional templates and the connector locus. }\DIFaddend The campaign used about 3,000 A100-hours\DIFdelbegin \DIFdel{. A chest block for }\DIFdelend \DIFaddbegin \DIFadd{; a chest block of }\DIFaddend a 7B checkpoint takes about ten\DIFdelbegin \DIFdel{GPU-hours}\DIFdelend .


\DIFdelbegin \begin{table}[htbp]
\DIFdeletedtable{%
{\small\textbf{Table (deleted).} \textbf{Coverage.} Modules completed per block at packaging, from the campaign manifest. Brackets name the primary template when it is not the \emph{is}/yes-no template.\par}\medskip
\centering
\setlength{\tabcolsep}{4pt}

\fitwidth{%
\fitwidth{%
\begin{tabular}{llll}
\toprule
Checkpoint & NIH ChestX-ray14 & CheXpert Plus & COCO\\
\midrule
Qwen2.5-VL-3B & C Cal P D R L Lc A Ad & C Cal P D R L Lc A Ad & C Cal P D R L Lc A Ad \\
Qwen2.5-VL-7B & C Cal P D R L Lc A Ad An E T Pr & C Cal P D R L Lc A Ad An E T & C Cal P D R L Lc A Ad An T Pr \\
Qwen2.5-VL-32B & C Cal P D R L Lc A Ad & C Cal P D R L Lc A Ad & C Cal P D R L Lc A Ad \\
Qwen2.5-VL-72B & C Cal P D R L Lc & C Cal D R Lc & C Cal P D R L Lc \\
Qwen3-VL-4B & C Cal P D R L Lc A Ad & C Cal P D R L Lc A Ad & C Cal P D R L Lc A Ad \\
Qwen3-VL-8B & C Cal P D R L Lc A Ad E T Pr & C Cal P D R L Lc A Ad & C Cal P D R L Lc A Ad T \\
Qwen3-VL-32B & C Cal P D R L Lc A Ad & C Cal P D R L Lc A Ad & C Cal P D R L Lc A Ad \\
InternVL3.5-8B & C Cal P D R L Lc A Ad E Pr & C Cal P D R L Lc A Ad & C Cal P D R L Lc A Ad \\
InternVL3.5-14B & C Cal P D R L Lc A Ad & C Cal P D R L Lc A Ad & C Cal P D R L Lc A Ad \\
InternVL3.5-38B & C Cal P D R L Lc A Ad & C Cal P D R L Lc A Ad & C Cal P D R L Lc A Ad \\
Gemma 3 4B & C Cal P D R L Lc A Ad Pr & C Cal P D R L Lc A Ad & C Cal P D R L Lc A Ad \\
Gemma 3 12B & C Cal P D R L Lc A Ad An E T Pr & C Cal P D R L Lc A Ad An E & C Cal P D R L Lc A Ad An \\
Gemma 3 27B & C Cal P D R L Lc A Ad Pr & C Cal P D R L Lc A Ad & C Cal P D R L Lc A Ad \\
MedGemma 4B & C Cal P D R L Lc A Ad Pr & C Cal P D R L Lc A Ad E & C Cal P D R L Lc A Ad \\
MedGemma 27B & C Cal P D R L Lc A Ad & C Cal P D R L Lc A Ad & C Cal P D R L Lc A Ad \\
Lingshu 7B & C Cal P D R L Lc A Ad & C Cal P D R L Lc A Ad & C Cal P D R L Lc A Ad \\
Lingshu 32B & C Cal P D R L Lc A Ad An E T Pr & C Cal P D R L Lc A Ad An E T & C Cal P D R L Lc A Ad An \\
LLaVA-1.5-7B & C Cal P D R L Lc A Ad & C Cal P D R L Lc A Ad & C Cal P D R L Lc A Ad \\
LLaVA-1.5-13B & C Cal P D R L Lc A Ad & C Cal P D R L Lc A Ad & C Cal P D R L Lc A Ad \\
LLaVA-Med 1.5 7B & C Cal P D R L Lc A Ad & C Cal P D R L Lc A Ad & C Cal P D R L Lc A Ad \\
CheXagent-2-3B & C Cal P D R L Lc & C Cal P D R L Lc & C Cal P D R L Lc \\
CheXagent-8B & C Cal P D R L Lc & C Cal P D R L Lc & C Cal P D R L Lc \\
LLaVA-Rad-7B & C Cal P D R L Lc & C Cal P D R L Lc & C Cal P D R L Lc \\
HuatuoGPT-Vision-7B & C Cal P D R L Lc & C Cal P D R L Lc & C Cal P D R L Lc \\
Llama 3.2 Vision 11B & C Cal P D R L Lc A [IB] & C Cal P D R L Lc A [IB] & C Cal P D R L Lc A [IB] \\
\midrule
\multicolumn{4}{p{1.357\textwidth}}{\textit{Legend:} C = write matrix, Cal = calibration pass, P = five additional templates, D = dose sweep, R = two probe refits, L = connector locus, Lc = connector-locus calibration, A = alternative direction constructions, Ad = displacement lifts, An = answer direction, E = extended competitor family, T = token-weighted writes, Pr = fp32 and batch-size-one rescoring.} \\
\bottomrule
\end{tabular}}}
}
\end{table}
%DIFDELCMD < 

%DIFDELCMD < %%%
\DIFdelend \subsection{Ownership of every concept}
\label{app:cf-tables}

Tables~\ref{tab:cf-own-nih}--\ref{tab:cf-own-coco} report the ownership contrast \DIFaddbegin \DIFadd{$O_q$ }\DIFaddend of every concept cell, one table per dataset\DIFdelbegin \DIFdel{. Its }\DIFdelend \DIFaddbegin \DIFadd{, with the checkpoints in the order of Table~\ref{tab:cf-main}. Their }\DIFaddend cells are shaded by magnitude at \DIFdelbegin \DIFdel{the thresholds }\DIFdelend $0.02$, $0.10$, and $0.25$, terracotta for positive and \DIFdelbegin \DIFdel{slate }\DIFdelend \DIFaddbegin \DIFadd{blue }\DIFaddend for negative ownership\DIFdelbegin \DIFdel{, and its checkpoints follow the order of Table~\ref{tab:cf-main}. Figure~\ref{fig:cf-w-structure} summarises the structure of the underlying write matrices.
Figure~\ref{fig:cf-examples} shows grades for individual images.
}\DIFdelend \DIFaddbegin \DIFadd{.
}\DIFaddend 

\DIFaddbegin \begin{DIFnomarkup}\cfnote{写作}{原稿 A.5 只有开头一段就接三张表，没有对各表的说明，也没有收尾。现在每张表前加一小段，给出该数据集被 own 的 cell 数、stronger competitor 数、$O_q>0$ 的比例和值得看的行；三张表之后加一段收尾，对比胸片与 COCO，并指向写入矩阵图和单图示例（这两个引用原来在开头段）。数字全部用宏。开头段的 Its 改为 Their，颜色说法与表注统一为 blue。}\end{DIFnomarkup}

\paragraph{\DIFadd{NIH ChestX-ray14.}} \cfNihOwned{} \DIFadd{of the }\cfNihCells{} \DIFadd{cells are owned, a stronger competitor is identified in }\cfNihCompetitor{}\DIFadd{, and the contrast is positive in }\cfGeoOwnWinsNihPct\DIFadd{\% of cells. The owned cells concentrate in a few rows: Qwen2.5-VL-72B owns Atelectasis ($O_q=\cfOQwenSeventyTwoAtel$), Cardiomegaly, and Mass, while in Gemma~3 12B the competing directions win most answers, down to $O_{\mathrm{Effusion}}=\cfOGemmaTwelveEff$.
}

\DIFaddend \begin{table}[htbp]
\centering
\setlength{\tabcolsep}{4pt}
\caption{\textbf{Ownership of every concept, NIH ChestX-ray14.} $O_q$ at $\alpha=+0.25$: bold is owned, $\dagger$ a stronger competitor, $\ddagger$ unresolved; the last row counts owned cells. Shading darkens with $|O_q|$ (terracotta positive, blue negative).}
\label{tab:cf-own-nih}
\DIFdelbeginFL %DIFDELCMD < \fitwidth{%
%DIFDELCMD < \begin{tabular}{lrrrrrr}
%DIFDELCMD < \toprule
%DIFDELCMD < Checkpoint & \multicolumn{1}{c}{Effusion} & \multicolumn{1}{c}{Atelectasis} & \multicolumn{1}{c}{Pneumothorax} & \multicolumn{1}{c}{Cardiomegaly} & \multicolumn{1}{c}{Mass} & \multicolumn{1}{c}{Nodule}\\
%DIFDELCMD < \midrule
%DIFDELCMD < Lingshu 32B & \cellcolor{cfSlate1}-0.01 & \cellcolor{cfSlate2}-0.08 & \cellcolor{cfTerra3}\textbf{+0.14} & \cellcolor{cfTerra2}\textbf{+0.05} & \cellcolor{cfSlate2}-0.08 & \cellcolor{cfSlate3}-0.18 \\
%DIFDELCMD < CheXagent-2-3B & \cellcolor{cfSlate1}-0.00$^{\ddagger}$ & \cellcolor{cfTerra2}\textbf{+0.04} & \cellcolor{cfSlate1}-0.00$^{\dagger}$ & \cellcolor{cfSlate1}-0.01 & \cellcolor{cfTerra1}\textbf{+0.01} & \cellcolor{cfSlate1}-0.01 \\
%DIFDELCMD < Qwen2.5-VL-72B & \cellcolor{cfSlate4}-0.32 & \cellcolor{cfTerra4}\textbf{+0.54} & \cellcolor{cfSlate4}-0.50 & \cellcolor{cfTerra4}\textbf{+0.26} & \cellcolor{cfTerra2}\textbf{+0.07} & \cellcolor{cfSlate3}-0.23 \\
%DIFDELCMD < Lingshu 7B & \cellcolor{cfTerra1}+0.02 & \cellcolor{cfTerra1}+0.01 & \cellcolor{cfTerra1}+0.01 & \cellcolor{cfTerra3}\textbf{+0.17} & \cellcolor{cfSlate2}-0.09 & \cellcolor{cfTerra1}+0.02 \\
%DIFDELCMD < Qwen3-VL-4B & \cellcolor{cfSlate2}-0.03$^{\dagger}$ & \cellcolor{cfSlate2}-0.07 & \cellcolor{cfSlate4}-0.29 & \cellcolor{cfSlate2}-0.08 & \cellcolor{cfSlate3}-0.13 & \cellcolor{cfSlate2}-0.09 \\
%DIFDELCMD < Qwen3-VL-32B & \cellcolor{cfSlate3}-0.14 & \cellcolor{cfSlate2}-0.04 & \cellcolor{cfTerra1}+0.00$^{\ddagger}$ & \cellcolor{cfSlate2}-0.04 & \cellcolor{cfSlate3}-0.11 & \cellcolor{cfTerra3}\textbf{+0.17} \\
%DIFDELCMD < InternVL3.5-38B & \cellcolor{cfTerra1}+0.00$^{\ddagger}$ & \cellcolor{cfSlate1}-0.02 & \cellcolor{cfSlate1}-0.00 & \cellcolor{cfSlate2}-0.03 & \cellcolor{cfTerra1}+0.02 & \cellcolor{cfTerra1}+0.00 \\
%DIFDELCMD < Gemma 3 12B & \cellcolor{cfSlate4}-0.83 & \cellcolor{cfSlate4}-0.57 & \cellcolor{cfSlate3}-0.17 & \cellcolor{cfSlate3}-0.19 & \cellcolor{cfSlate4}-0.65 & \cellcolor{cfTerra3}\textbf{+0.17} \\
%DIFDELCMD < MedGemma 4B & \cellcolor{cfTerra1}+0.01 & \cellcolor{cfSlate4}-0.56 & \cellcolor{cfTerra1}+0.01 & \cellcolor{cfSlate2}-0.07 & \cellcolor{cfSlate4}-0.27 & \cellcolor{cfTerra2}+0.09 \\
%DIFDELCMD < InternVL3.5-14B & \cellcolor{cfTerra1}\textbf{+0.01} & \cellcolor{cfSlate1}-0.01 & \cellcolor{cfSlate1}-0.00 & \cellcolor{cfSlate1}-0.00 & \cellcolor{cfTerra1}+0.00 & \cellcolor{cfSlate1}-0.01 \\
%DIFDELCMD < Qwen2.5-VL-32B & \cellcolor{cfSlate3}-0.22 & \cellcolor{cfSlate2}-0.06 & \cellcolor{cfSlate2}-0.06 & \cellcolor{cfTerra3}\textbf{+0.11} & \cellcolor{cfSlate2}-0.08 & \cellcolor{cfSlate2}-0.05 \\
%DIFDELCMD < InternVL3.5-8B & \cellcolor{cfTerra1}+0.00 & \cellcolor{cfSlate1}-0.02 & \cellcolor{cfSlate1}-0.00 & \cellcolor{cfTerra1}+0.00 & \cellcolor{cfSlate1}-0.01 & \cellcolor{cfSlate1}-0.01 \\
%DIFDELCMD < HuatuoGPT-Vision-7B & \cellcolor{cfSlate2}-0.04 & \cellcolor{cfSlate2}-0.03 & \cellcolor{cfSlate2}-0.03 & \cellcolor{cfSlate2}-0.08 & \cellcolor{cfSlate2}-0.04 & \cellcolor{cfTerra1}\textbf{+0.01} \\
%DIFDELCMD < LLaVA-Med 1.5 7B & \cellcolor{cfSlate1}-0.00 & \cellcolor{cfSlate1}-0.00 & \cellcolor{cfSlate2}-0.04 & \cellcolor{cfSlate1}-0.01 & \cellcolor{cfSlate1}-0.01 & \cellcolor{cfSlate2}-0.02 \\
%DIFDELCMD < CheXagent-8B & \cellcolor{cfTerra1}+0.01 & \cellcolor{cfSlate1}-0.01 & \cellcolor{cfSlate1}-0.00 & \cellcolor{cfSlate1}-0.00 & \cellcolor{cfTerra1}+0.00 & \cellcolor{cfSlate1}-0.01 \\
%DIFDELCMD < Qwen2.5-VL-3B & \cellcolor{cfSlate2}-0.09 & \cellcolor{cfSlate2}-0.04 & \cellcolor{cfSlate3}-0.18 & \cellcolor{cfSlate3}-0.11 & \cellcolor{cfSlate2}-0.04 & \cellcolor{cfSlate2}-0.04 \\
%DIFDELCMD < Qwen2.5-VL-7B & \cellcolor{cfSlate2}-0.07$^{\dagger}$ & \cellcolor{cfSlate3}-0.23 & \cellcolor{cfSlate3}-0.10 & \cellcolor{cfSlate3}-0.22 & \cellcolor{cfSlate3}-0.17 & \cellcolor{cfSlate2}-0.07 \\
%DIFDELCMD < Qwen3-VL-8B & \cellcolor{cfTerra2}+0.07 & \cellcolor{cfSlate1}-0.01 & \cellcolor{cfSlate2}-0.07 & \cellcolor{cfTerra1}+0.00 & \cellcolor{cfSlate3}-0.11 & \cellcolor{cfTerra2}+0.03 \\
%DIFDELCMD < Gemma 3 27B & \cellcolor{cfSlate2}-0.04 & \cellcolor{cfSlate2}-0.03 & \cellcolor{cfSlate1}-0.01 & \cellcolor{cfTerra1}+0.00$^{\ddagger}$ & \cellcolor{cfSlate2}-0.08 & \cellcolor{cfSlate1}-0.02 \\
%DIFDELCMD < Gemma 3 4B & \cellcolor{cfSlate1}-0.00 & \cellcolor{cfSlate1}-0.00 & \cellcolor{cfSlate1}-0.01 & \cellcolor{cfSlate3}-0.24 & \cellcolor{cfSlate4}-0.28 & \cellcolor{cfSlate3}-0.12 \\
%DIFDELCMD < LLaVA-1.5-7B & \cellcolor{cfSlate2}-0.03 & \cellcolor{cfSlate1}-0.01 & \cellcolor{cfSlate2}-0.08 & \cellcolor{cfSlate3}-0.13 & \cellcolor{cfTerra2}+0.02 & \cellcolor{cfSlate1}-0.02 \\
%DIFDELCMD < MedGemma 27B & \cellcolor{cfSlate2}-0.06 & \cellcolor{cfSlate4}-0.45 & \cellcolor{cfSlate1}-0.01 & \cellcolor{cfTerra1}+0.00 & \cellcolor{cfSlate3}-0.13 & \cellcolor{cfSlate3}-0.20 \\
%DIFDELCMD < LLaVA-1.5-13B & \cellcolor{cfSlate1}-0.01 & \cellcolor{cfSlate2}-0.05 & \cellcolor{cfSlate2}-0.04 & \cellcolor{cfSlate2}-0.03 & \cellcolor{cfTerra1}+0.01 & \cellcolor{cfSlate1}-0.02 \\
%DIFDELCMD < Llama 3.2 Vision 11B & \cellcolor{cfSlate2}-0.04 & \cellcolor{cfSlate2}-0.07 & \cellcolor{cfSlate2}-0.05 & \cellcolor{cfTerra2}+0.02 & \cellcolor{cfSlate2}-0.06 & \cellcolor{cfSlate1}-0.01 \\
%DIFDELCMD < LLaVA-Rad-7B & \cellcolor{cfSlate1}-0.02 & \cellcolor{cfSlate1}-0.01 & \cellcolor{cfSlate1}-0.00 & \cellcolor{cfTerra1}+0.00 & \cellcolor{cfSlate1}-0.01 & \cellcolor{cfSlate2}-0.03 \\
%DIFDELCMD < \midrule
%DIFDELCMD < \textbf{Owned} & 1 & 2 & 1 & 4 & 2 & 3 \\
%DIFDELCMD < \bottomrule
%DIFDELCMD < \end{tabular}}
%DIFDELCMD < %%%
\DIFdelendFL \DIFaddbeginFL \resizebox{\linewidth}{!}{%
\begin{tabular}{lrrrrrr}
\toprule
Checkpoint & \multicolumn{1}{c}{Effusion} & \multicolumn{1}{c}{Atelectasis} & \multicolumn{1}{c}{Pneumothorax} & \multicolumn{1}{c}{Cardiomegaly} & \multicolumn{1}{c}{Mass} & \multicolumn{1}{c}{Nodule}\\
\midrule
Lingshu 32B & \cellcolor{cfSlate1}-0.01 & \cellcolor{cfSlate2}-0.08 & \cellcolor{cfTerra3}\textbf{+0.14} & \cellcolor{cfTerra2}\textbf{+0.05} & \cellcolor{cfSlate2}-0.08 & \cellcolor{cfSlate3}-0.18 \\
CheXagent-2-3B & \cellcolor{cfSlate1}-0.00$^{\ddagger}$ & \cellcolor{cfTerra2}\textbf{+0.04} & \cellcolor{cfSlate1}-0.00$^{\dagger}$ & \cellcolor{cfSlate1}-0.01 & \cellcolor{cfTerra1}\textbf{+0.01} & \cellcolor{cfSlate1}-0.01 \\
Qwen2.5-VL-72B & \cellcolor{cfSlate4}-0.32 & \cellcolor{cfTerra4}\textbf{+0.54} & \cellcolor{cfSlate4}-0.50 & \cellcolor{cfTerra4}\textbf{+0.26} & \cellcolor{cfTerra2}\textbf{+0.07} & \cellcolor{cfSlate3}-0.23 \\
Lingshu 7B & \cellcolor{cfTerra1}+0.02 & \cellcolor{cfTerra1}+0.01 & \cellcolor{cfTerra1}+0.01 & \cellcolor{cfTerra3}\textbf{+0.17} & \cellcolor{cfSlate2}-0.09 & \cellcolor{cfTerra1}+0.02 \\
Qwen3-VL-4B & \cellcolor{cfSlate2}-0.03$^{\dagger}$ & \cellcolor{cfSlate2}-0.07 & \cellcolor{cfSlate4}-0.29 & \cellcolor{cfSlate2}-0.08 & \cellcolor{cfSlate3}-0.13 & \cellcolor{cfSlate2}-0.09 \\
Qwen3-VL-32B & \cellcolor{cfSlate3}-0.14 & \cellcolor{cfSlate2}-0.04 & \cellcolor{cfTerra1}+0.00$^{\ddagger}$ & \cellcolor{cfSlate2}-0.04 & \cellcolor{cfSlate3}-0.11 & \cellcolor{cfTerra3}\textbf{+0.17} \\
InternVL3.5-38B & \cellcolor{cfTerra1}+0.00$^{\ddagger}$ & \cellcolor{cfSlate1}-0.02 & \cellcolor{cfSlate1}-0.00 & \cellcolor{cfSlate2}-0.03 & \cellcolor{cfTerra1}+0.02 & \cellcolor{cfTerra1}+0.00 \\
Gemma 3 12B & \cellcolor{cfSlate4}-0.83 & \cellcolor{cfSlate4}-0.57 & \cellcolor{cfSlate3}-0.17 & \cellcolor{cfSlate3}-0.19 & \cellcolor{cfSlate4}-0.65 & \cellcolor{cfTerra3}\textbf{+0.17} \\
MedGemma 4B & \cellcolor{cfTerra1}+0.01 & \cellcolor{cfSlate4}-0.56 & \cellcolor{cfTerra1}+0.01 & \cellcolor{cfSlate2}-0.07 & \cellcolor{cfSlate4}-0.27 & \cellcolor{cfTerra2}+0.09 \\
InternVL3.5-14B & \cellcolor{cfTerra1}\textbf{+0.01} & \cellcolor{cfSlate1}-0.01 & \cellcolor{cfSlate1}-0.00 & \cellcolor{cfSlate1}-0.00 & \cellcolor{cfTerra1}+0.00 & \cellcolor{cfSlate1}-0.01 \\
Qwen2.5-VL-32B & \cellcolor{cfSlate3}-0.22 & \cellcolor{cfSlate2}-0.06 & \cellcolor{cfSlate2}-0.06 & \cellcolor{cfTerra3}\textbf{+0.11} & \cellcolor{cfSlate2}-0.08 & \cellcolor{cfSlate2}-0.05 \\
InternVL3.5-8B & \cellcolor{cfTerra1}+0.00 & \cellcolor{cfSlate1}-0.02 & \cellcolor{cfSlate1}-0.00 & \cellcolor{cfTerra1}+0.00 & \cellcolor{cfSlate1}-0.01 & \cellcolor{cfSlate1}-0.01 \\
HuatuoGPT-Vision-7B & \cellcolor{cfSlate2}-0.04 & \cellcolor{cfSlate2}-0.03 & \cellcolor{cfSlate2}-0.03 & \cellcolor{cfSlate2}-0.08 & \cellcolor{cfSlate2}-0.04 & \cellcolor{cfTerra1}\textbf{+0.01} \\
LLaVA-Med 1.5 7B & \cellcolor{cfSlate1}-0.00 & \cellcolor{cfSlate1}-0.00 & \cellcolor{cfSlate2}-0.04 & \cellcolor{cfSlate1}-0.01 & \cellcolor{cfSlate1}-0.01 & \cellcolor{cfSlate2}-0.02 \\
CheXagent-8B & \cellcolor{cfTerra1}+0.01 & \cellcolor{cfSlate1}-0.01 & \cellcolor{cfSlate1}-0.00 & \cellcolor{cfSlate1}-0.00 & \cellcolor{cfTerra1}+0.00 & \cellcolor{cfSlate1}-0.01 \\
Qwen2.5-VL-3B & \cellcolor{cfSlate2}-0.09 & \cellcolor{cfSlate2}-0.04 & \cellcolor{cfSlate3}-0.18 & \cellcolor{cfSlate3}-0.11 & \cellcolor{cfSlate2}-0.04 & \cellcolor{cfSlate2}-0.04 \\
Qwen2.5-VL-7B & \cellcolor{cfSlate2}-0.07$^{\dagger}$ & \cellcolor{cfSlate3}-0.23 & \cellcolor{cfSlate3}-0.10 & \cellcolor{cfSlate3}-0.22 & \cellcolor{cfSlate3}-0.17 & \cellcolor{cfSlate2}-0.07 \\
Qwen3-VL-8B & \cellcolor{cfTerra2}+0.07 & \cellcolor{cfSlate1}-0.01 & \cellcolor{cfSlate2}-0.07 & \cellcolor{cfTerra1}+0.00 & \cellcolor{cfSlate3}-0.11 & \cellcolor{cfTerra2}+0.03 \\
Gemma 3 27B & \cellcolor{cfSlate2}-0.04 & \cellcolor{cfSlate2}-0.03 & \cellcolor{cfSlate1}-0.01 & \cellcolor{cfTerra1}+0.00$^{\ddagger}$ & \cellcolor{cfSlate2}-0.08 & \cellcolor{cfSlate1}-0.02 \\
Gemma 3 4B & \cellcolor{cfSlate1}-0.00 & \cellcolor{cfSlate1}-0.00 & \cellcolor{cfSlate1}-0.01 & \cellcolor{cfSlate3}-0.24 & \cellcolor{cfSlate4}-0.28 & \cellcolor{cfSlate3}-0.12 \\
LLaVA-1.5-7B & \cellcolor{cfSlate2}-0.03 & \cellcolor{cfSlate1}-0.01 & \cellcolor{cfSlate2}-0.08 & \cellcolor{cfSlate3}-0.13 & \cellcolor{cfTerra2}+0.02 & \cellcolor{cfSlate1}-0.02 \\
MedGemma 27B & \cellcolor{cfSlate2}-0.06 & \cellcolor{cfSlate4}-0.45 & \cellcolor{cfSlate1}-0.01 & \cellcolor{cfTerra1}+0.00 & \cellcolor{cfSlate3}-0.13 & \cellcolor{cfSlate3}-0.20 \\
LLaVA-1.5-13B & \cellcolor{cfSlate1}-0.01 & \cellcolor{cfSlate2}-0.05 & \cellcolor{cfSlate2}-0.04 & \cellcolor{cfSlate2}-0.03 & \cellcolor{cfTerra1}+0.01 & \cellcolor{cfSlate1}-0.02 \\
Llama 3.2 Vision 11B & \cellcolor{cfSlate2}-0.04 & \cellcolor{cfSlate2}-0.07 & \cellcolor{cfSlate2}-0.05 & \cellcolor{cfTerra2}+0.02 & \cellcolor{cfSlate2}-0.06 & \cellcolor{cfSlate1}-0.01 \\
LLaVA-Rad-7B & \cellcolor{cfSlate1}-0.02 & \cellcolor{cfSlate1}-0.01 & \cellcolor{cfSlate1}-0.00 & \cellcolor{cfTerra1}+0.00 & \cellcolor{cfSlate1}-0.01 & \cellcolor{cfSlate2}-0.03 \\
\midrule
\textbf{Owned} & 1 & 2 & 1 & 4 & 2 & 3 \\
\bottomrule
\end{tabular}}
\DIFaddendFL \end{table}

\DIFaddbegin \paragraph{\DIFadd{CheXpert Plus.}} \cfChexOwned{} \DIFadd{of the }\cfChexCells{} \DIFadd{cells are owned, }\cfChexCompetitor{} \DIFadd{have a stronger competitor, and the contrast is positive in }\cfGeoOwnWinsChexPct\DIFadd{\% of cells. Lingshu contributes the largest share of the owned cells, with three concepts at 7B and five at 32B.
}

\DIFaddend \begin{table}[htbp]
\centering
\setlength{\tabcolsep}{4pt}
\caption{\textbf{Ownership of every concept, CheXpert Plus.} As Table~\ref{tab:cf-own-nih}.}
\label{tab:cf-own-chexpert}
\DIFdelbeginFL %DIFDELCMD < \fitwidth{%
%DIFDELCMD < \begin{tabular}{lrrrrrr}
%DIFDELCMD < \toprule
%DIFDELCMD < Checkpoint & \multicolumn{1}{c}{Effusion} & \multicolumn{1}{c}{Atelectasis} & \multicolumn{1}{c}{Pneumothorax} & \multicolumn{1}{c}{Cardiomegaly} & \multicolumn{1}{c}{Consolidation} & \multicolumn{1}{c}{Edema}\\
%DIFDELCMD < \midrule
%DIFDELCMD < Lingshu 32B & \cellcolor{cfTerra2}\textbf{+0.03} & \cellcolor{cfSlate2}-0.04$^{\dagger}$ & \cellcolor{cfTerra2}\textbf{+0.02} & \cellcolor{cfTerra3}\textbf{+0.11} & \cellcolor{cfTerra2}\textbf{+0.03} & \cellcolor{cfTerra3}\textbf{+0.11} \\
%DIFDELCMD < CheXagent-2-3B & \cellcolor{cfTerra2}\textbf{+0.03} & \cellcolor{cfSlate2}-0.02 & \cellcolor{cfTerra1}\textbf{+0.01} & \cellcolor{cfTerra1}\textbf{+0.02} & \cellcolor{cfTerra1}\textbf{+0.01} & \cellcolor{cfSlate1}-0.00$^{\ddagger}$ \\
%DIFDELCMD < Qwen2.5-VL-72B & \cellcolor{cfSlate3}-0.12 & \cellcolor{cfTerra2}+0.09 & \cellcolor{cfTerra2}\textbf{+0.03} & \cellcolor{cfSlate3}-0.11 & \cellcolor{cfSlate4}-0.29 & \cellcolor{cfSlate3}-0.13 \\
%DIFDELCMD < Lingshu 7B & \cellcolor{cfTerra2}\textbf{+0.10} & \cellcolor{cfSlate2}-0.09 & \cellcolor{cfTerra2}\textbf{+0.09} & \cellcolor{cfTerra1}+0.01$^{\ddagger}$ & \cellcolor{cfSlate2}-0.05 & \cellcolor{cfTerra3}\textbf{+0.21} \\
%DIFDELCMD < Qwen3-VL-4B & \cellcolor{cfSlate1}-0.01$^{\dagger}$ & \cellcolor{cfSlate1}-0.02 & \cellcolor{cfSlate3}-0.23 & \cellcolor{cfSlate1}-0.01 & \cellcolor{cfTerra2}\textbf{+0.09} & \cellcolor{cfTerra2}\textbf{+0.09} \\
%DIFDELCMD < Qwen3-VL-32B & \cellcolor{cfSlate2}-0.04 & \cellcolor{cfTerra1}+0.00$^{\ddagger}$ & \cellcolor{cfTerra1}+0.00 & \cellcolor{cfSlate2}-0.03 & \cellcolor{cfSlate3}-0.12 & \cellcolor{cfTerra2}\textbf{+0.05} \\
%DIFDELCMD < InternVL3.5-38B & \cellcolor{cfTerra1}\textbf{+0.02} & \cellcolor{cfSlate1}-0.01 & \cellcolor{cfSlate1}-0.00 & \cellcolor{cfSlate1}-0.00$^{\ddagger}$ & \cellcolor{cfTerra1}\textbf{+0.01} & \cellcolor{cfSlate2}-0.02 \\
%DIFDELCMD < Gemma 3 12B & \cellcolor{cfSlate4}-0.83 & \cellcolor{cfSlate3}-0.13 & \cellcolor{cfSlate4}-0.71 & \cellcolor{cfSlate4}-0.47 & \cellcolor{cfSlate4}-0.26 & \cellcolor{cfTerra3}\textbf{+0.15} \\
%DIFDELCMD < MedGemma 4B & \cellcolor{cfSlate3}-0.20 & \cellcolor{cfSlate2}-0.08 & \cellcolor{cfSlate4}-0.60 & \cellcolor{cfTerra2}\textbf{+0.04} & \cellcolor{cfSlate3}-0.20 & \cellcolor{cfTerra3}\textbf{+0.10} \\
%DIFDELCMD < InternVL3.5-14B & \cellcolor{cfTerra1}\textbf{+0.01} & \cellcolor{cfSlate1}-0.00 & \cellcolor{cfTerra1}+0.00 & \cellcolor{cfSlate1}-0.01 & \cellcolor{cfTerra1}+0.01 & \cellcolor{cfSlate1}-0.00 \\
%DIFDELCMD < Qwen2.5-VL-32B & \cellcolor{cfSlate2}-0.07 & \cellcolor{cfSlate3}-0.19 & \cellcolor{cfSlate3}-0.12 & \cellcolor{cfSlate3}-0.19 & \cellcolor{cfSlate2}-0.05 & \cellcolor{cfSlate2}-0.02 \\
%DIFDELCMD < InternVL3.5-8B & \cellcolor{cfSlate1}-0.01 & \cellcolor{cfSlate1}-0.01 & \cellcolor{cfTerra1}\textbf{+0.00} & \cellcolor{cfSlate1}-0.00 & \cellcolor{cfTerra1}+0.01 & \cellcolor{cfSlate1}-0.00 \\
%DIFDELCMD < HuatuoGPT-Vision-7B & \cellcolor{cfSlate2}-0.02 & \cellcolor{cfSlate1}-0.02 & \cellcolor{cfSlate1}-0.00 & \cellcolor{cfTerra1}+0.01 & \cellcolor{cfSlate2}-0.03 & \cellcolor{cfSlate1}-0.01 \\
%DIFDELCMD < LLaVA-Med 1.5 7B & \cellcolor{cfTerra1}\textbf{+0.00} & \cellcolor{cfSlate1}-0.00 & \cellcolor{cfSlate2}-0.02 & \cellcolor{cfSlate2}-0.03 & \cellcolor{cfSlate1}-0.00$^{\dagger}$ & \cellcolor{cfSlate1}-0.02 \\
%DIFDELCMD < CheXagent-8B & \cellcolor{cfSlate2}-0.02 & \cellcolor{cfSlate1}-0.02 & \cellcolor{cfSlate2}-0.02 & \cellcolor{cfSlate2}-0.07 & \cellcolor{cfTerra2}\textbf{+0.04} & \cellcolor{cfSlate2}-0.02 \\
%DIFDELCMD < Qwen2.5-VL-3B & \cellcolor{cfSlate3}-0.18 & \cellcolor{cfSlate1}-0.02 & \cellcolor{cfSlate2}-0.04 & \cellcolor{cfSlate2}-0.05 & \cellcolor{cfSlate2}-0.03 & \cellcolor{cfSlate2}-0.08 \\
%DIFDELCMD < Qwen2.5-VL-7B & \cellcolor{cfTerra1}+0.01 & \cellcolor{cfSlate2}-0.05 & \cellcolor{cfSlate2}-0.10 & \cellcolor{cfSlate1}-0.00 & \cellcolor{cfSlate3}-0.15 & \cellcolor{cfSlate3}-0.18 \\
%DIFDELCMD < Qwen3-VL-8B & \cellcolor{cfSlate2}-0.07$^{\dagger}$ & \cellcolor{cfSlate1}-0.00$^{\dagger}$ & \cellcolor{cfSlate2}-0.10 & \cellcolor{cfSlate1}-0.00 & \cellcolor{cfSlate2}-0.04$^{\dagger}$ & \cellcolor{cfSlate2}-0.06 \\
%DIFDELCMD < Gemma 3 27B & \cellcolor{cfSlate2}-0.03 & \cellcolor{cfSlate1}-0.00 & \cellcolor{cfSlate1}-0.01 & \cellcolor{cfSlate1}-0.00 & \cellcolor{cfSlate2}-0.03 & \cellcolor{cfTerra1}+0.00$^{\ddagger}$ \\
%DIFDELCMD < Gemma 3 4B & \cellcolor{cfSlate1}-0.00 & \cellcolor{cfSlate1}-0.00 & \cellcolor{cfSlate4}-0.65 & \cellcolor{cfSlate1}-0.02 & \cellcolor{cfSlate4}-0.35 & \cellcolor{cfSlate4}-0.58 \\
%DIFDELCMD < LLaVA-1.5-7B & \cellcolor{cfTerra1}+0.00 & \cellcolor{cfSlate2}-0.09 & \cellcolor{cfSlate2}-0.10 & \cellcolor{cfTerra2}+0.03 & \cellcolor{cfSlate2}-0.03 & \cellcolor{cfSlate1}-0.00 \\
%DIFDELCMD < MedGemma 27B & \cellcolor{cfSlate1}-0.01 & \cellcolor{cfTerra1}+0.00 & \cellcolor{cfSlate2}-0.06 & \cellcolor{cfTerra2}+0.03 & \cellcolor{cfSlate1}-0.02 & \cellcolor{cfSlate4}-0.28 \\
%DIFDELCMD < LLaVA-1.5-13B & \cellcolor{cfSlate1}-0.00 & \cellcolor{cfSlate2}-0.03 & \cellcolor{cfSlate1}-0.01 & \cellcolor{cfTerra1}+0.00 & \cellcolor{cfTerra1}+0.01 & \cellcolor{cfSlate1}-0.01 \\
%DIFDELCMD < Llama 3.2 Vision 11B & \cellcolor{cfSlate1}-0.01 & \cellcolor{cfSlate2}-0.05 & \cellcolor{cfSlate2}-0.08 & \cellcolor{cfSlate2}-0.07 & \cellcolor{cfSlate2}-0.05 & \cellcolor{cfSlate1}-0.01 \\
%DIFDELCMD < LLaVA-Rad-7B & \cellcolor{cfSlate2}-0.04 & \cellcolor{cfSlate1}-0.01 & \cellcolor{cfTerra1}+0.00 & \cellcolor{cfSlate2}-0.07 & \cellcolor{cfTerra2}+0.02 & \cellcolor{cfSlate1}-0.00 \\
%DIFDELCMD < \midrule
%DIFDELCMD < \textbf{Owned} & 6 & 0 & 5 & 3 & 5 & 6 \\
%DIFDELCMD < \bottomrule
%DIFDELCMD < \end{tabular}}
%DIFDELCMD < %%%
\DIFdelendFL \DIFaddbeginFL \resizebox{\linewidth}{!}{%
\begin{tabular}{lrrrrrr}
\toprule
Checkpoint & \multicolumn{1}{c}{Effusion} & \multicolumn{1}{c}{Atelectasis} & \multicolumn{1}{c}{Pneumothorax} & \multicolumn{1}{c}{Cardiomegaly} & \multicolumn{1}{c}{Consolidation} & \multicolumn{1}{c}{Edema}\\
\midrule
Lingshu 32B & \cellcolor{cfTerra2}\textbf{+0.03} & \cellcolor{cfSlate2}-0.04$^{\dagger}$ & \cellcolor{cfTerra2}\textbf{+0.02} & \cellcolor{cfTerra3}\textbf{+0.11} & \cellcolor{cfTerra2}\textbf{+0.03} & \cellcolor{cfTerra3}\textbf{+0.11} \\
CheXagent-2-3B & \cellcolor{cfTerra2}\textbf{+0.03} & \cellcolor{cfSlate2}-0.02 & \cellcolor{cfTerra1}\textbf{+0.01} & \cellcolor{cfTerra1}\textbf{+0.02} & \cellcolor{cfTerra1}\textbf{+0.01} & \cellcolor{cfSlate1}-0.00$^{\ddagger}$ \\
Qwen2.5-VL-72B & \cellcolor{cfSlate3}-0.12 & \cellcolor{cfTerra2}+0.09 & \cellcolor{cfTerra2}\textbf{+0.03} & \cellcolor{cfSlate3}-0.11 & \cellcolor{cfSlate4}-0.29 & \cellcolor{cfSlate3}-0.13 \\
Lingshu 7B & \cellcolor{cfTerra2}\textbf{+0.10} & \cellcolor{cfSlate2}-0.09 & \cellcolor{cfTerra2}\textbf{+0.09} & \cellcolor{cfTerra1}+0.01$^{\ddagger}$ & \cellcolor{cfSlate2}-0.05 & \cellcolor{cfTerra3}\textbf{+0.21} \\
Qwen3-VL-4B & \cellcolor{cfSlate1}-0.01$^{\dagger}$ & \cellcolor{cfSlate1}-0.02 & \cellcolor{cfSlate3}-0.23 & \cellcolor{cfSlate1}-0.01 & \cellcolor{cfTerra2}\textbf{+0.09} & \cellcolor{cfTerra2}\textbf{+0.09} \\
Qwen3-VL-32B & \cellcolor{cfSlate2}-0.04 & \cellcolor{cfTerra1}+0.00$^{\ddagger}$ & \cellcolor{cfTerra1}+0.00 & \cellcolor{cfSlate2}-0.03 & \cellcolor{cfSlate3}-0.12 & \cellcolor{cfTerra2}\textbf{+0.05} \\
InternVL3.5-38B & \cellcolor{cfTerra1}\textbf{+0.02} & \cellcolor{cfSlate1}-0.01 & \cellcolor{cfSlate1}-0.00 & \cellcolor{cfSlate1}-0.00$^{\ddagger}$ & \cellcolor{cfTerra1}\textbf{+0.01} & \cellcolor{cfSlate2}-0.02 \\
Gemma 3 12B & \cellcolor{cfSlate4}-0.83 & \cellcolor{cfSlate3}-0.13 & \cellcolor{cfSlate4}-0.71 & \cellcolor{cfSlate4}-0.47 & \cellcolor{cfSlate4}-0.26 & \cellcolor{cfTerra3}\textbf{+0.15} \\
MedGemma 4B & \cellcolor{cfSlate3}-0.20 & \cellcolor{cfSlate2}-0.08 & \cellcolor{cfSlate4}-0.60 & \cellcolor{cfTerra2}\textbf{+0.04} & \cellcolor{cfSlate3}-0.20 & \cellcolor{cfTerra3}\textbf{+0.10} \\
InternVL3.5-14B & \cellcolor{cfTerra1}\textbf{+0.01} & \cellcolor{cfSlate1}-0.00 & \cellcolor{cfTerra1}+0.00 & \cellcolor{cfSlate1}-0.01 & \cellcolor{cfTerra1}+0.01 & \cellcolor{cfSlate1}-0.00 \\
Qwen2.5-VL-32B & \cellcolor{cfSlate2}-0.07 & \cellcolor{cfSlate3}-0.19 & \cellcolor{cfSlate3}-0.12 & \cellcolor{cfSlate3}-0.19 & \cellcolor{cfSlate2}-0.05 & \cellcolor{cfSlate2}-0.02 \\
InternVL3.5-8B & \cellcolor{cfSlate1}-0.01 & \cellcolor{cfSlate1}-0.01 & \cellcolor{cfTerra1}\textbf{+0.00} & \cellcolor{cfSlate1}-0.00 & \cellcolor{cfTerra1}+0.01 & \cellcolor{cfSlate1}-0.00 \\
HuatuoGPT-Vision-7B & \cellcolor{cfSlate2}-0.02 & \cellcolor{cfSlate1}-0.02 & \cellcolor{cfSlate1}-0.00 & \cellcolor{cfTerra1}+0.01 & \cellcolor{cfSlate2}-0.03 & \cellcolor{cfSlate1}-0.01 \\
LLaVA-Med 1.5 7B & \cellcolor{cfTerra1}\textbf{+0.00} & \cellcolor{cfSlate1}-0.00 & \cellcolor{cfSlate2}-0.02 & \cellcolor{cfSlate2}-0.03 & \cellcolor{cfSlate1}-0.00$^{\dagger}$ & \cellcolor{cfSlate1}-0.02 \\
CheXagent-8B & \cellcolor{cfSlate2}-0.02 & \cellcolor{cfSlate1}-0.02 & \cellcolor{cfSlate2}-0.02 & \cellcolor{cfSlate2}-0.07 & \cellcolor{cfTerra2}\textbf{+0.04} & \cellcolor{cfSlate2}-0.02 \\
Qwen2.5-VL-3B & \cellcolor{cfSlate3}-0.18 & \cellcolor{cfSlate1}-0.02 & \cellcolor{cfSlate2}-0.04 & \cellcolor{cfSlate2}-0.05 & \cellcolor{cfSlate2}-0.03 & \cellcolor{cfSlate2}-0.08 \\
Qwen2.5-VL-7B & \cellcolor{cfTerra1}+0.01 & \cellcolor{cfSlate2}-0.05 & \cellcolor{cfSlate2}-0.10 & \cellcolor{cfSlate1}-0.00 & \cellcolor{cfSlate3}-0.15 & \cellcolor{cfSlate3}-0.18 \\
Qwen3-VL-8B & \cellcolor{cfSlate2}-0.07$^{\dagger}$ & \cellcolor{cfSlate1}-0.00$^{\dagger}$ & \cellcolor{cfSlate2}-0.10 & \cellcolor{cfSlate1}-0.00 & \cellcolor{cfSlate2}-0.04$^{\dagger}$ & \cellcolor{cfSlate2}-0.06 \\
Gemma 3 27B & \cellcolor{cfSlate2}-0.03 & \cellcolor{cfSlate1}-0.00 & \cellcolor{cfSlate1}-0.01 & \cellcolor{cfSlate1}-0.00 & \cellcolor{cfSlate2}-0.03 & \cellcolor{cfTerra1}+0.00$^{\ddagger}$ \\
Gemma 3 4B & \cellcolor{cfSlate1}-0.00 & \cellcolor{cfSlate1}-0.00 & \cellcolor{cfSlate4}-0.65 & \cellcolor{cfSlate1}-0.02 & \cellcolor{cfSlate4}-0.35 & \cellcolor{cfSlate4}-0.58 \\
LLaVA-1.5-7B & \cellcolor{cfTerra1}+0.00 & \cellcolor{cfSlate2}-0.09 & \cellcolor{cfSlate2}-0.10 & \cellcolor{cfTerra2}+0.03 & \cellcolor{cfSlate2}-0.03 & \cellcolor{cfSlate1}-0.00 \\
MedGemma 27B & \cellcolor{cfSlate1}-0.01 & \cellcolor{cfTerra1}+0.00 & \cellcolor{cfSlate2}-0.06 & \cellcolor{cfTerra2}+0.03 & \cellcolor{cfSlate1}-0.02 & \cellcolor{cfSlate4}-0.28 \\
LLaVA-1.5-13B & \cellcolor{cfSlate1}-0.00 & \cellcolor{cfSlate2}-0.03 & \cellcolor{cfSlate1}-0.01 & \cellcolor{cfTerra1}+0.00 & \cellcolor{cfTerra1}+0.01 & \cellcolor{cfSlate1}-0.01 \\
Llama 3.2 Vision 11B & \cellcolor{cfSlate1}-0.01 & \cellcolor{cfSlate2}-0.05 & \cellcolor{cfSlate2}-0.08 & \cellcolor{cfSlate2}-0.07 & \cellcolor{cfSlate2}-0.05 & \cellcolor{cfSlate1}-0.01 \\
LLaVA-Rad-7B & \cellcolor{cfSlate2}-0.04 & \cellcolor{cfSlate1}-0.01 & \cellcolor{cfTerra1}+0.00 & \cellcolor{cfSlate2}-0.07 & \cellcolor{cfTerra2}+0.02 & \cellcolor{cfSlate1}-0.00 \\
\midrule
\textbf{Owned} & 6 & 0 & 5 & 3 & 5 & 6 \\
\bottomrule
\end{tabular}}
\DIFaddendFL \end{table}

\DIFaddbegin \paragraph{\DIFadd{COCO.}} \cfCocoOwned{} \DIFadd{of the }\cfCocoOwnedN{} \DIFadd{cells are owned and }\cfCocoCompetitor{} \DIFadd{have a stronger competitor, and }\cfCocoBlocksFourPlus{} \DIFadd{of the }\cfCocoBlocks{} \DIFadd{checkpoints own at least four of the six objects. The rows without owned objects are CheXagent-8B and LLaVA-Rad-7B, whose clean answers to the COCO questions are near chance (Table~\ref{tab:cf-main}).
}

\DIFaddend \begin{table}[htbp]
\centering
\DIFdelbeginFL %DIFDELCMD < \setlength{\tabcolsep}{4pt}
%DIFDELCMD < %%%
\DIFdelendFL \DIFaddbeginFL \setlength{\tabcolsep}{11pt}
\begin{DIFnomarkup}\cfnote{表宽}{全部表格由 \texttt{\textbackslash fitwidth}（只缩不放，窄表会留白）改为 \texttt{\textbackslash resizebox\{\textbackslash linewidth\}\{!\}}，一律撑满版心宽度。本表自然宽度只有版心的 75\%（物体名短），直接放大会让字号比 NIH、CheXpert 两张表大三分之一，所以把列间距从 4pt 加到 11pt，使自然宽度接近版心，再缩放到满宽，三张热力表字号一致。其余表格原本就接近或超过版心宽度，视觉上几乎不变。}\end{DIFnomarkup}

\DIFaddendFL \caption{\textbf{Ownership of every concept, COCO.} As Table~\ref{tab:cf-own-nih}.}
\label{tab:cf-own-coco}
\DIFdelbeginFL %DIFDELCMD < \fitwidth{%
%DIFDELCMD < \begin{tabular}{lrrrrrr}
%DIFDELCMD < \toprule
%DIFDELCMD < Checkpoint & \multicolumn{1}{c}{person} & \multicolumn{1}{c}{dog} & \multicolumn{1}{c}{car} & \multicolumn{1}{c}{chair} & \multicolumn{1}{c}{bottle} & \multicolumn{1}{c}{bicycle}\\
%DIFDELCMD < \midrule
%DIFDELCMD < Lingshu 32B & \cellcolor{cfTerra4}\textbf{+0.29} & \cellcolor{cfTerra3}\textbf{+0.23} & \cellcolor{cfTerra2}\textbf{+0.08} & \cellcolor{cfTerra4}\textbf{+0.53} & \cellcolor{cfTerra4}\textbf{+0.34} & \cellcolor{cfTerra3}\textbf{+0.22} \\
%DIFDELCMD < CheXagent-2-3B & \cellcolor{cfTerra2}\textbf{+0.07} & \cellcolor{cfTerra2}\textbf{+0.07} & \cellcolor{cfTerra2}+0.03 & \cellcolor{cfSlate1}-0.01 & \cellcolor{cfTerra2}+0.04 & \cellcolor{cfTerra1}+0.00 \\
%DIFDELCMD < Qwen2.5-VL-72B & \cellcolor{cfTerra4}\textbf{+0.66} & \cellcolor{cfTerra4}\textbf{+0.44} & \cellcolor{cfTerra4}\textbf{+0.95} & \cellcolor{cfTerra4}\textbf{+0.89} & \cellcolor{cfTerra4}\textbf{+0.62} & \cellcolor{cfTerra4}\textbf{+0.40} \\
%DIFDELCMD < Lingshu 7B & \cellcolor{cfTerra4}\textbf{+0.39} & \cellcolor{cfTerra4}\textbf{+0.74} & \cellcolor{cfTerra4}\textbf{+0.70} & \cellcolor{cfTerra4}\textbf{+0.76} & \cellcolor{cfTerra4}\textbf{+0.64} & \cellcolor{cfTerra4}\textbf{+0.85} \\
%DIFDELCMD < Qwen3-VL-4B & \cellcolor{cfTerra4}\textbf{+0.46} & \cellcolor{cfTerra3}\textbf{+0.13} & \cellcolor{cfTerra4}\textbf{+0.33} & \cellcolor{cfTerra4}\textbf{+0.71} & \cellcolor{cfTerra4}\textbf{+0.34} & \cellcolor{cfTerra2}\textbf{+0.07} \\
%DIFDELCMD < Qwen3-VL-32B & \cellcolor{cfTerra3}\textbf{+0.13} & \cellcolor{cfTerra3}\textbf{+0.12} & \cellcolor{cfTerra3}\textbf{+0.11} & \cellcolor{cfTerra3}\textbf{+0.23} & \cellcolor{cfTerra3}\textbf{+0.24} & \cellcolor{cfTerra2}\textbf{+0.07} \\
%DIFDELCMD < InternVL3.5-38B & \cellcolor{cfTerra1}\textbf{+0.01} & \cellcolor{cfTerra1}\textbf{+0.00} & \cellcolor{cfTerra1}\textbf{+0.01} & \cellcolor{cfTerra1}\textbf{+0.02} & \cellcolor{cfTerra1}\textbf{+0.01} & \cellcolor{cfTerra1}\textbf{+0.01} \\
%DIFDELCMD < Gemma 3 12B & \cellcolor{cfTerra3}\textbf{+0.19} & \cellcolor{cfTerra4}\textbf{+0.76} & \cellcolor{cfTerra4}\textbf{+0.48} & \cellcolor{cfTerra4}\textbf{+0.45} & \cellcolor{cfTerra3}\textbf{+0.18} & \cellcolor{cfTerra4}\textbf{+0.61} \\
%DIFDELCMD < MedGemma 4B & \cellcolor{cfTerra4}\textbf{+0.35} & \cellcolor{cfTerra2}\textbf{+0.03} & \cellcolor{cfTerra4}\textbf{+0.43} & \cellcolor{cfTerra3}\textbf{+0.19} & \cellcolor{cfTerra3}\textbf{+0.20} & \cellcolor{cfTerra3}\textbf{+0.14} \\
%DIFDELCMD < InternVL3.5-14B & \cellcolor{cfTerra1}\textbf{+0.02} & \cellcolor{cfTerra1}\textbf{+0.01} & \cellcolor{cfTerra1}\textbf{+0.01} & \cellcolor{cfTerra2}\textbf{+0.04} & \cellcolor{cfTerra2}\textbf{+0.03} & \cellcolor{cfTerra1}+0.00$^{\ddagger}$ \\
%DIFDELCMD < Qwen2.5-VL-32B & \cellcolor{cfTerra4}\textbf{+0.44} & \cellcolor{cfTerra3}\textbf{+0.12} & \cellcolor{cfTerra4}\textbf{+0.27} & \cellcolor{cfTerra4}\textbf{+0.83} & \cellcolor{cfTerra4}\textbf{+0.41} & \cellcolor{cfTerra3}\textbf{+0.19} \\
%DIFDELCMD < InternVL3.5-8B & \cellcolor{cfTerra2}\textbf{+0.03} & \cellcolor{cfTerra1}\textbf{+0.01} & \cellcolor{cfTerra2}\textbf{+0.04} & \cellcolor{cfTerra2}\textbf{+0.04} & \cellcolor{cfTerra2}\textbf{+0.04} & \cellcolor{cfTerra1}\textbf{+0.02} \\
%DIFDELCMD < HuatuoGPT-Vision-7B & \cellcolor{cfTerra1}\textbf{+0.02} & \cellcolor{cfTerra1}+0.00$^{\ddagger}$ & \cellcolor{cfTerra1}\textbf{+0.00} & \cellcolor{cfTerra2}\textbf{+0.03} & \cellcolor{cfTerra2}\textbf{+0.03} & \cellcolor{cfTerra1}+0.00$^{\ddagger}$ \\
%DIFDELCMD < LLaVA-Med 1.5 7B & \cellcolor{cfTerra1}+0.02 & \cellcolor{cfTerra1}\textbf{+0.01} & \cellcolor{cfTerra1}+0.01 & \cellcolor{cfSlate2}-0.02 & \cellcolor{cfTerra1}+0.00 & \cellcolor{cfSlate1}-0.01 \\
%DIFDELCMD < CheXagent-8B & \cellcolor{cfSlate1}-0.00 & \cellcolor{cfTerra1}+0.01 & \cellcolor{cfSlate1}-0.01 & \cellcolor{cfSlate2}-0.03 & \cellcolor{cfSlate3}-0.13 & \cellcolor{cfSlate2}-0.03 \\
%DIFDELCMD < Qwen2.5-VL-3B & \cellcolor{cfTerra4}\textbf{+0.46} & \cellcolor{cfTerra4}\textbf{+0.42} & \cellcolor{cfTerra4}\textbf{+0.53} & \cellcolor{cfTerra4}\textbf{+0.68} & \cellcolor{cfTerra4}\textbf{+0.61} & \cellcolor{cfTerra4}\textbf{+0.39} \\
%DIFDELCMD < Qwen2.5-VL-7B & \cellcolor{cfTerra4}\textbf{+0.48} & \cellcolor{cfTerra4}\textbf{+0.75} & \cellcolor{cfTerra4}\textbf{+0.61} & \cellcolor{cfTerra4}\textbf{+0.87} & \cellcolor{cfTerra4}\textbf{+0.51} & \cellcolor{cfTerra4}\textbf{+0.81} \\
%DIFDELCMD < Qwen3-VL-8B & \cellcolor{cfTerra4}\textbf{+0.42} & \cellcolor{cfTerra4}\textbf{+0.51} & \cellcolor{cfTerra4}\textbf{+0.74} & \cellcolor{cfTerra4}\textbf{+0.79} & \cellcolor{cfTerra4}\textbf{+0.59} & \cellcolor{cfTerra4}\textbf{+0.43} \\
%DIFDELCMD < Gemma 3 27B & \cellcolor{cfTerra3}\textbf{+0.24} & \cellcolor{cfTerra4}\textbf{+0.36} & \cellcolor{cfTerra2}\textbf{+0.09} & \cellcolor{cfTerra4}\textbf{+0.28} & \cellcolor{cfTerra3}\textbf{+0.15} & \cellcolor{cfTerra4}\textbf{+0.75} \\
%DIFDELCMD < Gemma 3 4B & \cellcolor{cfTerra3}\textbf{+0.16} & \cellcolor{cfTerra4}\textbf{+0.54} & \cellcolor{cfTerra3}\textbf{+0.24} & \cellcolor{cfTerra4}\textbf{+0.34} & \cellcolor{cfTerra4}\textbf{+0.44} & \cellcolor{cfSlate4}-0.27 \\
%DIFDELCMD < LLaVA-1.5-7B & \cellcolor{cfTerra2}\textbf{+0.06} & \cellcolor{cfTerra1}\textbf{+0.00} & \cellcolor{cfTerra1}\textbf{+0.01} & \cellcolor{cfTerra2}\textbf{+0.07} & \cellcolor{cfTerra1}+0.01 & \cellcolor{cfTerra1}\textbf{+0.00} \\
%DIFDELCMD < MedGemma 27B & \cellcolor{cfTerra4}\textbf{+0.33} & \cellcolor{cfSlate2}-0.04 & \cellcolor{cfTerra2}\textbf{+0.08} & \cellcolor{cfTerra3}\textbf{+0.18} & \cellcolor{cfTerra3}\textbf{+0.21} & \cellcolor{cfSlate2}-0.04 \\
%DIFDELCMD < LLaVA-1.5-13B & \cellcolor{cfTerra2}\textbf{+0.05} & \cellcolor{cfTerra1}+0.00$^{\ddagger}$ & \cellcolor{cfTerra1}+0.02 & \cellcolor{cfTerra2}\textbf{+0.08} & \cellcolor{cfTerra1}+0.02 & \cellcolor{cfSlate2}-0.03 \\
%DIFDELCMD < Llama 3.2 Vision 11B & \cellcolor{cfSlate2}-0.03 & \cellcolor{cfSlate2}-0.02 & \cellcolor{cfTerra2}+0.02 & \cellcolor{cfSlate1}-0.02 & \cellcolor{cfSlate2}-0.04 & \cellcolor{cfTerra1}\textbf{+0.01} \\
%DIFDELCMD < LLaVA-Rad-7B & \cellcolor{cfSlate2}-0.06 & \cellcolor{cfTerra1}+0.01 & \cellcolor{cfSlate2}-0.02 & \cellcolor{cfSlate2}-0.06 & \cellcolor{cfSlate2}-0.04 & \cellcolor{cfSlate1}-0.00 \\
%DIFDELCMD < \midrule
%DIFDELCMD < \textbf{Owned} & 21 & 19 & 19 & 20 & 18 & 16 \\
%DIFDELCMD < \bottomrule
%DIFDELCMD < \end{tabular}}
%DIFDELCMD < %%%
\DIFdelendFL \DIFaddbeginFL \resizebox{\linewidth}{!}{%
\begin{tabular}{lrrrrrr}
\toprule
Checkpoint & \multicolumn{1}{c}{person} & \multicolumn{1}{c}{dog} & \multicolumn{1}{c}{car} & \multicolumn{1}{c}{chair} & \multicolumn{1}{c}{bottle} & \multicolumn{1}{c}{bicycle}\\
\midrule
Lingshu 32B & \cellcolor{cfTerra4}\textbf{+0.29} & \cellcolor{cfTerra3}\textbf{+0.23} & \cellcolor{cfTerra2}\textbf{+0.08} & \cellcolor{cfTerra4}\textbf{+0.53} & \cellcolor{cfTerra4}\textbf{+0.34} & \cellcolor{cfTerra3}\textbf{+0.22} \\
CheXagent-2-3B & \cellcolor{cfTerra2}\textbf{+0.07} & \cellcolor{cfTerra2}\textbf{+0.07} & \cellcolor{cfTerra2}+0.03 & \cellcolor{cfSlate1}-0.01 & \cellcolor{cfTerra2}+0.04 & \cellcolor{cfTerra1}+0.00 \\
Qwen2.5-VL-72B & \cellcolor{cfTerra4}\textbf{+0.66} & \cellcolor{cfTerra4}\textbf{+0.44} & \cellcolor{cfTerra4}\textbf{+0.95} & \cellcolor{cfTerra4}\textbf{+0.89} & \cellcolor{cfTerra4}\textbf{+0.62} & \cellcolor{cfTerra4}\textbf{+0.40} \\
Lingshu 7B & \cellcolor{cfTerra4}\textbf{+0.39} & \cellcolor{cfTerra4}\textbf{+0.74} & \cellcolor{cfTerra4}\textbf{+0.70} & \cellcolor{cfTerra4}\textbf{+0.76} & \cellcolor{cfTerra4}\textbf{+0.64} & \cellcolor{cfTerra4}\textbf{+0.85} \\
Qwen3-VL-4B & \cellcolor{cfTerra4}\textbf{+0.46} & \cellcolor{cfTerra3}\textbf{+0.13} & \cellcolor{cfTerra4}\textbf{+0.33} & \cellcolor{cfTerra4}\textbf{+0.71} & \cellcolor{cfTerra4}\textbf{+0.34} & \cellcolor{cfTerra2}\textbf{+0.07} \\
Qwen3-VL-32B & \cellcolor{cfTerra3}\textbf{+0.13} & \cellcolor{cfTerra3}\textbf{+0.12} & \cellcolor{cfTerra3}\textbf{+0.11} & \cellcolor{cfTerra3}\textbf{+0.23} & \cellcolor{cfTerra3}\textbf{+0.24} & \cellcolor{cfTerra2}\textbf{+0.07} \\
InternVL3.5-38B & \cellcolor{cfTerra1}\textbf{+0.01} & \cellcolor{cfTerra1}\textbf{+0.00} & \cellcolor{cfTerra1}\textbf{+0.01} & \cellcolor{cfTerra1}\textbf{+0.02} & \cellcolor{cfTerra1}\textbf{+0.01} & \cellcolor{cfTerra1}\textbf{+0.01} \\
Gemma 3 12B & \cellcolor{cfTerra3}\textbf{+0.19} & \cellcolor{cfTerra4}\textbf{+0.76} & \cellcolor{cfTerra4}\textbf{+0.48} & \cellcolor{cfTerra4}\textbf{+0.45} & \cellcolor{cfTerra3}\textbf{+0.18} & \cellcolor{cfTerra4}\textbf{+0.61} \\
MedGemma 4B & \cellcolor{cfTerra4}\textbf{+0.35} & \cellcolor{cfTerra2}\textbf{+0.03} & \cellcolor{cfTerra4}\textbf{+0.43} & \cellcolor{cfTerra3}\textbf{+0.19} & \cellcolor{cfTerra3}\textbf{+0.20} & \cellcolor{cfTerra3}\textbf{+0.14} \\
InternVL3.5-14B & \cellcolor{cfTerra1}\textbf{+0.02} & \cellcolor{cfTerra1}\textbf{+0.01} & \cellcolor{cfTerra1}\textbf{+0.01} & \cellcolor{cfTerra2}\textbf{+0.04} & \cellcolor{cfTerra2}\textbf{+0.03} & \cellcolor{cfTerra1}+0.00$^{\ddagger}$ \\
Qwen2.5-VL-32B & \cellcolor{cfTerra4}\textbf{+0.44} & \cellcolor{cfTerra3}\textbf{+0.12} & \cellcolor{cfTerra4}\textbf{+0.27} & \cellcolor{cfTerra4}\textbf{+0.83} & \cellcolor{cfTerra4}\textbf{+0.41} & \cellcolor{cfTerra3}\textbf{+0.19} \\
InternVL3.5-8B & \cellcolor{cfTerra2}\textbf{+0.03} & \cellcolor{cfTerra1}\textbf{+0.01} & \cellcolor{cfTerra2}\textbf{+0.04} & \cellcolor{cfTerra2}\textbf{+0.04} & \cellcolor{cfTerra2}\textbf{+0.04} & \cellcolor{cfTerra1}\textbf{+0.02} \\
HuatuoGPT-Vision-7B & \cellcolor{cfTerra1}\textbf{+0.02} & \cellcolor{cfTerra1}+0.00$^{\ddagger}$ & \cellcolor{cfTerra1}\textbf{+0.00} & \cellcolor{cfTerra2}\textbf{+0.03} & \cellcolor{cfTerra2}\textbf{+0.03} & \cellcolor{cfTerra1}+0.00$^{\ddagger}$ \\
LLaVA-Med 1.5 7B & \cellcolor{cfTerra1}+0.02 & \cellcolor{cfTerra1}\textbf{+0.01} & \cellcolor{cfTerra1}+0.01 & \cellcolor{cfSlate2}-0.02 & \cellcolor{cfTerra1}+0.00 & \cellcolor{cfSlate1}-0.01 \\
CheXagent-8B & \cellcolor{cfSlate1}-0.00 & \cellcolor{cfTerra1}+0.01 & \cellcolor{cfSlate1}-0.01 & \cellcolor{cfSlate2}-0.03 & \cellcolor{cfSlate3}-0.13 & \cellcolor{cfSlate2}-0.03 \\
Qwen2.5-VL-3B & \cellcolor{cfTerra4}\textbf{+0.46} & \cellcolor{cfTerra4}\textbf{+0.42} & \cellcolor{cfTerra4}\textbf{+0.53} & \cellcolor{cfTerra4}\textbf{+0.68} & \cellcolor{cfTerra4}\textbf{+0.61} & \cellcolor{cfTerra4}\textbf{+0.39} \\
Qwen2.5-VL-7B & \cellcolor{cfTerra4}\textbf{+0.48} & \cellcolor{cfTerra4}\textbf{+0.75} & \cellcolor{cfTerra4}\textbf{+0.61} & \cellcolor{cfTerra4}\textbf{+0.87} & \cellcolor{cfTerra4}\textbf{+0.51} & \cellcolor{cfTerra4}\textbf{+0.81} \\
Qwen3-VL-8B & \cellcolor{cfTerra4}\textbf{+0.42} & \cellcolor{cfTerra4}\textbf{+0.51} & \cellcolor{cfTerra4}\textbf{+0.74} & \cellcolor{cfTerra4}\textbf{+0.79} & \cellcolor{cfTerra4}\textbf{+0.59} & \cellcolor{cfTerra4}\textbf{+0.43} \\
Gemma 3 27B & \cellcolor{cfTerra3}\textbf{+0.24} & \cellcolor{cfTerra4}\textbf{+0.36} & \cellcolor{cfTerra2}\textbf{+0.09} & \cellcolor{cfTerra4}\textbf{+0.28} & \cellcolor{cfTerra3}\textbf{+0.15} & \cellcolor{cfTerra4}\textbf{+0.75} \\
Gemma 3 4B & \cellcolor{cfTerra3}\textbf{+0.16} & \cellcolor{cfTerra4}\textbf{+0.54} & \cellcolor{cfTerra3}\textbf{+0.24} & \cellcolor{cfTerra4}\textbf{+0.34} & \cellcolor{cfTerra4}\textbf{+0.44} & \cellcolor{cfSlate4}-0.27 \\
LLaVA-1.5-7B & \cellcolor{cfTerra2}\textbf{+0.06} & \cellcolor{cfTerra1}\textbf{+0.00} & \cellcolor{cfTerra1}\textbf{+0.01} & \cellcolor{cfTerra2}\textbf{+0.07} & \cellcolor{cfTerra1}+0.01 & \cellcolor{cfTerra1}\textbf{+0.00} \\
MedGemma 27B & \cellcolor{cfTerra4}\textbf{+0.33} & \cellcolor{cfSlate2}-0.04 & \cellcolor{cfTerra2}\textbf{+0.08} & \cellcolor{cfTerra3}\textbf{+0.18} & \cellcolor{cfTerra3}\textbf{+0.21} & \cellcolor{cfSlate2}-0.04 \\
LLaVA-1.5-13B & \cellcolor{cfTerra2}\textbf{+0.05} & \cellcolor{cfTerra1}+0.00$^{\ddagger}$ & \cellcolor{cfTerra1}+0.02 & \cellcolor{cfTerra2}\textbf{+0.08} & \cellcolor{cfTerra1}+0.02 & \cellcolor{cfSlate2}-0.03 \\
Llama 3.2 Vision 11B & \cellcolor{cfSlate2}-0.03 & \cellcolor{cfSlate2}-0.02 & \cellcolor{cfTerra2}+0.02 & \cellcolor{cfSlate1}-0.02 & \cellcolor{cfSlate2}-0.04 & \cellcolor{cfTerra1}\textbf{+0.01} \\
LLaVA-Rad-7B & \cellcolor{cfSlate2}-0.06 & \cellcolor{cfTerra1}+0.01 & \cellcolor{cfSlate2}-0.02 & \cellcolor{cfSlate2}-0.06 & \cellcolor{cfSlate2}-0.04 & \cellcolor{cfSlate1}-0.00 \\
\midrule
\textbf{Owned} & 21 & 19 & 19 & 20 & 18 & 16 \\
\bottomrule
\end{tabular}}
\DIFaddendFL \end{table}

\FloatBarrier

\DIFdelbegin \subsection{\DIFdel{Controls}}
%DIFAUXCMD
\addtocounter{subsection}{-1}%DIFAUXCMD
%DIFDELCMD < \label{app:cf-controls-sec}
%DIFDELCMD < 

%DIFDELCMD < %%%
\DIFdel{Table~\ref{tab:cf-controls} summarises the reference family and the per-block controls of Section~\ref{sec:controls} across the grid}\DIFdelend \DIFaddbegin \DIFadd{Across the three tables, }\cfChestOwned{} \DIFadd{of the }\cfChestCellsN{} \DIFadd{chest cells are owned and most chest contrasts are negative, against }\cfCocoOwned{} \DIFadd{of the }\cfCocoOwnedN{} \DIFadd{COCO cells for the same checkpoints under the same procedure. Figure~\ref{fig:cf-w-structure} summarises the write matrices behind these contrasts (Appendix~\ref{app:cf-figures}), and Figure~\ref{fig:cf-examples} shows the grades on single images (Appendix~\ref{app:cf-examples})}\DIFaddend .

\DIFdelbegin \begin{table}[htbp]
\DIFdeletedtable{%
{\small\textbf{Table (deleted).} \textbf{Steering reference against verdict.} Primary write-matrix cells of the included blocks, cross-classified by the steering reference and by the simultaneous-bound verdict against the five other clinical directions.\par}\medskip
\centering
\setlength{\tabcolsep}{4pt}

\fitwidth{%
\fitwidth{%
\begin{tabular}{llrrrr}
\toprule
Cells & Steering reference & \multicolumn{1}{c}{advantage} & \multicolumn{1}{c}{stronger competitor} & \multicolumn{1}{c}{unresolved} & \multicolumn{1}{c}{total}\\
\midrule
Chest, all cells & met & 38 & 9 & 9 & 56 \\
 & not met & 24 & 187 & 33 & 244 \\
Chest, readable and answerable & met & 28 & 5 & 5 & 38 \\
 & not met & 13 & 74 & 16 & 103 \\
COCO, all cells & met & 113 & 0 & 4 & 117 \\
 & not met & 10 & 17 & 6 & 33 \\
\bottomrule
\end{tabular}}}
}
\end{table}
%DIFDELCMD < %%%
\DIFdelend \DIFaddbegin \begin{DIFnomarkup}\cfnote{删表}{原 A.6 Controls 小节在删去 Table 8、9 后只剩一句“Table 9 summarises…”，整节删去。本小节删去 ledger 表，保留说明文字（它已按参照 / 规则 / sham 列出全部例外），首句不再指向表格。}\end{DIFnomarkup}
\DIFaddend 

\DIFdelbegin \begin{table}[htbp]
\DIFdeletedtable{%
{\small\textbf{Table (deleted).} \textbf{Reference family and controls.} Median and interquartile range per dataset, over all cells in the first four rows and over blocks with the module in the next six; the last three rows count cells.\par}\medskip
\centering
\setlength{\tabcolsep}{4pt}

\fitwidth{%
\fitwidth{%
\begin{tabular}{lrrrrrr}
\toprule
\multirow{2}{*}{Quantity} & \multicolumn{2}{c}{NIH ChestX-ray14} & \multicolumn{2}{c}{CheXpert Plus} & \multicolumn{2}{c}{COCO (control)}\\
\cmidrule(lr){2-3}\cmidrule(lr){4-5}\cmidrule(lr){6-7}
& \multicolumn{1}{c}{median} & \multicolumn{1}{c}{IQR} & \multicolumn{1}{c}{median} & \multicolumn{1}{c}{IQR} & \multicolumn{1}{c}{median} & \multicolumn{1}{c}{IQR}\\
\midrule
Random-family p95 of $W$ & +0.063 & [+0.021, +0.165] & +0.066 & [+0.027, +0.157] & +0.023 & [+0.008, +0.055] \\
$|$Sham write$|$ & +0.028 & [+0.009, +0.075] & +0.027 & [+0.008, +0.076] & +0.011 & [+0.003, +0.034] \\
Concept write $W_{q,q}$ & +0.010 & [-0.006, +0.067] & +0.020 & [-0.001, +0.082] & +0.127 & [+0.020, +0.482] \\
Ownership $O_q$ & -0.017 & [-0.078, +0.000] & -0.014 & [-0.065, +0.001] & +0.102 & [+0.009, +0.433] \\
Signed dose range of $O_q$ & +0.266 & [+0.061, +0.444] & +0.162 & [+0.034, +0.359] & +0.399 & [+0.050, +0.612] \\
Refit SD of $O_q$ (3 seeds) & +0.039 & [+0.015, +0.072] & +0.037 & [+0.017, +0.052] & +0.043 & [+0.006, +0.074] \\
Connector-locus median $O_q$ & -0.013 & [-0.023, -0.005] & -0.009 & [-0.020, -0.004] & +0.003 & [+0.001, +0.007] \\
Label gap (logits) & -0.1 & [-0.9, -0.0] & -0.3 & [-1.3, +0.0] & -5.3 & [-14.8, -0.3] \\
$|$Wording contrast$|$ (IY$-$WY) & +0.033 & [+0.008, +0.078] & +0.021 & [+0.007, +0.035] & +0.028 & [+0.007, +0.065] \\
$|$Mapping contrast$|$ (IA$-$IB) & +0.037 & [+0.008, +0.073] & +0.019 & [+0.003, +0.055] & +0.030 & [+0.007, +0.051] \\
\midrule
Reference met & \multicolumn{2}{c}{20 of 150} & \multicolumn{2}{c}{36 of 150} & \multicolumn{2}{c}{117 of 150} \\
Owned & \multicolumn{2}{c}{13 of 150} & \multicolumn{2}{c}{25 of 150} & \multicolumn{2}{c}{113 of 150} \\
Reference met but not owned & \multicolumn{2}{c}{7 of 150} & \multicolumn{2}{c}{11 of 150} & \multicolumn{2}{c}{4 of 150} \\
\bottomrule
\end{tabular}}}
}
\end{table}
%DIFDELCMD < 

%DIFDELCMD < %%%
\DIFdelend \subsection{What each analysis scores its writes against}
\label{app:cf-ledger}

\DIFdelbegin \DIFdel{Table~\ref{tab:cf-ledger} lists the }%DIFDELCMD < \cfLedgerAnalyses{} %%%
\DIFdel{graded analyses (}%DIFDELCMD < \cfLedgerCells{} %%%
\DIFdel{cells); the packaged per-block statistics populate each row and check its declared reference and sham. The write matrixscores its own write against }\DIFdelend \DIFaddbegin \begin{DIFnomarkup}\cfnote{附录精简}{说明文字压成先给默认规则、再列三类例外，约 160 词减到 130 词；不再逐项罗列 sham 的计数。}\end{DIFnomarkup}

\DIFadd{Unless stated otherwise, an analysis uses the write matrix's rule: the own write must exceed }\DIFaddend the 95th percentile of the 119-direction random family and the absolute sham, \DIFdelbegin \DIFdel{decides by }\DIFdelend \DIFaddbegin \DIFadd{and the verdict comes from }\DIFaddend simultaneous max-$T$ bounds \DIFdelbegin \DIFdel{, and permutes the direction it writes}\DIFdelend \DIFaddbegin \DIFadd{on the contrasts with the competitors}\DIFaddend . \cfLedgerFullReference{} of the \DIFdelbegin \DIFdel{analyses share this reference , }\DIFdelend \DIFaddbegin \cfLedgerAnalyses{} \DIFadd{graded analyses (}\cfLedgerCells{} \DIFadd{cells) keep this reference and }\DIFaddend \cfLedgerMaxT{} \DIFdelbegin \DIFdel{this rule, and }%DIFDELCMD < \cfLedgerOwnSham{} %%%
\DIFdel{this sham. The semantic endpoints instead }\DIFdelend \DIFaddbegin \DIFadd{keep this rule. The exceptions are the semantic endpoints, which }\DIFaddend take the maximum of \DIFdelbegin \DIFdel{an }\DIFdelend \cfSemLedgerRandom{} \DIFdelbegin \DIFdel{-direction random family }\DIFdelend \DIFaddbegin \DIFadd{random directions }\DIFaddend of their own\DIFdelbegin \DIFdel{, and }\DIFdelend \DIFaddbegin \DIFadd{; }\DIFaddend the answer direction under \DIFdelbegin \DIFdel{the }\DIFdelend held-out templates\DIFaddbegin \DIFadd{, which }\DIFaddend has the random 95th percentile \DIFdelbegin \DIFdel{for }\DIFdelend \DIFaddbegin \DIFadd{in }\DIFaddend \cfAnsdirtRefPairs{} of its \cfAnsdirtLedgerPairs{} (cell, template) pairs and the sham alone \DIFdelbegin \DIFdel{for the rest. The }\DIFdelend \DIFaddbegin \DIFadd{in the rest; and the }\DIFaddend logit-scale grade\DIFaddbegin \DIFadd{, which }\DIFaddend decides by a percentile interval on the \DIFdelbegin \DIFdel{ownership contrast, whose maximum it recomputes in every draw}\DIFdelend \DIFaddbegin \DIFadd{margin-scale ownership contrast}\DIFaddend . The probe refits, alternative constructions, displacement lifts, and expert-label refits (\cfLedgerSeedShamCells{} cells) \DIFdelbegin \DIFdel{permute the }\DIFdelend \DIFaddbegin \DIFadd{take the sham of the }\DIFaddend seed-0 logistic normal\DIFdelbegin \DIFdel{rather than the direction they write, and the token-weighted writes are referenced against the sham of the uniformly written normal. }\DIFdelend \DIFaddbegin \DIFadd{.
}\DIFaddend 

\DIFdelbegin \begin{table}[htbp]
\DIFdeletedtable{%
{\small\textbf{Table (deleted).} \textbf{What each analysis scores its writes against.} Per graded analysis: scored rows, steering reference, decision rule, the direction its sham permutes, and graded (question, arm) cells.\par}\medskip
\centering
\setlength{\tabcolsep}{3pt}

\fitwidth{%
\fitwidth{%
\begin{tabular}{>{\raggedright\arraybackslash}p{0.221\textwidth}>{\raggedright\arraybackslash}p{0.126\textwidth}>{\raggedright\arraybackslash}p{0.307\textwidth}>{\raggedright\arraybackslash}p{0.250\textwidth}>{\raggedright\arraybackslash}p{0.179\textwidth}r}
\toprule
Analysis & Rows & Reference & Rule & Sham & \multicolumn{1}{c}{Cells}\\
\midrule
Write matrix, primary template & 600 test & 119 random $p_{95}$, sham & simultaneous max-$T$, 30 contrasts & the direction written & 450 \\
Probe refits, seeds 1 and 2 & 600 test & 119 random $p_{95}$, sham, the write matrix's & simultaneous max-$T$, 30 contrasts & the seed-0 normal & 900 \\
Logit-scale grade & 600 test & 119 random $p_{95}$, sham, on the margin scale & percentile interval on $O^m_q$ & the direction written & 450 \\
Unsaturated-row grade & unsaturated test & 119 random $p_{95}$, sham, on those rows & simultaneous max-$T$, 30 contrasts & the direction written & 450 \\
Alternative direction constructions & 600 test & 119 random $p_{95}$, sham, the write matrix's & simultaneous max-$T$, 30 contrasts & the seed-0 normal & 1{,}440 \\
Displacement lifts & 600 test & 119 random $p_{95}$, sham, the write matrix's & simultaneous max-$T$, 30 contrasts & the seed-0 normal & 684 \\
Extended competitor family & 600 test & 119 random $p_{95}$, sham & simultaneous max-$T$, 8 or 11 contrasts & the direction written & 54 \\
Token-weighted writes & 600 test & 119 random $p_{95}$, sham, the sham unweighted & simultaneous max-$T$, 30 contrasts & the direction written & 96 \\
fp32 and batch-size-one rescoring & 200 test & 119 random $p_{95}$, sham & simultaneous max-$T$, 30 contrasts & the direction written & 108 \\
Answer direction, primary template & 600 test & 119 random $p_{95}$, sham & simultaneous max-$T$, 30 contrasts & the direction written & 54 \\
Answer direction, held-out templates & 600 test & 119 random $p_{95}$, sham in 80 pairs; sham alone in 190 & simultaneous max-$T$, 30 contrasts & the direction written & 270 \\
Non-clinical attributes & 600 test & 119 random $p_{95}$, sham & simultaneous max-$T$, 72 contrasts & the direction written & 114 \\
Findings in the attribute family & 600 test & 119 random $p_{95}$, sham & simultaneous max-$T$, 72 contrasts & the direction written & 228 \\
Radiologist-label regrade & 200 valid & 119 random $p_{95}$, sham & simultaneous max-$T$, 30 contrasts & the direction written & 114 \\
Expert-label refits & 600 test & 119 random $p_{95}$, sham, the write matrix's & simultaneous max-$T$, 30 contrasts & the seed-0 normal & 36 \\
Projection seeds 1 and 2 & 600 test & 119 random $p_{95}$, sham, that projection's & simultaneous max-$T$, 30 contrasts & the direction written & 36 \\
Crossed tower and reader & 200 test & 119 random $p_{95}$, sham, of the tower in use & simultaneous max-$T$, 30 contrasts & the direction written & 132 \\
Identical-tensor replay & 200 test & 119 random $p_{95}$, sham, the source block's & simultaneous max-$T$, 30 contrasts & the direction written & 48 \\
Semantic endpoints & 600 test & 18 random, maximum; sham, the endpoint's own & simultaneous max-$T$, 30 contrasts & the direction written & 144 \\
Fine-grained object family & 600 test & 119 random $p_{95}$, sham, the family's own & simultaneous max-$T$, 30 contrasts & the direction written & 36 \\
\bottomrule
\end{tabular}}}
}
\end{table}
%DIFDELCMD < %%%
\DIFdelend \DIFaddbegin \begin{DIFnomarkup}\cfnote{附录精简}{原稿 A.8 Robustness、A.9 Further results、A.10 Crossing 按“稳健性 / 更多结果 / 交叉实验”切分，同一话题散在三处（answer direction 4 处、放射科标签 2 处、属性 2 处、语义终点 2 处、替代构造 3 处）。现改为按正文主线的话题分节：A.7 控制与稳健性、A.8 reader、A.9 非临床属性、A.10 answer direction、A.11 难度，每节“一句方法 + 结果”，合并后去掉重复的数字。这一块的文字由约 5,400 词压到约 3,200 词；删去的是出处与流程细节（换塔的张量与收据、ridge 惩罚各取值被选中的次数、基线逐位一致率、NIH 最小正例数等），对论点不利的结果全部保留。定义改用公式：位移 lift、token 加权写入、标签来源与样本量的拆分、reader / tower 效应、answer direction 的 ridge 拟合、白化 cosine、column selectivity、语义终点的判定规则、两个匹配估计量。新符号都在首次出现处定义，没有复用已有字母。正文中指向原 A.9、A.10 的引用改指新小节。}\end{DIFnomarkup}
\DIFaddend 

\DIFdelbegin \subsection{\DIFdel{Robustness of the deficit}}
%DIFAUXCMD
\addtocounter{subsection}{-1}%DIFAUXCMD
\DIFdelend \DIFaddbegin \subsection{\DIFadd{Controls and robustness}}
\DIFaddend \label{app:cf-robustness}
\DIFdelbegin %DIFDELCMD < 

%DIFDELCMD < %%%
\DIFdel{The projected features are $x=R^{\top}\hvec$, where $R\in\R^{D\times512}$ is the fixed Gaussian projection. Two lifts map vectors from the projected, train-scaled space back to the consumed space.
}\DIFdelend \DIFaddbegin 

\begin{DIFnomarkup}\cfnote{附录精简}{原 Further results 里的读取细节和原 Scale 段末尾的样本量支持规则合并为一段；只保留决定资格的计数（CheXpert Atelectasis、COCO bicycle），删去 NIH 最小正例数和测试集计数。}\end{DIFnomarkup}

\paragraph{\DIFadd{Reading and support.}} \DIFadd{Median probe selectivity over the chest cells is $\cfSelChestMed$ (quartiles $\cfSelChestQOne$ and $\cfSelChestQThree$), and }\cfReadAns{} \DIFadd{of the }\cfChestAnswerable{} \DIFadd{answerable chest cells are also readable. Effusion and Cardiomegaly are the most readable of the four concepts both chest sets share, with mean selectivity $\cfSelEffusion$ and $\cfSelCardiomegaly$ over }\cfChestProbeBlocks{} \DIFadd{chest blocks. Both grades need ten known positives and ten known negatives among the 400 calibration rows. CheXpert Atelectasis lacks them in each of the }\cfSupportInelChexBlocks{} \DIFadd{CheXpert blocks (}\cfCalChexInelPos{} \DIFadd{positives, }\cfCalChexInelNeg{} \DIFadd{negatives), which makes }\cfSupportInelChest{} \DIFadd{of the }\cfChestCellsN{} \DIFadd{chest cells ineligible, and COCO bicycle lacks them with }\cfCalCocoInelPos{} \DIFadd{positives. Among eligible chest cells, }\cfSupportReadablePctChest\DIFadd{\% are readable and }\cfSupportAnswerablePctChest\DIFadd{\% answerable.
}

\paragraph{\DIFadd{Dose and locus.}} \DIFadd{The dose sweep ran on }\cfDoseBlocksNih{} \DIFadd{NIH and }\cfDoseBlocksCoco{} \DIFadd{COCO blocks, and on COCO the ownership margins span $\cfCocoMarginMin$--$\cfCocoMarginMax$ in the Qwen, Lingshu, and Gemma families. Writes at the connector locus are an order of magnitude smaller than at the consumed block: median $|W_{q,q}|$ is $\cfConnWqqNih$ against $\cfPrimWqqNih$ on NIH and $\cfConnWqqCoco$ against $\cfPrimWqqCoco$ on COCO.
}

\paragraph{\DIFadd{Alternative directions.}} \DIFadd{Besides the logistic normal, we write four constructions per concept (Table~\ref{tab:cf-altdir}): the difference of class means, the Haufe pattern, the normal orthogonalised against the other five normals, and the normal refitted on features residualised against the other five labels. }\DIFaddend A classifier coefficient \DIFdelbegin \DIFdel{$\boldsymbol\beta$ lifts as $\hat\wvec\propto R\,\mathrm{diag}(1/\boldsymbol s)\boldsymbol\beta$. The resulting written direction scores the concept as the probe does (}\DIFdelend \DIFaddbegin \DIFadd{lifts by }\DIFaddend Equation~\ref{eq:lift}\DIFdelbegin \DIFdel{). A }\DIFdelend \DIFaddbegin \DIFadd{. A difference of means or a pattern is instead a }\DIFaddend displacement $\boldsymbol u$ \DIFdelbegin \DIFdel{lifts via }\DIFdelend \DIFaddbegin \DIFadd{in the projected, train-scaled space, and }\DIFaddend its minimum-norm preimage \DIFdelbegin \DIFdel{as $\hat\wvec\propto R(R^{\top}R)^{-1}\mathrm{diag}(\boldsymbol s)\boldsymbol u$. This lift reproduces the displacement exactly: $R^{\top}\hat\wvec\propto\mathrm{diag}(\boldsymbol s)\boldsymbol u$. The }\DIFdelend \DIFaddbegin \DIFadd{gives the displacement lift
}\begin{equation}
\hat\wvec\propto R(R^{\top}R)^{-1}\mathrm{diag}(\boldsymbol s)\,\boldsymbol u,\qquad\text{so that}\qquad R^{\top}\hat\wvec\propto\mathrm{diag}(\boldsymbol s)\,\boldsymbol u .
\label{eq:displacement-lift}
\end{equation}
\DIFadd{The two }\DIFaddend lifts differ because $R^{\top}R$ is not a multiple of the identity \DIFdelbegin \DIFdel{. Measurements on every block with the displacement lifts show that $R^{\top}R/D$ is full rank, with eigenvalues $\cfAltdirdEigMin$--$\cfAltdirdEigMax$ and condition number }%DIFDELCMD < \cfAltdirdCondMin{}%%%
\DIFdel{--}%DIFDELCMD < \cfAltdirdCondMax{}%%%
\DIFdel{. We compute the }\DIFdelend \DIFaddbegin \DIFadd{(Appendix~\ref{app:repro-probe}). The }\DIFaddend difference-of-means and pattern \DIFdelbegin \DIFdel{vectors in the same projected, train-scaled space as the logistic normal and report both lifts. The alternative-construction module scores the coefficient-lift family. The displacement lift gives the geometrically matched comparison for these two families: a difference of class means and a Haufe pattern are displacements in the projected space rather than classifier coefficients. We form the orthogonalised normal by removing the span of the other five normals }\DIFdelend \DIFaddbegin \DIFadd{directions lie at a median absolute raw model-space cosine of $\cfAltdirCosMin$ to $\cfAltdirCosMax$ }\DIFaddend from the logistic normal.
\DIFdelbegin \DIFdel{We refit the residualised normal on features residualised against the other five labels.
}\DIFdelend 

\paragraph{\DIFdelbegin \DIFdel{Alternative directions, }\DIFdelend \DIFaddbegin \DIFadd{Token }\DIFaddend weights \DIFdelbegin \DIFdel{, }\DIFdelend and \DIFdelbegin \DIFdel{competitor families}\DIFdelend \DIFaddbegin \DIFadd{competitors}\DIFaddend .} \DIFdelbegin \DIFdel{Token-weighted variants multiply the per-token dose $\alpha\lVert\hvec_{it}\rVert$ by a weight $w_t$ whose mean over the consumed tokens is one}\DIFdelend \DIFaddbegin \DIFadd{A token-weighted write replaces the uniform dose of Equation~\ref{eq:intervention} with weights $\omega_t$ that average one over the $n_{\mathrm{tok}}$ consumed tokens,
}\begin{equation}
\hvec'_{it}=\hvec_{it}+\alpha\,\omega_t\lVert\hvec_{it}\rVert_2\,\hat\wvec_c,\qquad
\omega_t=\begin{cases}n_{\mathrm{tok}}\,e^{z_t}\big/\sum_{t'}e^{z_{t'}} & \text{softmax,}\\ 4\cdot\mathbf 1[\,z_t\text{ in the top quarter}\,] & \text{top quarter,}\end{cases}
\label{eq:token-weights}
\end{equation}
\DIFadd{where $z_t$ is the probe's score of token $t$}\DIFaddend . The softmax variant \DIFdelbegin \DIFdel{weights tokens by the softmax of their probe scores. It preserves the }\DIFdelend \DIFaddbegin \DIFadd{keeps the }\DIFaddend uniform write's L1 dose \DIFdelbegin \DIFdel{. The }\DIFdelend \DIFaddbegin \DIFadd{and the }\DIFaddend top-quarter variant \DIFdelbegin \DIFdel{($w_t=4$ on the 25\% highest-scoring tokens) has twice the uniform write's }\DIFdelend \DIFaddbegin \DIFadd{doubles its }\DIFaddend Frobenius norm. The extended competitor family\DIFaddbegin \DIFadd{, scored across }\cfExtcompBlocks{} \DIFadd{blocks, }\DIFaddend adds a direction for every other \DIFdelbegin \DIFdel{dataset }\DIFdelend label with at least 100 known positives and 100 known negatives (NIH: Consolidation, Edema, Infiltration; CheXpert: Enlarged Cardiomediastinum, Fracture, Lung Lesion, Lung Opacity, Pneumonia, Support Devices)\DIFdelbegin \DIFdel{. We fit and write each added direction }\DIFdelend \DIFaddbegin \DIFadd{, fitted and written }\DIFaddend like the protocol normals\DIFdelbegin \DIFdel{. Under this family, }\DIFdelend \DIFaddbegin \DIFadd{; }\DIFaddend a cell is \DIFdelbegin \DIFdel{owned when its write meets the steering reference and holds the simultaneous advantage over the five protocol competitors and every extra direction.
Table~\ref{tab:cf-altdir} reports every construction, lift, token weighting, and competitor family per dataset. Each panel includes the protocol's logistic normal on the same blocks.
}\DIFdelend \DIFaddbegin \DIFadd{then owned only if its write also beats every added direction.
}\DIFaddend 

\DIFaddbegin \begin{table}[h]
\centering
\setlength{\tabcolsep}{4pt}
\begin{DIFnomarkup}\cfnote{表格分布}{原稿的稳健性与补充结果表集中排在 A.8 末尾，与讨论它们的文字相隔很远。现在按话题小节放置：A.7 放 Alternative directions、Geometry、Refit（单一结构、支撑正文 5.2 与 5.7），A.8 放 Shared-tower pairs（支撑 5.4），A.9 放属性表中属性 probe 那一块，A.10 放 answer direction 表与终点表的计数部分（各去掉与 Table 12 重复或拼接的部分），A.11 把原 fine-grained 表中的匹配估计量做成一张 4 行小表。表里有的数字从正文删去（例如 geometry 段的随机率、Spearman、NIH 余弦，替代构造的逐项计数），正文改为引用表格；限定论点的匹配结果仍在文字中保留。}\end{DIFnomarkup}

\caption{\textbf{\DIFaddFL{Alternative directions and competitor families.}} \DIFaddFL{Cells owned under the paper's rule within each family, and the median $O_q$, per dataset. Each panel repeats the reference write on its own blocks.}}
\label{tab:cf-altdir}
\resizebox{\linewidth}{!}{%
\begin{tabular}{llrrrrrr}
\toprule
\multirow{2}{*}{Direction} & \multirow{2}{*}{Lift} & \multicolumn{2}{c}{NIH ChestX-ray14} & \multicolumn{2}{c}{CheXpert Plus} & \multicolumn{2}{c}{COCO}\\
\cmidrule(lr){3-4}\cmidrule(lr){5-6}\cmidrule(lr){7-8}
&  & \multicolumn{1}{c}{owned} & \multicolumn{1}{c}{med.\ $O_q$} & \multicolumn{1}{c}{owned} & \multicolumn{1}{c}{med.\ $O_q$} & \multicolumn{1}{c}{owned} & \multicolumn{1}{c}{med.\ $O_q$}\\
\midrule
\multicolumn{8}{l}{\textit{Direction construction (20 blocks per dataset)}} \\
\quad logistic normal & coefficient & 7/120 & -0.03 & 19/120 & -0.01 & 101/120 & +0.18 \\
\quad difference of means & coefficient & 16/120 & -0.02 & 18/120 & -0.01 & 62/120 & +0.02 \\
\quad Haufe pattern & coefficient & 17/120 & -0.02 & 17/120 & -0.01 & 53/120 & +0.01 \\
\quad orthogonalised normal & coefficient & 8/120 & -0.02 & 14/120 & -0.01 & 98/120 & +0.13 \\
\quad residualised normal & coefficient & 8/120 & -0.02 & 13/120 & -0.01 & 102/120 & +0.20 \\
\quad owned under at least one &  & 25/120 & -- & 41/120 & -- & 108/120 & -- \\
\midrule
\multicolumn{8}{l}{\textit{Displacement lift (19 blocks per dataset)}} \\
\quad logistic normal & coefficient & 7/114 & -0.03 & 19/114 & -0.01 & 100/114 & +0.19 \\
\quad difference of means & displacement & 21/114 & -0.03 & 19/114 & -0.00 & 58/114 & +0.02 \\
\quad Haufe pattern & displacement & 17/114 & -0.02 & 17/114 & -0.00 & 49/114 & +0.01 \\
\midrule
\multicolumn{8}{l}{\textit{Token weights of the logistic normal (4 NIH, 2 CheXpert, 2 COCO blocks)}} \\
\quad uniform & coefficient & 3/24 & -0.08 & 5/12 & +0.00 & 12/12 & +0.60 \\
\quad softmax of the probe score & coefficient & 0/24 & -0.01 & 0/12 & -0.00 & 12/12 & +0.16 \\
\quad top quarter of tokens & coefficient & 2/24 & -0.12 & 3/12 & -0.05 & 12/12 & +0.57 \\
\midrule
\multicolumn{8}{l}{\textit{Competitor family of the logistic normal (5 NIH, 4 CheXpert blocks)}} \\
\quad six clinical directions & coefficient & 3/30 & -0.07 & 8/24 & -0.07 & -- & -- \\
\quad with every extra dataset label & coefficient & 3/30 & -0.07 & 5/24 & -0.10 & -- & -- \\
\bottomrule
\end{tabular}}
\end{table}

\DIFaddend \paragraph{\DIFdelbegin \DIFdel{Non-clinical attributes}\DIFdelend \DIFaddbegin \DIFadd{Geometry of the competitors}\DIFaddend .} The \DIFdelbegin \DIFdel{attribute directions use three labels from each dataset's metadata: the view field (anteroposterior or portable against posteroanterior), recorded sex, and age at least 60. We fit these directions using the same images, features, projection, and probe settings as the clinical probes. We then lift and write them with the six clinical directions as one nine-direction family. Attribute and clinical questions share the same steering reference in all }%DIFDELCMD < \cfAttrRefBlocks{} %%%
\DIFdel{blocks of Table~\ref{tab:cf-attr}: the 119-direction random family written at }\DIFdelend \DIFaddbegin \DIFadd{strongest competitor coincides with the most similar direction or }\DIFaddend the \DIFdelbegin \DIFdel{same dose and the coordinate-permutation sham. We score each attribute question in }%DIFDELCMD < \cfAttrqPhrasings{} %%%
\DIFdel{phrasings on the calibration rows. The phrasing used for the write reaches a median clean-answer AUROC of }%DIFDELCMD < \cfAttrqWrittenAuroc{}%%%
\DIFdel{, compared with }%DIFDELCMD < \cfAttrqSelectedAuroc{} %%%
\DIFdel{for the phrasing in the same cell with the highest one-sided lower bound. Across the }%DIFDELCMD < \cfAttrqCells{} %%%
\DIFdel{attribute cells, the median gap is }%DIFDELCMD < \cfAttrqGain{} %%%
\DIFdel{and the largest gap is }%DIFDELCMD < \cfAttrqGainMax{}%%%
\DIFdel{. }%DIFDELCMD < \cfAttrqCapable{} %%%
\DIFdel{of those cells clear the answerable rule as written. We assess the comparison three ways: over every cell, over answerable cells alone, and over attribute cells paired with clinical cells in the same block within }%DIFDELCMD < \cfAttrPairWindow{} %%%
\DIFdel{clean-answer AUROC. This matching places }%DIFDELCMD < \cfAttrPairMatched{} %%%
\DIFdel{of the }%DIFDELCMD < \cfAttrPairCellsN{} %%%
\DIFdel{attribute cells into }%DIFDELCMD < \cfAttrPairAttrN{} %%%
\DIFdel{pairs and leaves }%DIFDELCMD < \cfAttrPairUnmatched{} %%%
\DIFdel{without a partner in their own block. Attribute probes reach a median AUROC of }%DIFDELCMD < \cfAttrAurocMin{}%%%
\DIFdel{--}%DIFDELCMD < \cfAttrAurocMax{} %%%
\DIFdel{against the attribute label and }%DIFDELCMD < \cfAttrAnswerAurocMin{}%%%
\DIFdel{--}%DIFDELCMD < \cfAttrAnswerAurocMax{} %%%
\DIFdel{against the model's own clean answer. We summarise each attribute--finding pair by its median over blocks; over the }%DIFDELCMD < \cfAttrCosPairs{} %%%
\DIFdel{pairs the largest absolute raw model-space cosine between an attribute direction and a clinical normal is }%DIFDELCMD < \cfAttrCosMax{}%%%
\DIFdel{, for }%DIFDELCMD < \cfAttrCosPairLabel{}%%%
\DIFdel{. Across the }%DIFDELCMD < \cfAttrBlocks{} %%%
\DIFdel{blocks with the module, ownership counts are }%DIFDELCMD < \cfAttrChestAttrOwned{} %%%
\DIFdel{of }%DIFDELCMD < \cfAttrChestAttrN{} %%%
\DIFdel{cells for attributes and }%DIFDELCMD < \cfAttrChestClinOwned{} %%%
\DIFdel{of }%DIFDELCMD < \cfAttrChestClinN{} %%%
\DIFdel{for findings in the same family. The corresponding attribute and finding counts are }%DIFDELCMD < \cfAttrChexAttrOwned{}%%%
\DIFdel{/}%DIFDELCMD < \cfAttrChexAttrN{} %%%
\DIFdelend \DIFaddbegin \DIFadd{most co-occurring label at about the chance rate of one cell in five, }\DIFaddend and \DIFdelbegin %DIFDELCMD < \cfAttrChexClinOwned{}%%%
\DIFdel{/}%DIFDELCMD < \cfAttrChexClinN{} %%%
\DIFdel{on CheXpert, and }%DIFDELCMD < \cfAttrNihAttrOwned{}%%%
\DIFdel{/}%DIFDELCMD < \cfAttrNihAttrN{} %%%
\DIFdel{and }%DIFDELCMD < \cfAttrNihClinOwned{}%%%
\DIFdel{/}%DIFDELCMD < \cfAttrNihClinN{} %%%
\DIFdel{on NIH.
}%DIFDELCMD < 

%DIFDELCMD < %%%
\paragraph{\DIFdel{Answer direction: rescue, loss, and templates.}} %DIFAUXCMD
\addtocounter{paragraph}{-1}%DIFAUXCMD
\DIFdel{We lift the answer direction $a_q$ of Section~\ref{sec:ansdir} as we lift the probe normal. We write it at the same dose, against the write matrix's own 119-direction random family and a coordinate-permutation sham of the $a_q$ under test, using the other five $a_d$ as clinical competitors. A concept is }\emph{\DIFdel{kept}} %DIFAUXCMD
\DIFdel{when both $a_q$ and the label direction $\hat w_q$ are owned, }\emph{\DIFdel{rescued}} %DIFAUXCMD
\DIFdel{when only $a_q$ is owned, and }\emph{\DIFdel{lost}} %DIFAUXCMD
\DIFdel{when only $\hat w_q$ is owned (Table~\ref{tab:cf-ansdir}, top). The release reports both ownership contrasts, column selectivity, and the whitened cosine for every cell. We also write the $a_q$ fitted on the primary template IY unchanged under the rewording WY and the answer-mapped templates IA, IB, WA, and WB. Each template uses its own clean baseline, the $a_d$ of the other five concepts, and a coordinate-permutation sham of the $a_q$ under test. The held-out templates score no random family of their own. The random 95th percentile is available where the additional-template module scored one under that template, which is }%DIFDELCMD < \cfAnsdirtRefPairs{} %%%
\DIFdel{of the }%DIFDELCMD < \cfAnsdirtLedgerPairs{} %%%
\DIFdel{(cell, template) pairs; the remaining }%DIFDELCMD < \cfAnsdirtShamPairs{} %%%
\DIFdel{pairs meet the steering reference against their own sham alone. Every pair takes the same $6\times5$ max-$T$ verdict (Table~\ref{tab:cf-ledger}). A block enters when every held-out template is scored on all 600 test rows (Table~\ref{tab:cf-ansdir}, bottom).
}%DIFDELCMD < 

%DIFDELCMD < %%%
\paragraph{\DIFdel{Direction and label geometry.}} %DIFAUXCMD
\addtocounter{paragraph}{-1}%DIFAUXCMD
\DIFdel{Neither correlations among the six directions nor co-occurrence of the six report labels explains why a competing direction wins }\DIFdelend \DIFaddbegin \DIFadd{the competitor's advantage $W_{q,d}-W_{q,q}$ is uncorrelated with the raw cosine between the two directions and with the correlation $\phi$ between the two binary labels }\DIFaddend (Table~\ref{tab:cf-geometry}).
\DIFdelbegin \DIFdel{The direction cosine is the raw model-space cosine between two of a block's six write directions, averaged over all blocks' off-diagonal pairs. The label $\phi$ measures the correlation between two of the six test labels. On NIH , the }\DIFdelend \DIFaddbegin 

\DIFadd{NIH shows the deficit although its }\DIFaddend six directions are \DIFdelbegin \DIFdel{close to orthogonal : their mean absolute off-diagonal cosine is $\cfGeoAbsCosNih$, against $\cfGeoCosRandom$ between pairs of random unit directions, and their signed mean is $\cfGeoCosNih$. The labels are nearly uncorrelated(mean $\phi=\cfGeoPhiNih$). Yet NIH shows the same off-diagonal dominance as CheXpert, where }\DIFdelend \DIFaddbegin \DIFadd{nearly orthogonal and its labels nearly uncorrelated, while on CheXpert }\DIFaddend directions and labels are correlated (\DIFdelbegin \DIFdel{mean absolute off-diagonal cosine $\cfGeoAbsCosChex$, signed mean $\cfGeoCosChex$; }\DIFdelend Effusion--Consolidation \DIFaddbegin \DIFadd{cosine }\DIFaddend $\cfGeoCosEffCons$ averaged over blocks, $\phi=\cfGeoPhiEffCons$) \DIFdelbegin \DIFdel{. The ownership contrast is positive, $O_q>0$, in }%DIFDELCMD < \cfGeoOwnWinsNihPct%%%
\DIFdel{\% of the }%DIFDELCMD < \cfNihCells{} %%%
\DIFdel{NIH cells and }%DIFDELCMD < \cfGeoOwnWinsChexPct%%%
\DIFdel{\% of the }%DIFDELCMD < \cfChexCells{} %%%
\DIFdel{CheXpert cells. Each cell has five competing directions, so its strongest competitor matches any fixed competitor in one cell in five by chance. The strongest competitor matches the most similar direction at this rate (}%DIFDELCMD < \cfGeoChanceCosMin{}%%%
\DIFdel{--}%DIFDELCMD < \cfGeoChanceCosMax%%%
\DIFdel{\% of chest cells per dataset). It also corresponds to the label with the largest $\phi$ or the most co-occurring label at this rate (}%DIFDELCMD < \cfGeoChanceLabelMin{}%%%
\DIFdel{--}%DIFDELCMD < \cfGeoChanceLabelMax%%%
\DIFdel{\%). Across the }%DIFDELCMD < \cfGeoOffdiagCells{} %%%
\DIFdel{off-diagonal chest cells, the competitor's advantage $W_{q,d}-W_{q,q}$ is uncorrelated with the direction cosine (Spearman over cells $\cfGeoRhoAdvCos$, $p=\cfGeoRhoAdvCosP$) and with the label $\phi$ ($\cfGeoRhoAdvPhi$)}\DIFdelend \DIFaddbegin \DIFadd{and the dominance is the same}\DIFaddend . At least two competitors exceed the own write in \DIFdelbegin %DIFDELCMD < \cfGeoTwoAbove{} %%%
\DIFdel{of }%DIFDELCMD < \cfGeoTwoAboveN{} %%%
\DIFdel{chest cells(}\DIFdelend \cfGeoTwoAbovePct\DIFdelbegin \DIFdel{\%). Removing }\DIFdelend \DIFaddbegin \DIFadd{\% of chest cells, so removing }\DIFaddend the strongest competitor \DIFdelbegin \DIFdel{therefore }\DIFdelend leaves $O_q$ negative in those cells. \DIFdelbegin \DIFdel{The last three rows of Table~\ref{tab:cf-geometry} count the sign of $O_q$ after removal, not the owned grade. }\DIFdelend On COCO, \DIFdelbegin \DIFdel{we restrict each cell's competitors to a }\DIFdelend \DIFaddbegin \DIFadd{competitors restricted to }\DIFaddend co-occurring \DIFdelbegin \DIFdel{family of objects (person with bicycle, $\phi=\cfGeoFamPhiPersonBicycle$; person with }\DIFdelend \DIFaddbegin \DIFadd{objects (for example person and }\DIFaddend dog, $\phi=\cfGeoFamPhiPersonDog$\DIFdelbegin \DIFdel{; chair with bottle, $\phi=\cfGeoFamPhiChairBottle$; car with bicycle, $\phi=\cfGeoFamPhiCarBicycle$; person with chair and bottle, mean $\phi=\cfGeoFamPhiPersonChairBottle$) . This leaves }\DIFdelend \DIFaddbegin \DIFadd{) still leave }\DIFaddend $O_q>0$ in \cfGeoCocoFamMinPct{}--\cfGeoCocoFamMaxPct\% of each family's cells, \DIFdelbegin \DIFdel{compared with }\DIFdelend \DIFaddbegin \DIFadd{against }\DIFaddend \cfGeoCocoFamFullMinPct{}--\cfGeoCocoFamFullMaxPct\% \DIFdelbegin \DIFdel{of the same cells }\DIFdelend with all five\DIFdelbegin \DIFdel{competitors. Of the five directions in the }\DIFdelend \DIFaddbegin \DIFadd{. In the }\DIFaddend Effusion case, Nodule is the \DIFdelbegin \emph{\DIFdel{least}} %DIFAUXCMD
\DIFdel{similar }\DIFdelend \DIFaddbegin \DIFadd{least similar direction }\DIFaddend to Effusion (raw \DIFdelbegin \DIFdel{model-space }\DIFdelend cosine $\cfSeedCosNodule$\DIFdelbegin \DIFdel{in this block). It still }\DIFdelend \DIFaddbegin \DIFadd{), yet it }\DIFaddend moves the Effusion answer \DIFdelbegin \DIFdel{more than }\DIFdelend \DIFaddbegin \DIFadd{by $\cfSeedWNodule$ against $\cfSeedWEff$ for }\DIFaddend Effusion's own direction\DIFdelbegin \DIFdel{($\cfSeedWNodule$ against $\cfSeedWEff$).
The most similar direction, Atelectasis (raw model-space cosine $\cfSeedCosAtel$), moves it by $\cfSeedWAtel$}\DIFdelend .

\DIFaddbegin \begin{table}[h]
\centering
\setlength{\tabcolsep}{5pt}
\caption{\textbf{\DIFaddFL{Direction and label geometry against off-diagonal dominance.}} \DIFaddFL{One column per dataset, with the two chest sets pooled in the third. A ``--'' marks a row that is defined per dataset only.}}
\label{tab:cf-geometry}
\resizebox{\linewidth}{!}{%
\begin{tabular}{lrrrr}
\toprule
& \multicolumn{1}{c}{NIH ChestX-ray14} & \multicolumn{1}{c}{CheXpert Plus} & \multicolumn{1}{c}{both chest sets} & \multicolumn{1}{c}{COCO}\\
\midrule
\multicolumn{5}{l}{\textit{Geometry of the six directions and the six labels}} \\
\quad mean off-diagonal cosine between the directions & 0.019 & 0.215 & -- & 0.043 \\
\quad mean off-diagonal $|\cos|$ between the directions & 0.093 & 0.218 & -- & 0.058 \\
\quad mean $|\cos|$ between random unit directions & 0.024 & 0.024 & -- & 0.024 \\
\quad mean off-diagonal $\phi$ between the labels & 0.026 & 0.381 & -- & 0.022 \\
\midrule
\multicolumn{5}{l}{\textit{Cells whose strongest competitor is (chance: one in five)}} \\
\quad the most similar direction & 31/150 & 33/150 & 64/300 & 31/150 \\
\quad the label with the largest $\phi$ & 35/150 & 36/150 & 71/300 & 29/150 \\
\quad the most co-occurring label & 24/150 & 32/150 & 56/300 & 22/150 \\
\midrule
\multicolumn{5}{l}{\textit{Spearman $\rho$ over off-diagonal cells of $W_{q,d}-W_{q,q}$ with}} \\
\quad the cosine between the two directions & 0.085 & 0.048 & 0.037 & 0.042 \\
\quad the $\phi$ between the two labels & 0.036 & 0.037 & -0.001 & 0.007 \\
\midrule
\multicolumn{5}{l}{\textit{Sign of $O_q$}} \\
\quad cells with $O_q>0$ & 38/150 & 43/150 & 81/300 & 129/150 \\
\quad cells with $O_q>0$ without the strongest competitor & 62/150 & 75/150 & 137/300 & 136/150 \\
\quad cells with at least two competitors above the own write & 88/150 & 75/150 & 163/300 & 14/150 \\
\bottomrule
\end{tabular}}
\end{table}

\DIFaddend \paragraph{Scale \DIFdelbegin \DIFdel{, saturation, refits, }\DIFdelend and \DIFdelbegin \DIFdel{numerics}\DIFdelend \DIFaddbegin \DIFadd{saturation}\DIFaddend .} \DIFdelbegin \DIFdel{The grades hold on the logit scale. Recomputing the grade from }\DIFdelend \DIFaddbegin \DIFadd{Let $O^m_q$ denote ownership recomputed from changes in }\DIFaddend the margin $\ell_{\mathrm{yes}}-\ell_{\mathrm{no}}$ \DIFdelbegin \DIFdel{reproduces the probability-scale }\DIFdelend \DIFaddbegin \DIFadd{instead of $P(\mathrm{yes})$. The margin-scale grade reproduces the }\DIFaddend owned grade in \cfScaleAgree{} of \cfScaleAgreeN{} cells\DIFdelbegin \DIFdel{. Every owned cell retains a positive ownership contrast on that scale, }\DIFdelend \DIFaddbegin \DIFadd{, every owned cell keeps }\DIFaddend $O^m_q>0$\DIFdelbegin \DIFdel{. The full owned grade persists in }\DIFdelend \DIFaddbegin \DIFadd{, and the full margin-scale grade keeps }\DIFaddend \cfScaleOwnedKeepFull{} of the \cfScaleOwned{} owned cells\DIFdelbegin \DIFdel{, with the steering reference recomputed on the margin scale and a positive percentile interval for $O^m_q$ (Table~\ref{tab:cf-scale}). Saturation counters a ceiling explanation: COCO blocks contain far more rows at $P(\mathrm{yes})>0.99$ or $<0.01$ than chest }\DIFdelend \DIFaddbegin \DIFadd{. Rows with $P(\mathrm{yes})\notin[0.01,0.99]$ are far more common in COCO }\DIFaddend blocks (median block share $\cfScaleSatCoco$\DIFdelbegin \DIFdel{against }\DIFdelend \DIFaddbegin \DIFadd{) than in NIH (}\DIFaddend $\cfScaleSatNih$\DIFdelbegin \DIFdel{and }\DIFdelend \DIFaddbegin \DIFadd{) and CheXpert (}\DIFaddend $\cfScaleSatChex$) \DIFdelbegin \DIFdel{. Restricting each cell to its unsaturated rows preserves $O_q>0$ in }\DIFdelend \DIFaddbegin \DIFadd{blocks. On unsaturated rows alone, }\DIFaddend all \cfScaleUnsatOwned{} owned chest cells \DIFdelbegin \DIFdel{and in }\DIFdelend \DIFaddbegin \DIFadd{keep $O_q>0$ and }\DIFaddend \cfScaleUnsatOtherPos{} of the \cfScaleUnsatOther{} others \DIFdelbegin \DIFdel{. Full regrading uses a steering reference recomputed on these rows and the simultaneous max-$T$ advantage over the same six-direction family. It preserves the owned grade in }\DIFdelend \DIFaddbegin \DIFadd{have $O_q>0$; the full regrade keeps }\DIFaddend \cfScaleUnsatGradeKept{} of \cfScaleUnsatGradeN{} owned chest cells and raises \cfScaleUnsatGradeGained{} of the \cfScaleUnsatGradeOtherN{} others to owned. The \DIFdelbegin \DIFdel{same pattern holds in the }\DIFdelend clean-No, clean-Yes, label-present, and label-absent strata \DIFdelbegin \DIFdel{. }\DIFdelend \DIFaddbegin \DIFadd{give the same pattern.
}

\paragraph{\DIFadd{Refits and numerics.}} \DIFaddend Refitting the probes on resampled training patients \DIFdelbegin \DIFdel{, }\DIFdelend and grading each \DIFdelbegin \DIFdel{refitted direction }\DIFdelend \DIFaddbegin \DIFadd{refit }\DIFaddend against the seed-0 \DIFdelbegin \DIFdel{random 95th percentile and the seed-0 sham of the same block and question, }\DIFdelend \DIFaddbegin \DIFadd{references }\DIFaddend changes $O_q$ by a median of $\cfScaleRefitMedian$ (Table~\DIFdelbegin \DIFdel{\ref{tab:cf-ledger}); the last column of Table~\ref{tab:cf-refit}regrades seed 0 with the same draws and counts agreement with the paper's grade}\DIFdelend \DIFaddbegin \DIFadd{\ref{tab:cf-refit})}\DIFaddend . Rescoring the \DIFdelbegin \DIFdel{write grid on the }\DIFdelend first 200 test rows of each block \DIFdelbegin \DIFdel{, }\DIFdelend in fp32 and at batch size one \DIFdelbegin \DIFdel{, moves the $6\times126$ written cells of a block }\DIFdelend \DIFaddbegin \DIFadd{moves a written cell }\DIFaddend by at most \cfPrecisionMaxDW{} \DIFdelbegin \DIFdel{. Table~\ref{tab:cf-scale} gives that largest change and the largest in a clinical contrast $C_{q,d}$, the difference between the own write and a competing write on the same question. Regraded under either setting, verdicts change in }%DIFDELCMD < \cfPrecisionVerdictChanges{} %%%
\DIFdelend and \DIFdelbegin \DIFdel{steering references in }%DIFDELCMD < \cfPrecisionRefChanges{} %%%
\DIFdel{of the }\DIFdelend \DIFaddbegin \DIFadd{changes }\cfPrecisionVerdictChanges{} \DIFadd{verdicts and }\cfPrecisionRefChanges{} \DIFadd{steering references among }\DIFaddend \cfPrecisionGradeCells{} grades\DIFdelbegin \DIFdel{, and the two point verdicts change in }%DIFDELCMD < \cfPrecisionPointVerdictChanges{} %%%
\DIFdel{(Table~\ref{tab:cf-scale})}\DIFdelend . The bf16 batch effects declared at the gates (\DIFdelbegin \DIFdel{candidate-logit differences of }\DIFdelend $\cfScaleBatchMin$--$\cfScaleBatchMax$ \DIFdelbegin \DIFdel{) affect no owned cell. The margin sign agrees }\DIFdelend \DIFaddbegin \DIFadd{in candidate logits) leave the margin sign unchanged }\DIFaddend in \cfScaleBatchAgree{} of \cfScaleBatchN{} blocks\DIFdelbegin \DIFdel{; }\DIFdelend \DIFaddbegin \DIFadd{, and }\DIFaddend the \cfScaleBatchDisagree{} \DIFdelbegin \DIFdel{disagreeing blocks }\DIFdelend \DIFaddbegin \DIFadd{others }\DIFaddend contain no owned cell\DIFdelbegin \DIFdel{s}\DIFdelend  . Re-scoring the clean baseline in a second module \DIFdelbegin \DIFdel{gives bitwise agreement in 100\% of rows in }\DIFdelend \DIFaddbegin \DIFadd{agrees bit for bit on every row in }\DIFaddend \cfScaleBaselineIdentical{} of \cfScaleBaselineN{} blocks\DIFdelbegin \DIFdel{and in at least }%DIFDELCMD < \cfScaleBaselineRestPct%%%
\DIFdel{\% of rows in the rest.
Differences satisfying $|\Delta P(\mathrm{yes})|>0.1$ occur in at most }%DIFDELCMD < \cfScaleBaselineDpPct%%%
\DIFdel{\% of rows. Writes at the connector locus are an order of magnitude smaller than those at the consumed block (median $|W_{q,q}|$ $\cfConnWqqNih$ against $\cfPrimWqqNih$ on NIH and $\cfConnWqqCoco$ against $\cfPrimWqqCoco$ on COCO). The readable and answerable grades use the 400 calibration rows for support. }%DIFDELCMD < \cfSupportInelChest{} %%%
\DIFdel{of the }%DIFDELCMD < \cfChestCellsN{} %%%
\DIFdel{chest cells fail that cohort's requirement of ten known positives and ten known negatives: CheXpert Atelectasis in each of the }%DIFDELCMD < \cfSupportInelChexBlocks{} %%%
\DIFdel{CheXpert blocks, with }%DIFDELCMD < \cfCalChexInelPos{} %%%
\DIFdel{known positives and }%DIFDELCMD < \cfCalChexInelNeg{} %%%
\DIFdel{known negatives among the 400 calibration rows. Every NIH concept meets the requirement on that cohort; Pneumothorax has the smallest count, at }%DIFDELCMD < \cfCalNihMinPos{} %%%
\DIFdel{positives. On COCO, bicycle is ineligible in each of the }%DIFDELCMD < \cfSupportInelCocoBlocks{} %%%
\DIFdel{COCO blocks, with }%DIFDELCMD < \cfCalCocoInelPos{} %%%
\DIFdel{positives. Support counts for these two grades come solely from calibration. The test cohort carrying the writes has its own counts (}%DIFDELCMD < \cfTestChexInelPos{} %%%
\DIFdel{positives and }%DIFDELCMD < \cfTestChexInelNeg{} %%%
\DIFdel{negatives for CheXpert Atelectasis, }%DIFDELCMD < \cfTestCocoInelPos{} %%%
\DIFdel{positives for COCO bicycle). Among eligible cells, }%DIFDELCMD < \cfSupportReadableChest{} %%%
\DIFdel{of }%DIFDELCMD < \cfSupportReadableNChest{} %%%
\DIFdel{are readable (}%DIFDELCMD < \cfSupportReadablePctChest%%%
\DIFdel{\%) and }%DIFDELCMD < \cfSupportAnswerableChest{} %%%
\DIFdel{of }%DIFDELCMD < \cfSupportAnswerableNChest{} %%%
\DIFdel{are answerable (}%DIFDELCMD < \cfSupportAnswerablePctChest%%%
\DIFdel{\%).
}\DIFdelend \DIFaddbegin \DIFadd{.
}\DIFaddend 

\DIFaddbegin \begin{table}[h]
\centering
\setlength{\tabcolsep}{3pt}
\caption{\textbf{\DIFaddFL{The full grade under refitted probes.}} \DIFaddFL{Cells per seed-0 grade (owned, stronger competitor, unresolved, fixed-family advantage without the reference) and how many keep it under both refits, one, or none.}}
\label{tab:cf-refit}
\resizebox{\linewidth}{!}{%
\begin{tabular}{lrrrrrrrrrrrrrrr}
\toprule
\multirow{2}{*}{Dataset} & \multirow{2}{*}{blocks} & \multirow{2}{*}{cells} & \multicolumn{4}{c}{owned} & \multicolumn{3}{c}{stronger competitor} & \multicolumn{3}{c}{unresolved} & \multicolumn{2}{c}{adv.\ no ref.} & \multicolumn{1}{c}{seed 0}\\
\cmidrule(lr){4-7}\cmidrule(lr){8-10}\cmidrule(lr){11-13}\cmidrule(lr){14-15}
 &  &  & \multicolumn{1}{c}{$n$} & \multicolumn{1}{c}{both} & \multicolumn{1}{c}{$\geq1$} & \multicolumn{1}{c}{none} & \multicolumn{1}{c}{$n$} & \multicolumn{1}{c}{both} & \multicolumn{1}{c}{$\geq1$} & \multicolumn{1}{c}{$n$} & \multicolumn{1}{c}{both} & \multicolumn{1}{c}{$\geq1$} & \multicolumn{1}{c}{$n$} & \multicolumn{1}{c}{both} & \multicolumn{1}{c}{agrees}\\
\midrule
NIH ChestX-ray14 & 25 & 150 & 13 & 4 & 8 & 5 & 101 & 69 & 97 & 21 & 5 & 12 & 15 & 0 & 150/150 \\
CheXpert Plus & 25 & 150 & 25 & 6 & 17 & 8 & 95 & 63 & 89 & 21 & 3 & 10 & 9 & 3 & 150/150 \\
COCO & 25 & 150 & 113 & 99 & 110 & 3 & 17 & 13 & 17 & 10 & 4 & 7 & 10 & 2 & 149/150 \\
\textit{chest} & 50 & 300 & 38 & 10 & 25 & 13 & 196 & 132 & 186 & 42 & 8 & 22 & 24 & 3 & 300/300 \\
\textit{all} & 75 & 450 & 151 & 109 & 135 & 16 & 213 & 145 & 203 & 52 & 12 & 29 & 34 & 5 & 449/450 \\
\bottomrule
\end{tabular}}
\end{table}

\begin{DIFnomarkup}\cfnote{附录精简}{原 Robustness 小节的 label source 段（584 词）与原 Further results 小节的 Labels 段内容重复（专家 refit 的 AUROC、样本量与标签来源效应出现两次），合并为一处；三组对照写成两个差值公式 $\Delta_{\text{label}}$、$\Delta_{\text{size}}$。answer direction 的拟合细节移到 A.10。}\end{DIFnomarkup}

\DIFaddend \paragraph{\DIFdelbegin \DIFdel{The answer direction and the label source}\DIFdelend \DIFaddbegin \DIFadd{Radiologist labels}\DIFaddend .} We \DIFdelbegin \DIFdel{fit the answer direction of Section~\ref{sec:ansdir} to one quantity: the model's own clean answer margin $\ell_{\mathrm{yes}}-\ell_{\mathrm{no}}$ on the block's primary template, read from clean forward passes with the write disabled. We use the block's own projected, train-scaled features from the first 3}%DIFDELCMD < {%%%
\DIFdel{,}%DIFDELCMD < }%%%
\DIFdel{000 training-cohort rows in the cohort's fixed manifest order. A seeded shuffle determines this order on NIH, and a hash of patient or image identity determines it on CheXpert and COCO. Neither labels nor answers enter the ordering. For each question, we fit the margin from these features by ridge regression. We select the penalty by five-fold cross-validated $R^2$ over $\{0.1,1,10,100\}$, then refit on all 3}%DIFDELCMD < {%%%
\DIFdel{,}%DIFDELCMD < }%%%
\DIFdel{000 rows with the chosen penalty. Across the }%DIFDELCMD < \cfAnsAlphaCells{} %%%
\DIFdel{fitted questions, the selected penalty is $1$ for }%DIFDELCMD < \cfAnsAlphaOne{}%%%
\DIFdel{, $10$ for }%DIFDELCMD < \cfAnsAlphaTen{}%%%
\DIFdel{, and $100$ for }%DIFDELCMD < \cfAnsAlphaHundred{}%%%
\DIFdel{. The folds split rows: one per patient on CheXpert, one per image on COCO, and rows from }%DIFDELCMD < \cfAnsTrainUnitsNih{} %%%
\DIFdel{patients on NIH. We lift and normalise the fitted coefficient vector exactly as we do a probe normal. The radiologist-label cohort provides two distinct tests. First, we }\DIFdelend re-grade the \DIFdelbegin \DIFdel{unchanged directions fitted on report-derived labels on }\DIFdelend \DIFaddbegin \DIFadd{report-label directions on }\DIFaddend \cfValidRows{} films labelled by radiologist consensus\DIFdelbegin \DIFdel{. These patients appear }\DIFdelend \DIFaddbegin \DIFadd{, from patients }\DIFaddend in no campaign cohort\DIFdelbegin \DIFdel{(Table~\ref{tab:cf-valid}). The absolute difference }\DIFdelend \DIFaddbegin \DIFadd{. }\DIFaddend $|O_q(\text{valid})-O_q(\text{test})|$ has a 90th percentile of \cfValidNinetyDO{}\DIFdelbegin \DIFdel{over all cells. On }\DIFdelend \DIFaddbegin \DIFadd{, and on }\DIFaddend the \cfValidSupported{} \DIFdelbegin \DIFdel{cells with at least ten known positives and ten known negatives on both cohorts, }\DIFdelend \DIFaddbegin \DIFadd{supported cells }\DIFaddend readability agrees in \cfValidReadableAgree{} of \cfValidReadableN{} and answerability in \cfValidAnswerAgree{} of \cfValidAnswerN{}. \DIFdelbegin \DIFdel{Second, we refit the }\DIFdelend \DIFaddbegin \DIFadd{On CheXpert, restricting to rows with a known label keeps the sign of $O_q$ in }\cfValKnownSignAgree{} \DIFadd{of }\cfValKnownN{} \DIFadd{cells.
}

\DIFadd{Refitting the }\DIFaddend six directions on \DIFdelbegin \DIFdel{those radiologist labelsusing the block's own projection, its train-only scaler, and the protocol's probe settings. We cross-fit }\DIFdelend \DIFaddbegin \DIFadd{the radiologist labels, cross-fitted }\DIFaddend over patients in \cfValidfitFolds{} folds \DIFdelbegin \DIFdel{, scoring every film with a direction not fitted on it. We write the refitted directions on the }\DIFdelend \DIFaddbegin \DIFadd{and written on the }\DIFaddend 600 test rows\DIFdelbegin \DIFdel{at the primary dose and template, against the write matrix's own 119-direction random family and its seed-0 sham, and take the same $6\times5$ max-$T$ verdict; the sham here is the permutation of the report-label direction rather than of the refitted one (Table~\ref{tab:cf-ledger}). Across the }%DIFDELCMD < \cfValidfitBlocks{} %%%
\DIFdel{CheXpert blocks with this refit, the expert-label and report-label directions have }\DIFdelend \DIFaddbegin \DIFadd{, gives directions at }\DIFaddend a median raw \DIFdelbegin \DIFdel{model-space }\DIFdelend cosine of \cfValidfitCosRaw{} \DIFdelbegin \DIFdel{and a median covariance-whitened cosine of }\DIFdelend \DIFaddbegin \DIFadd{(whitened }\DIFaddend \cfValidfitCosWhitened{}\DIFdelbegin \DIFdel{. The expert-label refit achieves a median }\DIFdelend \DIFaddbegin \DIFadd{) from the report-label ones, owned in }\cfValidfitOwnedExpert{} \DIFadd{of }\cfValidfitCells{} \DIFadd{cells against }\cfValidfitOwnedReport{}\DIFadd{. The refit reaches a }\DIFaddend held-out AUROC of \cfValidfitAurocExpert{} on the radiologist labels, against \cfValidfitAurocReportDir{} for the report-label direction\DIFdelbegin \DIFdel{on the same films. Those two directions differ in two ways at once: the labels they were fitted to, and the number of rows available to fit them (}\DIFdelend \DIFaddbegin \DIFadd{, and is fitted on }\DIFaddend \cfVfArmExpertRows{} \DIFaddbegin \DIFadd{rows }\DIFaddend per fold against a median of \cfVfArmReportFullRows{}\DIFdelbegin \DIFdel{labelled training rows per concept)}\DIFdelend . Three further arms\DIFdelbegin \DIFdel{separate the two, all }\DIFdelend \DIFaddbegin \DIFadd{, }\DIFaddend scored on the same films \DIFdelbegin \DIFdel{, with the block's own projection, its stored train-only scaler, the protocol's probe settings and the same }%DIFDELCMD < \cfVfArmFolds{}%%%
\DIFdel{-fold }\DIFdelend \DIFaddbegin \DIFadd{with the same }\DIFaddend cross-fitting\DIFdelbegin \DIFdel{over patients (Table~\ref{tab:cf-valid}, lower panel). Report-derived }\DIFdelend \DIFaddbegin \DIFadd{, separate the label source from the sample size: report-derived }\DIFaddend labels of those \DIFdelbegin \DIFdel{same }\DIFdelend \cfVfArmRows{} films reach \cfVfArmReportSameFilmsAuroc{}; report-derived labels of a training subsample \DIFdelbegin \DIFdel{drawn to leave each concept the same }\DIFdelend \DIFaddbegin \DIFadd{with the refit's }\DIFaddend \cfVfArmReportSubRows{} fitting rows\DIFdelbegin \DIFdel{as the refit}\DIFdelend , redrawn \cfVfArmDraws{} times, reach \cfVfArmReportSubAuroc{}; \DIFdelbegin \DIFdel{report-derived labels of every training row reach }\DIFdelend \DIFaddbegin \DIFadd{and every labelled training row reaches }\DIFaddend \cfVfArmReportFullAuroc{}. \DIFdelbegin \DIFdel{Holding the number of fitting rows fixed, the label source moves the result by $\cfVfArmDeltaLabel$; holding the label source fixed, the sample size moves it by $\cfVfArmDeltaSize$. }\DIFdelend \DIFaddbegin \DIFadd{With $\mathrm{AUROC}_{L,n}$ the held-out AUROC against the radiologist labels of a direction fitted on $n$ rows carrying labels $L$, $n=\cfVfArmReportSubRows$ the refit's sample and $N$ every labelled training row,
}\begin{equation}
\begin{aligned}
\Delta_{\text{label}}&=\mathrm{AUROC}_{\text{radiologist},n}-\mathrm{AUROC}_{\text{report},n}=\cfVfArmDeltaLabel,\\
\Delta_{\text{size}}&=\mathrm{AUROC}_{\text{report},N}-\mathrm{AUROC}_{\text{report},n}=\cfVfArmDeltaSize .
\end{aligned}
\label{eq:label-size}
\end{equation}
\DIFaddend The size-matched report refit also \DIFdelbegin \DIFdel{sits at a raw model-space cosine of }%DIFDELCMD < \cfVfArmReportSubCos{} %%%
\DIFdel{to }\DIFdelend \DIFaddbegin \DIFadd{lies as far from }\DIFaddend the shipped direction \DIFdelbegin \DIFdel{and a whitened cosine of }%DIFDELCMD < \cfVfArmReportSubCosWhite{}%%%
\DIFdel{, against }\DIFdelend \DIFaddbegin \DIFadd{as the radiologist refit (raw cosine }\cfVfArmReportSubCos{} \DIFadd{against }\DIFaddend \cfVfArmExpertCos{}\DIFdelbegin \DIFdel{and }%DIFDELCMD < \cfVfArmExpertCosWhite{} %%%
\DIFdel{for the expert refit. The distance between the refit and the shipped direction is the estimation sample, not the label source. The expert-label directions are owned in }%DIFDELCMD < \cfValidfitOwnedExpert{} %%%
\DIFdel{of }%DIFDELCMD < \cfValidfitCells{} %%%
\DIFdel{cells, against }%DIFDELCMD < \cfValidfitOwnedReport{} %%%
\DIFdel{for the report-label directions of the same blocks.
}%DIFDELCMD < 

%DIFDELCMD < %%%
\DIFdel{Table~\ref{tab:cf-geometry} reports direction and label geometry. Table~\ref{tab:cf-scale} reports scale, saturation, refit, connector, batch, and numerics statistics supporting Section~\ref{sec:robust-geometry}. Table~\ref{tab:cf-refit} reports the full grade under refitted probes. Table~\ref{tab:cf-pairs} reports paired shared-tower differences supporting Section~\ref{sec:reader}. Table~\ref{tab:cf-altdir} reports alternative directions, lifts, token weights, and competitor families. Table~\ref{tab:cf-ansdir} reports the answer direction on the primary and held-out templates. Table~\ref{tab:cf-attr} reports non-clinical attribute directions. Table~\ref{tab:cf-validation} reports what the answer direction reads and the column selectivity of $W$. Table~\ref{tab:cf-valid} reports the grade on radiologist labels and separates the label source of a refitted direction from the number of rows it was fitted on. Table~\ref{tab:cf-towerswap} reports the crossed tower and reader. Table~\ref{tab:cf-semend} reports the four semantic endpoints. Table~\ref{tab:cf-fgobj} reports the fine-grained object family with the difficulty-matched comparison (Appendix~\ref{app:cf-crossed}}\DIFdelend ).

\DIFdelbegin %DIFDELCMD < \begin{table}[h]
%DIFDELCMD < \centering
%DIFDELCMD < \setlength{\tabcolsep}{5pt}
%DIFDELCMD < %%%
%DIFDELCMD < \caption{%
{%DIFAUXCMD
\textbf{\DIFdelFL{Direction and label geometry against off-diagonal dominance.}} %DIFAUXCMD
\DIFdelFL{One column per dataset, with the two chest sets pooled in the third. A ``--'' marks a row that is defined per dataset only.}}
%DIFAUXCMD
%DIFDELCMD < \label{tab:cf-geometry}
%DIFDELCMD < \fitwidth{%
%DIFDELCMD < \begin{tabular}{lrrrr}
%DIFDELCMD < \toprule
%DIFDELCMD < & \multicolumn{1}{c}{NIH ChestX-ray14} & \multicolumn{1}{c}{CheXpert Plus} & \multicolumn{1}{c}{both chest sets} & \multicolumn{1}{c}{COCO}\\
%DIFDELCMD < \midrule
%DIFDELCMD < \multicolumn{5}{l}{\textit{Geometry of the six directions and the six labels}} \\
%DIFDELCMD < \quad mean off-diagonal cosine between the directions & 0.019 & 0.215 & -- & 0.043 \\
%DIFDELCMD < \quad mean off-diagonal $|\cos|$ between the directions & 0.093 & 0.218 & -- & 0.058 \\
%DIFDELCMD < \quad mean $|\cos|$ between random unit directions & 0.024 & 0.024 & -- & 0.024 \\
%DIFDELCMD < \quad mean off-diagonal $\phi$ between the labels & 0.026 & 0.381 & -- & 0.022 \\
%DIFDELCMD < \midrule
%DIFDELCMD < \multicolumn{5}{l}{\textit{Cells whose strongest competitor is (chance: one in five)}} \\
%DIFDELCMD < \quad the most similar direction & 31/150 & 33/150 & 64/300 & 31/150 \\
%DIFDELCMD < \quad the label with the largest $\phi$ & 35/150 & 36/150 & 71/300 & 29/150 \\
%DIFDELCMD < \quad the most co-occurring label & 24/150 & 32/150 & 56/300 & 22/150 \\
%DIFDELCMD < \midrule
%DIFDELCMD < \multicolumn{5}{l}{\textit{Spearman $\rho$ over off-diagonal cells of $W_{q,d}-W_{q,q}$ with}} \\
%DIFDELCMD < \quad the cosine between the two directions & 0.085 & 0.048 & 0.037 & 0.042 \\
%DIFDELCMD < \quad the $\phi$ between the two labels & 0.036 & 0.037 & -0.001 & 0.007 \\
%DIFDELCMD < \midrule
%DIFDELCMD < \multicolumn{5}{l}{\textit{Sign of $O_q$}} \\
%DIFDELCMD < \quad cells with $O_q>0$ & 38/150 & 43/150 & 81/300 & 129/150 \\
%DIFDELCMD < \quad cells with $O_q>0$ without the strongest competitor & 62/150 & 75/150 & 137/300 & 136/150 \\
%DIFDELCMD < \quad cells with at least two competitors above the own write & 88/150 & 75/150 & 163/300 & 14/150 \\
%DIFDELCMD < \bottomrule
%DIFDELCMD < \end{tabular}}
%DIFDELCMD < \end{table}
\begin{table}[htbp]
\DIFdeletedtable{%
{\small\textbf{Table (deleted).} \textbf{Robustness of the owned grade.} Owned cells per dataset under each alternative scale, stratum, refit, and locus, beside the grade on the probability ($p$) and logit-margin ($m$) scales; the lower panel rescores in fp32 at batch size one.\par}\medskip
\centering
\setlength{\tabcolsep}{4pt}

\fitwidth{%
\fitwidth{%
\begin{tabular}{lrrrrrrrrrr}
\toprule
\multirow{2}{*}{Dataset} & \multicolumn{3}{c}{owned cells} & \multicolumn{2}{c}{unsaturated, $O_q>0$} & \multicolumn{2}{c}{refits} & \multicolumn{2}{c}{median $|W_{q,q}|$} & \multicolumn{1}{c}{connector}\\
\cmidrule(lr){2-4}\cmidrule(lr){5-6}\cmidrule(lr){7-8}\cmidrule(lr){9-10}
& \multicolumn{1}{c}{$p$} & \multicolumn{1}{c}{$m$} & \multicolumn{1}{c}{both} & \multicolumn{1}{c}{owned} & \multicolumn{1}{c}{other} & \multicolumn{1}{c}{med.\ $|\Delta O|$} & \multicolumn{1}{c}{stable} & \multicolumn{1}{c}{primary} & \multicolumn{1}{c}{connector} & \multicolumn{1}{c}{$>$ p95}\\
\midrule
NIH ChestX-ray14 & 13/150 & 12/150 & 11 & 13/13 & 25/131 & 0.020 & 8/13 & 0.032 & 0.008 & 15/150 \\
CheXpert Plus & 25/150 & 27/150 & 25 & 25/25 & 16/121 & 0.023 & 10/25 & 0.032 & 0.010 & 16/144 \\
COCO & 113/150 & 116/150 & 112 & 105/106 & 16/37 & 0.020 & 107/113 & 0.127 & 0.008 & 106/150 \\
\midrule
\multicolumn{11}{p{1.286\linewidth}}{\textit{Batch effects at the gates (61 blocks): batched-vs-single candidate-logit difference 0.00--1.84 (median 0.28); margin sign agreement in 55 of 61 blocks; owned cells in disagreeing blocks: 0.}} \\
\bottomrule
\end{tabular}}}

\medskip
\setlength{\tabcolsep}{5pt}
\fitwidth{%
\fitwidth{%
\begin{tabular}{llrrrr}
\toprule
\multicolumn{6}{l}{\textit{Numerics: the same grid rescored in fp32 and at batch size one}} \\
\midrule
 & & \multicolumn{2}{c}{fp32 weights and forward} & \multicolumn{2}{c}{bf16, batch size one} \\
\cmidrule(lr){3-4}\cmidrule(lr){5-6}
Checkpoint & Dataset & $\max|\Delta W|$ & $\max|\Delta C|$ & $\max|\Delta W|$ & $\max|\Delta C|$ \\
\midrule
Qwen2.5-VL-7B & NIH ChestX-ray14 & 0.0087 & 0.0063 & 0.0065 & 0.0049 \\
Qwen2.5-VL-7B & COCO & 0.0060 & 0.0032 & 0.0028 & 0.0020 \\
Qwen3-VL-8B & NIH ChestX-ray14 & 0.0177 & 0.0089 & 0.0204 & 0.0089 \\
InternVL3.5-8B & NIH ChestX-ray14 & 0.0018 & 0.0014 & 0.0018 & 0.0014 \\
Gemma 3 4B & NIH ChestX-ray14 & 0.0544 & 0.0493 & 0.0261 & 0.0377 \\
Gemma 3 12B & NIH ChestX-ray14 & 0.0558 & 0.0359 & 0.0111 & 0.0109 \\
Gemma 3 27B & NIH ChestX-ray14 & 0.0413 & 0.0186 & 0.0149 & 0.0053 \\
MedGemma 4B & NIH ChestX-ray14 & 0.0125 & 0.0069 & 0.0056 & 0.0054 \\
Lingshu 32B & NIH ChestX-ray14 & 0.0032 & 0.0023 & 0.0009 & 0.0007 \\
\bottomrule
\end{tabular}}}
}
\end{table}
%DIFDELCMD < \begin{table}[h]
%DIFDELCMD < \centering
%DIFDELCMD < \setlength{\tabcolsep}{3pt}
%DIFDELCMD < %%%
%DIFDELCMD < \caption{%
{%DIFAUXCMD
\textbf{\DIFdelFL{The full grade under refitted probes.}} %DIFAUXCMD
\DIFdelFL{Cells per seed-0 grade (owned, stronger competitor, unresolved, fixed-family advantage without the reference) and how many keep it under both refits, one, or none.}}
%DIFAUXCMD
%DIFDELCMD < \label{tab:cf-refit}
%DIFDELCMD < \fitwidth{%
%DIFDELCMD < \begin{tabular}{lrrrrrrrrrrrrrrr}
%DIFDELCMD < \toprule
%DIFDELCMD < \multirow{2}{*}{Dataset} & \multirow{2}{*}{blocks} & \multirow{2}{*}{cells} & \multicolumn{4}{c}{owned} & \multicolumn{3}{c}{stronger competitor} & \multicolumn{3}{c}{unresolved} & \multicolumn{2}{c}{adv.\ no ref.} & \multicolumn{1}{c}{seed 0}\\
%DIFDELCMD < \cmidrule(lr){4-7}\cmidrule(lr){8-10}\cmidrule(lr){11-13}\cmidrule(lr){14-15}
%DIFDELCMD <  &  &  & \multicolumn{1}{c}{$n$} & \multicolumn{1}{c}{both} & \multicolumn{1}{c}{$\geq1$} & \multicolumn{1}{c}{none} & \multicolumn{1}{c}{$n$} & \multicolumn{1}{c}{both} & \multicolumn{1}{c}{$\geq1$} & \multicolumn{1}{c}{$n$} & \multicolumn{1}{c}{both} & \multicolumn{1}{c}{$\geq1$} & \multicolumn{1}{c}{$n$} & \multicolumn{1}{c}{both} & \multicolumn{1}{c}{agrees}\\
%DIFDELCMD < \midrule
%DIFDELCMD < NIH ChestX-ray14 & 25 & 150 & 13 & 4 & 8 & 5 & 101 & 69 & 97 & 21 & 5 & 12 & 15 & 0 & 150/150 \\
%DIFDELCMD < CheXpert Plus & 25 & 150 & 25 & 6 & 17 & 8 & 95 & 63 & 89 & 21 & 3 & 10 & 9 & 3 & 150/150 \\
%DIFDELCMD < COCO & 25 & 150 & 113 & 99 & 110 & 3 & 17 & 13 & 17 & 10 & 4 & 7 & 10 & 2 & 149/150 \\
%DIFDELCMD < \textit{chest} & 50 & 300 & 38 & 10 & 25 & 13 & 196 & 132 & 186 & 42 & 8 & 22 & 24 & 3 & 300/300 \\
%DIFDELCMD < \textit{all} & 75 & 450 & 151 & 109 & 135 & 16 & 213 & 145 & 203 & 52 & 12 & 29 & 34 & 5 & 449/450 \\
%DIFDELCMD < \bottomrule
%DIFDELCMD < \end{tabular}}
%DIFDELCMD < \end{table}
%DIFDELCMD < \begin{table}[h]
%DIFDELCMD < \centering
%DIFDELCMD < \setlength{\tabcolsep}{4pt}
%DIFDELCMD < %%%
%DIFDELCMD < \caption{%
{%DIFAUXCMD
\textbf{\DIFdelFL{Same write, different reader: paired differences.}} %DIFAUXCMD
\DIFdelFL{$\Delta O_q=O_q(\text{larger})-O_q(\text{smaller})$ per shared-tower pair on the same 600 test rows, with paired patient-bootstrap 95\% intervals; the last two columns divide $|\Delta O_q|$ by the two blocks' mean refit standard deviation.}}
%DIFAUXCMD
%DIFDELCMD < \label{tab:cf-pairs}
%DIFDELCMD < \fitwidth{%
%DIFDELCMD < \begin{tabular}{llcrrrrr}
%DIFDELCMD < \toprule
%DIFDELCMD < \multirow{2}{*}{Pair} & \multirow{2}{*}{Dataset} & \multirow{2}{*}{relation} & \multicolumn{1}{c}{$\Delta O_q$ CI} & \multicolumn{2}{c}{$|\Delta O_q|$} & \multicolumn{2}{c}{ratio to refit SD}\\
%DIFDELCMD < \cmidrule(lr){5-6}\cmidrule(lr){7-8}
%DIFDELCMD < &  &  & \multicolumn{1}{c}{$\not\ni 0$} & \multicolumn{1}{c}{median} & \multicolumn{1}{c}{max} & \multicolumn{1}{c}{median} & \multicolumn{1}{c}{max}\\
%DIFDELCMD < \midrule
%DIFDELCMD < Gemma 3 4B $\to$ Gemma 3 12B & NIH ChestX-ray14 & bit-identical & 5/6 & 0.33 & 0.83 & 2.7 & 6.9 \\
%DIFDELCMD < Gemma 3 4B $\to$ Gemma 3 12B & CheXpert Plus & bit-identical & 6/6 & 0.29 & 0.83 & 2.1 & 5.9 \\
%DIFDELCMD < Gemma 3 4B $\to$ Gemma 3 12B & COCO & bit-identical & 5/6 & 0.23 & 0.87 & 1.6 & 6.0 \\
%DIFDELCMD < Gemma 3 12B $\to$ Gemma 3 27B & NIH ChestX-ray14 & bf16-equal & 6/6 & 0.37 & 0.79 & 3.7 & 8.0 \\
%DIFDELCMD < Gemma 3 12B $\to$ Gemma 3 27B & CheXpert Plus & bf16-equal & 6/6 & 0.35 & 0.81 & 1.4 & 3.3 \\
%DIFDELCMD < Gemma 3 12B $\to$ Gemma 3 27B & COCO & bf16-equal & 5/6 & 0.15 & 0.40 & 1.7 & 4.3 \\
%DIFDELCMD < Gemma 3 4B $\to$ Gemma 3 27B & NIH ChestX-ray14 & bf16-equal & 5/6 & 0.07 & 0.24 & 1.8 & 6.0 \\
%DIFDELCMD < Gemma 3 4B $\to$ Gemma 3 27B & CheXpert Plus & bf16-equal & 6/6 & 0.17 & 0.65 & 4.6 & 17.3 \\
%DIFDELCMD < Gemma 3 4B $\to$ Gemma 3 27B & COCO & bf16-equal & 6/6 & 0.16 & 1.02 & 1.4 & 8.8 \\
%DIFDELCMD < MedGemma 4B $\to$ MedGemma 27B & NIH ChestX-ray14 & bf16-equal & 5/6 & 0.09 & 0.30 & 1.5 & 5.0 \\
%DIFDELCMD < MedGemma 4B $\to$ MedGemma 27B & CheXpert Plus & bf16-equal & 5/6 & 0.19 & 0.54 & 1.8 & 5.0 \\
%DIFDELCMD < MedGemma 4B $\to$ MedGemma 27B & COCO & bf16-equal & 3/6 & 0.05 & 0.35 & 0.6 & 4.6 \\
%DIFDELCMD < \bottomrule
%DIFDELCMD < \end{tabular}}
%DIFDELCMD < \end{table}
%DIFDELCMD < \begin{table}[h]
%DIFDELCMD < \centering
%DIFDELCMD < \setlength{\tabcolsep}{4pt}
%DIFDELCMD < %%%
%DIFDELCMD < \caption{%
{%DIFAUXCMD
\textbf{\DIFdelFL{Alternative directions and competitor families.}} %DIFAUXCMD
\DIFdelFL{Cells owned under the paper's rule within each family, and the median $O_q$, per dataset. Each panel repeats the reference write on its own blocks.}}
%DIFAUXCMD
%DIFDELCMD < \label{tab:cf-altdir}
%DIFDELCMD < \fitwidth{%
%DIFDELCMD < \begin{tabular}{llrrrrrr}
%DIFDELCMD < \toprule
%DIFDELCMD < \multirow{2}{*}{Direction} & \multirow{2}{*}{Lift} & \multicolumn{2}{c}{NIH ChestX-ray14} & \multicolumn{2}{c}{CheXpert Plus} & \multicolumn{2}{c}{COCO}\\
%DIFDELCMD < \cmidrule(lr){3-4}\cmidrule(lr){5-6}\cmidrule(lr){7-8}
%DIFDELCMD < &  & \multicolumn{1}{c}{owned} & \multicolumn{1}{c}{med.\ $O_q$} & \multicolumn{1}{c}{owned} & \multicolumn{1}{c}{med.\ $O_q$} & \multicolumn{1}{c}{owned} & \multicolumn{1}{c}{med.\ $O_q$}\\
%DIFDELCMD < \midrule
%DIFDELCMD < \multicolumn{8}{l}{\textit{Direction construction (20 blocks per dataset)}} \\
%DIFDELCMD < \quad logistic normal & coefficient & 7/120 & -0.03 & 19/120 & -0.01 & 101/120 & +0.18 \\
%DIFDELCMD < \quad difference of means & coefficient & 16/120 & -0.02 & 18/120 & -0.01 & 62/120 & +0.02 \\
%DIFDELCMD < \quad Haufe pattern & coefficient & 17/120 & -0.02 & 17/120 & -0.01 & 53/120 & +0.01 \\
%DIFDELCMD < \quad orthogonalised normal & coefficient & 8/120 & -0.02 & 14/120 & -0.01 & 98/120 & +0.13 \\
%DIFDELCMD < \quad residualised normal & coefficient & 8/120 & -0.02 & 13/120 & -0.01 & 102/120 & +0.20 \\
%DIFDELCMD < \quad owned under at least one &  & 25/120 & -- & 41/120 & -- & 108/120 & -- \\
%DIFDELCMD < \midrule
%DIFDELCMD < \multicolumn{8}{l}{\textit{Displacement lift (19 blocks per dataset)}} \\
%DIFDELCMD < \quad logistic normal & coefficient & 7/114 & -0.03 & 19/114 & -0.01 & 100/114 & +0.19 \\
%DIFDELCMD < \quad difference of means & displacement & 21/114 & -0.03 & 19/114 & -0.00 & 58/114 & +0.02 \\
%DIFDELCMD < \quad Haufe pattern & displacement & 17/114 & -0.02 & 17/114 & -0.00 & 49/114 & +0.01 \\
%DIFDELCMD < \midrule
%DIFDELCMD < \multicolumn{8}{l}{\textit{Token weights of the logistic normal (4 NIH, 2 CheXpert, 2 COCO blocks)}} \\
%DIFDELCMD < \quad uniform & coefficient & 3/24 & -0.08 & 5/12 & +0.00 & 12/12 & +0.60 \\
%DIFDELCMD < \quad softmax of the probe score & coefficient & 0/24 & -0.01 & 0/12 & -0.00 & 12/12 & +0.16 \\
%DIFDELCMD < \quad top quarter of tokens & coefficient & 2/24 & -0.12 & 3/12 & -0.05 & 12/12 & +0.57 \\
%DIFDELCMD < \midrule
%DIFDELCMD < \multicolumn{8}{l}{\textit{Competitor family of the logistic normal (5 NIH, 4 CheXpert blocks)}} \\
%DIFDELCMD < \quad six clinical directions & coefficient & 3/30 & -0.07 & 8/24 & -0.07 & -- & -- \\
%DIFDELCMD < \quad with every extra dataset label & coefficient & 3/30 & -0.07 & 5/24 & -0.10 & -- & -- \\
%DIFDELCMD < \bottomrule
%DIFDELCMD < \end{tabular}}
%DIFDELCMD < \end{table}
%DIFDELCMD < \begin{table}[h]
%DIFDELCMD < \centering
%DIFDELCMD < \setlength{\tabcolsep}{3.5pt}
%DIFDELCMD < %%%
%DIFDELCMD < \caption{%
{%DIFAUXCMD
\textbf{\DIFdelFL{The answer direction.}} %DIFAUXCMD
\DIFdelFL{Top, primary template: concepts owned by the label direction $\hat w_q$ and the answer direction $a_q$ (kept, rescued, lost), and medians over the six questions. Bottom: $a_q$ under each held-out template, with owned concepts, in parentheses those owned under IY, and kept (concept, template) pairs.}}
%DIFAUXCMD
%DIFDELCMD < \label{tab:cf-ansdir}
%DIFDELCMD < \fitwidth{%
%DIFDELCMD < \begin{tabular}{llrrrrrrrrrr}
%DIFDELCMD < \toprule
%DIFDELCMD < \multirow{2}{*}{Checkpoint} & \multirow{2}{*}{Dataset} & \multicolumn{5}{c}{owned concepts} & \multirow{2}{*}{$a_q>$ comp.} & \multicolumn{4}{c}{median}\\
%DIFDELCMD < \cmidrule(lr){3-7}\cmidrule(lr){9-12}
%DIFDELCMD < &  & \multicolumn{1}{c}{$\hat w_q$} & \multicolumn{1}{c}{$a_q$} & \multicolumn{1}{c}{kept} & \multicolumn{1}{c}{rescued} & \multicolumn{1}{c}{lost} &  & \multicolumn{1}{c}{$W^a_{q,q}$} & \multicolumn{1}{c}{$\cos$} & \multicolumn{1}{c}{$\cos_\Sigma$} & \multicolumn{1}{c}{$R^2$}\\
%DIFDELCMD < \midrule
%DIFDELCMD < Qwen2.5-VL-7B & NIH ChestX-ray14 & 0/6 & 3/6 & 0 & 3 & 0 & 6/6 & 0.61 & 0.07 & 0.35 & 0.54 \\
%DIFDELCMD < Qwen2.5-VL-7B & CheXpert Plus & 0/6 & 6/6 & 0 & 6 & 0 & 6/6 & 0.55 & 0.06 & 0.36 & 0.48 \\
%DIFDELCMD < Qwen2.5-VL-7B & COCO & 6/6 & 6/6 & 6 & 0 & 0 & 6/6 & 0.78 & 0.64 & 0.84 & 0.58 \\
%DIFDELCMD < Gemma 3 12B & NIH ChestX-ray14 & 1/6 & 2/6 & 1 & 1 & 0 & 5/6 & 0.47 & 0.18 & 0.33 & 0.17 \\
%DIFDELCMD < Gemma 3 12B & CheXpert Plus & 1/6 & 0/6 & 0 & 0 & 1 & 1/6 & 0.37 & 0.17 & 0.10 & 0.19 \\
%DIFDELCMD < Gemma 3 12B & COCO & 6/6 & 6/6 & 6 & 0 & 0 & 6/6 & 0.61 & 0.59 & 0.79 & 0.37 \\
%DIFDELCMD < Lingshu 32B & NIH ChestX-ray14 & 2/6 & 3/6 & 2 & 1 & 0 & 6/6 & 0.24 & 0.30 & 0.60 & 0.77 \\
%DIFDELCMD < Lingshu 32B & CheXpert Plus & 5/6 & 5/6 & 5 & 0 & 0 & 6/6 & 0.22 & 0.33 & 0.87 & 0.80 \\
%DIFDELCMD < Lingshu 32B & COCO & 6/6 & 6/6 & 6 & 0 & 0 & 6/6 & 0.36 & 0.54 & 0.73 & 0.55 \\
%DIFDELCMD < \midrule
%DIFDELCMD < NIH ChestX-ray14 & 3 blocks & 3/18 & 8/18 & 3 & 5 & 0 & 17/18 & 0.47 & 0.18 & 0.43 & 0.54 \\
%DIFDELCMD < CheXpert Plus & 3 blocks & 6/18 & 11/18 & 5 & 6 & 1 & 13/18 & 0.35 & 0.16 & 0.36 & 0.48 \\
%DIFDELCMD < \textit{chest} & 6 blocks & 9/36 & 19/36 & 8 & 11 & 1 & 30/36 & 0.41 & 0.18 & 0.39 & 0.50 \\
%DIFDELCMD < COCO & 3 blocks & 18/18 & 18/18 & 18 & 0 & 0 & 18/18 & 0.61 & 0.61 & 0.81 & 0.54 \\
%DIFDELCMD < \bottomrule
%DIFDELCMD < \end{tabular}}
%DIFDELCMD < %%%
\DIFdelendFL \DIFaddbeginFL \begin{DIFnomarkup}\cfnote{附录精简}{共享塔配对、规模与家族、胸片训练塔（原 Further results）与换塔、tensor replay（原 Crossing）合为一节。换塔三段（约 680 词）压成两段约 400 词：删去收据的逐项清单，只留逐位验证的结论；reader / tower 效应写成公式；保留原生组合在小样本上的核对、三个数据集与三个量的效应，以及“混搭组合回答变差”这一不利结果。replay 由约 290 词改为两段约 230 词，先写做法与漂移、再写自身回放的对照与结果。}\end{DIFnomarkup}
\DIFaddendFL 

\DIFdelbeginFL %DIFDELCMD < \medskip
%DIFDELCMD < \fitwidth{%
%DIFDELCMD < \begin{tabular}{llrrrrrr}
%DIFDELCMD < \toprule
%DIFDELCMD < \multirow{2}{*}{Checkpoint} & \multirow{2}{*}{Dataset} & \multicolumn{5}{c}{held-out template: owned (IY-owned kept)} & \multirow{2}{*}{kept pairs}\\
%DIFDELCMD < \cmidrule(lr){3-7}
%DIFDELCMD < &  & \multicolumn{1}{c}{WY} & \multicolumn{1}{c}{IA} & \multicolumn{1}{c}{IB} & \multicolumn{1}{c}{WA} & \multicolumn{1}{c}{WB} & \\
%DIFDELCMD < \midrule
%DIFDELCMD < Qwen2.5-VL-7B & NIH ChestX-ray14 & 4 (3) & 3 (3) & 4 (3) & 4 (3) & 4 (3) & 15/15 \\
%DIFDELCMD < Qwen2.5-VL-7B & CheXpert Plus & 4 (4) & 6 (6) & 5 (5) & 5 (5) & 5 (5) & 25/30 \\
%DIFDELCMD < Qwen2.5-VL-7B & COCO & 6 (6) & 6 (6) & 6 (6) & 6 (6) & 6 (6) & 30/30 \\
%DIFDELCMD < Gemma 3 12B & NIH ChestX-ray14 & 2 (1) & 1 (1) & 0 (0) & 1 (0) & 1 (0) & 2/10 \\
%DIFDELCMD < Gemma 3 12B & CheXpert Plus & 0 (0) & 1 (0) & 1 (0) & 1 (0) & 1 (0) & 0/0 \\
%DIFDELCMD < Gemma 3 12B & COCO & 5 (5) & 6 (6) & 6 (6) & 6 (6) & 6 (6) & 29/30 \\
%DIFDELCMD < Lingshu 32B & NIH ChestX-ray14 & 3 (3) & 3 (3) & 3 (3) & 3 (3) & 3 (3) & 15/15 \\
%DIFDELCMD < Lingshu 32B & CheXpert Plus & 5 (5) & 3 (3) & 3 (3) & 4 (4) & 4 (4) & 19/25 \\
%DIFDELCMD < Lingshu 32B & COCO & 6 (6) & 6 (6) & 6 (6) & 6 (6) & 6 (6) & 30/30 \\
%DIFDELCMD < \midrule
%DIFDELCMD < \textit{chest} & 6 blocks & 18 (16) & 17 (16) & 16 (14) & 18 (15) & 18 (15) & 76/95 \\
%DIFDELCMD < COCO & 3 blocks & 17 (17) & 18 (18) & 18 (18) & 18 (18) & 18 (18) & 89/90 \\
%DIFDELCMD < \bottomrule
%DIFDELCMD < \end{tabular}}
%DIFDELCMD < \end{table}
\begin{table}[htbp]
\DIFoldtable{%
{\small\textbf{Table (before revision).} \textbf{The answer direction.} Top, primary template: concepts owned by the label direction $\hat w_q$ and the answer direction $a_q$ (kept, rescued, lost), and medians over the six questions. Bottom: $a_q$ under each held-out template, with owned concepts, in parentheses those owned under IY, and kept (concept, template) pairs.\par}\medskip
\centering
\setlength{\tabcolsep}{3.5pt}

\fitwidth{%
\fitwidth{%
\begin{tabular}{llrrrrrrrrrr}
\toprule
\multirow{2}{*}{Checkpoint} & \multirow{2}{*}{Dataset} & \multicolumn{5}{c}{owned concepts} & \multirow{2}{*}{$a_q>$ comp.} & \multicolumn{4}{c}{median}\\
\cmidrule(lr){3-7}\cmidrule(lr){9-12}
&  & \multicolumn{1}{c}{$\hat w_q$} & \multicolumn{1}{c}{$a_q$} & \multicolumn{1}{c}{kept} & \multicolumn{1}{c}{rescued} & \multicolumn{1}{c}{lost} &  & \multicolumn{1}{c}{$W^a_{q,q}$} & \multicolumn{1}{c}{$\cos$} & \multicolumn{1}{c}{$\cos_\Sigma$} & \multicolumn{1}{c}{$R^2$}\\
\midrule
Qwen2.5-VL-7B & NIH ChestX-ray14 & 0/6 & 3/6 & 0 & 3 & 0 & 6/6 & 0.61 & 0.07 & 0.35 & 0.54 \\
Qwen2.5-VL-7B & CheXpert Plus & 0/6 & 6/6 & 0 & 6 & 0 & 6/6 & 0.55 & 0.06 & 0.36 & 0.48 \\
Qwen2.5-VL-7B & COCO & 6/6 & 6/6 & 6 & 0 & 0 & 6/6 & 0.78 & 0.64 & 0.84 & 0.58 \\
Gemma 3 12B & NIH ChestX-ray14 & 1/6 & 2/6 & 1 & 1 & 0 & 5/6 & 0.47 & 0.18 & 0.33 & 0.17 \\
Gemma 3 12B & CheXpert Plus & 1/6 & 0/6 & 0 & 0 & 1 & 1/6 & 0.37 & 0.17 & 0.10 & 0.19 \\
Gemma 3 12B & COCO & 6/6 & 6/6 & 6 & 0 & 0 & 6/6 & 0.61 & 0.59 & 0.79 & 0.37 \\
Lingshu 32B & NIH ChestX-ray14 & 2/6 & 3/6 & 2 & 1 & 0 & 6/6 & 0.24 & 0.30 & 0.60 & 0.77 \\
Lingshu 32B & CheXpert Plus & 5/6 & 5/6 & 5 & 0 & 0 & 6/6 & 0.22 & 0.33 & 0.87 & 0.80 \\
Lingshu 32B & COCO & 6/6 & 6/6 & 6 & 0 & 0 & 6/6 & 0.36 & 0.54 & 0.73 & 0.55 \\
\midrule
NIH ChestX-ray14 & 3 blocks & 3/18 & 8/18 & 3 & 5 & 0 & 17/18 & 0.47 & 0.18 & 0.43 & 0.54 \\
CheXpert Plus & 3 blocks & 6/18 & 11/18 & 5 & 6 & 1 & 13/18 & 0.35 & 0.16 & 0.36 & 0.48 \\
\textit{chest} & 6 blocks & 9/36 & 19/36 & 8 & 11 & 1 & 30/36 & 0.41 & 0.18 & 0.39 & 0.50 \\
COCO & 3 blocks & 18/18 & 18/18 & 18 & 0 & 0 & 18/18 & 0.61 & 0.61 & 0.81 & 0.54 \\
\bottomrule
\end{tabular}}}

\medskip
\fitwidth{%
\fitwidth{%
\begin{tabular}{llrrrrrr}
\toprule
\multirow{2}{*}{Checkpoint} & \multirow{2}{*}{Dataset} & \multicolumn{5}{c}{held-out template: owned (IY-owned kept)} & \multirow{2}{*}{kept pairs}\\
\cmidrule(lr){3-7}
&  & \multicolumn{1}{c}{WY} & \multicolumn{1}{c}{IA} & \multicolumn{1}{c}{IB} & \multicolumn{1}{c}{WA} & \multicolumn{1}{c}{WB} & \\
\midrule
Qwen2.5-VL-7B & NIH ChestX-ray14 & 4 (3) & 3 (3) & 4 (3) & 4 (3) & 4 (3) & 15/15 \\
Qwen2.5-VL-7B & CheXpert Plus & 4 (4) & 6 (6) & 5 (5) & 5 (5) & 5 (5) & 25/30 \\
Qwen2.5-VL-7B & COCO & 6 (6) & 6 (6) & 6 (6) & 6 (6) & 6 (6) & 30/30 \\
Gemma 3 12B & NIH ChestX-ray14 & 2 (1) & 1 (1) & 0 (0) & 1 (0) & 1 (0) & 2/10 \\
Gemma 3 12B & CheXpert Plus & 0 (0) & 1 (0) & 1 (0) & 1 (0) & 1 (0) & 0/0 \\
Gemma 3 12B & COCO & 5 (5) & 6 (6) & 6 (6) & 6 (6) & 6 (6) & 29/30 \\
Lingshu 32B & NIH ChestX-ray14 & 3 (3) & 3 (3) & 3 (3) & 3 (3) & 3 (3) & 15/15 \\
Lingshu 32B & CheXpert Plus & 5 (5) & 3 (3) & 3 (3) & 4 (4) & 4 (4) & 19/25 \\
Lingshu 32B & COCO & 6 (6) & 6 (6) & 6 (6) & 6 (6) & 6 (6) & 30/30 \\
\midrule
\textit{chest} & 6 blocks & 18 (16) & 17 (16) & 16 (14) & 18 (15) & 18 (15) & 76/95 \\
COCO & 3 blocks & 17 (17) & 18 (18) & 18 (18) & 18 (18) & 18 (18) & 89/90 \\
\bottomrule
\end{tabular}}}
}
\end{table}
%DIFDELCMD < \begin{table}[htbp]
%DIFDELCMD < \centering
%DIFDELCMD < \setlength{\tabcolsep}{4pt}
%DIFDELCMD < %%%
%DIFDELCMD < \caption{%
{%DIFAUXCMD
\textbf{\DIFdelFL{Non-clinical attributes of the same radiographs.}} %DIFAUXCMD
\DIFdelFL{Top: owned attribute and clinical cells in the nine-direction family per checkpoint and dataset, naming the owned attributes. Bottom: medians over blocks per attribute of probe AUROC, selectivity, clean-answer AUROC, and cosine to the clinical normals.}}
%DIFAUXCMD
%DIFDELCMD < \label{tab:cf-attr}
%DIFDELCMD < \fitwidth{%
%DIFDELCMD < \begin{tabular}{llrrccc}
%DIFDELCMD < \toprule
%DIFDELCMD < Checkpoint & Dataset & \multicolumn{1}{c}{attributes owned} & \multicolumn{1}{c}{findings owned} & \multicolumn{1}{c}{view AP} & \multicolumn{1}{c}{sex F} & \multicolumn{1}{c}{age $\geq60$}\\
%DIFDELCMD < \midrule
%DIFDELCMD < Qwen2.5-VL-3B & NIH ChestX-ray14 & 0/3 & 0/6 & -- & -- & -- \\
%DIFDELCMD < Qwen2.5-VL-3B & CheXpert Plus & 0/3 & 0/6 & -- & -- & -- \\
%DIFDELCMD < Qwen2.5-VL-7B & NIH ChestX-ray14 & 0/3 & 0/6 & -- & -- & -- \\
%DIFDELCMD < Qwen2.5-VL-7B & CheXpert Plus & 0/3 & 0/6 & -- & -- & -- \\
%DIFDELCMD < Qwen2.5-VL-32B & NIH ChestX-ray14 & 0/3 & 0/6 & -- & -- & -- \\
%DIFDELCMD < Qwen2.5-VL-32B & CheXpert Plus & 0/3 & 0/6 & -- & -- & -- \\
%DIFDELCMD < Qwen3-VL-4B & NIH ChestX-ray14 & 1/3 & 0/6 & -- & yes & -- \\
%DIFDELCMD < Qwen3-VL-4B & CheXpert Plus & 1/3 & 2/6 & -- & yes & -- \\
%DIFDELCMD < Qwen3-VL-8B & NIH ChestX-ray14 & 0/3 & 0/6 & -- & -- & -- \\
%DIFDELCMD < Qwen3-VL-8B & CheXpert Plus & 1/3 & 0/6 & -- & -- & yes \\
%DIFDELCMD < Qwen3-VL-32B & NIH ChestX-ray14 & 0/3 & 1/6 & -- & -- & -- \\
%DIFDELCMD < Qwen3-VL-32B & CheXpert Plus & 1/3 & 1/6 & -- & -- & yes \\
%DIFDELCMD < InternVL3.5-8B & NIH ChestX-ray14 & 0/3 & 0/6 & -- & -- & -- \\
%DIFDELCMD < InternVL3.5-8B & CheXpert Plus & 0/3 & 1/6 & -- & -- & -- \\
%DIFDELCMD < InternVL3.5-14B & NIH ChestX-ray14 & 0/3 & 1/6 & -- & -- & -- \\
%DIFDELCMD < InternVL3.5-14B & CheXpert Plus & 0/3 & 1/6 & -- & -- & -- \\
%DIFDELCMD < InternVL3.5-38B & NIH ChestX-ray14 & 0/3 & 0/6 & -- & -- & -- \\
%DIFDELCMD < InternVL3.5-38B & CheXpert Plus & 0/3 & 2/6 & -- & -- & -- \\
%DIFDELCMD < Gemma 3 4B & NIH ChestX-ray14 & 0/3 & 0/6 & -- & -- & -- \\
%DIFDELCMD < Gemma 3 4B & CheXpert Plus & 0/3 & 0/6 & -- & -- & -- \\
%DIFDELCMD < Gemma 3 12B & NIH ChestX-ray14 & 0/3 & 1/6 & -- & -- & -- \\
%DIFDELCMD < Gemma 3 12B & CheXpert Plus & 0/3 & 1/6 & -- & -- & -- \\
%DIFDELCMD < Gemma 3 27B & NIH ChestX-ray14 & 0/3 & 0/6 & -- & -- & -- \\
%DIFDELCMD < Gemma 3 27B & CheXpert Plus & 0/3 & 0/6 & -- & -- & -- \\
%DIFDELCMD < MedGemma 4B & NIH ChestX-ray14 & 0/3 & 0/6 & -- & -- & -- \\
%DIFDELCMD < MedGemma 4B & CheXpert Plus & 1/3 & 1/6 & -- & -- & yes \\
%DIFDELCMD < MedGemma 27B & NIH ChestX-ray14 & 1/3 & 0/6 & -- & -- & yes \\
%DIFDELCMD < MedGemma 27B & CheXpert Plus & 0/3 & 0/6 & -- & -- & -- \\
%DIFDELCMD < Lingshu 7B & NIH ChestX-ray14 & 2/3 & 1/6 & -- & yes & yes \\
%DIFDELCMD < Lingshu 7B & CheXpert Plus & 0/3 & 3/6 & -- & -- & -- \\
%DIFDELCMD < Lingshu 32B & NIH ChestX-ray14 & 2/3 & 2/6 & -- & yes & yes \\
%DIFDELCMD < Lingshu 32B & CheXpert Plus & 1/3 & 4/6 & -- & -- & yes \\
%DIFDELCMD < LLaVA-1.5-7B & NIH ChestX-ray14 & 0/3 & 0/6 & -- & -- & -- \\
%DIFDELCMD < LLaVA-1.5-7B & CheXpert Plus & 0/3 & 0/6 & -- & -- & -- \\
%DIFDELCMD < LLaVA-1.5-13B & NIH ChestX-ray14 & 1/3 & 0/6 & yes & -- & -- \\
%DIFDELCMD < LLaVA-1.5-13B & CheXpert Plus & 0/3 & 0/6 & -- & -- & -- \\
%DIFDELCMD < LLaVA-Med 1.5 7B & NIH ChestX-ray14 & 0/3 & 0/6 & -- & -- & -- \\
%DIFDELCMD < LLaVA-Med 1.5 7B & CheXpert Plus & 0/3 & 1/6 & -- & -- & -- \\
%DIFDELCMD < \midrule
%DIFDELCMD < NIH ChestX-ray14 & 19 blocks & 7/57 & 6/114 & 1/19 & 3/19 & 3/19 \\
%DIFDELCMD < CheXpert Plus & 19 blocks & 5/57 & 17/114 & 0/19 & 1/19 & 4/19 \\
%DIFDELCMD < \textit{chest} & 38 blocks & 12/114 & 23/228 & 1/38 & 4/38 & 7/38 \\
%DIFDELCMD < \midrule
%DIFDELCMD < \multicolumn{7}{l}{\textit{Attribute probes and directions: medians over the blocks of each dataset}} \\
%DIFDELCMD < Attribute & Dataset & probe AUROC & selectivity & answer AUROC & med.\ $|\cos|$ & owned \\
%DIFDELCMD < \midrule
%DIFDELCMD < AP (portable) projection & NIH ChestX-ray14 & 0.999 & 0.32 & 0.76 & 0.03 & 1/19 \\
%DIFDELCMD < AP (portable) projection & CheXpert Plus & 0.998 & 0.31 & 0.73 & 0.05 & 0/19 \\
%DIFDELCMD < AP (portable) projection & \textit{chest} & 0.999 & 0.31 & 0.74 & 0.04 & 1/38 \\
%DIFDELCMD < female sex & NIH ChestX-ray14 & 0.984 & 0.30 & 0.73 & 0.05 & 3/19 \\
%DIFDELCMD < female sex & CheXpert Plus & 0.957 & 0.27 & 0.65 & 0.06 & 1/19 \\
%DIFDELCMD < female sex & \textit{chest} & 0.971 & 0.28 & 0.71 & 0.05 & 4/38 \\
%DIFDELCMD < age $\geq60$ & NIH ChestX-ray14 & 0.896 & 0.21 & 0.62 & 0.05 & 3/19 \\
%DIFDELCMD < age $\geq60$ & CheXpert Plus & 0.898 & 0.21 & 0.64 & 0.06 & 4/19 \\
%DIFDELCMD < age $\geq60$ & \textit{chest} & 0.896 & 0.21 & 0.63 & 0.06 & 7/38 \\
%DIFDELCMD < \midrule
%DIFDELCMD < \multicolumn{7}{l}{\textit{The same comparison read three ways, pooled over the 38 chest blocks}} \\
%DIFDELCMD < Comparison & unit & attribute cells & owned & finding cells & owned & difference \\
%DIFDELCMD < \midrule
%DIFDELCMD < every cell & cell & 114 & 12 & 228 & 23 & 0.004 \\
%DIFDELCMD < answerable cells only & cell & 100 & 11 & 141 & 22 & -0.046 \\
%DIFDELCMD < answerability-matched pairs & pair & 188 & 14 & 188 & 17 & -0.016 \\
%DIFDELCMD < \bottomrule
%DIFDELCMD < \end{tabular}}
%DIFDELCMD < \end{table}
\begin{table}[htbp]
\DIFoldtable{%
{\small\textbf{Table (before revision).} \textbf{Non-clinical attributes of the same radiographs.} Top: owned attribute and clinical cells in the nine-direction family per checkpoint and dataset, naming the owned attributes. Bottom: medians over blocks per attribute of probe AUROC, selectivity, clean-answer AUROC, and cosine to the clinical normals.\par}\medskip
\centering
\setlength{\tabcolsep}{4pt}

\fitwidth{%
\fitwidth{%
\begin{tabular}{llrrccc}
\toprule
Checkpoint & Dataset & \multicolumn{1}{c}{attributes owned} & \multicolumn{1}{c}{findings owned} & \multicolumn{1}{c}{view AP} & \multicolumn{1}{c}{sex F} & \multicolumn{1}{c}{age $\geq60$}\\
\midrule
Qwen2.5-VL-3B & NIH ChestX-ray14 & 0/3 & 0/6 & -- & -- & -- \\
Qwen2.5-VL-3B & CheXpert Plus & 0/3 & 0/6 & -- & -- & -- \\
Qwen2.5-VL-7B & NIH ChestX-ray14 & 0/3 & 0/6 & -- & -- & -- \\
Qwen2.5-VL-7B & CheXpert Plus & 0/3 & 0/6 & -- & -- & -- \\
Qwen2.5-VL-32B & NIH ChestX-ray14 & 0/3 & 0/6 & -- & -- & -- \\
Qwen2.5-VL-32B & CheXpert Plus & 0/3 & 0/6 & -- & -- & -- \\
Qwen3-VL-4B & NIH ChestX-ray14 & 1/3 & 0/6 & -- & yes & -- \\
Qwen3-VL-4B & CheXpert Plus & 1/3 & 2/6 & -- & yes & -- \\
Qwen3-VL-8B & NIH ChestX-ray14 & 0/3 & 0/6 & -- & -- & -- \\
Qwen3-VL-8B & CheXpert Plus & 1/3 & 0/6 & -- & -- & yes \\
Qwen3-VL-32B & NIH ChestX-ray14 & 0/3 & 1/6 & -- & -- & -- \\
Qwen3-VL-32B & CheXpert Plus & 1/3 & 1/6 & -- & -- & yes \\
InternVL3.5-8B & NIH ChestX-ray14 & 0/3 & 0/6 & -- & -- & -- \\
InternVL3.5-8B & CheXpert Plus & 0/3 & 1/6 & -- & -- & -- \\
InternVL3.5-14B & NIH ChestX-ray14 & 0/3 & 1/6 & -- & -- & -- \\
InternVL3.5-14B & CheXpert Plus & 0/3 & 1/6 & -- & -- & -- \\
InternVL3.5-38B & NIH ChestX-ray14 & 0/3 & 0/6 & -- & -- & -- \\
InternVL3.5-38B & CheXpert Plus & 0/3 & 2/6 & -- & -- & -- \\
Gemma 3 4B & NIH ChestX-ray14 & 0/3 & 0/6 & -- & -- & -- \\
Gemma 3 4B & CheXpert Plus & 0/3 & 0/6 & -- & -- & -- \\
Gemma 3 12B & NIH ChestX-ray14 & 0/3 & 1/6 & -- & -- & -- \\
Gemma 3 12B & CheXpert Plus & 0/3 & 1/6 & -- & -- & -- \\
Gemma 3 27B & NIH ChestX-ray14 & 0/3 & 0/6 & -- & -- & -- \\
Gemma 3 27B & CheXpert Plus & 0/3 & 0/6 & -- & -- & -- \\
MedGemma 4B & NIH ChestX-ray14 & 0/3 & 0/6 & -- & -- & -- \\
MedGemma 4B & CheXpert Plus & 1/3 & 1/6 & -- & -- & yes \\
MedGemma 27B & NIH ChestX-ray14 & 1/3 & 0/6 & -- & -- & yes \\
MedGemma 27B & CheXpert Plus & 0/3 & 0/6 & -- & -- & -- \\
Lingshu 7B & NIH ChestX-ray14 & 2/3 & 1/6 & -- & yes & yes \\
Lingshu 7B & CheXpert Plus & 0/3 & 3/6 & -- & -- & -- \\
Lingshu 32B & NIH ChestX-ray14 & 2/3 & 2/6 & -- & yes & yes \\
Lingshu 32B & CheXpert Plus & 1/3 & 4/6 & -- & -- & yes \\
LLaVA-1.5-7B & NIH ChestX-ray14 & 0/3 & 0/6 & -- & -- & -- \\
LLaVA-1.5-7B & CheXpert Plus & 0/3 & 0/6 & -- & -- & -- \\
LLaVA-1.5-13B & NIH ChestX-ray14 & 1/3 & 0/6 & yes & -- & -- \\
LLaVA-1.5-13B & CheXpert Plus & 0/3 & 0/6 & -- & -- & -- \\
LLaVA-Med 1.5 7B & NIH ChestX-ray14 & 0/3 & 0/6 & -- & -- & -- \\
LLaVA-Med 1.5 7B & CheXpert Plus & 0/3 & 1/6 & -- & -- & -- \\
\midrule
NIH ChestX-ray14 & 19 blocks & 7/57 & 6/114 & 1/19 & 3/19 & 3/19 \\
CheXpert Plus & 19 blocks & 5/57 & 17/114 & 0/19 & 1/19 & 4/19 \\
\textit{chest} & 38 blocks & 12/114 & 23/228 & 1/38 & 4/38 & 7/38 \\
\midrule
\multicolumn{7}{l}{\textit{Attribute probes and directions: medians over the blocks of each dataset}} \\
Attribute & Dataset & probe AUROC & selectivity & answer AUROC & med.\ $|\cos|$ & owned \\
\midrule
AP (portable) projection & NIH ChestX-ray14 & 0.999 & 0.32 & 0.76 & 0.03 & 1/19 \\
AP (portable) projection & CheXpert Plus & 0.998 & 0.31 & 0.73 & 0.05 & 0/19 \\
AP (portable) projection & \textit{chest} & 0.999 & 0.31 & 0.74 & 0.04 & 1/38 \\
female sex & NIH ChestX-ray14 & 0.984 & 0.30 & 0.73 & 0.05 & 3/19 \\
female sex & CheXpert Plus & 0.957 & 0.27 & 0.65 & 0.06 & 1/19 \\
female sex & \textit{chest} & 0.971 & 0.28 & 0.71 & 0.05 & 4/38 \\
age $\geq60$ & NIH ChestX-ray14 & 0.896 & 0.21 & 0.62 & 0.05 & 3/19 \\
age $\geq60$ & CheXpert Plus & 0.898 & 0.21 & 0.64 & 0.06 & 4/19 \\
age $\geq60$ & \textit{chest} & 0.896 & 0.21 & 0.63 & 0.06 & 7/38 \\
\midrule
\multicolumn{7}{l}{\textit{The same comparison read three ways, pooled over the 38 chest blocks}} \\
Comparison & unit & attribute cells & owned & finding cells & owned & difference \\
\midrule
every cell & cell & 114 & 12 & 228 & 23 & 0.004 \\
answerable cells only & cell & 100 & 11 & 141 & 22 & -0.046 \\
answerability-matched pairs & pair & 188 & 14 & 188 & 17 & -0.016 \\
\bottomrule
\end{tabular}}}
}
\end{table}
\begin{table}[htbp]
\DIFdeletedtable{%
{\small\textbf{Table (deleted).} \textbf{Attributes and findings by clean-answer AUROC.} The nine-direction family's cells binned by the clean answer's AUROC against the question's own label; each entry is owned cells of cells in the bin.\par}\medskip
\centering
\setlength{\tabcolsep}{4pt}

\fitwidth{%
\fitwidth{%
\begin{tabular}{lrrrrrr}
\toprule
Cells & \multicolumn{1}{c}{$<$ 0.50} & \multicolumn{1}{c}{0.50--0.70} & \multicolumn{1}{c}{0.70--0.80} & \multicolumn{1}{c}{0.80--0.90} & \multicolumn{1}{c}{$\geq$ 0.90} & \multicolumn{1}{c}{all cells}\\
\midrule
attributes, NIH ChestX-ray14 & 1/4 & 1/26 & 3/14 & 2/6 & 0/7 & 7/57 \\
attributes, CheXpert Plus & 0/5 & 3/29 & 0/12 & 2/5 & 0/6 & 5/57 \\
\textit{attributes, chest} & 1/9 & 4/55 & 3/26 & 4/11 & 0/13 & 12/114 \\
findings, NIH ChestX-ray14 & 0/6 & 2/63 & 2/32 & 2/13 & 0/0 & 6/114 \\
findings, CheXpert Plus & 0/13 & 4/47 & 2/9 & 5/21 & 6/24 & 17/114 \\
\textit{findings, chest} & 0/19 & 6/110 & 4/41 & 7/34 & 6/24 & 23/228 \\
\bottomrule
\end{tabular}}}
}
\end{table}
%DIFDELCMD < \begin{table}[h]
%DIFDELCMD < \centering
%DIFDELCMD < \setlength{\tabcolsep}{4pt}
%DIFDELCMD < %%%
%DIFDELCMD < \caption{%
{%DIFAUXCMD
\textbf{\DIFdelFL{The answer direction and column selectivity.}} %DIFAUXCMD
\DIFdelFL{Top: per block, medians over the six questions of the AUROCs, the Pearson correlation, and the raw and whitened cosines between $a_q$ and $\hat w_q$. Bottom: column selectivity $\kappa_d=W_{d,d}/\sum_q|W_{q,d}|$ per dataset, and ownership on the rows whose label for $q$ is known.}}
%DIFAUXCMD
%DIFDELCMD < \label{tab:cf-validation}
%DIFDELCMD < \fitwidth{%
%DIFDELCMD < \begin{tabular}{llrrrrrrr}
%DIFDELCMD < \toprule
%DIFDELCMD < \multirow{2}{*}{Checkpoint} & \multirow{2}{*}{Dataset} & \multicolumn{2}{c}{AUROC, label} & \multicolumn{2}{c}{AUROC, clean answer} & \multirow{2}{*}{Pearson} & \multicolumn{2}{c}{cosine}\\
%DIFDELCMD < \cmidrule(lr){3-4}\cmidrule(lr){5-6}\cmidrule(lr){8-9}
%DIFDELCMD < &  & \multicolumn{1}{c}{$a_q$} & \multicolumn{1}{c}{$\hat w_q$} & \multicolumn{1}{c}{$a_q$} & \multicolumn{1}{c}{$\hat w_q$} &  & \multicolumn{1}{c}{raw} & \multicolumn{1}{c}{$\Sigma$}\\
%DIFDELCMD < \midrule
%DIFDELCMD < Qwen2.5-VL-7B & NIH ChestX-ray14 & 0.60 & 0.73 & 0.93 & 0.66 & 0.36 & 0.07 & 0.35 \\
%DIFDELCMD < Qwen2.5-VL-7B & CheXpert Plus & 0.67 & 0.86 & 0.85 & 0.66 & 0.34 & 0.06 & 0.36 \\
%DIFDELCMD < Qwen2.5-VL-7B & COCO & 0.97 & 0.97 & 0.98 & 0.98 & 0.84 & 0.64 & 0.84 \\
%DIFDELCMD < Lingshu 32B & NIH ChestX-ray14 & 0.75 & 0.76 & 0.95 & 0.79 & 0.62 & 0.30 & 0.60 \\
%DIFDELCMD < Lingshu 32B & CheXpert Plus & 0.85 & 0.86 & 0.95 & 0.89 & 0.88 & 0.33 & 0.87 \\
%DIFDELCMD < Lingshu 32B & COCO & 0.94 & 0.97 & 0.93 & 0.92 & 0.75 & 0.54 & 0.73 \\
%DIFDELCMD < Gemma 3 12B & NIH ChestX-ray14 & 0.59 & 0.65 & 0.72 & 0.61 & 0.31 & 0.18 & 0.33 \\
%DIFDELCMD < Gemma 3 12B & CheXpert Plus & 0.61 & 0.78 & 0.73 & 0.56 & 0.13 & 0.17 & 0.10 \\
%DIFDELCMD < Gemma 3 12B & COCO & 0.93 & 0.96 & 0.80 & 0.76 & 0.81 & 0.59 & 0.79 \\
%DIFDELCMD < \midrule
%DIFDELCMD < \multicolumn{9}{p{1.314\linewidth}}{\textit{Alignment against ownership over the 54 answer-direction cells: Spearman correlation between $\cos_\Sigma(a_q,\hat w_q)$ and the label direction's $O_q$ $\rho=0.75$ ($p=8\times10^{-11}$); with the label direction's owned indicator $\rho=0.74$ ($p=1\times10^{-10}$); median $\cos_\Sigma$ 0.80 in owned against 0.34 in not-owned cells; within the chest sets alone $\rho=0.59$ ($p=2\times10^{-4}$, 36 cells).}} \\
%DIFDELCMD < \midrule
%DIFDELCMD < \multicolumn{9}{l}{\textit{Column selectivity and known-label ownership, per dataset}} \\
%DIFDELCMD < Dataset & blocks & cells & med.\ $\kappa_d$ & $\kappa_d\geq0.5$ & known rows & sign kept & owned kept & \\
%DIFDELCMD < \midrule
%DIFDELCMD < NIH ChestX-ray14 & 25 & 150 & 0.12 & 2/150 & 600 & 150/150 & 13/13 & \\
%DIFDELCMD < CheXpert Plus & 25 & 150 & 0.12 & 2/150 & 165 & 143/150 & 23/25 & \\
%DIFDELCMD < COCO & 25 & 150 & 0.57 & 92/150 & 600 & 150/150 & 113/113 & \\
%DIFDELCMD < \bottomrule
%DIFDELCMD < \end{tabular}}
%DIFDELCMD < \end{table}
\begin{table}[htbp]
\DIFoldtable{%
{\small\textbf{Table (before revision).} \textbf{The answer direction and column selectivity.} Top: per block, medians over the six questions of the AUROCs, the Pearson correlation, and the raw and whitened cosines between $a_q$ and $\hat w_q$. Bottom: column selectivity $\kappa_d=W_{d,d}/\sum_q|W_{q,d}|$ per dataset, and ownership on the rows whose label for $q$ is known.\par}\medskip
\centering
\setlength{\tabcolsep}{4pt}

\fitwidth{%
\fitwidth{%
\begin{tabular}{llrrrrrrr}
\toprule
\multirow{2}{*}{Checkpoint} & \multirow{2}{*}{Dataset} & \multicolumn{2}{c}{AUROC, label} & \multicolumn{2}{c}{AUROC, clean answer} & \multirow{2}{*}{Pearson} & \multicolumn{2}{c}{cosine}\\
\cmidrule(lr){3-4}\cmidrule(lr){5-6}\cmidrule(lr){8-9}
&  & \multicolumn{1}{c}{$a_q$} & \multicolumn{1}{c}{$\hat w_q$} & \multicolumn{1}{c}{$a_q$} & \multicolumn{1}{c}{$\hat w_q$} &  & \multicolumn{1}{c}{raw} & \multicolumn{1}{c}{$\Sigma$}\\
\midrule
Qwen2.5-VL-7B & NIH ChestX-ray14 & 0.60 & 0.73 & 0.93 & 0.66 & 0.36 & 0.07 & 0.35 \\
Qwen2.5-VL-7B & CheXpert Plus & 0.67 & 0.86 & 0.85 & 0.66 & 0.34 & 0.06 & 0.36 \\
Qwen2.5-VL-7B & COCO & 0.97 & 0.97 & 0.98 & 0.98 & 0.84 & 0.64 & 0.84 \\
Lingshu 32B & NIH ChestX-ray14 & 0.75 & 0.76 & 0.95 & 0.79 & 0.62 & 0.30 & 0.60 \\
Lingshu 32B & CheXpert Plus & 0.85 & 0.86 & 0.95 & 0.89 & 0.88 & 0.33 & 0.87 \\
Lingshu 32B & COCO & 0.94 & 0.97 & 0.93 & 0.92 & 0.75 & 0.54 & 0.73 \\
Gemma 3 12B & NIH ChestX-ray14 & 0.59 & 0.65 & 0.72 & 0.61 & 0.31 & 0.18 & 0.33 \\
Gemma 3 12B & CheXpert Plus & 0.61 & 0.78 & 0.73 & 0.56 & 0.13 & 0.17 & 0.10 \\
Gemma 3 12B & COCO & 0.93 & 0.96 & 0.80 & 0.76 & 0.81 & 0.59 & 0.79 \\
\midrule
\multicolumn{9}{p{1.314\linewidth}}{\textit{Alignment against ownership over the 54 answer-direction cells: Spearman correlation between $\cos_\Sigma(a_q,\hat w_q)$ and the label direction's $O_q$ $\rho=0.75$ ($p=8\times10^{-11}$); with the label direction's owned indicator $\rho=0.74$ ($p=1\times10^{-10}$); median $\cos_\Sigma$ 0.80 in owned against 0.34 in not-owned cells; within the chest sets alone $\rho=0.59$ ($p=2\times10^{-4}$, 36 cells).}} \\
\midrule
\multicolumn{9}{l}{\textit{Column selectivity and known-label ownership, per dataset}} \\
Dataset & blocks & cells & med.\ $\kappa_d$ & $\kappa_d\geq0.5$ & known rows & sign kept & owned kept & \\
\midrule
NIH ChestX-ray14 & 25 & 150 & 0.12 & 2/150 & 600 & 150/150 & 13/13 & \\
CheXpert Plus & 25 & 150 & 0.12 & 2/150 & 165 & 143/150 & 23/25 & \\
COCO & 25 & 150 & 0.57 & 92/150 & 600 & 150/150 & 113/113 & \\
\bottomrule
\end{tabular}}}
}
\end{table}
\begin{table}[htbp]
\DIFdeletedtable{%
{\small\textbf{Table (deleted).} \textbf{Radiologist labels.} Top: the complete grade per CheXpert block on 200 radiologist-labelled films, against the same grade on the 600 labeler-labelled test rows. Bottom: the six directions refitted on the radiologist labels.\par}\medskip
\centering
\setlength{\tabcolsep}{4pt}

\fitwidth{%
\fitwidth{%
\begin{tabular}{lrrrrr}
\toprule
Checkpoint & \multicolumn{1}{c}{owned (test)} & \multicolumn{1}{c}{owned (valid)} & \multicolumn{1}{c}{verdict agrees} & \multicolumn{1}{c}{owned grade agrees} & \multicolumn{1}{c}{median $|\Delta O_q|$}\\
\midrule
Qwen2.5-VL-3B & 0 & 0 & 5/6 & 6/6 & 0.005 \\
Qwen2.5-VL-7B & 0 & 0 & 6/6 & 6/6 & 0.005 \\
Qwen2.5-VL-32B & 0 & 0 & 6/6 & 6/6 & 0.004 \\
Qwen3-VL-4B & 2 & 2 & 4/6 & 6/6 & 0.009 \\
Qwen3-VL-8B & 0 & 0 & 5/6 & 6/6 & 0.003 \\
Qwen3-VL-32B & 1 & 1 & 6/6 & 6/6 & 0.004 \\
InternVL3.5-8B & 1 & 0 & 5/6 & 5/6 & 0.001 \\
InternVL3.5-14B & 1 & 1 & 3/6 & 6/6 & 0.002 \\
InternVL3.5-38B & 2 & 2 & 6/6 & 6/6 & 0.001 \\
Gemma 3 4B & 0 & 0 & 4/6 & 6/6 & 0.007 \\
Gemma 3 12B & 1 & 1 & 6/6 & 6/6 & 0.016 \\
Gemma 3 27B & 0 & 0 & 3/6 & 6/6 & 0.002 \\
MedGemma 4B & 2 & 1 & 5/6 & 5/6 & 0.025 \\
MedGemma 27B & 0 & 0 & 5/6 & 6/6 & 0.010 \\
Lingshu 7B & 3 & 3 & 6/6 & 6/6 & 0.007 \\
Lingshu 32B & 5 & 5 & 6/6 & 6/6 & 0.005 \\
LLaVA-1.5-7B & 0 & 0 & 5/6 & 6/6 & 0.003 \\
LLaVA-1.5-13B & 0 & 0 & 5/6 & 6/6 & 0.000 \\
LLaVA-Med 1.5 7B & 1 & 1 & 6/6 & 6/6 & 0.000 \\
\midrule
\textit{all} (19 blocks) & 19 & 17 & 97/114 & 112/114 & 0.003 \\
\bottomrule
\end{tabular}}}

\medskip
\setlength{\tabcolsep}{4pt}
\fitwidth{%
\fitwidth{%
\begin{tabular}{llrrrrr}
\toprule
\multicolumn{7}{p{1.357\linewidth}}{\textit{The label source of a refitted direction against the number of rows it was fitted on. The expert-label refit is estimated on 200 radiologist-labelled films and the shipped direction on about 20{,}000 report-labelled rows, so comparing only those two confounds the two. Every arm keeps the block's own projection, its stored train-only scaler, the protocol's probe settings and the same 5-fold cross-fitting over patients, and is scored on the same films by a direction not fitted on them. The two training subsample arms redraw the fitting rows 20 times and report the median; the second draws enough rows to leave the concept with as many labelled rows as the expert refit has. Cosines are to the shipped direction. Medians over the 36 cells of 6 CheXpert blocks.}} \\
\midrule
 & & & \multicolumn{2}{c}{cross-fitted AUROC against} & \multicolumn{2}{c}{cosine to shipped} \\
\cmidrule(lr){4-5}\cmidrule(lr){6-7}
Labels & Fitted on & rows & radiologist & report & raw & whitened \\
\midrule
radiologist & 200 valid films & 160 & 0.69 & 0.71 & 0.12 & 0.56 \\
report-derived & the same 200 valid films & 44 & 0.67 & 0.71 & 0.12 & 0.59 \\
report-derived & 200 training rows & 42 & 0.66 & 0.70 & 0.12 & 0.58 \\
report-derived & 200 labelled training rows & 160 & 0.70 & 0.74 & 0.22 & 0.69 \\
report-derived & every training row & 5332 & 0.79 & 0.86 & 1.00 & 1.00 \\
\midrule
\multicolumn{7}{l}{\textit{What each contrast isolates, in AUROC against the radiologist labels}} \\
\midrule
\multicolumn{6}{l}{label source, both arms fitting the same number of rows} & -0.01 \\
\multicolumn{6}{l}{sample size, both arms on report-derived labels} & 0.09 \\
\multicolumn{6}{l}{the expert refit against the shipped direction, which mixes the two} & -0.11 \\
\bottomrule
\end{tabular}}}
}
\end{table}
%DIFDELCMD < 

%DIFDELCMD < %%%
\subsection{\DIFdel{Further results}}
%DIFAUXCMD
\addtocounter{subsection}{-1}%DIFAUXCMD
%DIFDELCMD < \label{app:cf-further}
%DIFDELCMD < 

%DIFDELCMD < %%%
\DIFdel{This subsection reports the detail behind Sections~\ref{sec:headline}--\ref{sec:controls}.
}%DIFDELCMD < 

%DIFDELCMD < %%%
\paragraph{\DIFdel{Reading and answering.}} %DIFAUXCMD
\addtocounter{paragraph}{-1}%DIFAUXCMD
\DIFdel{Median probe selectivity over the chest cells is $\cfSelChestMed$, with quartiles $\cfSelChestQOne$ and $\cfSelChestQThree$; of the }%DIFDELCMD < \cfChestAnswerable{} %%%
\DIFdel{answerable chest cells, }%DIFDELCMD < \cfReadAns{} %%%
\DIFdel{are also readable. Effusion and Cardiomegaly are the most readable of the four concepts both chest sets share, with mean selectivity $\cfSelEffusion$ and $\cfSelCardiomegaly$ over }%DIFDELCMD < \cfChestProbeBlocks{} %%%
\DIFdel{chest blocks. Table~\ref{tab:cf-contingency} cross-classifies every chest and COCO cell by the steering reference and the verdict.
}%DIFDELCMD < 

%DIFDELCMD < %%%
\paragraph{\DIFdel{Natural objects.}} %DIFAUXCMD
\addtocounter{paragraph}{-1}%DIFAUXCMD
\DIFdel{Ownership margins on COCO span $\cfCocoMarginMin$--$\cfCocoMarginMax$ in the Qwen, Lingshu, and Gemma families. The dose sweep ran on }%DIFDELCMD < \cfDoseBlocksNih{} %%%
\DIFdel{NIH and }%DIFDELCMD < \cfDoseBlocksCoco{} %%%
\DIFdel{COCO blocks. Within a checkpoint, the counts of owned fine-grained and easy objects differ by $\cfFgSameDiffMin$ to $\cfFgSameDiffMax$ of six, and median probe selectivity is }%DIFDELCMD < \cfFgFineSel{} %%%
\DIFdel{for the fine-grained objects against }%DIFDELCMD < \cfFgEasySameSel{} %%%
\DIFdel{for the easy ones. Table~\ref{tab:cf-fgobj} gives the remaining difficulty bins, the per-dataset counts, and the matched estimator.
}%DIFDELCMD < 

%DIFDELCMD < %%%
\DIFdelend \paragraph{Shared-tower pairs.} \DIFdelbegin \DIFdel{Table~\ref{tab:cf-pairs}records the bitwise relation of every pair. The median paired difference in ownership is }%DIFDELCMD < \cfPairsRatioMin{}%%%
\DIFdel{--}%DIFDELCMD < \cfPairsRatioMax{} %%%
\DIFdel{times the refit standard deviation, up to }%DIFDELCMD < \cfPairsRatioSingleMax{} %%%
\DIFdel{times for single concepts, and it exceeds that standard deviation in }%DIFDELCMD < \cfPairsRatioAboveOne{} %%%
\DIFdel{of }%DIFDELCMD < \cfPairsRatioN{} %%%
\DIFdel{pairs. }\DIFdelend \DIFaddbegin \DIFadd{Gemma~3 4B and 12B write bit-identical vectors, and the other Gemma~3 and MedGemma pairs agree to bf16 precision (Table~\ref{tab:cf-pairs}). }\DIFaddend On NIH Effusion, the identical write \DIFdelbegin \DIFdel{yields an $O_q$ difference of $\cfPairDeltaEffTwelveTwentySeven$ $[\cfPairDeltaEffTwelveTwentySevenLo,\cfPairDeltaEffTwelveTwentySevenHi]$ between }\DIFdelend \DIFaddbegin \DIFadd{gives }\DIFaddend Gemma~3 12B and 27B \DIFaddbegin \DIFadd{an $O_q$ difference of $\cfPairDeltaEffTwelveTwentySeven$ $[\cfPairDeltaEffTwelveTwentySevenLo,\cfPairDeltaEffTwelveTwentySevenHi]$}\DIFaddend . In 12B, competing directions win the chest answers outright ($O_{\mathrm{Effusion}}=\cfOGemmaTwelveEff$ on both chest sets, $O_{\mathrm{Mass}}=\cfOGemmaTwelveMass$ on NIH, $O_{\mathrm{Pneumothorax}}=\cfOGemmaTwelvePneu$ on CheXpert); in 27B, no direction moves the same \DIFdelbegin \DIFdel{three }\DIFdelend answers ($\cfOGemmaTwentySevenEff$, $\cfOGemmaTwentySevenMass$, $\cfOGemmaTwentySevenPneu$). \DIFdelbegin \DIFdel{In }\DIFdelend \DIFaddbegin \DIFadd{Gemma~3 }\DIFaddend 4B \DIFdelbegin \DIFdel{, the clean answer is }\DIFdelend \DIFaddbegin \DIFadd{answers }\DIFaddend yes on \cfGemmaFourYesMin--\cfGemmaFourYesMax\% of chest rows, so many \DIFaddbegin \DIFadd{of its }\DIFaddend cells sit at ceiling\DIFdelbegin \DIFdel{(clean answer above $0.99$ or below $0.01$ on more than 90\% of rows) and $P(\mathrm{yes})$ has little room to move. The logit margin $\ell_{\mathrm{yes}}-\ell_{\mathrm{no}}$ does not saturate, and ownership recomputed with margin changes in place of probability changes, }\DIFdelend \DIFaddbegin \DIFadd{; on the margin scale its }\DIFaddend $O^m_q$ \DIFdelbegin \DIFdel{, }\DIFdelend runs from $\cfPairsGemmaFourOmNear$ to $\cfPairsGemmaFourOmFar$ with every interval below zero\DIFdelbegin \DIFdel{: }\DIFdelend \DIFaddbegin \DIFadd{, so }\DIFaddend every own write lowers its answer's margin \DIFdelbegin \DIFdel{by }\DIFdelend more than any \DIFdelbegin \DIFdel{competing write }\DIFdelend \DIFaddbegin \DIFadd{competitor }\DIFaddend raises it. MedGemma 4B and 27B differ in the same way (\DIFaddbegin \DIFadd{CheXpert }\DIFaddend Pneumothorax $\cfOMedFourPneu$ \DIFdelbegin \DIFdel{versus }\DIFdelend \DIFaddbegin \DIFadd{against }\DIFaddend $\cfOMedTwentySevenPneu$\DIFdelbegin \DIFdel{and }\DIFdelend \DIFaddbegin \DIFadd{, }\DIFaddend Edema $\cfOMedFourEdema$ \DIFdelbegin \DIFdel{versus }\DIFdelend \DIFaddbegin \DIFadd{against }\DIFaddend $\cfOMedTwentySevenEdema$\DIFdelbegin \DIFdel{on CheXpert).
The reader effect spans the whole matrix, including cells whose verdict stays the same. In the crossed swap of the Gemma~3 4B and MedGemma 4B towers and readers (Appendix~\ref{app:cf-crossed}), changing the reader at a fixed tower changes $O_q$ on NIH by a mean absolute }%DIFDELCMD < \cfSwapReaderNihOMin{}%%%
\DIFdel{, against }%DIFDELCMD < \cfSwapTowerNihOMin{} %%%
\DIFdel{for changing the tower; the same appendix replays identical stored tensors to every reader of a shared tower.
}\DIFdelend \DIFaddbegin \DIFadd{).
}\DIFaddend 

\DIFaddbegin \begin{table}[h]
\centering
\setlength{\tabcolsep}{4pt}
\caption{\textbf{\DIFaddFL{Same write, different reader: paired differences.}} \DIFaddFL{$\Delta O_q=O_q(\text{larger})-O_q(\text{smaller})$ per shared-tower pair on the same 600 test rows, with paired patient-bootstrap 95\% intervals; the last two columns divide $|\Delta O_q|$ by the two blocks' mean refit standard deviation.}}
\label{tab:cf-pairs}
\resizebox{\linewidth}{!}{%
\begin{tabular}{llcrrrrr}
\toprule
\multirow{2}{*}{Pair} & \multirow{2}{*}{Dataset} & \multirow{2}{*}{relation} & \multicolumn{1}{c}{$\Delta O_q$ CI} & \multicolumn{2}{c}{$|\Delta O_q|$} & \multicolumn{2}{c}{ratio to refit SD}\\
\cmidrule(lr){5-6}\cmidrule(lr){7-8}
&  &  & \multicolumn{1}{c}{$\not\ni 0$} & \multicolumn{1}{c}{median} & \multicolumn{1}{c}{max} & \multicolumn{1}{c}{median} & \multicolumn{1}{c}{max}\\
\midrule
Gemma 3 4B $\to$ Gemma 3 12B & NIH ChestX-ray14 & bit-identical & 5/6 & 0.33 & 0.83 & 2.7 & 6.9 \\
Gemma 3 4B $\to$ Gemma 3 12B & CheXpert Plus & bit-identical & 6/6 & 0.29 & 0.83 & 2.1 & 5.9 \\
Gemma 3 4B $\to$ Gemma 3 12B & COCO & bit-identical & 5/6 & 0.23 & 0.87 & 1.6 & 6.0 \\
Gemma 3 12B $\to$ Gemma 3 27B & NIH ChestX-ray14 & bf16-equal & 6/6 & 0.37 & 0.79 & 3.7 & 8.0 \\
Gemma 3 12B $\to$ Gemma 3 27B & CheXpert Plus & bf16-equal & 6/6 & 0.35 & 0.81 & 1.4 & 3.3 \\
Gemma 3 12B $\to$ Gemma 3 27B & COCO & bf16-equal & 5/6 & 0.15 & 0.40 & 1.7 & 4.3 \\
Gemma 3 4B $\to$ Gemma 3 27B & NIH ChestX-ray14 & bf16-equal & 5/6 & 0.07 & 0.24 & 1.8 & 6.0 \\
Gemma 3 4B $\to$ Gemma 3 27B & CheXpert Plus & bf16-equal & 6/6 & 0.17 & 0.65 & 4.6 & 17.3 \\
Gemma 3 4B $\to$ Gemma 3 27B & COCO & bf16-equal & 6/6 & 0.16 & 1.02 & 1.4 & 8.8 \\
MedGemma 4B $\to$ MedGemma 27B & NIH ChestX-ray14 & bf16-equal & 5/6 & 0.09 & 0.30 & 1.5 & 5.0 \\
MedGemma 4B $\to$ MedGemma 27B & CheXpert Plus & bf16-equal & 5/6 & 0.19 & 0.54 & 1.8 & 5.0 \\
MedGemma 4B $\to$ MedGemma 27B & COCO & bf16-equal & 3/6 & 0.05 & 0.35 & 0.6 & 4.6 \\
\bottomrule
\end{tabular}}
\end{table}

\DIFaddend \paragraph{Size and family.} \DIFdelbegin \DIFdel{Ownership on the chest sets is concentrated among }\DIFdelend \DIFaddbegin \DIFadd{Chest ownership concentrates in }\DIFaddend a few readers and varies by concept and dataset (Figure~\ref{fig:ladders}; Table~\ref{tab:cf-models}). Qwen2.5-VL owns no concept\DIFdelbegin \DIFdel{s}\DIFdelend  at 3B or 7B\DIFdelbegin \DIFdel{. At }\DIFdelend \DIFaddbegin \DIFadd{, Cardiomegaly at }\DIFaddend 32B, \DIFdelbegin \DIFdel{it owns Cardiomegaly ; at 72B, }\DIFdelend \DIFaddbegin \DIFadd{and }\DIFaddend Atelectasis ($O_{\mathrm{Atelectasis}}=\cfOQwenSeventyTwoAtel$), Cardiomegaly, and Mass \DIFdelbegin \DIFdel{. }\DIFdelend \DIFaddbegin \DIFadd{at 72B; }\DIFaddend Effusion has negative ownership at every size. Lingshu, the medical Qwen2.5-VL derivative, owns three CheXpert concepts at 7B and five at 32B, the \DIFdelbegin \DIFdel{highest clinical ownership }\DIFdelend \DIFaddbegin \DIFadd{most }\DIFaddend in the grid. LLaVA-1.5 reads the chest concepts but answers none above chance and owns none.

\paragraph{\DIFdelbegin \DIFdel{Medical and chest-specialised checkpoints}\DIFdelend \DIFaddbegin \DIFadd{Chest-trained towers}\DIFaddend .} \DIFdelbegin \DIFdel{We divide the grid into three groups: checkpoints whose towers are trained on chest radiographs (}\DIFdelend \DIFaddbegin \DIFadd{The towers of }\DIFaddend CheXagent-2-3B, CheXagent-8B, \DIFaddbegin \DIFadd{and }\DIFaddend LLaVA-Rad-7B \DIFdelbegin \DIFdel{), medically tuned checkpoints on general towers, and general checkpoints}\DIFdelend \DIFaddbegin \DIFadd{are trained on chest radiographs}\DIFaddend ; HuatuoGPT-Vision-7B is tuned on medical data over a stock CLIP tower. The chest-trained towers \DIFdelbegin \DIFdel{' }\DIFdelend \DIFaddbegin \DIFadd{reach a }\DIFaddend median probe AUROC \DIFaddbegin \DIFadd{of }\cfSpecTowerProbeAuroc{} \DIFaddend across \cfSpecTowerProbeN{} NIH cells\DIFdelbegin \DIFdel{is }%DIFDELCMD < \cfSpecTowerProbeAuroc{}%%%
\DIFdelend , against \cfSpecStockProbeAuroc{} for HuatuoGPT-Vision \DIFdelbegin \DIFdel{'s stock CLIP tower }\DIFdelend and \cfSpecRestProbeAuroc{} \DIFdelbegin \DIFdel{across the }%DIFDELCMD < \cfSpecRestProbeN{} %%%
\DIFdel{remaining NIH cells. Across their }%DIFDELCMD < \cfSpecTowerN{} %%%
\DIFdel{chest cells, }\DIFdelend \DIFaddbegin \DIFadd{for the remaining }\cfSpecRestProbeN{}\DIFadd{, and }\DIFaddend they read \cfSpecTowerReadable{} and answer \cfSpecTowerAnswerable{} \DIFaddbegin \DIFadd{of their }\cfSpecTowerN{} \DIFadd{chest cells}\DIFaddend . On NIH, \DIFdelbegin \DIFdel{the three groups own }\DIFdelend \cfSpecNihTowerOwned{} of \cfSpecNihTowerN{} \DIFaddbegin \DIFadd{cells are owned for chest-trained towers}\DIFaddend , \cfSpecNihMedOwned{} of \cfSpecNihMedN{} \DIFaddbegin \DIFadd{for medically tuned checkpoints}\DIFaddend , and \cfSpecNihGenOwned{} of \cfSpecNihGenN{} \DIFdelbegin \DIFdel{cells. The natural-image control of Section~\ref{sec:coco} }\DIFdelend \DIFaddbegin \DIFadd{for general checkpoints. The COCO control }\DIFaddend rests on checkpoints that answer natural-image questions\DIFdelbegin \DIFdel{. The general group answers }\DIFdelend \DIFaddbegin \DIFadd{: of the COCO cells, }\DIFaddend \cfSpecGenCocoAnswerable{} of \DIFdelbegin %DIFDELCMD < \cfSpecGenCocoN{} %%%
\DIFdel{COCO cells at a median }\DIFdelend \DIFaddbegin \cfSpecGenCocoN{} \DIFadd{are answerable and }\cfSpecGenCocoOwned{} \DIFadd{owned for the general group (median }\DIFaddend clean-answer \DIFdelbegin \DIFdel{AUROC of }\DIFdelend \DIFaddbegin \DIFadd{AUROC }\DIFaddend \cfSpecGenCocoAnsAuroc{}\DIFdelbegin \DIFdel{and owns }%DIFDELCMD < \cfSpecGenCocoOwned{}%%%
\DIFdel{; the medically tuned group answers }%DIFDELCMD < \cfSpecMedCocoAnswerable{} %%%
\DIFdel{of }%DIFDELCMD < \cfSpecMedCocoN{} %%%
\DIFdel{at }%DIFDELCMD < \cfSpecMedCocoAnsAuroc{} %%%
\DIFdel{and owns }%DIFDELCMD < \cfSpecMedCocoOwned{}%%%
\DIFdel{. The }\DIFdelend \DIFaddbegin \DIFadd{), }\cfSpecMedCocoAnswerable{} \DIFadd{of }\cfSpecMedCocoN{} \DIFadd{and }\cfSpecMedCocoOwned{} \DIFadd{for the medically tuned group (}\cfSpecMedCocoAnsAuroc{}\DIFadd{), and }\cfSpecTowerCocoAnswerable{} \DIFadd{of }\cfSpecTowerCocoN{} \DIFadd{and }\cfSpecTowerCocoOwned{} \DIFadd{for the }\DIFaddend chest-trained \DIFdelbegin \DIFdel{towers answer }%DIFDELCMD < \cfSpecTowerCocoAnswerable{} %%%
\DIFdel{of }%DIFDELCMD < \cfSpecTowerCocoN{} %%%
\DIFdel{COCO cellsat }%DIFDELCMD < \cfSpecTowerCocoAnsAuroc{} %%%
\DIFdel{and own }%DIFDELCMD < \cfSpecTowerCocoOwned{}%%%
\DIFdel{, so their COCO rows give the grade little to assess.
}\DIFdelend \DIFaddbegin \DIFadd{towers (}\cfSpecTowerCocoAnsAuroc{}\DIFadd{).
}\DIFaddend 

\DIFdelbegin \paragraph{\DIFdel{Non-clinical attributes.}} %DIFAUXCMD
\addtocounter{paragraph}{-1}%DIFAUXCMD
\DIFdel{All }%DIFDELCMD < \cfAttrReadable{} %%%
\DIFdel{of }%DIFDELCMD < \cfAttrChestAttrN{} %%%
\DIFdel{attribute cells are readable. Among answerable cells, attributes are owned in }%DIFDELCMD < \cfAttrCapAttrOwned{} %%%
\DIFdel{of }%DIFDELCMD < \cfAttrCapAttrN{} %%%
\DIFdel{and findings in }%DIFDELCMD < \cfAttrCapClinOwned{} %%%
\DIFdel{of }%DIFDELCMD < \cfAttrCapClinN{}%%%
\DIFdel{. Pairing each attribute cell with clinical cells from its own block within }%DIFDELCMD < \cfAttrPairWindow{} %%%
\DIFdel{clean-answer AUROC yields }%DIFDELCMD < \cfAttrPairAttrN{} %%%
\DIFdel{pairs, with attributes owned in }%DIFDELCMD < \cfAttrPairAttrOwned{} %%%
\DIFdel{and findings in }%DIFDELCMD < \cfAttrPairClinOwned{}%%%
\DIFdel{. On NIH, ownership is }%DIFDELCMD < \cfAttrNihAttrOwned{}%%%
\DIFdel{/}%DIFDELCMD < \cfAttrNihAttrN{} %%%
\DIFdel{for attributes and }%DIFDELCMD < \cfAttrNihClinOwned{}%%%
\DIFdel{/}%DIFDELCMD < \cfAttrNihClinN{} %%%
\DIFdel{for findings.
}%DIFDELCMD < 

%DIFDELCMD < %%%
\paragraph{\DIFdel{Answer direction.}} %DIFAUXCMD
\addtocounter{paragraph}{-1}%DIFAUXCMD
\DIFdel{The answer direction's write beats the strongest logistic competitor in }%DIFDELCMD < \cfAnsChestBeats{} %%%
\DIFdel{chest cells. Under the five held-out templates, $a_q$ remains owned in }%DIFDELCMD < \cfAnsdirtCocoKept{} %%%
\DIFdel{of }%DIFDELCMD < \cfAnsdirtCocoKeptN{} %%%
\DIFdel{COCO (cell, template) pairs, and over all pairs it is owned in }%DIFDELCMD < \cfAnsdirtChestOwnedPairs{} %%%
\DIFdel{of }%DIFDELCMD < \cfAnsdirtChestPairsAll{} %%%
\DIFdel{chest and }%DIFDELCMD < \cfAnsdirtCocoOwnedPairs{} %%%
\DIFdel{of }%DIFDELCMD < \cfAnsdirtCocoPairsAll{} %%%
\DIFdel{COCO pairs. In Qwen2.5-VL-7B, it is owned for }%DIFDELCMD < \cfAnsQwenOwnedNih{} %%%
\DIFdel{of the six NIH concepts (for Effusion, its ownership against the other five answer directions is $\cfAnsQwenEffO$, with an own write of $\cfAnsQwenEffW$) and all }%DIFDELCMD < \cfAnsQwenOwnedChex{} %%%
\DIFdel{CheXpert concepts; in Lingshu 32B, for }%DIFDELCMD < \cfAnsLingOwnedNih{} %%%
\DIFdel{NIH and }%DIFDELCMD < \cfAnsLingOwnedChex{} %%%
\DIFdel{CheXpert concepts; in Gemma~3 12B, for }%DIFDELCMD < \cfAnsGemOwnedNih{} %%%
\DIFdel{NIH and }%DIFDELCMD < \cfAnsGemOwnedChex{} %%%
\DIFdel{CheXpert concepts. In Qwen2.5-VL-7B on NIH, the label AUROC is $\cfValAurocAnsQwenNih$ for the answer direction and $\cfValAurocProbeQwenNih$ for the probe, and the clean-answer AUROC is $\cfValAurocAnsCleanQwenNih$ and $\cfValAurocProbeCleanQwenNih$.
}%DIFDELCMD < 

%DIFDELCMD < %%%
\DIFdel{Cosine and projection depend on the inner product chosen for the representation \mbox{%DIFAUXCMD
\citep{park2024linear}}\hskip0pt%DIFAUXCMD
, so each cosine names its metric. The median raw model-space cosine between the answer and label directions spans $\cfAnsCosQwenMin$--$\cfAnsCosQwenMax$ across the two chest blocks of Qwen2.5-VL-7B and $\cfAnsCosLingMin$--$\cfAnsCosLingMax$ across those of Lingshu 32B, compared with $\cfAnsCosCocoMin$--$\cfAnsCosCocoMax$ across the COCO blocks. For two coefficient vectors $\boldsymbol u,\boldsymbol v$ in the projected, train-scaled feature space, the whitened cosine $\cos_\Sigma(\boldsymbol u,\boldsymbol v)=\boldsymbol u^{\top}\Sigma\boldsymbol v/\sqrt{\boldsymbol u^{\top}\Sigma\boldsymbol u\;\boldsymbol v^{\top}\Sigma\boldsymbol v}$, with $\Sigma$ the covariance of the training features, is the correlation, over training images, between the scores that $\boldsymbol u$ and $\boldsymbol v$ assign to the train-scaled features. Between the answer direction's and the probe's coefficients, its median over each dataset's cells is $\cfValCosSigmaChestMin$--$\cfValCosSigmaChestMax$ on the two chest sets and $\cfValCosSigmaCoco$ on COCO. The metrics differ most on the CheXpert block of Lingshu 32B, where the median over its six concepts is $\cfValCosRawLingChex$ raw and $\cfValCosSigmaLingChex$ whitened. The column-selectivity index $\kappa_d=W_{d,d}/\sum_q|W_{q,d}|$ divides a direction's signed effect on its own question by its summed absolute effect on all six questions, so $\kappa_d=1$ when a write moves only its own answer. Across the }%DIFDELCMD < \cfAnsColSelChestN{} %%%
\DIFdel{chest cells, its median is $\cfAnsColSelChest$ for answer directions and $\cfAnsColSelLabelChest$ for label directions in the same blocks, and $\cfAnsColSelCoco$ and $\cfAnsColSelLabelCoco$ over the }%DIFDELCMD < \cfAnsColSelCocoN{} %%%
\DIFdel{COCO cells: on the chest sets, both answer and label directions move the other five answers along with their own. Over all cells, $\kappa_d$ has median }%DIFDELCMD < \cfValColSelNih{} %%%
\DIFdel{on NIH and }%DIFDELCMD < \cfValColSelChex{} %%%
\DIFdel{on CheXpert, versus }%DIFDELCMD < \cfValColSelCoco{} %%%
\DIFdel{on COCO (Table~\ref{tab:cf-validation}).
}%DIFDELCMD < 

%DIFDELCMD < %%%
\paragraph{\DIFdel{Semantic endpoints.}} %DIFAUXCMD
\addtocounter{paragraph}{-1}%DIFAUXCMD
\DIFdel{The endpoint cells come from }%DIFDELCMD < \cfSemBlocks{} %%%
\DIFdel{scored blocks. The negated question moves opposite to the affirmative one in }%DIFDELCMD < \cfSemNegOpposite{} %%%
\DIFdel{of the }%DIFDELCMD < \cfSemNegN{} %%%
\DIFdel{concept--block pairs. The forced choice is owned in both presentation orders in }%DIFDELCMD < \cfSemForcedOwned{} %%%
\DIFdel{of }%DIFDELCMD < \cfSemForcedN{} %%%
\DIFdel{pairs, with a mean absolute gap of }%DIFDELCMD < \cfSemGapMin{}%%%
\DIFdel{--}%DIFDELCMD < \cfSemGapMax{} %%%
\DIFdel{between orders. The report continuation is owned in }%DIFDELCMD < \cfSemReportOwned{} %%%
\DIFdel{of }%DIFDELCMD < \cfSemReportCells{} %%%
\DIFdel{cells.
}%DIFDELCMD < 

%DIFDELCMD < %%%
\paragraph{\DIFdel{Labels and alternative constructions.}} %DIFAUXCMD
\addtocounter{paragraph}{-1}%DIFAUXCMD
\DIFdel{On CheXpert, restricting evaluation to rows with a known label preserves the sign of $O_q$ in }%DIFDELCMD < \cfValKnownSignAgree{} %%%
\DIFdel{of }%DIFDELCMD < \cfValKnownN{} %%%
\DIFdel{cells; NIH and COCO have every test label known, so their known-row numbers reproduce the all-row numbers. Refitting the six directions on the radiologist labels themselves leaves them a median raw model-space cosine of }%DIFDELCMD < \cfValidfitCosRaw{} %%%
\DIFdel{and a median whitened cosine of }%DIFDELCMD < \cfValidfitCosWhitened{} %%%
\DIFdel{from the report-label directions, across the }%DIFDELCMD < \cfValidfitBlocks{} %%%
\DIFdel{CheXpert blocks with the refit. The expert-label refit reads the radiologist labels with a median cross-fitted AUROC of }%DIFDELCMD < \cfValidfitAurocExpert{}%%%
\DIFdel{, versus }%DIFDELCMD < \cfValidfitAurocReportDir{} %%%
\DIFdel{for the report-label direction on the same films, and a report-label direction fitted on as many rows as the refit reaches }%DIFDELCMD < \cfVfArmReportSubAuroc{}%%%
\DIFdel{: the estimation sample moves the result by $\cfVfArmDeltaSize$ and the label source by $\cfVfArmDeltaLabel$. The difference-of-means and pattern directions sit at a median absolute raw model-space cosine of $\cfAltdirCosMin$ to $\cfAltdirCosMax$ from the logistic normal, and each is owned in more NIH cells and fewer COCO cells than the logistic normal. Lifted as displacements, they are owned in }%DIFDELCMD < \cfAltdirdChestDom{} %%%
\DIFdel{and }%DIFDELCMD < \cfAltdirdChestPattern{} %%%
\DIFdel{of }%DIFDELCMD < \cfAltdirdChestN{} %%%
\DIFdel{chest cells and in }%DIFDELCMD < \cfAltdirdCocoDom{} %%%
\DIFdel{and }%DIFDELCMD < \cfAltdirdCocoPattern{} %%%
\DIFdel{of }%DIFDELCMD < \cfAltdirdCocoN{} %%%
\DIFdel{COCO cells, respectively (Table~\ref{tab:cf-altdir}).
Every construction preserves the chest-to-COCO gap: at least one construction is owned in }%DIFDELCMD < \cfAltdirNihAny{} %%%
\DIFdel{of }%DIFDELCMD < \cfAltdirNihN{} %%%
\DIFdel{cells on NIH and }%DIFDELCMD < \cfAltdirChexAny{} %%%
\DIFdel{of }%DIFDELCMD < \cfAltdirChexN{} %%%
\DIFdel{on CheXpert, compared with }%DIFDELCMD < \cfAltdirCocoAny{} %%%
\DIFdel{of }%DIFDELCMD < \cfAltdirCocoN{} %%%
\DIFdel{on COCO. The alternative constructions are scored across }%DIFDELCMD < \cfAltdirBlocks{} %%%
\DIFdel{blocks, the extended competitor family across }%DIFDELCMD < \cfExtcompBlocks{} %%%
\DIFdel{blocks, and the token-weighted writes across }%DIFDELCMD < \cfTokenwChestBlocks{} %%%
\DIFdel{chest blocks; on COCO, the softmax and top-quarter writes are owned in }%DIFDELCMD < \cfTokenwSoftCocoOwned{} %%%
\DIFdel{and }%DIFDELCMD < \cfTokenwTopqCocoOwned{} %%%
\DIFdel{of }%DIFDELCMD < \cfTokenwCocoCells{} %%%
\DIFdel{cells.
}%DIFDELCMD < 

%DIFDELCMD < %%%
\subsection{\DIFdel{Crossing the factors, replaying the tensors, and changing the endpoint}}
%DIFAUXCMD
\addtocounter{subsection}{-1}%DIFAUXCMD
%DIFDELCMD < \label{app:cf-crossed}
%DIFDELCMD < 

%DIFDELCMD < %%%
\DIFdelend \paragraph{Crossed tower and reader.} \DIFdelbegin \DIFdel{The swap replaces the host checkpoint's vision tower with the donor checkpoint's tower, leaving the host's }\DIFdelend \DIFaddbegin \DIFadd{We load the vision tower of Gemma~3 4B (G) behind the }\DIFaddend connector and language model \DIFdelbegin \DIFdel{untouched. For every swap, the module records a receipt listing the replaced }\DIFdelend \DIFaddbegin \DIFadd{of MedGemma 4B (M), and the reverse. All }\cfSwapTensors{} \DIFadd{weight }\DIFaddend tensors of the \DIFdelbegin \DIFdel{tower root, their parameter count, a digest of the donor tower, and tensor-by-tensor comparisons of the loaded tower against the donorand of everything outside the tower against the host's own weights. Of the host tower's }%DIFDELCMD < \cfSwapTowerTensors{} %%%
\DIFdel{tensors, }%DIFDELCMD < \cfSwapTensors{} %%%
\DIFdel{carry weights; the module replaces each and verifies bitwise equality with the donor's (}%DIFDELCMD < \cfSwapParamsM{} %%%
\DIFdel{million parameters). The module rebuilds the remaining tensor, a position buffer. None }\DIFdelend \DIFaddbegin \DIFadd{loaded tower equal the donor's bit for bit, and none }\DIFaddend of the \cfSwapOutside{} tensors outside \DIFdelbegin \DIFdel{the tower changes. Write directions are fitted in the representation each towerproduces, so a swapped arm writes the donor's directions. The two factors therefore carry different loads: at a fixed tower the swap changes the connector and the language model and nothing else, while replacing the tower replaces the representation and the six directions fitted on it together. The tower factor is the joint effect of both; the reader factor is the effect of neither. Each of the }\DIFdelend \DIFaddbegin \DIFadd{it changes. Directions travel with their tower, so changing the tower changes the representation and its six directions together, while changing the reader changes only the connector and the language model. All }\DIFaddend \cfSwapCombos{} combinations \DIFdelbegin \DIFdel{of tower and reader is }\DIFdelend \DIFaddbegin \DIFadd{are }\DIFaddend graded on the same \cfSwapRows{} test rows of the block, at the primary dose and template, against the 119-direction random family and the sham of the tower in use\DIFdelbegin \DIFdel{, with the $6\times5$ max-$T$ verdict (Table~\ref{tab:cf-ledger})}\DIFdelend . The two native combinations \DIFdelbegin \DIFdel{give the two }\DIFdelend \DIFaddbegin \DIFadd{reproduce the }\DIFaddend checkpoints' own grades on \DIFdelbegin \DIFdel{that cohort. Across }\DIFdelend \DIFaddbegin \DIFadd{this smaller cohort: across }\DIFaddend the \cfSwapNativeCells{} native cells, the ownership contrast \DIFdelbegin \DIFdel{changes }\DIFdelend \DIFaddbegin \DIFadd{moves }\DIFaddend by at most \cfSwapNativeMaxDO{}\DIFdelbegin \DIFdel{relative to the grade on the block's whole test cohort. On the smaller cohort, }\DIFdelend \DIFaddbegin \DIFadd{, and }\DIFaddend \cfSwapNativeUnresolved{} of the \cfSwapNativeOwnedFull{} \DIFdelbegin \DIFdel{cells owned in that grade }\DIFdelend \DIFaddbegin \DIFadd{owned cells }\DIFaddend become unresolved. \DIFdelbegin \DIFdel{The crossover contrasts hold one factor fixed and change the other, and report three quantities of the same six questions: the own write $W_{q,q}$, }\DIFdelend \DIFaddbegin \DIFadd{Writing $Y_q(g/h)$ for a quantity of question $q$ with tower $g$ and reader $h$, }\DIFaddend the \DIFdelbegin \DIFdel{ownership contrast $O_q$, and the margin of the own writeover the strongest competitor in the }\DIFdelend \DIFaddbegin \DIFadd{mean reader and tower effects are
}\begin{equation}
\delta_{\text{reader}}=\frac{1}{12}\sum_{g}\sum_{q}\bigl|Y_q(g/\mathrm G)-Y_q(g/\mathrm M)\bigr|,\qquad
\delta_{\text{tower}}=\frac{1}{12}\sum_{h}\sum_{q}\bigl|Y_q(\mathrm G/h)-Y_q(\mathrm M/h)\bigr|,
\label{eq:crossover}
\end{equation}
\DIFadd{for $Y_q\in\{W_{q,q},\,O_q,\,M_q\}$, where $M_q=W_{q,q}-\max(\max_{d\neq q}W_{q,d},\,p_{95},\,|W^{\mathrm{sham}}_q|)$ is the own write's margin over the }\DIFaddend whole reference family, \DIFdelbegin \DIFdel{$M_q=W_{q,q}-\max(\max_{d\neq q}W_{q,d},\,\text{random }p_{95},\,|\text{sham}|)$, each as the mean over the six questions of the absolute paired change, with a percentile interval over the same 2}%DIFDELCMD < {%%%
\DIFdel{,}%DIFDELCMD < }%%%
\DIFdel{000 draws of the rows and , in brackets, the number of questions of six whose change is simultaneously nonzero.
The last two rows of each block average each factor over its two levels. The three come from one pass over the four arms' packaged outcomes, so a row's columns are the same arms read three ways. The four arms of a dataset are the same four arms whichever of the pair's blocks indexes them, and the two host blocks give the same crossover: the middle panel counts each dataset's four combinationsonce, and an arm graded in both host blocks agrees in all }%DIFDELCMD < \cfSwapArmsTwice{} %%%
\DIFdel{cases where it appears twice. A dash marks an arm that is not graded in that host block. Table~\ref{tab:cf-towerswap} reports every combination and contrast.
}\DIFdelend \DIFaddbegin \DIFadd{with $p_{95}$ the random 95th percentile and $W^{\mathrm{sham}}_q$ the sham write.
}\DIFaddend 

\DIFdelbegin \paragraph{\DIFdel{What the crossing owns.}} %DIFAUXCMD
\addtocounter{paragraph}{-1}%DIFAUXCMD
\DIFdel{Gemma~3 4B and MedGemma 4B give }%DIFDELCMD < \cfSwapCombos{} %%%
\DIFdel{combinations of tower and reader on each dataset, and written directions travel with the tower that produced them. On NIH ChestX-ray14 and CheXpert Plus, no combination owns a finding: }\DIFdelend \DIFaddbegin \DIFadd{Under the four combinations, }\DIFaddend \cfSwapChestOwned{} of \cfSwapChestCells{} \DIFdelbegin \DIFdel{cells across the }%DIFDELCMD < \cfSwapChestCombos{} %%%
\DIFdel{combinations. On COCO, the same crossing owns }\DIFdelend \DIFaddbegin \DIFadd{chest cells are owned, against }\DIFaddend \cfSwapCocoOwned{} of \cfSwapCocoCells{} \DIFdelbegin \DIFdel{cells. Each native combination owns }\DIFdelend \DIFaddbegin \DIFadd{COCO cells: on COCO, }\DIFaddend \cfSwapCocoNativeMin{} of the six objects \DIFdelbegin \DIFdel{, }\DIFdelend \DIFaddbegin \DIFadd{are owned under each native combination, }\cfSwapCocoGemmaTowerCross{} \DIFadd{with }\DIFaddend the Gemma tower behind the MedGemma reader\DIFdelbegin \DIFdel{owns }%DIFDELCMD < \cfSwapCocoGemmaTowerCross{}%%%
\DIFdel{, and }\DIFdelend \DIFaddbegin \DIFadd{, and }\cfSwapCocoMedTowerCross{} \DIFadd{with }\DIFaddend the MedGemma tower behind the Gemma~3 reader\DIFdelbegin \DIFdel{owns }%DIFDELCMD < \cfSwapCocoMedTowerCross{}%%%
\DIFdel{. The crossover holds one factor fixed and changes the other on three quantities of the same six questions. }\DIFdelend \DIFaddbegin \DIFadd{. }\DIFaddend On NIH, the \DIFdelbegin \DIFdel{mean absolute }\DIFdelend reader effect is \cfSwapReaderNihMin{} $[\cfSwapReaderNihLo,\cfSwapReaderNihHi]$ on \DIFdelbegin \DIFdel{the own write }\DIFdelend \DIFaddbegin \DIFadd{$W_{q,q}$ }\DIFaddend and \cfSwapReaderNihOMin{} $[\cfSwapReaderNihOLo,\cfSwapReaderNihOHi]$ on \DIFdelbegin \DIFdel{the ownership contrast, against a tower effect of }\DIFdelend \DIFaddbegin \DIFadd{$O_q$, against }\DIFaddend \cfSwapTowerNihMin{} \DIFdelbegin \DIFdel{$[\cfSwapTowerNihLo,\cfSwapTowerNihHi]$ }\DIFdelend and \cfSwapTowerNihOMin{} \DIFdelbegin \DIFdel{$[\cfSwapTowerNihOLo,\cfSwapTowerNihOHi]$. On CheXpert the reader effect }\DIFdelend \DIFaddbegin \DIFadd{for the tower; on CheXpert it }\DIFaddend is \cfSwapReaderChexMin{} \DIFdelbegin \DIFdel{on the own write }\DIFdelend and \cfSwapReaderChexOMin{}\DIFdelbegin \DIFdel{on the ownership contrast}\DIFdelend , against \cfSwapTowerChexMin{} and \cfSwapTowerChexOMin{}\DIFdelbegin \DIFdel{for the tower. On the family margin }\DIFdelend \DIFaddbegin \DIFadd{; and on the family margin $M_q$ }\DIFaddend the two chest sets give \cfSwapReaderNihMMin{} and \cfSwapReaderChexMMin{} for the reader against \cfSwapTowerNihMMin{} and \cfSwapTowerChexMMin{} for the tower. On COCO the tower effect is the larger \DIFdelbegin \DIFdel{of the two }\DIFdelend on all three quantities \DIFdelbegin \DIFdel{: }\DIFdelend \DIFaddbegin \DIFadd{(}\DIFaddend \cfSwapTowerCocoMin{}, \cfSwapTowerCocoOMin{}, and \cfSwapTowerCocoMMin{} against \cfSwapReaderCocoMin{}, \cfSwapReaderCocoOMin{}, and \cfSwapReaderCocoMMin{}\DIFaddbegin \DIFadd{)}\DIFaddend .

\DIFdelbegin \paragraph{\DIFdel{How well the crossed arms answer.}} %DIFAUXCMD
\addtocounter{paragraph}{-1}%DIFAUXCMD
\DIFdel{The hybrid arms }\DIFdelend \DIFaddbegin \DIFadd{The hybrid combinations }\DIFaddend answer worse than the native \DIFdelbegin \DIFdel{arms. On the labelled rows the crossing grades, with }\DIFdelend \DIFaddbegin \DIFadd{ones. With }\DIFaddend no write applied, \DIFdelbegin \DIFdel{the }\DIFdelend \DIFaddbegin \DIFadd{all }\cfSwapAnsWorse{} \DIFadd{of the }\cfSwapAnsArms{} \DIFadd{hybrid arms have a lower }\DIFaddend median clean-answer AUROC \DIFdelbegin \DIFdel{over a block's six concepts }\DIFdelend \DIFaddbegin \DIFadd{than the native arm of their block, by }\cfSwapAnsDropMin{} \DIFadd{to }\cfSwapAnsDropMax{}\DIFadd{: it }\DIFaddend falls from \cfSwapAnsNativeGemCoco{} \DIFdelbegin \DIFdel{in the native Gemma~3 4B arm on COCO }\DIFdelend to \cfSwapAnsHybridGemCoco{} \DIFdelbegin \DIFdel{in }\DIFdelend \DIFaddbegin \DIFadd{when }\DIFaddend the MedGemma tower \DIFdelbegin \DIFdel{behind that reader, from }%DIFDELCMD < \cfSwapAnsNativeMedNih{} %%%
\DIFdel{in the native MedGemma 4B arm on NIH ChestX-ray14 to }%DIFDELCMD < \cfSwapAnsHybridMedNih{} %%%
\DIFdel{in the }\DIFdelend \DIFaddbegin \DIFadd{serves the }\DIFaddend Gemma~3 \DIFdelbegin \DIFdel{tower behind that reader }\DIFdelend \DIFaddbegin \DIFadd{reader on COCO}\DIFaddend , and from \DIFaddbegin \cfSwapAnsNativeMedNih{} \DIFadd{to }\cfSwapAnsHybridMedNih{} \DIFadd{and from }\DIFaddend \cfSwapAnsNativeMedChex{} to \cfSwapAnsHybridMedChex{} \DIFdelbegin \DIFdel{on CheXpert Plus. The remaining three arms give }%DIFDELCMD < \cfSwapAnsNativeGemNih{} %%%
\DIFdel{against }%DIFDELCMD < \cfSwapAnsHybridGemNih{} %%%
\DIFdel{and }%DIFDELCMD < \cfSwapAnsNativeGemChex{} %%%
\DIFdel{against }%DIFDELCMD < \cfSwapAnsHybridGemChex{} %%%
\DIFdel{in }\DIFdelend \DIFaddbegin \DIFadd{when }\DIFaddend the Gemma~3 \DIFdelbegin \DIFdel{4B blocks, and }%DIFDELCMD < \cfSwapAnsNativeMedCoco{} %%%
\DIFdel{against }%DIFDELCMD < \cfSwapAnsHybridMedCoco{} %%%
\DIFdel{in the MedGemma 4B block on COCO.
All }%DIFDELCMD < \cfSwapAnsWorse{} %%%
\DIFdel{of the }%DIFDELCMD < \cfSwapAnsArms{} %%%
\DIFdel{hybrid arms fall below the native arm of their own block, by }%DIFDELCMD < \cfSwapAnsDropMin{} %%%
\DIFdel{to }%DIFDELCMD < \cfSwapAnsDropMax{}%%%
\DIFdel{.
}\DIFdelend \DIFaddbegin \DIFadd{tower serves the MedGemma reader on NIH and CheXpert.
}\DIFaddend 

\DIFdelbegin \begin{table}[htbp]
\DIFdeletedtable{%
{\small\textbf{Table (deleted).} \textbf{Crossing the tower and the reader.} Each column of the top two panels is tower/reader, with G for Gemma~3 4B and M for MedGemma 4B. The bottom panel changes one factor at a time and gives the mean absolute paired change over the six questions.\par}\medskip
\centering
\setlength{\tabcolsep}{4pt}

\fitwidth{%
\fitwidth{%
\begin{tabular}{llcccccc}
\toprule
\multirow{2}{*}{Host checkpoint} & \multirow{2}{*}{Dataset} & \multicolumn{2}{c}{native} & \multicolumn{2}{c}{swapped} & \multirow{2}{*}{concepts} & \multirow{2}{*}{crossover}\\
\cmidrule(lr){3-4}\cmidrule(lr){5-6}
&  & \multicolumn{1}{c}{G/G} & \multicolumn{1}{c}{M/M} & \multicolumn{1}{c}{M/G} & \multicolumn{1}{c}{G/M} &  & \\
\midrule
Gemma 3 4B & NIH ChestX-ray14 & 0 & 0 & 0 & 0 & 6 & complete \\
Gemma 3 4B & CheXpert Plus & 0 & 0 & 0 & 0 & 6 & complete \\
Gemma 3 4B & COCO & 5 & 5 & 0 & 4 & 6 & complete \\
MedGemma 4B & NIH ChestX-ray14 & 0 & 0 & 0 & 0 & 6 & complete \\
MedGemma 4B & CheXpert Plus & 0 & 0 & -- & 0 & 6 & -- \\
MedGemma 4B & COCO & 5 & 5 & -- & 4 & 6 & -- \\
\midrule
\multicolumn{8}{l}{\textit{The four combinations of each dataset, shared arms counted once}} \\
Combinations & Dataset & G/G & M/M & M/G & G/M & concepts & owned of cells \\
\midrule
all four & NIH ChestX-ray14 & 0 & 0 & 0 & 0 & 6 & 0/24 \\
all four & CheXpert Plus & 0 & 0 & 0 & 0 & 6 & 0/24 \\
all four & COCO & 5 & 5 & 0 & 4 & 6 & 14/24 \\
\bottomrule
\end{tabular}}}\par
\vspace{4pt}
\setlength{\tabcolsep}{2.5pt}
\fitwidth{%
\fitwidth{%
\begin{tabular}{llccc}
\toprule
\multicolumn{5}{l}{\textit{Crossover: the effect of replacing one factor with the other held fixed}} \\
\midrule
Block & Changed, held & own write $W_{q,q}$ & ownership $O_q$ & family margin $M_q$ \\
\midrule
Gemma 3 4B, NIH & reader, own tower & 0.133 [0.106, 0.169] (1) & 0.158 [0.136, 0.187] (4) & 0.425 [0.393, 0.458] (5) \\
 & reader, partner tower & 0.240 [0.221, 0.259] (5) & 0.210 [0.192, 0.233] (4) & 0.286 [0.268, 0.309] (6) \\
 & tower, own reader & 0.164 [0.140, 0.189] (4) & 0.061 [0.046, 0.080] (1) & 0.137 [0.120, 0.164] (4) \\
 & tower, partner reader & 0.160 [0.125, 0.193] (6) & 0.265 [0.239, 0.299] (5) & 0.496 [0.462, 0.525] (6) \\
 & \quad mean reader effect & 0.186 [0.170, 0.207] & 0.184 [0.169, 0.203] & 0.356 [0.339, 0.376] \\
 & \quad mean tower effect & 0.162 [0.140, 0.182] & 0.163 [0.149, 0.185] & 0.316 [0.299, 0.338] \\
\addlinespace
Gemma 3 4B, CheXpert & reader, own tower & 0.382 [0.353, 0.412] (5) & 0.303 [0.273, 0.332] (6) & 0.490 [0.465, 0.515] (5) \\
 & reader, partner tower & 0.125 [0.102, 0.151] (4) & 0.214 [0.194, 0.240] (5) & 0.240 [0.214, 0.263] (6) \\
 & tower, own reader & 0.269 [0.245, 0.293] (5) & 0.223 [0.205, 0.241] (6) & 0.372 [0.352, 0.390] (5) \\
 & tower, partner reader & 0.183 [0.153, 0.217] (4) & 0.297 [0.267, 0.332] (5) & 0.517 [0.485, 0.553] (6) \\
 & \quad mean reader effect & 0.254 [0.235, 0.274] & 0.259 [0.240, 0.280] & 0.365 [0.347, 0.383] \\
 & \quad mean tower effect & 0.226 [0.208, 0.246] & 0.260 [0.243, 0.279] & 0.444 [0.425, 0.466] \\
\addlinespace
Gemma 3 4B, COCO & reader, own tower & 0.072 [0.049, 0.104] (2) & 0.229 [0.187, 0.271] (5) & 0.143 [0.114, 0.185] (3) \\
 & reader, partner tower & 0.202 [0.170, 0.239] (3) & 0.425 [0.390, 0.459] (6) & 0.432 [0.398, 0.464] (6) \\
 & tower, own reader & 0.295 [0.261, 0.337] (5) & 0.497 [0.453, 0.540] (6) & 0.408 [0.368, 0.450] (6) \\
 & tower, partner reader & 0.200 [0.175, 0.227] (4) & 0.335 [0.307, 0.364] (5) & 0.298 [0.273, 0.326] (5) \\
 & \quad mean reader effect & 0.137 [0.119, 0.161] & 0.327 [0.298, 0.354] & 0.288 [0.265, 0.315] \\
 & \quad mean tower effect & 0.248 [0.228, 0.271] & 0.416 [0.392, 0.441] & 0.353 [0.331, 0.379] \\
\addlinespace
MedGemma 4B, NIH & reader, own tower & 0.240 [0.221, 0.259] (5) & 0.210 [0.192, 0.233] (4) & 0.286 [0.268, 0.309] (6) \\
 & reader, partner tower & 0.133 [0.106, 0.169] (1) & 0.158 [0.136, 0.187] (4) & 0.425 [0.393, 0.458] (5) \\
 & tower, own reader & 0.160 [0.125, 0.193] (6) & 0.265 [0.239, 0.299] (5) & 0.496 [0.462, 0.525] (6) \\
 & tower, partner reader & 0.164 [0.140, 0.189] (4) & 0.061 [0.046, 0.080] (1) & 0.137 [0.120, 0.164] (4) \\
 & \quad mean reader effect & 0.186 [0.170, 0.207] & 0.184 [0.169, 0.203] & 0.356 [0.339, 0.376] \\
 & \quad mean tower effect & 0.162 [0.140, 0.182] & 0.163 [0.149, 0.185] & 0.316 [0.299, 0.338] \\
\bottomrule
\end{tabular}}}
}
\end{table}
%DIFDELCMD < 

%DIFDELCMD < %%%
\DIFdelend \paragraph{Identical-tensor replay.} Checkpoints \DIFdelbegin \DIFdel{within }\DIFdelend \DIFaddbegin \DIFadd{of }\DIFaddend one family share a vision tower, but two readers' inputs agree only up to bf16 rounding. The replay module \DIFdelbegin \DIFdel{removes this rounding difference. It }\DIFdelend stores the consumed block's token tensors once from the group's source checkpoint and feeds them to every reader \DIFdelbegin \DIFdel{in the group , replacing each }\DIFdelend \DIFaddbegin \DIFadd{of the group in place of that }\DIFaddend reader's own tower output\DIFdelbegin \DIFdel{so the two }\DIFdelend \DIFaddbegin \DIFadd{, so that the }\DIFaddend readers consume the same bits. Before \DIFdelbegin \DIFdel{replacement, it measures drift as }\DIFdelend \DIFaddbegin \DIFadd{the replacement it measures the drift, }\DIFaddend the largest absolute difference between the reader's own tower output and the stored tensor on the same rows. \DIFdelbegin \DIFdel{A reader replaying the tensors of its own block controls for the machinery: its replayed grade is its own grade computed through the replacement path. The module also recomputes each pair's between-reader difference twice: once with each reader running its own tower and once with both reading the stored tensors. The two computations differ only in their input. Across the }%DIFDELCMD < \cfReplaySelf{} %%%
\DIFdel{replays in which a reader reads }\DIFdelend \DIFaddbegin \DIFadd{In }\cfReplayDriftZero{} \DIFadd{of the }\cfReplayBlocks{} \DIFadd{replays that output already matches the stored tensor bit for bit; the exception is LLaVA-Med, whose fine-tuned tower departs by up to }\cfReplayDriftMax{} \DIFadd{against a mean absolute tensor entry of }\cfReplayTensorMin{}\DIFadd{--}\cfReplayTensorMax{}\DIFadd{.
}

\DIFadd{A reader that replays }\DIFaddend its own block's tensors \DIFdelbegin \DIFdel{, }\DIFdelend \DIFaddbegin \DIFadd{tests the replacement path itself: across the }\cfReplaySelf{} \DIFadd{such replays }\DIFaddend no verdict and no ownership grade changes\DIFdelbegin \DIFdel{; }\DIFdelend \DIFaddbegin \DIFadd{, }\DIFaddend \cfReplaySelfExact{} reproduce the block's \DIFdelbegin \DIFdel{own grade exactly}\DIFdelend \DIFaddbegin \DIFadd{grade exactly, }\DIFaddend and the third differs by at most $\cfReplayMaxDWSelf$ in a write magnitude. \DIFdelbegin \DIFdel{In }%DIFDELCMD < \cfReplayDriftZero{} %%%
\DIFdel{of the }\DIFdelend \DIFaddbegin \DIFadd{Across all }\DIFaddend \cfReplayBlocks{} \DIFdelbegin \DIFdel{replays the reader's own tower output already matches the stored tensor bit for bit. The exception is LLaVA-Med, whose fine-tuned tower departs from it by up to }%DIFDELCMD < \cfReplayDriftMax{}%%%
\DIFdel{, against a mean absolute tensor entry of }%DIFDELCMD < \cfReplayTensorMin{}%%%
\DIFdel{--}%DIFDELCMD < \cfReplayTensorMax{}%%%
\DIFdel{. Across the }%DIFDELCMD < \cfReplayBlocks{} %%%
\DIFdelend replays, the change is distinguishable from zero in \cfReplayMoved{} of \cfReplayCells{} concept cells, the largest change in a write magnitude is $\cfReplayMaxDW$, and \cfReplayVerdictChanges{} verdicts and \cfReplayOwnedChanges{} ownership grades change. \DIFdelbegin \DIFdel{Across }\DIFdelend \DIFaddbegin \DIFadd{For }\DIFaddend the \cfReplayPairs{} reader pairs, the mean absolute difference between the \DIFaddbegin \DIFadd{two }\DIFaddend readers' own-question \DIFdelbegin \DIFdel{write magnitudes }\DIFdelend \DIFaddbegin \DIFadd{writes }\DIFaddend is \cfReplayPairOwnMin{}--\cfReplayPairOwnMax{} when each \DIFdelbegin \DIFdel{reader runs }\DIFdelend \DIFaddbegin \DIFadd{reads }\DIFaddend its own tower, and feeding both \DIFdelbegin \DIFdel{readers }\DIFdelend the stored tensors changes that difference by at most \cfReplayChangeMax{}, or \cfReplayRelMax\% of its size.

\DIFdelbegin \paragraph{\DIFdel{Semantic endpoints.}} %DIFAUXCMD
\addtocounter{paragraph}{-1}%DIFAUXCMD
\DIFdel{The four endpoints change what is read after the write}\DIFdelend \DIFaddbegin \subsection{\DIFadd{Non-clinical attributes}}
\label{app:cf-attributes}

\begin{DIFnomarkup}\cfnote{附录精简}{原 Robustness 小节的属性段（方法 + 结果）与原 Further results 小节的属性段（结果）合并，三种比较写在同一句里；删去每个数据集的分项计数和措辞差距的最大值。}\end{DIFnomarkup}

\DIFadd{The three attribute labels come from each dataset's metadata: anteroposterior or portable view against posteroanterior, recorded sex, and age at least 60. Their directions are fitted, lifted, and written with the six clinical directions as one nine-direction family, against the same random family and sham in all }\cfAttrRefBlocks{} \DIFadd{blocks. All }\cfAttrReadable{} \DIFadd{of the }\cfAttrChestAttrN{} \DIFadd{attribute cells are readable, with median probe AUROC }\cfAttrAurocMin{}\DIFadd{--}\cfAttrAurocMax{} \DIFadd{against the label (Table~\ref{tab:cf-attr}). Each attribute question is scored in }\cfAttrqPhrasings{} \DIFadd{phrasings on the calibration rows; the written phrasing reaches a median clean-answer AUROC of }\cfAttrqWrittenAuroc{}\DIFadd{, against }\cfAttrqSelectedAuroc{} \DIFadd{for the best phrasing of the cell, and }\cfAttrqCapable{} \DIFadd{of the }\cfAttrqCells{} \DIFadd{attribute cells are answerable. Attributes are owned no more often than findings under three comparisons: over every cell (}\cfAttrChestAttrOwned{} \DIFadd{of }\cfAttrChestAttrN{} \DIFadd{against }\cfAttrChestClinOwned{} \DIFadd{of }\cfAttrChestClinN{}\DIFadd{), over answerable cells (}\cfAttrCapAttrOwned{} \DIFadd{of }\cfAttrCapAttrN{} \DIFadd{against }\cfAttrCapClinOwned{} \DIFadd{of }\cfAttrCapClinN{}\DIFadd{), and over }\cfAttrPairAttrN{} \DIFadd{pairs matched within }\cfAttrPairWindow{} \DIFadd{clean-answer AUROC in the same block (}\cfAttrPairAttrOwned{} \DIFadd{against }\cfAttrPairClinOwned{}\DIFadd{). Over the }\cfAttrCosPairs{} \DIFadd{attribute--finding pairs, the largest absolute raw cosine between an attribute direction and a clinical normal is }\cfAttrCosMax{}\DIFadd{, for }\cfAttrCosPairLabel{}\DIFadd{.
}

\begin{table}[htbp]
\centering
\setlength{\tabcolsep}{6pt}
\caption{\textbf{\DIFaddFL{Non-clinical attributes of the same radiographs.}} \DIFaddFL{Per attribute and dataset, medians over blocks of the attribute probe's AUROC and selectivity, the clean-answer AUROC of the attribute question, and the absolute raw cosine to the clinical normals, with the owned attribute cells.}}
\label{tab:cf-attr}
\resizebox{\linewidth}{!}{%
\begin{tabular}{llrrrrr}
\toprule
Attribute & Dataset & \multicolumn{1}{c}{probe AUROC} & \multicolumn{1}{c}{selectivity} & \multicolumn{1}{c}{answer AUROC} & \multicolumn{1}{c}{med.\ $|\cos|$} & \multicolumn{1}{c}{owned} \\
\midrule
AP (portable) projection & NIH ChestX-ray14 & 0.999 & 0.32 & 0.76 & 0.03 & 1/19 \\
AP (portable) projection & CheXpert Plus & 0.998 & 0.31 & 0.73 & 0.05 & 0/19 \\
AP (portable) projection & \textit{chest} & 0.999 & 0.31 & 0.74 & 0.04 & 1/38 \\
\addlinespace
female sex & NIH ChestX-ray14 & 0.984 & 0.30 & 0.73 & 0.05 & 3/19 \\
female sex & CheXpert Plus & 0.957 & 0.27 & 0.65 & 0.06 & 1/19 \\
female sex & \textit{chest} & 0.971 & 0.28 & 0.71 & 0.05 & 4/38 \\
\addlinespace
age $\geq60$ & NIH ChestX-ray14 & 0.896 & 0.21 & 0.62 & 0.05 & 3/19 \\
age $\geq60$ & CheXpert Plus & 0.898 & 0.21 & 0.64 & 0.06 & 4/19 \\
age $\geq60$ & \textit{chest} & 0.896 & 0.21 & 0.63 & 0.06 & 7/38 \\
\bottomrule
\end{tabular}}
\end{table}

\subsection{\DIFadd{The answer direction}}
\label{app:cf-ansdir}

\begin{DIFnomarkup}\cfnote{附录精简}{原来分在 4 处的 answer direction 内容合为一节：拟合（ridge 写成公式，删去惩罚取值的分布与行序的实现细节）、拥有情况、读到了什么（白化 cosine、$\kappa_d$ 写成公式）、语义终点（判定规则写成不等式；增量效度的零结果保留）。}\end{DIFnomarkup}

\paragraph{\DIFadd{Fitting and writing.}} \DIFadd{For question $q$, let $m_{iq}=\ell_{\mathrm{yes}}-\ell_{\mathrm{no}}$ be the clean answer margin of training image $i$ on the block's primary template, and $\boldsymbol x_i$ its projected, train-scaled features. The answer direction $a_q$ lifts, by Equation~\ref{eq:lift}, the ridge coefficients
}\begin{equation}
\boldsymbol\gamma_q=\arg\min_{\gamma_0,\boldsymbol\gamma}\;\sum_{i}\bigl(m_{iq}-\gamma_0-\boldsymbol x_i^{\top}\boldsymbol\gamma\bigr)^2+\lambda\lVert\boldsymbol\gamma\rVert_2^2 ,
\label{eq:answer-direction}
\end{equation}
\DIFadd{fitted on the first 3}{\DIFaddend ,\DIFdelbegin \DIFdel{not what is written }\DIFdelend \DIFaddbegin }\DIFadd{000 training rows in the manifest's label-blind order, with $\lambda\in\{0.1,1,10,100\}$ chosen by five-fold cross-validated $R^2$ over patients or images. We write $a_q$ at the same dose against the write matrix's random family and a sham of the $a_q$ under test, with the other five $a_d$ as competitors. Under the five held-out templates, $a_q$ is written unchanged against each template's own clean baseline and sham.
}

\paragraph{\DIFadd{Where it is owned.}} \DIFadd{The answer direction beats the strongest logistic competitor in }\cfAnsChestBeats{} \DIFadd{chest cells (Table~\ref{tab:cf-ansdir}). Over all held-out (cell, template) pairs it is owned in }\cfAnsdirtChestOwnedPairs{} \DIFadd{of }\cfAnsdirtChestPairsAll{} \DIFadd{chest and }\cfAnsdirtCocoOwnedPairs{} \DIFadd{of }\cfAnsdirtCocoPairsAll{} \DIFadd{COCO pairs, and it stays owned in }\cfAnsdirtCocoKept{} \DIFadd{of }\cfAnsdirtCocoKeptN{} \DIFadd{COCO pairs owned under the primary template. In Qwen2.5-VL-7B on NIH, the Effusion answer direction reaches $O_q=\cfAnsQwenEffO$ with an own write of $\cfAnsQwenEffW$}\DIFaddend .
\DIFdelbegin \DIFdel{The negated question asks whether there is no evidence of the finding. A direction that carries the concept raises the affirmative answer}\DIFdelend \DIFaddbegin 

\begin{table}[htbp]
\centering
\setlength{\tabcolsep}{5pt}
\caption{\textbf{\DIFaddFL{Ownership under the answer direction.}} \DIFaddFL{Per block, on the primary template: concepts owned by the label direction $\hat w_q$ and by the answer direction $a_q$, split into kept (both), rescued ($a_q$ only), and lost ($\hat w_q$ only); cells where the write of $a_q$ beats the strongest logistic competitor; and medians over the six questions of the own write of $a_q$, $W^a_{q,q}$, and of the ridge fit's $R^2$.}}
\label{tab:cf-ansdir}
\resizebox{\linewidth}{!}{%
\begin{tabular}{llrrrrrrrr}
\toprule
\multirow{2}{*}{Checkpoint} & \multirow{2}{*}{Dataset} & \multicolumn{5}{c}{owned concepts} & \multirow{2}{*}{$a_q>$ comp.} & \multicolumn{2}{c}{median}\\
\cmidrule(lr){3-7}\cmidrule(lr){9-10}
&  & \multicolumn{1}{c}{$\hat w_q$} & \multicolumn{1}{c}{$a_q$} & \multicolumn{1}{c}{kept} & \multicolumn{1}{c}{rescued} & \multicolumn{1}{c}{lost} &  & \multicolumn{1}{c}{$W^a_{q,q}$} & \multicolumn{1}{c}{$R^2$}\\
\midrule
Qwen2.5-VL-7B & NIH ChestX-ray14 & 0/6 & 3/6 & 0 & 3 & 0 & 6/6 & 0.61 & 0.54 \\
Qwen2.5-VL-7B & CheXpert Plus & 0/6 & 6/6 & 0 & 6 & 0 & 6/6 & 0.55 & 0.48 \\
Qwen2.5-VL-7B & COCO & 6/6 & 6/6 & 6 & 0 & 0 & 6/6 & 0.78 & 0.58 \\
Gemma 3 12B & NIH ChestX-ray14 & 1/6 & 2/6 & 1 & 1 & 0 & 5/6 & 0.47 & 0.17 \\
Gemma 3 12B & CheXpert Plus & 1/6 & 0/6 & 0 & 0 & 1 & 1/6 & 0.37 & 0.19 \\
Gemma 3 12B & COCO & 6/6 & 6/6 & 6 & 0 & 0 & 6/6 & 0.61 & 0.37 \\
Lingshu 32B & NIH ChestX-ray14 & 2/6 & 3/6 & 2 & 1 & 0 & 6/6 & 0.24 & 0.77 \\
Lingshu 32B & CheXpert Plus & 5/6 & 5/6 & 5 & 0 & 0 & 6/6 & 0.22 & 0.80 \\
Lingshu 32B & COCO & 6/6 & 6/6 & 6 & 0 & 0 & 6/6 & 0.36 & 0.55 \\
\midrule
NIH ChestX-ray14 & 3 blocks & 3/18 & 8/18 & 3 & 5 & 0 & 17/18 & 0.47 & 0.54 \\
CheXpert Plus & 3 blocks & 6/18 & 11/18 & 5 & 6 & 1 & 13/18 & 0.35 & 0.48 \\
\textit{chest} & 6 blocks & 9/36 & 19/36 & 8 & 11 & 1 & 30/36 & 0.41 & 0.50 \\
COCO & 3 blocks & 18/18 & 18/18 & 18 & 0 & 0 & 18/18 & 0.61 & 0.54 \\
\bottomrule
\end{tabular}}
\end{table}

\paragraph{\DIFadd{What it reads.}} \DIFadd{In Qwen2.5-VL-7B on NIH (Table~\ref{tab:cf-validation}), the answer direction's AUROC is $\cfValAurocAnsQwenNih$ against the label and $\cfValAurocAnsCleanQwenNih$ against the clean answer, against $\cfValAurocProbeQwenNih$ and $\cfValAurocProbeCleanQwenNih$ for the probe. Cosines depend on the inner product chosen for the representation \mbox{%DIFAUXCMD
\citep{park2024linear}}\hskip0pt%DIFAUXCMD
, so besides the raw cosine we report, for two coefficient vectors $\boldsymbol\beta,\boldsymbol\beta'$ in the projected space,
}\begin{equation}
\cos_\Sigma(\boldsymbol\beta,\boldsymbol\beta')=\frac{\boldsymbol\beta^{\top}\Sigma\boldsymbol\beta'}{\sqrt{\boldsymbol\beta^{\top}\Sigma\boldsymbol\beta\;\boldsymbol\beta'^{\top}\Sigma\boldsymbol\beta'}},
\label{eq:whitened-cos}
\end{equation}
\DIFadd{with $\Sigma$ the covariance of the training features: the correlation over training images between the two vectors' scores. Between answer and probe coefficients, the median raw cosine is $\cfAnsCosQwenMin$--$\cfAnsCosQwenMax$ on the chest blocks of Qwen2.5-VL-7B, $\cfAnsCosLingMin$--$\cfAnsCosLingMax$ on those of Lingshu 32B, and $\cfAnsCosCocoMin$--$\cfAnsCosCocoMax$ on COCO; the median whitened cosine is $\cfValCosSigmaChestMin$--$\cfValCosSigmaChestMax$ on the chest sets and $\cfValCosSigmaCoco$ on COCO, and the two metrics differ most for Lingshu 32B on CheXpert ($\cfValCosRawLingChex$ raw, $\cfValCosSigmaLingChex$ whitened).
}

\DIFadd{The column selectivity
}\begin{equation}
\kappa_d=\frac{W_{d,d}}{\sum_q\lvert W_{q,d}\rvert}
\label{eq:column-selectivity}
\end{equation}
\DIFadd{equals one when a write moves only its own answer. Over the }\cfAnsColSelChestN{} \DIFadd{chest cells its median is $\cfAnsColSelChest$ for answer directions and $\cfAnsColSelLabelChest$ for label directions, against $\cfAnsColSelCoco$ }\DIFaddend and \DIFdelbegin \DIFdel{lowers the negated one, so the two raw effects must have opposite signs.
The forced choice asks which is present: the finding or }\DIFdelend \DIFaddbegin \DIFadd{$\cfAnsColSelLabelCoco$ over }\cfAnsColSelCocoN{} \DIFadd{COCO cells; over all cells it is }\cfValColSelNih{} \DIFadd{on NIH, }\cfValColSelChex{} \DIFadd{on CheXpert, and }\cfValColSelCoco{} \DIFadd{on COCO. On the chest sets, both directions move the other five answers along with their own.
}

\begin{table}[htbp]
\centering
\setlength{\tabcolsep}{8pt}
\caption{\textbf{\DIFaddFL{The answer direction against the label direction.}} \DIFaddFL{Per block, medians over the six questions of the AUROCs of $a_q$ and $\hat w_q$ against the label and against the clean answer, the Pearson correlation, and the raw and whitened cosines between $a_q$ and $\hat w_q$.}}
\begin{DIFnomarkup}\cfnote{删表}{原表由三块拼成（逐 block 的 AUROC 与 cosine、一段 Spearman 脚注、按数据集的 column selectivity），列结构互不相同。只保留上半部分：9 行同一套列，直接显示胸片上 $a_q$ 与 $\hat w_q$ 的 cosine 低、COCO 上高，$a_q$ 对 clean answer 的 AUROC 高于 $\hat w_q$。Spearman 脚注没有宏、正文不依赖，删去；$\kappa_d$ 汇总已在 A.10 文字中。表注改为只描述保留的部分，正文 5.6 与 A.10 引用本表。}\end{DIFnomarkup}

\label{tab:cf-validation}
\resizebox{\linewidth}{!}{%
\begin{tabular}{llrrrrrrr}
\toprule
\multirow{3}{*}{Checkpoint} & \multirow{3}{*}{Dataset} & \multicolumn{4}{c}{AUROC} & \multirow{3}{*}{Pearson} & \multicolumn{2}{c}{\multirow{2}{*}{cosine}}\\
\cmidrule(lr){3-6}
&  & \multicolumn{2}{c}{label} & \multicolumn{2}{c}{clean answer} &  & & \\
\cmidrule(lr){3-4}\cmidrule(lr){5-6}\cmidrule(lr){8-9}
&  & \multicolumn{1}{c}{$a_q$} & \multicolumn{1}{c}{$\hat w_q$} & \multicolumn{1}{c}{$a_q$} & \multicolumn{1}{c}{$\hat w_q$} &  & \multicolumn{1}{c}{raw} & \multicolumn{1}{c}{$\Sigma$}\\
\midrule
Qwen2.5-VL-7B & NIH ChestX-ray14 & 0.60 & 0.73 & 0.93 & 0.66 & 0.36 & 0.07 & 0.35 \\
Qwen2.5-VL-7B & CheXpert Plus & 0.67 & 0.86 & 0.85 & 0.66 & 0.34 & 0.06 & 0.36 \\
Qwen2.5-VL-7B & COCO & 0.97 & 0.97 & 0.98 & 0.98 & 0.84 & 0.64 & 0.84 \\
Lingshu 32B & NIH ChestX-ray14 & 0.75 & 0.76 & 0.95 & 0.79 & 0.62 & 0.30 & 0.60 \\
Lingshu 32B & CheXpert Plus & 0.85 & 0.86 & 0.95 & 0.89 & 0.88 & 0.33 & 0.87 \\
Lingshu 32B & COCO & 0.94 & 0.97 & 0.93 & 0.92 & 0.75 & 0.54 & 0.73 \\
Gemma 3 12B & NIH ChestX-ray14 & 0.59 & 0.65 & 0.72 & 0.61 & 0.31 & 0.18 & 0.33 \\
Gemma 3 12B & CheXpert Plus & 0.61 & 0.78 & 0.73 & 0.56 & 0.13 & 0.17 & 0.10 \\
Gemma 3 12B & COCO & 0.93 & 0.96 & 0.80 & 0.76 & 0.81 & 0.59 & 0.79 \\
\bottomrule
\end{tabular}}
\end{table}

\paragraph{\DIFadd{Semantic endpoints.}} \DIFadd{Four endpoints keep the write and change what is read after it: the negated question, whose effect must take the opposite sign; the forced choice between the finding and }\DIFaddend its strongest competitor, \DIFdelbegin \DIFdel{taken from the block's own primary-template write matrix. We ask it in both presentation orders , so }\DIFdelend \DIFaddbegin \DIFadd{asked in both orders so that }\DIFaddend a position preference \DIFdelbegin \DIFdel{appears }\DIFdelend \DIFaddbegin \DIFadd{shows }\DIFaddend as a gap\DIFdelbegin \DIFdel{between them. The report continuation requests the findings section and reads }\DIFdelend \DIFaddbegin \DIFadd{; and }\DIFaddend the log-probability of the finding word \DIFdelbegin \DIFdel{. Each endpoint }\DIFdelend \DIFaddbegin \DIFadd{in a report continuation. Each }\DIFaddend carries its own clean baseline, \DIFdelbegin \DIFdel{a coordinate-permutation shamof each direction under test, and its own family of }\DIFdelend \DIFaddbegin \DIFadd{sham, and }\DIFaddend \cfSemLedgerRandom{} random directions\DIFdelbegin \DIFdel{written at the same dose. An endpoint }\DIFdelend \DIFaddbegin \DIFadd{. With $E_q$ the signed endpoint effect of the own write, a }\DIFaddend cell meets the \DIFdelbegin \DIFdel{steering reference when
its own signed effect is positive, above every one of those }\DIFdelend \DIFaddbegin \DIFadd{reference when
}\begin{equation}
E_q>\max\Bigl(0,\;\max_{k}E^{\mathrm{rand}}_{k},\;\bigl|E^{\mathrm{sham}}_q\bigr|\Bigr),
\label{eq:endpoint-reference}
\end{equation}
\DIFadd{the inner maximum running over the }\DIFaddend \cfSemLedgerRandom{} random \DIFdelbegin \DIFdel{writes, and above the absolute sham, and }\DIFdelend \DIFaddbegin \DIFadd{directions, and }\DIFaddend it is owned when it also holds the \DIFdelbegin \DIFdel{campaign's simultaneous }\DIFdelend $6\times5$ max-$T$ advantage.
\DIFdelbegin \DIFdel{The random bar here is the maximum of a family of }%DIFDELCMD < \cfSemLedgerRandom{} %%%
\DIFdel{rather than the 95th percentile of the campaign's 119; everything else is the campaign's rule }\DIFdelend \DIFaddbegin 

\DIFadd{Across }\cfSemBlocks{} \DIFadd{blocks, the negated question moves opposite to the affirmative one in }\cfSemNegOpposite{} \DIFadd{of }\cfSemNegN{} \DIFadd{concept--block pairs, the forced choice is owned in both orders in }\cfSemForcedOwned{} \DIFadd{of }\cfSemForcedN{} \DIFadd{pairs (mean order gap }\cfSemGapMin{}\DIFadd{--}\cfSemGapMax{}\DIFadd{), and the report continuation is owned in }\cfSemReportOwned{} \DIFadd{of }\cfSemReportCells{} \DIFadd{cells }\DIFaddend (Table~\DIFdelbegin \DIFdel{\ref{tab:cf-ledger}). The incremental-validity test regresses the }\DIFdelend \DIFaddbegin \DIFadd{\ref{tab:cf-semend}). Adding the ownership contrast to a regression of the }\DIFaddend endpoint effect on write magnitude, probe selectivity, and clean-answer AUROC \DIFdelbegin \DIFdel{, with endpoint dummies, then adds the ownership contrast and its owned indicator. The penalty and the standardisation are fitted inside each training fold, so no }\DIFdelend \DIFaddbegin \DIFadd{changes the }\DIFaddend held-out \DIFdelbegin \DIFdel{cell enters its own prediction. The predictors are per-concept block-level quantities, so the effective sample is the six concepts rather than the cells. The patient bootstrap preserves the dependence between cells that share rows. We cross-validate both models twice, leaving out one concept at a time inside a block and one block at a time across the }%DIFDELCMD < \cfSemCvBlocks{} %%%
\DIFdel{blocks. Held-out }\DIFdelend $R^2$ \DIFdelbegin \DIFdel{changes }\DIFdelend by $\cfSemCvLocoDelta$ \DIFdelbegin \DIFdel{under the concept split }\DIFdelend \DIFaddbegin \DIFadd{when one concept is left out }\DIFaddend and $\cfSemCvLoboDelta$ \DIFdelbegin \DIFdel{under the block split. A null that permutes ownership across the concepts of a block over }\DIFdelend \DIFaddbegin \DIFadd{when one of }\cfSemCvBlocks{} \DIFadd{blocks is left out; permuting ownership within blocks }\DIFaddend \cfSemCvPerms{} \DIFdelbegin \DIFdel{draws puts both at }\DIFdelend \DIFaddbegin \DIFadd{times gives }\DIFaddend $p=\cfSemCvLocoP$ and $p=\cfSemCvLoboP$\DIFdelbegin \DIFdel{, and no single split reaches $p<0.05$. Table~\ref{tab:cf-semend} reports every endpoint, test, and regression}\DIFdelend . \DIFaddbegin \DIFadd{The predictors vary by concept, so each block contributes an effective sample of six.
}\DIFaddend 

\begin{table}[htbp]
\centering
\DIFdelbeginFL %DIFDELCMD < \setlength{\tabcolsep}{3pt}
%DIFDELCMD < %%%
\DIFdelendFL \DIFaddbeginFL \setlength{\tabcolsep}{6pt}
\DIFaddendFL \caption{\textbf{Four endpoints no direction was fitted against.} \DIFdelbeginFL \DIFdelFL{Both direction families of Section~\ref{sec:ansdir}}\DIFdelendFL \DIFaddbeginFL \DIFaddFL{Per block and endpoint}\DIFaddendFL , \DIFdelbeginFL \DIFdelFL{written at }\DIFdelendFL the \DIFdelbeginFL \DIFdelFL{primary dose }\DIFdelendFL \DIFaddbeginFL \DIFaddFL{concepts of six whose own write is owned }\DIFaddendFL and \DIFdelbeginFL \DIFdelFL{read at }\DIFdelendFL \DIFaddbeginFL \DIFaddFL{whose own write meets }\DIFaddendFL the \DIFdelbeginFL \DIFdelFL{negated question}\DIFdelendFL \DIFaddbeginFL \DIFaddFL{steering reference}\DIFaddendFL , \DIFaddbeginFL \DIFaddFL{for }\DIFaddendFL the \DIFdelbeginFL \DIFdelFL{forced choice in both orders, }\DIFdelendFL \DIFaddbeginFL \DIFaddFL{label direction }\DIFaddendFL and \DIFdelbeginFL \DIFdelFL{a report continuation. Bottom: }\DIFdelendFL the \DIFdelbeginFL \DIFdelFL{ownership contrast added to a regression of the endpoint effect, cross-validated by concept and by block}\DIFdelendFL \DIFaddbeginFL \DIFaddFL{answer direction}\DIFaddendFL .}
\label{tab:cf-semend}
\DIFdelbeginFL %DIFDELCMD < \fitwidth{%
%DIFDELCMD < \begin{tabular}{lllcccc}
%DIFDELCMD < \toprule
%DIFDELCMD < \multirow{2}{*}{Checkpoint} & \multirow{2}{*}{Dataset} & \multirow{2}{*}{Endpoint} & \multicolumn{2}{c}{label direction} & \multicolumn{2}{c}{answer direction}\\
%DIFDELCMD < \cmidrule(lr){4-5}\cmidrule(lr){6-7}
%DIFDELCMD < &  &  & \multicolumn{1}{c}{owned} & \multicolumn{1}{c}{reference} & \multicolumn{1}{c}{owned} & \multicolumn{1}{c}{reference}\\
%DIFDELCMD < \midrule
%DIFDELCMD < Lingshu 32B & CheXpert Plus & negated question & 2/6 & 2/6 & 3/6 & 4/6 \\
%DIFDELCMD < Lingshu 32B & CheXpert Plus & forced choice, finding first & 3/6 & 3/6 & 4/6 & 4/6 \\
%DIFDELCMD < Lingshu 32B & CheXpert Plus & forced choice, competitor first & 4/6 & 4/6 & 5/6 & 5/6 \\
%DIFDELCMD < Lingshu 32B & CheXpert Plus & report continuation & 0/6 & 0/6 & 0/6 & 0/6 \\
%DIFDELCMD < Qwen2.5-VL-7B & NIH ChestX-ray14 & negated question & 0/6 & 1/6 & 3/6 & 6/6 \\
%DIFDELCMD < Qwen2.5-VL-7B & NIH ChestX-ray14 & forced choice, finding first & 0/6 & 0/6 & 3/6 & 3/6 \\
%DIFDELCMD < Qwen2.5-VL-7B & NIH ChestX-ray14 & forced choice, competitor first & 0/6 & 0/6 & 2/6 & 2/6 \\
%DIFDELCMD < Qwen2.5-VL-7B & NIH ChestX-ray14 & report continuation & 0/6 & 0/6 & 0/6 & 1/6 \\
%DIFDELCMD < Qwen2.5-VL-7B & CheXpert Plus & negated question & 1/6 & 1/6 & 5/6 & 6/6 \\
%DIFDELCMD < Qwen2.5-VL-7B & CheXpert Plus & forced choice, finding first & 0/6 & 0/6 & 3/6 & 4/6 \\
%DIFDELCMD < Qwen2.5-VL-7B & CheXpert Plus & forced choice, competitor first & 0/6 & 0/6 & 2/6 & 2/6 \\
%DIFDELCMD < Qwen2.5-VL-7B & CheXpert Plus & report continuation & 0/6 & 0/6 & 0/6 & 0/6 \\
%DIFDELCMD < \midrule
%DIFDELCMD < \textit{all} & \textit{3 blocks} & \textit{4 endpoints} & 10/72 & 11/72 & 30/72 & 37/72 \\
%DIFDELCMD < \midrule
%DIFDELCMD < \multicolumn{7}{l}{\textit{The three endpoint-specific tests}} \\
%DIFDELCMD < Checkpoint & Dataset & opposite on the negation & order gap & forced choice owned & report owned & core owned \\
%DIFDELCMD < \midrule
%DIFDELCMD < Lingshu 32B & CheXpert Plus & 4/6 & 0.026 & 3/6 & 0/6 & 5/6 \\
%DIFDELCMD < Qwen2.5-VL-7B & NIH ChestX-ray14 & 6/6 & 0.089 & 0/6 & 0/6 & 0/6 \\
%DIFDELCMD < Qwen2.5-VL-7B & CheXpert Plus & 6/6 & 0.111 & 0/6 & 0/6 & 0/6 \\
%DIFDELCMD < \midrule
%DIFDELCMD < \multicolumn{7}{l}{\textit{What ownership adds, held out: leaving out one concept at a time}} \\
%DIFDELCMD < Checkpoint & Dataset & Direction & held-out $R^2$ & with ownership & $\Delta R^2$ & permutation $p$ \\
%DIFDELCMD < \midrule
%DIFDELCMD < Lingshu 32B & CheXpert Plus & label direction & -0.034 & 0.147 & 0.181 & 0.39 \\
%DIFDELCMD < Lingshu 32B & CheXpert Plus & answer direction & 0.069 & -0.223 & -0.291 & 0.70 \\
%DIFDELCMD < Qwen2.5-VL-7B & CheXpert Plus & label direction & -0.139 & -0.149 & -0.010 & 0.55 \\
%DIFDELCMD < Qwen2.5-VL-7B & CheXpert Plus & answer direction & -3.168 & -0.112 & 3.057 & 0.18 \\
%DIFDELCMD < Qwen2.5-VL-7B & NIH ChestX-ray14 & label direction & -0.195 & -0.147 & 0.048 & 0.31 \\
%DIFDELCMD < Qwen2.5-VL-7B & NIH ChestX-ray14 & answer direction & -1.310 & -0.176 & 1.134 & 0.50 \\
%DIFDELCMD < \textit{pooled} & \textit{6 tests} & \textit{both} & -- & -- & 0.687 & 0.26 \\
%DIFDELCMD < \midrule
%DIFDELCMD < \multicolumn{7}{l}{\textit{The same, leaving out one block at a time}} \\
%DIFDELCMD < \midrule
%DIFDELCMD < \textit{all 3} & \textit{pooled} & label direction & -0.578 & -0.500 & 0.078 & 0.30 \\
%DIFDELCMD < \textit{all 3} & \textit{pooled} & answer direction & -0.461 & -0.398 & 0.063 & 0.26 \\
%DIFDELCMD < \textit{pooled} & \textit{2 tests} & \textit{both} & -- & -- & 0.070 & 0.26 \\
%DIFDELCMD < \midrule
%DIFDELCMD < \multicolumn{7}{l}{\textit{In-sample reference only, since adding a predictor cannot lower a training $R^2$}} \\
%DIFDELCMD < Checkpoint & Dataset & Direction & training $R^2$ & with ownership & $\Delta R^2$ & \\
%DIFDELCMD < \midrule
%DIFDELCMD < Lingshu 32B & CheXpert Plus & label direction & 0.769 & 0.788 & 0.019 & \\
%DIFDELCMD < Lingshu 32B & CheXpert Plus & answer direction & 0.368 & 0.423 & 0.055 & \\
%DIFDELCMD < Qwen2.5-VL-7B & NIH ChestX-ray14 & label direction & 0.241 & 0.246 & 0.004 & \\
%DIFDELCMD < Qwen2.5-VL-7B & NIH ChestX-ray14 & answer direction & 0.267 & 0.392 & 0.125 & \\
%DIFDELCMD < Qwen2.5-VL-7B & CheXpert Plus & label direction & 0.175 & 0.250 & 0.074 & \\
%DIFDELCMD < Qwen2.5-VL-7B & CheXpert Plus & answer direction & 0.473 & 0.509 & 0.036 & \\
%DIFDELCMD < \bottomrule
%DIFDELCMD < \end{tabular}}
%DIFDELCMD < %%%
\DIFdelendFL \DIFaddbeginFL \resizebox{\linewidth}{!}{%
\begin{tabular}{lllcccc}
\toprule
\multirow{2}{*}{Checkpoint} & \multirow{2}{*}{Dataset} & \multirow{2}{*}{Endpoint} & \multicolumn{2}{c}{label direction} & \multicolumn{2}{c}{answer direction}\\
\cmidrule(lr){4-5}\cmidrule(lr){6-7}
&  &  & \multicolumn{1}{c}{owned} & \multicolumn{1}{c}{reference} & \multicolumn{1}{c}{owned} & \multicolumn{1}{c}{reference}\\
\midrule
Lingshu 32B & CheXpert Plus & negated question & 2/6 & 2/6 & 3/6 & 4/6 \\
Lingshu 32B & CheXpert Plus & forced choice, finding first & 3/6 & 3/6 & 4/6 & 4/6 \\
Lingshu 32B & CheXpert Plus & forced choice, competitor first & 4/6 & 4/6 & 5/6 & 5/6 \\
Lingshu 32B & CheXpert Plus & report continuation & 0/6 & 0/6 & 0/6 & 0/6 \\
Qwen2.5-VL-7B & NIH ChestX-ray14 & negated question & 0/6 & 1/6 & 3/6 & 6/6 \\
Qwen2.5-VL-7B & NIH ChestX-ray14 & forced choice, finding first & 0/6 & 0/6 & 3/6 & 3/6 \\
Qwen2.5-VL-7B & NIH ChestX-ray14 & forced choice, competitor first & 0/6 & 0/6 & 2/6 & 2/6 \\
Qwen2.5-VL-7B & NIH ChestX-ray14 & report continuation & 0/6 & 0/6 & 0/6 & 1/6 \\
Qwen2.5-VL-7B & CheXpert Plus & negated question & 1/6 & 1/6 & 5/6 & 6/6 \\
Qwen2.5-VL-7B & CheXpert Plus & forced choice, finding first & 0/6 & 0/6 & 3/6 & 4/6 \\
Qwen2.5-VL-7B & CheXpert Plus & forced choice, competitor first & 0/6 & 0/6 & 2/6 & 2/6 \\
Qwen2.5-VL-7B & CheXpert Plus & report continuation & 0/6 & 0/6 & 0/6 & 0/6 \\
\midrule
\textit{all} & \textit{3 blocks} & \textit{4 endpoints} & 10/72 & 11/72 & 30/72 & 37/72 \\
\bottomrule
\end{tabular}}
\DIFaddendFL \end{table}

\DIFaddbegin \subsection{\DIFadd{Difficulty}}
\label{app:cf-difficulty}

\begin{DIFnomarkup}\cfnote{附录精简}{细粒度物体的选择规则只保留在 F.6 一处，这里只引用；两个匹配估计量写成公式，代替原来各两三句的文字定义；删去描述分档实现细节的句子。}\end{DIFnomarkup}

\DIFaddend \paragraph{Fine-grained objects\DIFdelbegin \DIFdel{and difficulty matching}\DIFdelend .} \DIFdelbegin \DIFdel{The rule that picks the six additional COCO categories , and the two matching windows, are }\DIFdelend \DIFaddbegin \DIFadd{Six further COCO categories test whether ownership on COCO depends on large, easy objects. A rule }\DIFaddend stated in full in Appendix~\ref{app:repro-protocol} \DIFdelbegin \DIFdel{. The rule reads }\DIFdelend \DIFaddbegin \DIFadd{picks them from }\DIFaddend the label counts and \DIFdelbegin \DIFdel{the }\DIFdelend mean instance areas of the fixed cohorts \DIFdelbegin \DIFdel{and nothing else; the registered protocol file carries it under the date of amendment A15, and the }\DIFdelend \DIFaddbegin \DIFadd{alone, and the protocol file registers the rule and the two matching windows before the }\DIFaddend module's first outcome\DIFdelbegin \DIFdel{shard is later than that date, so the categories and the windows were fixed before any of them was scored and before any ownership outcome of theirs existed. We fit, lift, write, and grade these categories exactly }\DIFdelend \DIFaddbegin \DIFadd{. The categories are fitted, written, and graded }\DIFaddend like the six easy objects, \DIFaddbegin \DIFadd{as their own family }\DIFaddend in the same blocks and \DIFdelbegin \DIFdel{on the same rows, as their own six-direction family, against a random family and a sham of their own (Table~\ref{tab:cf-ledger}). We measure their }\DIFdelend \DIFaddbegin \DIFadd{rows, with their }\DIFaddend clean-answer AUROC \DIFaddbegin \DIFadd{measured }\DIFaddend on the same calibration rows\DIFdelbegin \DIFdel{and by the same rule as every other cell, giving the matching variable the same meaning on both sides.
The matched comparison pairs each chest cell with every }\DIFdelend \DIFaddbegin \DIFadd{. Within a checkpoint the owned counts of fine-grained and easy objects differ by $\cfFgSameDiffMin$ to $\cfFgSameDiffMax$ of six, and median probe selectivity is }\cfFgFineSel{} \DIFadd{for the fine-grained objects against }\cfFgEasySameSel{} \DIFadd{for the easy ones.
}

\begin{DIFnomarkup}\cfnote{删表}{Fine-grained 表的匹配估计量此前只在表里：按 clean-answer AUROC 匹配后差异缩小到 $-0.066$，两个变量同时匹配时区间包含 0。这是限定“难度不是原因”这一论点的结果，现在既在 Table 14 中，也在文字里保留（全部用现成的宏）；两个估计量写成公式，代替原来各两三句的文字定义。}\end{DIFnomarkup}

\paragraph{\DIFadd{Difficulty matching.}} \DIFadd{Each chest cell $c$ is paired with the }\DIFaddend natural-image \DIFdelbegin \DIFdel{cell }\DIFdelend \DIFaddbegin \DIFadd{cells $\mathcal P(c)$ }\DIFaddend whose matching variables lie within \cfFgWindow{} of \DIFdelbegin \DIFdel{it. It reports the ownership rate on each side of the surviving pairs under two estimators. The }\emph{\DIFdel{pooled}} %DIFAUXCMD
\DIFdel{difference is the ownership rate among the matched chest cells minus the rate among the natural-image cells that served as a partner at least once. The }\emph{\DIFdel{matched}} %DIFAUXCMD
\DIFdel{difference is the mean over matched chest cells of that cell's own ownership minus the ownership rate of its own partner set, so a chest cell with many partners and one with few count equally. Neither estimator }\DIFdelend \DIFaddbegin \DIFadd{its own. With $o\in\{0,1\}$ a cell's owned indicator, $\mathcal C$ the chest cells with at least one partner, and $\mathcal P=\bigcup_{c\in\mathcal C}\mathcal P(c)$,
}\begin{equation}
\Delta_{\text{pool}}=\frac{1}{|\mathcal C|}\sum_{c\in\mathcal C}o_c-\frac{1}{|\mathcal P|}\sum_{p\in\mathcal P}o_p,\qquad
\Delta_{\text{match}}=\frac{1}{|\mathcal C|}\sum_{c\in\mathcal C}\Bigl(o_c-\frac{1}{|\mathcal P(c)|}\sum_{p\in\mathcal P(c)}o_p\Bigr);
\label{eq:matching}
\end{equation}
\DIFadd{$\Delta_{\text{pool}}$ compares the two matched pools, while $\Delta_{\text{match}}$ weights every chest cell equally however many partners it has; neither }\DIFaddend can be obtained from the other. \DIFdelbegin \DIFdel{The }\emph{\DIFdel{chest cells}} %DIFAUXCMD
\DIFdel{and }\emph{\DIFdel{partners}} %DIFAUXCMD
\DIFdel{columns give the effective sample of both: matched chest cells of all chest cells, and distinct natural-image cells used as a partner of all of them. Both intervals are percentile intervals over }\DIFdelend \DIFaddbegin \DIFadd{Chest cells without a partner leave the matching, which the stratified comparison below avoids. Intervals come from }\DIFaddend 5{,}000 draws of a cluster bootstrap \DIFdelbegin \DIFdel{that resamples models , keeping each model's chest and natural-image cellstogether. Chest cells without a partner leave the comparison. }\DIFdelend \DIFaddbegin \DIFadd{over models (Table~\ref{tab:cf-matching}). Matched on clean-answer AUROC, $\Delta_{\text{match}}$ shrinks to $\cfFgMatchAnsDiff$ $[\cfFgMatchAnsLo,\cfFgMatchAnsHi]$, and matched on both variables, which keeps }\cfFgMatchBothN{} \DIFadd{of the }\cfFgMatchBothCells{} \DIFadd{chest cells, its interval $[\cfFgMatchBothLo,\cfFgMatchBothHi]$ includes zero. }\DIFaddend The stratified comparison \DIFdelbegin \DIFdel{answers the same question without dropping a cell. It bins every cell of }\DIFdelend \DIFaddbegin \DIFadd{of Section~\ref{sec:coco} keeps every cell instead, binning }\DIFaddend both pools by \DIFdelbegin \DIFdel{its }\DIFdelend clean-answer AUROC at the edges \cfStratEdges{}\DIFdelbegin \DIFdel{, fixed in the released analysis code, and reports the ownership rate inside each bin with the same cluster bootstrap over models on the difference in the top bin. The last bin is closed on the right and every cell of each pool falls in exactly one bin. The }\DIFdelend \DIFaddbegin \DIFadd{; the }\DIFaddend \cfStratChestTopN{} chest cells of \DIFdelbegin \DIFdel{that bin are }\DIFdelend \DIFaddbegin \DIFadd{the top bin are all }\DIFaddend CheXpert Plus cells\DIFdelbegin \DIFdel{; }\DIFdelend \DIFaddbegin \DIFadd{, since }\DIFaddend no NIH ChestX-ray14 cell reaches \DIFdelbegin \DIFdel{a clean-answer AUROC of }\DIFdelend \cfStratTopEdge{}.
\DIFdelbegin \DIFdel{The same stratification of the nine-direction attribute family appears in Table~\ref{tab:cf-attr}. Table~\ref{tab:cf-fgobj} reports the per-block counts, the three pools, each matching, and every bin.
}\DIFdelend 

\begin{table}[htbp]
\centering
\DIFdelbeginFL %DIFDELCMD < \setlength{\tabcolsep}{4pt}
%DIFDELCMD < \setlength{\abovecaptionskip}{2pt}
%DIFDELCMD < \setlength{\belowcaptionskip}{0pt}
%DIFDELCMD < \renewcommand{\arraystretch}{0.95}
%DIFDELCMD < %%%
\DIFdelendFL \DIFaddbeginFL \setlength{\tabcolsep}{8pt}
\DIFaddendFL \caption{\DIFdelbeginFL \textbf{\DIFdelFL{Small and fine-grained objects, and difficulty matching.}} %DIFAUXCMD
\DIFdelFL{Top and middle}\DIFdelendFL \DIFaddbeginFL \textbf{\DIFaddFL{Difficulty matching.}} \DIFaddFL{Per matching variable}\DIFaddendFL : \DIFdelbeginFL \DIFdelFL{six fine-grained COCO categories graded like the six easy objects, against the easy objects of every COCO block and of the 6 blocks with fine-grained cells. Bottom: }\DIFdelendFL \DIFaddbeginFL \DIFaddFL{matched }\DIFaddendFL chest cells\DIFdelbeginFL \DIFdelFL{matched to }\DIFdelendFL \DIFaddbeginFL \DIFaddFL{, }\DIFaddendFL natural-image \DIFaddbeginFL \DIFaddFL{partner }\DIFaddendFL cells\DIFdelbeginFL \DIFdelFL{on difficulty}\DIFdelendFL , and the \DIFdelbeginFL \DIFdelFL{same cells stratified on clean-answer AUROC}\DIFdelendFL \DIFaddbeginFL \DIFaddFL{pooled and matched differences of Equation~\ref{eq:matching} with 95\% cluster-bootstrap intervals over models}\DIFaddendFL .}
\DIFdelbeginFL %DIFDELCMD < \label{tab:cf-fgobj}
%DIFDELCMD < \fitwidth{%
%DIFDELCMD < \begin{tabular}{lrrrrrr}
%DIFDELCMD < \toprule
%DIFDELCMD < \multirow{2}{*}{Checkpoint} & \multicolumn{2}{c}{owned of six} & \multicolumn{2}{c}{median answer AUROC} & \multicolumn{2}{c}{median selectivity}\\
%DIFDELCMD < \cmidrule(lr){2-3}\cmidrule(lr){4-5}\cmidrule(lr){6-7}
%DIFDELCMD < & \multicolumn{1}{c}{fine} & \multicolumn{1}{c}{easy} & \multicolumn{1}{c}{fine} & \multicolumn{1}{c}{easy} & \multicolumn{1}{c}{fine} & \multicolumn{1}{c}{easy}\\
%DIFDELCMD < \midrule
%DIFDELCMD < Qwen2.5-VL-3B & 6 & 6 & 0.947 & 0.971 & 0.024 & 0.037 \\
%DIFDELCMD < Qwen3-VL-4B & 5 & 6 & 0.984 & 0.990 & 0.036 & 0.033 \\
%DIFDELCMD < InternVL3.5-8B & 6 & 6 & 0.981 & 0.978 & 0.104 & 0.117 \\
%DIFDELCMD < Gemma 3 4B & 4 & 5 & 0.948 & 0.924 & 0.076 & 0.070 \\
%DIFDELCMD < Lingshu 7B & 6 & 6 & 0.950 & 0.968 & 0.025 & 0.023 \\
%DIFDELCMD < LLaVA-1.5-7B & 3 & 5 & 0.956 & 0.979 & 0.266 & 0.285 \\
%DIFDELCMD < \midrule
%DIFDELCMD < \multicolumn{7}{l}{\textit{Paired within these blocks:} fine-grained minus easy is +0, -1, +0, -1, -2, +0 cells of six; fine higher in 0,} \\
%DIFDELCMD < \multicolumn{7}{l}{easy higher in 3, tied in 3; exact two-sided sign test $p=0.25$} \\
%DIFDELCMD < \midrule
%DIFDELCMD < \multicolumn{7}{l}{\textit{The four pools of cells}} \\
%DIFDELCMD < Cells & blocks & owned & rate & median answer AUROC & median selectivity & readable \\
%DIFDELCMD < \midrule
%DIFDELCMD < chest findings & 50 & 38/300 & 0.127 & 0.695 & 0.137 & 178/300 \\
%DIFDELCMD < easy objects, every COCO block & 25 & 113/150 & 0.753 & 0.977 & 0.064 & 90/150 \\
%DIFDELCMD < easy objects, these blocks & 6 & 34/36 & 0.944 & 0.980 & 0.065 & 24/36 \\
%DIFDELCMD < fine-grained objects, these blocks & 6 & 30/36 & 0.833 & 0.963 & 0.058 & 23/36 \\
%DIFDELCMD < \midrule
%DIFDELCMD < \multicolumn{7}{l}{\textit{Chest cells matched to natural-image cells on difficulty, under both estimators}} \\
%DIFDELCMD < Matching & chest cells & partners & chest rate & natural rate & difference & 95\% interval \\
%DIFDELCMD < \midrule
%DIFDELCMD < probe selectivity alone & 300/300 & 186/186 & 0.127 & 0.769 &  &  \\
%DIFDELCMD < \quad pooled &  &  &  &  & -0.642 & [-0.763, -0.512] \\
%DIFDELCMD < \quad matched &  &  &  &  & -0.539 & [-0.678, -0.359] \\
%DIFDELCMD < clean-answer AUROC alone & 300/300 & 186/186 & 0.127 & 0.769 &  &  \\
%DIFDELCMD < \quad pooled &  &  &  &  & -0.642 & [-0.749, -0.505] \\
%DIFDELCMD < \quad matched &  &  &  &  & -0.066 & [-0.304, -0.004] \\
%DIFDELCMD < both variables & 212/300 & 101/186 & 0.137 & 0.624 &  &  \\
%DIFDELCMD < \quad pooled &  &  &  &  & -0.487 & [-0.653, -0.201] \\
%DIFDELCMD < \quad matched &  &  &  &  & -0.024 & [-0.198, 0.076] \\
%DIFDELCMD < both variables, easy partners only & 211/300 & 79/150 & 0.137 & 0.570 &  &  \\
%DIFDELCMD < \quad pooled &  &  &  &  & -0.432 & [-0.625, -0.150] \\
%DIFDELCMD < \quad matched &  &  &  &  & -0.008 & [-0.198, 0.085] \\
%DIFDELCMD < \midrule
%DIFDELCMD < \multicolumn{7}{l}{\textit{Chest cells and natural-image cells stratified on clean-answer AUROC, owned of cells}} \\
%DIFDELCMD < Cells & $<$ 0.50 & 0.50--0.70 & 0.70--0.80 & 0.80--0.90 & $\geq$ 0.90 & all cells \\
%DIFDELCMD < \midrule
%DIFDELCMD < NIH ChestX-ray14 & 0/7 & 6/80 & 3/43 & 4/19 & 0/1 & 13/150 \\
%DIFDELCMD < CheXpert Plus & 0/13 & 6/52 & 2/20 & 6/29 & 11/36 & 25/150 \\
%DIFDELCMD < \textit{chest findings} & 0/20 & 12/132 & 5/63 & 10/48 & 11/37 & 38/300 \\
%DIFDELCMD < easy objects, every COCO block & 0/6 & 0/6 & 0/7 & 1/4 & 112/127 & 113/150 \\
%DIFDELCMD < fine-grained objects, these blocks & 0/0 & 0/0 & 0/0 & 1/1 & 29/35 & 30/36 \\
%DIFDELCMD < \textit{natural-image objects} & 0/6 & 0/6 & 0/7 & 2/5 & 141/162 & 143/186 \\
%DIFDELCMD < \multicolumn{7}{l}{\quad chest minus natural-image in the $\geq$ 0.90 bin: -0.573, 95\% interval [-0.763, -0.388]} \\
%DIFDELCMD < \bottomrule
%DIFDELCMD < \end{tabular}}
%DIFDELCMD < %%%
\DIFdelendFL \DIFaddbeginFL \label{tab:cf-matching}
\resizebox{\linewidth}{!}{%
\begin{tabular}{lrrcc}
\toprule
Matching & \multicolumn{1}{c}{chest cells} & \multicolumn{1}{c}{partners} & $\Delta_{\text{pool}}$ [95\% interval] & $\Delta_{\text{match}}$ [95\% interval] \\
\midrule
probe selectivity alone & 300/300 & 186/186 & $-0.642$ $[-0.763, -0.512]$ & $-0.539$ $[-0.678, -0.359]$ \\
clean-answer AUROC alone & 300/300 & 186/186 & $-0.642$ $[-0.749, -0.505]$ & $-0.066$ $[-0.304, -0.004]$ \\
both variables & 212/300 & 101/186 & $-0.487$ $[-0.653, -0.201]$ & $-0.024$ $[-0.198, 0.076]$ \\
both, easy partners only & 211/300 & 79/150 & $-0.432$ $[-0.625, -0.150]$ & $-0.008$ $[-0.198, 0.085]$ \\
\bottomrule
\end{tabular}}
\DIFaddendFL \end{table}

\subsection{Write structure across the grid}
\label{app:cf-figures}

\DIFaddbegin \begin{DIFnomarkup}\cfnote{公式}{write-flow 份额写成公式并直接给出三个数；与 $\kappa_d$ 的区别从四个短句并成一句。}\end{DIFnomarkup}

\DIFaddend Figure~\ref{fig:cf-w-structure} summarises the $6\times6$ write matrices across all completed checkpoints as the cell-wise median matrix for each dataset. Figure~\ref{fig:write-flow} reads the same matrices as flow\DIFdelbegin \DIFdel{: band width is }\DIFdelend \DIFaddbegin \DIFadd{, with band width }\DIFaddend proportional to $F_{q,d}$, the mean \DIFdelbegin \DIFdel{positive answer change across checkpoints , coloured bands reach the writer's own question, and grey bands reach other questions. The write-flow share is the total width of own-question bands divided by the total width of all bands, $\sum_d F_{d,d}/\sum_{q,d}F_{q,d}$: }%DIFDELCMD < \cfOwnShareNih%%%
\DIFdel{\% on NIH and }%DIFDELCMD < \cfOwnShareChex%%%
\DIFdel{\% on CheXpert, compared with }%DIFDELCMD < \cfOwnShareCoco%%%
\DIFdel{\% on COCO. The column-selectivity index }\DIFdelend \DIFaddbegin \DIFadd{over checkpoints of $\max(W_{q,d},0)$. The own-question share of that flow is
}\begin{equation}
\frac{\sum_d F_{d,d}}{\sum_{q,d}F_{q,d}}=\cfOwnShareNih\%\ \text{(NIH)},\quad \cfOwnShareChex\%\ \text{(CheXpert)},\quad \cfOwnShareCoco\%\ \text{(COCO)}.
\label{eq:flow-share}
\end{equation}
\DIFadd{It pools positive changes over checkpoints, whereas the column selectivity }\DIFaddend $\kappa_d$ \DIFdelbegin \DIFdel{in Table~\ref{tab:cf-validation} measures a different quantity. It retains }\DIFdelend \DIFaddbegin \DIFadd{of Equation~\ref{eq:column-selectivity} keeps }\DIFaddend the sign of \DIFdelbegin \DIFdel{a direction's own effect. Its denominator includes negative changes. It gives every cell equal weight. We report its median over cells (}%DIFDELCMD < \cfValColSelNih{}%%%
\DIFdel{, }%DIFDELCMD < \cfValColSelChex{}%%%
\DIFdel{, and }%DIFDELCMD < \cfValColSelCoco{}%%%
\DIFdel{)}\DIFdelend \DIFaddbegin \DIFadd{the own effect, counts negative changes, and is summarised per cell}\DIFaddend . The largest sinks on the chest sets\DIFdelbegin \DIFdel{are the Cardiomegaly question }\DIFdelend \DIFaddbegin \DIFadd{, questions that every direction moves, are Cardiomegaly }\DIFaddend on NIH and \DIFdelbegin \DIFdel{the }\DIFdelend Edema and Effusion \DIFdelbegin \DIFdel{questions }\DIFdelend on CheXpert.
\DIFdelbegin \DIFdel{These questions receive writes from every direction.
}\DIFdelend 

\begin{figure}[p]
    \centering
    \resizebox{\textwidth}{!}{\includegraphics{figures/figA1_write_structure.pdf}}
    \caption{\textbf{Write structure across the grid.} Cell-wise median of the write matrix $W_{q,d}$ over completed checkpoints on NIH ChestX-ray14, CheXpert Plus, and COCO. Figure~\ref{fig:write-flow} reads the same matrices as flow.}
    \label{fig:cf-w-structure}
\end{figure}

\subsection{Per-image examples}
\label{app:cf-examples}

Figure~\ref{fig:cf-examples} shows grades for individual test images, with $P(\mathrm{yes})$ in parentheses on every answer line. The first four chest examples are deficit cells of Qwen2.5-VL-7B, the next two are owned cells (Lingshu 32B Edema, Qwen2.5-VL-72B Atelectasis), and the natural-image examples are owned cells of Qwen2.5-VL-7B, Lingshu 32B, and MedGemma 4B. For a chest example, we pool test rows where the finding is labelled present, the probe reads it (probability above $0.5$), and the clean answer is No. For a natural-image example, we pool rows where the probe reads the label correctly and the clean answer is No. We split each pool by whether the concept's own write or the strongest competing write raises the answer more. Among rows where either write raises the answer by at least $0.2$, we show the row with the gap closest to the median. Each example's last line reports how many pool rows share its ordering, giving the example its own base rate.

\begin{figure}[p]
    \centering
    \resizebox{\textwidth}{!}{\includegraphics{figures/figA2_examples.pdf}}
    \caption{\textbf{Reading, answering, and writing on single images.} Per image: the question, the report label, the probe's reading, the clean answer, and the answers under the own write and the strongest competing write at $\alpha=+0.25$. Beside the probe's reading, a check or cross marks whether it agrees with the label; beside the own write, whether that write moves the answer more than the competitor. The grey line counts, among the block's test rows the probe reads correctly and the model answers no when clean (on the chest sets, rows labelled present), those with the same ordering.}
    \label{fig:cf-examples}
\end{figure}
\DIFaddbegin \begin{DIFnomarkup}\cfnote{附录精简}{pilot 协议里的选择性、lift、写入、$O_q$ 四个公式与正文 Eq.~1–4 重复，且全文无引用，改为引用正文，只写 pilot 与正式评估不同的地方；原 decision rule 写成公式。max-$T$ 公式移到附录 E，作为唯一定义处。本节由约 1,430 词减到约 750 词。}\end{DIFnomarkup}

\DIFaddend \section{Pilot study: registered protocol}
\label{app:seed}

The two-model study that motivated the evaluation registered the procedures below; Appendix~\ref{app:seed-results} reports its results. \DIFaddbegin \DIFadd{It uses the construction of Section~\ref{sec:framework}, with the differences stated here.
}\DIFaddend 

\paragraph{Data\DIFaddbegin \DIFadd{, models, }\DIFaddend and \DIFdelbegin \DIFdel{split}\DIFdelend \DIFaddbegin \DIFadd{reading}\DIFaddend .}
We use NIH ChestX-ray14 \citep{wang2017chestxray}, whose findings are mined from radiology reports. A deterministic hash of patient identity assigns rows to disjoint train, validation, and test partitions\DIFdelbegin \DIFdel{. The registered manifest contains }\DIFdelend \DIFaddbegin \DIFadd{: }\DIFaddend 18,212 training images and 5,343 test images from 1,659 test patients. The three cells pair LLaVA-1.5-7B \DIFaddbegin \DIFadd{\mbox{%DIFAUXCMD
\citep{liu2024improved} }\hskip0pt%DIFAUXCMD
}\DIFaddend with Effusion and Edema, and Qwen2.5-VL-7B \DIFdelbegin \DIFdel{with Effusion. All reported uncertainty clusters by patient. Prompts appear in Appendix~\ref{app:protocol}.
}%DIFDELCMD < 

%DIFDELCMD < %%%
\paragraph{\DIFdel{Models and consumed loci.}}
%DIFAUXCMD
\addtocounter{paragraph}{-1}%DIFAUXCMD
\DIFdel{For LLaVA-1.5-7B \mbox{%DIFAUXCMD
\citep{liu2024improved} }\hskip0pt%DIFAUXCMD
and Qwen2.5-VL-7B \mbox{%DIFAUXCMD
\citep{bai2025qwen25vl}}\hskip0pt%DIFAUXCMD
, we use }\DIFdelend \DIFaddbegin \DIFadd{\mbox{%DIFAUXCMD
\citep{bai2025qwen25vl} }\hskip0pt%DIFAUXCMD
with Effusion, at }\DIFaddend the final visual block consumed by the connector or merger\DIFdelbegin \DIFdel{. Writing $\alpha$ for the signed intervention dose, synthetic hook tests establish downstream-path membership, a bitwise }\DIFdelend \DIFaddbegin \DIFadd{; hook tests confirm that this block lies on the downstream path, that }\DIFaddend $\alpha=0$ \DIFaddbegin \DIFadd{is a bitwise }\DIFaddend no-op, and \DIFdelbegin \DIFdel{changed downstream representations and logitsat nonzero dose. Probe and intervention therefore share one consumed representation.
}%DIFDELCMD < 

%DIFDELCMD < %%%
\paragraph{\DIFdel{Controlled decodability.}}
%DIFAUXCMD
\addtocounter{paragraph}{-1}%DIFAUXCMD
\DIFdel{For image $i$, let $\hvec_{it}\in\R^D$ denote the $D$-dimensional representation of visual token $t$ at the consumed final block; mean pooling over tokens yields $\hvec_i$. A fixed seed-0 Gaussian matrix $R\in\R^{D\times512}$ projects every model to a common readout dimension; featurewise scale $\boldsymbol{s}$ is fitted on training rows only. We then fit logistic regression with inverse regularisation strength $C=1$, maximum 2,000 iterations, and }\DIFdelend \DIFaddbegin \DIFadd{that nonzero doses change downstream logits. Projection, scaling, and probe follow Section~\ref{sec:directions} with }\DIFaddend seed 0. \DIFdelbegin \DIFdel{Each of }\DIFdelend \DIFaddbegin \DIFadd{Selectivity is Equation~\ref{eq:selectivity} on test rows, with }\DIFaddend 20 control seeds \DIFdelbegin \DIFdel{maps the recurring input type }\DIFdelend \DIFaddbegin \DIFadd{mapping the 39 recurring input types }\DIFaddend (AP view\DIFdelbegin \DIFdel{indicator, sexindicator}\DIFdelend , \DIFaddbegin \DIFadd{sex, }\DIFaddend and age decade\DIFdelbegin \DIFdel{, giving 39 strata) to a random label. Let $y$ be the real test labels and $r$ the fitted probe's continuous scores. For control seed $k\in\{0,\ldots,19\}$, $y_k^{\mathrm{ctrl}}$ denotes the random labels and $r_k^{\mathrm{ctrl}}$ the corresponding control probe scores. Selectivity compares their AUROCs:
}%DIFDELCMD < \begin{equation}
%DIFDELCMD < S = \operatorname{AUROC}(y,r)-\frac{1}{20}\sum_{k=0}^{19}\operatorname{AUROC}(y_k^{\mathrm{ctrl}},r_k^{\mathrm{ctrl}}).
%DIFDELCMD < \label{eq:seed-selectivity}
%DIFDELCMD < \end{equation}
%DIFDELCMD < %%%
\DIFdel{We report the real AUROC, control distribution, and $S$ with }\DIFdelend \DIFaddbegin \DIFadd{) to random labels and }\DIFaddend 2,000 patient-cluster bootstrap \DIFdelbegin \DIFdel{resamples. A positive interval for $S$ means that the target is more readable than same-capacity }\DIFdelend \DIFaddbegin \DIFadd{draws. Selectivity compares the concept with }\DIFaddend nuisance-type label maps \DIFdelbegin \DIFdel{. It does not yet imply downstream use.
Paired contrasts reuse identical ordered test rows and bootstrap draws.
}\DIFdelend \DIFaddbegin \DIFadd{of the same capacity; the write tests whether the model uses it. All uncertainty clusters by patient; prompts appear in Appendix~\ref{app:protocol}.
}\DIFaddend 

\paragraph{\DIFdelbegin \DIFdel{Probe-normal intervention}\DIFdelend \DIFaddbegin \DIFadd{Write and doses}\DIFaddend .}
\DIFdelbegin \DIFdel{Let $\boldsymbol\beta_c$ be the projected-space logistic coefficient for the positive class of concept $c$. We lift it to the consumed $D$-dimensional representation and normalize it exactly as
}%DIFDELCMD < \begin{equation}
%DIFDELCMD < \hat\wvec_c=\frac{\boldsymbol g_c}{\lVert\boldsymbol g_c\rVert_2},\qquad \boldsymbol g_c=R\,\mathrm{diag}(\boldsymbol s)^{-1}\boldsymbol\beta_c ,
%DIFDELCMD < \label{eq:seed-lift}
%DIFDELCMD < \end{equation}
%DIFDELCMD < %%%
\DIFdelend \DIFaddbegin \DIFadd{The write is Equation~\ref{eq:intervention} along the lifted normal of Equation~\ref{eq:lift}, }\DIFaddend with every entry of $\boldsymbol s$ floored at $10^{-8}$\DIFdelbegin \DIFdel{. The positive-class coefficient fixes the direction's clinical pole. At every visual token $t$, the registered intervention is
}%DIFDELCMD < \begin{equation}
%DIFDELCMD < \hvec'_{it}=\hvec_{it}+\alpha\lVert\hvec_{it}\rVert_2\hat\wvec_c.
%DIFDELCMD < \label{eq:seed-intervention}
%DIFDELCMD < \end{equation}
%DIFDELCMD < %%%
\DIFdel{Here the sign of $\alpha$ selects the positive or negative clinical pole, while $|\alpha|$ is the perturbation norm as a fraction of each token's activation norm. At the first answer position, $P(\mathrm{yes})=\sigma(\ell_{\mathrm{yes}}-\ell_{\mathrm{no}})$, where each log probability is the maximum over case and leading-space variants that tokenize to one token; }\DIFdelend \DIFaddbegin \DIFadd{, and }\DIFaddend the endpoint is \DIFdelbegin \DIFdel{its }\DIFdelend \DIFaddbegin \DIFadd{the }\DIFaddend paired mean change \DIFaddbegin \DIFadd{in $P(\mathrm{yes})$ }\DIFaddend from the common unsteered baseline \DIFaddbegin \DIFadd{(Section~\ref{sec:ownership})}\DIFaddend . Original cells use \DIFdelbegin \DIFdel{concept }\DIFdelend \DIFaddbegin \DIFadd{the }\DIFaddend doses $\{-1,-0.5,-0.25,-0.1,0,0.1,0.25,0.5,1\}$ on 200 held-out test images. \DIFdelbegin \DIFdel{Controls are evaluated at the six nonzero doses excluding $\pm0.1$: }\DIFdelend \DIFaddbegin \DIFadd{At the six controlled doses $\mathcal A=\{\pm1,\pm0.5,\pm0.25\}$, }\DIFaddend 20 isotropic random unit directions, one \DIFdelbegin \DIFdel{sham formed by a fixed coordinate permutation of $\hat\wvec_c$ and renormalization, and }\DIFdelend \DIFaddbegin \DIFadd{coordinate-permutation sham, and the }\DIFaddend fixed alternative clinical \DIFdelbegin \DIFdel{probe normals . Every control receives }\DIFdelend \DIFaddbegin \DIFadd{normals receive }\DIFaddend the same signed dose on the same rows. Dose monotonicity over $[-0.5,0.5]$ is descriptive.

\paragraph{Original-cell decision rule.}
\DIFdelbegin \DIFdel{Let }\DIFdelend \DIFaddbegin \DIFadd{With }\DIFaddend $\Delta_c(\alpha)$ \DIFdelbegin \DIFdel{be }\DIFdelend the target direction's paired mean \DIFdelbegin \DIFdel{probability change , and let $\mathcal A=\{-1,-0.5,-0.25,0.25,0.5,1\}$ be the six controlled doses. The primary effect is $\max_{\alpha\in\mathcal A}\operatorname{sign}(\alpha)\Delta_c(\alpha)$, attained at dose $\alpha^\star$. }\DIFdelend \DIFaddbegin \DIFadd{change at dose $\alpha$, the primary effect and the dose that attains it are
}\begin{equation}
\Delta^\star_c=\max_{\alpha\in\mathcal A}\operatorname{sign}(\alpha)\,\Delta_c(\alpha),\qquad
\alpha^\star=\operatorname*{arg\,max}_{\alpha\in\mathcal A}\operatorname{sign}(\alpha)\,\Delta_c(\alpha).
\label{eq:seed-rule}
\end{equation}
\DIFaddend A cell clears the generic \DIFdelbegin \DIFdel{intervention gate when this selected effect exceeds }\DIFdelend \DIFaddbegin \DIFadd{gate when $\Delta^\star_c$ exceeds both }\DIFaddend the 95th percentile of the 20 random directions' \DIFdelbegin \DIFdel{mean probability changes at that same }\DIFdelend \DIFaddbegin \DIFadd{changes at }\DIFaddend $\alpha^\star$, each multiplied by $\operatorname{sign}(\alpha^\star)$, and the \DIFdelbegin \DIFdel{maximum }\DIFdelend \DIFaddbegin \DIFadd{largest }\DIFaddend absolute sham effect over $\mathcal A$. The random directions are \DIFdelbegin \DIFdel{not separately maximized over doses. This asymmetric selection defines }\DIFdelend \DIFaddbegin \DIFadd{evaluated at $\alpha^\star$ alone, so the gate is }\DIFaddend a registered reference \DIFdelbegin \DIFdel{gate, not a cross-dose-corrected 5\% test.
The conjunction is fixed across the three cells. The independent supplement below instead evaluates a dose locked before observing its patients.
}\DIFdelend \DIFaddbegin \DIFadd{without cross-dose correction.
}\DIFaddend 

\paragraph{Direction-specificity supplement.}
The \DIFdelbegin \DIFdel{initial experiment prespecified the random and sham comparisons; its clinical-direction analysis was descriptive. The }\DIFdelend Nodule response motivated a separately registered supplement at the locked Qwen dose \DIFdelbegin \DIFdel{, }\DIFdelend $\alpha=+0.25$\DIFdelbegin \DIFdel{. We select one row }\DIFdelend \DIFaddbegin \DIFadd{, on one row per patient, chosen }\DIFaddend by label-blind hash\DIFdelbegin \DIFdel{for each of }\DIFdelend \DIFaddbegin \DIFadd{, for }\DIFaddend 400 ChestX-ray14 test patients absent from the original \DIFdelbegin \DIFdel{200-image intervention cohort. The semantic family contains all five non-target normals from the }\DIFdelend \DIFaddbegin \DIFadd{cohort. It estimates the Effusion row of the write matrix (Section~\ref{sec:ownership}) against the five other }\DIFaddend pre-existing Qwen \DIFdelbegin \DIFdel{direction bundle }\DIFdelend \DIFaddbegin \DIFadd{normals }\DIFaddend (Atelectasis, Pneumothorax, Cardiomegaly, Mass, and Nodule), \DIFdelbegin \DIFdel{rather than directions chosen from the supplement outcome. Each uses the same training rows, projection, scaling, classifier, lifting rule, and normalization as Effusion. Together with a common baseline, }\DIFdelend \DIFaddbegin \DIFadd{fitted like Effusion, together with }\DIFaddend 20 random directions \DIFdelbegin \DIFdel{, }\DIFdelend and the sham\DIFdelbegin \DIFdel{, the grid produces }\DIFdelend \DIFaddbegin \DIFadd{: }\DIFaddend 28 configurations and 11,200 per-image outcomes\DIFdelbegin \DIFdel{; exact row selection appears in Appendix~\ref{app:seed-results}. }%DIFDELCMD < 

%DIFDELCMD < %%%
\DIFdel{Semantic ownership is a contrast on a direction--question response surface, not a property of a direction in isolation. Let $W_{q,d}$ denote the paired mean change in question $q$'s yes probability after intervening with direction $d$ at the locked dose. A large $W_{q,d}$ establishes that $d$ can write $q$'s answer; it does not establish that $d$ owns that answer. Within a fixed clinical family, the ownership contrast is
}%DIFDELCMD < \begin{equation}
%DIFDELCMD < O_q=W_{q,q}-\max_{d\neq q}W_{q,d}.
%DIFDELCMD < \label{eq:seed-margin}
%DIFDELCMD < \end{equation}
%DIFDELCMD < %%%
\DIFdel{The registered supplement estimates the Effusion row of this surface, so the reported ownership contrast }\DIFdelend \DIFaddbegin \DIFadd{. Direction specificity requires an Effusion effect above the random 95th percentile and the absolute sham, and a one-sided 95\% lower bound of }\DIFaddend $O_{\mathrm{Effusion}}$ \DIFdelbegin \DIFdel{is Effusion minus its maximum alternative. The maximum is recomputed inside }\DIFdelend \DIFaddbegin \DIFadd{(Equation~\ref{eq:margin}) above zero, with the maximum recomputed in }\DIFaddend each of 5,000 patient-bootstrap draws. \DIFdelbegin \DIFdel{Thus $O_{\mathrm{Effusion}}>0$ means Effusion beats every member of the fixed clinical family, while $O_{\mathrm{Effusion}}<0$ means at least one alternative is stronger. }\DIFdelend Adding competitors can only \DIFdelbegin \DIFdel{maintain or decrease }\DIFdelend \DIFaddbegin \DIFadd{lower }\DIFaddend $O_q$\DIFdelbegin \DIFdel{: }\DIFdelend \DIFaddbegin \DIFadd{, so }\DIFaddend a positive advantage \DIFdelbegin \DIFdel{is conditional on the candidate family , whereas a negative counterexample does not require an }\DIFdelend \DIFaddbegin \DIFadd{holds for the family tested, while a negative one needs no }\DIFaddend exhaustive family.
\DIFdelbegin \DIFdel{This operational ownership criterion measures relative advantage, not a necessary condition for a concept to participate in computation. Direction specificity requires a positive Effusion effect above random p95 and absolute sham, together with a one-sided 95\% lower bound for $O_{\mathrm{Effusion}}$ above zero.
}\DIFdelend 

\paragraph{Causal-ownership matrix.}
\DIFdelbegin \DIFdel{To test whether another clinical normal better matches the answer pathway, we prospectively cross }\DIFdelend \DIFaddbegin \DIFadd{A prospective grid crosses }\DIFaddend six Qwen directions (Effusion, Atelectasis, Pneumothorax, Cardiomegaly, Mass, and Nodule) with their six \DIFdelbegin \DIFdel{matching }\DIFdelend questions on 400 new patients \DIFdelbegin \DIFdel{. Every question uses the same }\DIFdelend \DIFaddbegin \DIFadd{at }\DIFaddend $\alpha=+0.25$, \DIFaddbegin \DIFadd{with }\DIFaddend 119 random directions \DIFdelbegin \DIFdel{, }\DIFdelend and a sham. \DIFdelbegin \DIFdel{For question }\DIFdelend \DIFaddbegin \DIFadd{Question }\DIFaddend $q$ \DIFdelbegin \DIFdel{, ownership requires the diagonal effect }\DIFdelend \DIFaddbegin \DIFadd{is owned when }\DIFaddend $W_{q,q}$ \DIFdelbegin \DIFdel{to exceed }\DIFdelend \DIFaddbegin \DIFadd{exceeds }\DIFaddend every other clinical direction, every random direction, and the absolute sham\DIFdelbegin \DIFdel{effect. A centered studentized max-$T$ bootstrap gives simultaneous 95\% lower bounds over all }\DIFdelend \DIFaddbegin \DIFadd{, with the }\DIFaddend 30 clinical contrasts \DIFdelbegin \DIFdel{. The shared-alias family contains two statistics per direction: mean excess over the matching diagonals }\DIFdelend \DIFaddbegin \DIFadd{bounded simultaneously by Equation~\ref{eq:max-t}. A separate family of 12 statistics tests shared aliasing: for each direction, its signed mean effect }\DIFaddend on its five off-diagonal questions and \DIFdelbegin \DIFdel{signed mean effect on those questions. These give 12 max-$T$ bounds; the effect must also exceed the random alias maximum and maximum absolute sham effect, detailed in Appendix~\ref{app:seed-results}}\DIFdelend \DIFaddbegin \DIFadd{its mean excess over the matching diagonals there}\DIFaddend .

\DIFdelbegin \DIFdel{For a family of $m$ statistics, let $\widehat C_j$ be statistic $j$'s estimate and $C^*_{bj}$ its value in joint patient-bootstrap draw $b$, with $B=5{,}000$ draws. Write $\tau_j=\operatorname{SD}_{b}(C^*_{bj})$ for the bootstrap standard deviation with denominator $B-1$, held fixed across draws, and $Q_{.95}$ for the empirical 95th percentile with linear interpolation. For $\tau_j>10^{-12}$, the centered standardized draw $Z_{bj}$, common critical value $c_{.95}$, and simultaneous lower bound $L_j$ are
}%DIFDELCMD < \begin{equation}
%DIFDELCMD < Z_{bj}=\frac{C^*_{bj}-\widehat C_j}{\tau_j},\qquad
%DIFDELCMD < c_{.95}=Q_{.95}\!\left(\left\{\max_{1\leq j\leq m}Z_{bj}\right\}_{b=1}^{B}\right),\qquad
%DIFDELCMD < L_j=\widehat C_j-c_{.95}\tau_j.
%DIFDELCMD < \label{eq:max-t}
%DIFDELCMD < \end{equation}
%DIFDELCMD < %%%
\DIFdel{Joint resampling preserves paired conditions; no standard error is re-estimated within a draw. Appendix~\ref{app:simultaneous-bounds} specifies the families, degenerate components, and the separate percentile intervals for target-minus-maximum margins.
}%DIFDELCMD < 

%DIFDELCMD < %%%
\DIFdelend \paragraph{Answer-encoding crossover and Mass confirmation.} A registered crossover on 200 patients of the causal-ownership cohort, with the other 200 as a patient-disjoint calibration set, crosses the original yes/no prompt with A-present and B-present prompts that exchange the finding-to-letter mapping, at $\alpha=\pm0.25$ with six clinical directions, 20 random directions, and a sham; a question enters only if both coded mappings have clean-answer AUROC lower bounds above 0.5 with at least ten patients per class. An independent confirmation then tests Mass on 200 patients absent from every earlier Qwen cohort, crossing two wordings (\textit{show}, \textit{is}) with explicit yes/no, A-present, and B-present instructions at $\alpha=+0.25$, with 119 random directions, a sham, and 20 clinical comparisons under simultaneous 95\% lower bounds. Confirmation requires positive bounds and Mass effects above every random effect and the absolute sham, across all four coded cells or both original-wording cells.

\section{Pilot study: results}
\label{app:seed-results}

\begin{table}[t]
\centering
\setlength{\tabcolsep}{4pt}
\input{tables/cf_colors}
\caption{\textbf{Pilot replication on 600 new patients.} Qwen2.5-VL-7B on NIH ChestX-ray14. Per concept: the concept write $W_{q,q}$, the random 95th percentile (p95), the absolute sham, and its rank among the 120 compared effects; blue shades a negative ownership contrast.}
\label{tab:cf-seed}
\DIFdelbeginFL %DIFDELCMD < \fitwidth{%
%DIFDELCMD < \begin{tabular}{lrrrrrrrll}
%DIFDELCMD < \toprule
%DIFDELCMD < Concept & \multicolumn{1}{c}{$S$} & \multicolumn{1}{c}{AUROC} & \multicolumn{1}{c}{$W_{q,q}$} & \multicolumn{1}{c}{p95} & \multicolumn{1}{c}{$|$sham$|$} & \multicolumn{1}{c}{rank} & \multicolumn{1}{c}{$O_q$ [95\% CI]} & Competitor & Verdict\\
%DIFDELCMD < \midrule
%DIFDELCMD < Effusion & 0.102 & 0.593 & +0.190 & 0.137 & 0.017 & 2/120 & \cellcolor{cfSlate2}-0.071 [-0.080, -0.062] & Nodule & competitor \\
%DIFDELCMD < Atelectasis & 0.062 & 0.566 & +0.101 & 0.297 & 0.096 & 55/120 & \cellcolor{cfSlate3}-0.226 [-0.233, -0.219] & Nodule & competitor \\
%DIFDELCMD < Pneumothorax & 0.152 & 0.508 & +0.089 & 0.320 & 0.113 & 52/120 & \cellcolor{cfSlate3}-0.102 [-0.108, -0.096] & Effusion & competitor \\
%DIFDELCMD < Cardiomegaly & 0.205 & 0.742 & +0.085 & 0.311 & 0.158 & 33/120 & \cellcolor{cfSlate3}-0.219 [-0.227, -0.211] & Atelectasis & competitor \\
%DIFDELCMD < Mass & 0.021 & 0.678 & +0.359 & 0.372 & 0.092 & 10/120 & \cellcolor{cfSlate3}-0.169 [-0.178, -0.160] & Effusion & competitor \\
%DIFDELCMD < Nodule & -0.041 & 0.642 & +0.118 & 0.158 & 0.055 & 15/120 & \cellcolor{cfSlate2}-0.072 [-0.081, -0.063] & Effusion & competitor \\
%DIFDELCMD < \bottomrule
%DIFDELCMD < \end{tabular}}
%DIFDELCMD < %%%
\DIFdelendFL \DIFaddbeginFL \resizebox{\linewidth}{!}{%
\begin{tabular}{lrrrrrrrll}
\toprule
Concept & \multicolumn{1}{c}{$S$} & \multicolumn{1}{c}{AUROC} & \multicolumn{1}{c}{$W_{q,q}$} & \multicolumn{1}{c}{p95} & \multicolumn{1}{c}{$|$sham$|$} & \multicolumn{1}{c}{rank} & \multicolumn{1}{c}{$O_q$ [95\% CI]} & Competitor & Verdict\\
\midrule
Effusion & 0.102 & 0.593 & +0.190 & 0.137 & 0.017 & 2/120 & \cellcolor{cfSlate2}-0.071 [-0.080, -0.062] & Nodule & competitor \\
Atelectasis & 0.062 & 0.566 & +0.101 & 0.297 & 0.096 & 55/120 & \cellcolor{cfSlate3}-0.226 [-0.233, -0.219] & Nodule & competitor \\
Pneumothorax & 0.152 & 0.508 & +0.089 & 0.320 & 0.113 & 52/120 & \cellcolor{cfSlate3}-0.102 [-0.108, -0.096] & Effusion & competitor \\
Cardiomegaly & 0.205 & 0.742 & +0.085 & 0.311 & 0.158 & 33/120 & \cellcolor{cfSlate3}-0.219 [-0.227, -0.211] & Atelectasis & competitor \\
Mass & 0.021 & 0.678 & +0.359 & 0.372 & 0.092 & 10/120 & \cellcolor{cfSlate3}-0.169 [-0.178, -0.160] & Effusion & competitor \\
Nodule & -0.041 & 0.642 & +0.118 & 0.158 & 0.055 & 15/120 & \cellcolor{cfSlate2}-0.072 [-0.081, -0.063] & Effusion & competitor \\
\bottomrule
\end{tabular}}
\DIFaddendFL \end{table}

\subsection{A non-generic write effect reproduces}
\label{sec:qwen-replication}

The first finding starts with a successful intervention: Qwen Effusion reproduces beyond random and sham controls before clinical alternatives change its attribution.

Controlled decodability is evaluated at the final visual block each connector consumes: the penultimate CLIP encoder layer for LLaVA and the last of Qwen's 32 visual blocks. Intervals are 95\% patient-cluster bootstrap intervals over 2,000 draws; the control AUROC is the mean over the 20 recurring-type random-label maps, with its 5th--95th percentile spread in parentheses.

Qwen Effusion first satisfies the reader test. Its probe reaches AUROC \cfSeedDecQwenEffAuroc{} $[\cfSeedDecQwenEffAurocLo,\cfSeedDecQwenEffAurocHi]$ against a control AUROC of \cfSeedDecQwenEffCtrl{} (\cfSeedDecQwenEffCtrlLo--\cfSeedDecQwenEffCtrlHi), a controlled selectivity of \cfSeedDecQwenEffSel{} (95\% CI $[\cfSeedDecQwenEffSelLo,\cfSeedDecQwenEffSelHi]$). LLaVA-1.5-7B reads both of its concepts at its consumed block: Effusion at AUROC \cfSeedDecLlavaEffAuroc{} $[\cfSeedDecLlavaEffAurocLo,\cfSeedDecLlavaEffAurocHi]$ with selectivity \cfSeedDecLlavaEffSel{} $[\cfSeedDecLlavaEffSelLo,\cfSeedDecLlavaEffSelHi]$, and Edema at AUROC \cfSeedDecLlavaEdemaAuroc{} $[\cfSeedDecLlavaEdemaAurocLo,\cfSeedDecLlavaEdemaAurocHi]$ with selectivity \cfSeedDecLlavaEdemaSel{} $[\cfSeedDecLlavaEdemaSelLo,\cfSeedDecLlavaEdemaSelHi]$, against a control AUROC of \cfSeedDecLlavaEffCtrl{} (\cfSeedDecLlavaEffCtrlLo--\cfSeedDecLlavaEffCtrlHi). Original-cohort intervention outcomes maximize the concept-consistent effect over six fixed nonzero doses. Qwen makes a strong write-direction case in the original 200-image cohort: the maximum concept-consistent change is 0.2343 at $\alpha=+0.25$, above the random-direction 95th percentile of 0.0963 and the maximum absolute sham effect of 0.0852.

The response at the locked dose $\alpha=+0.25$ reproduces on the patient-disjoint 400-patient cohort. Effusion changes mean $P(\mathrm{yes})$ by 0.1945, again above the random-direction 95th percentile of 0.0878 and absolute sham effect of 0.0157. None of these patients appears in the original intervention cohort, and the dose was not reselected, supporting a large and reproducible non-generic effect within ChestX-ray14.


\subsection{Semantic competition changes attribution, not efficacy}
\label{sec:specificity-results}

The fixed clinical family provides a harder identification test than non-semantic controls. Nodule changes the Effusion answer by 0.2519 and realizes the point-estimate maximum among the five alternatives. The Effusion ownership contrast $O_{\mathrm{Effusion}}$ is $-0.0573$. The registered fixed-family endpoint recomputes the maximum inside each of 5,000 patient-bootstrap draws, giving a two-sided 95\% CI of $[-0.0694,-0.0450]$ and a one-sided 95\% lower bound of $-0.0678$. The entire interval is below zero. Descriptively, all five alternative clinical normals move the Effusion endpoint in the same positive direction; Effusion exceeds four by 0.1172--0.1778, while Nodule alone has a larger point estimate. This shared-sign pattern motivates an answer-sensitive-alias hypothesis but cannot distinguish it from question-specific ownership.

Qwen Effusion therefore remains a reproducible direction that changes the answer beyond random and sham perturbations, while its specific semantic attribution fails. The familywise endpoint establishes that at least one fixed clinical alternative is stronger; Nodule is the point-estimate maximizer, not a separately familywise-identified semantic owner.

\subsection{Wording changes clinical competition for Mass}
\label{sec:encoding-results}

In the registered crossover, the four questions that pass capability calibration (Atelectasis, Pneumothorax, Mass, and Nodule) respond more in the component that reverses the raw A-minus-B sign when the finding-to-letter mapping is exchanged, but so do all 20 random directions and the sham, so aggregate mapping response does not separate clinical normals from generic perturbations. Mass is the exception: its clinical margin, the Mass effect minus the largest of the five other clinical effects in finding-present logit units, changes from $-0.7181$ under the original yes/no prompt to $0.2825$ and $0.2900$ under A-present and B-present (Table~\ref{tab:seed-mass}, top). The independent confirmation on 200 new patients reproduces the clinical advantage under both \textit{show} coded prompts, with simultaneous lower bounds of 0.1675 and 0.1462, while both \textit{is} margins are negative and the Mass effects stay below the maximum of the 119 random directions (rank 3 of 120), so neither registered specificity conjunction is met (Table~\ref{tab:seed-mass}, bottom). Mass steering increases Brier error in every cell, and every AUROC-change interval spans zero.

\begin{table}[H]
\centering
\setlength{\tabcolsep}{5pt}
\caption{\textbf{Mass clinical competition across prompts.} Clinical margin per prompt: the Mass effect minus the largest of the five other clinical effects, in finding-present logit units. Top, the discovery crossover; bottom, the registered confirmation on new patients.}
\label{tab:seed-mass}
\DIFdelbeginFL %DIFDELCMD < \fitwidth{%
%DIFDELCMD < \begin{tabular}{lrrrrr}
%DIFDELCMD < \toprule
%DIFDELCMD < Prompt & Clinical margin & 95\% CI & Mass effect & Random max & $|$Sham$|$ \\
%DIFDELCMD < \midrule
%DIFDELCMD < Standard yes/no & $-0.7181$ & $[-0.7956,-0.6412]$ & 2.8625 & 3.0412 & 0.5869 \\
%DIFDELCMD < A-present & $0.2825$ & $[0.2269,0.3406]$ & 1.6125 & 1.1650 & 0.7550 \\
%DIFDELCMD < B-present & $0.2900$ & $[0.2300,0.3506]$ & 1.7469 & 1.2581 & 1.1475 \\
%DIFDELCMD < \midrule
%DIFDELCMD < Condition & Clinical margin & Minimum LCB & Mass effect & Random max & \\
%DIFDELCMD < \midrule
%DIFDELCMD < \textit{show}/A-present & 0.2525 & 0.1675 & 1.7019 & 2.2100 & \\
%DIFDELCMD < \textit{show}/B-present & 0.2338 & 0.1462 & 1.8188 & 2.2169 & \\
%DIFDELCMD < \textit{is}/A-present & $-0.1087$ & $-0.1751$ & 1.9225 & 2.4750 & \\
%DIFDELCMD < \textit{is}/B-present & $-0.2006$ & $-0.2686$ & 2.0019 & 2.4944 & \\
%DIFDELCMD < \bottomrule
%DIFDELCMD < \end{tabular}}
%DIFDELCMD < %%%
\DIFdelendFL \DIFaddbeginFL \resizebox{\linewidth}{!}{%
\begin{tabular}{lrrrrr}
\toprule
Prompt & Clinical margin & 95\% CI & Mass effect & Random max & $|$Sham$|$ \\
\midrule
Standard yes/no & $-0.7181$ & $[-0.7956,-0.6412]$ & 2.8625 & 3.0412 & 0.5869 \\
A-present & $0.2825$ & $[0.2269,0.3406]$ & 1.6125 & 1.1650 & 0.7550 \\
B-present & $0.2900$ & $[0.2300,0.3506]$ & 1.7469 & 1.2581 & 1.1475 \\
\midrule
Condition & Clinical margin & Minimum LCB & Mass effect & Random max & \\
\midrule
\textit{show}/A-present & 0.2525 & 0.1675 & 1.7019 & 2.2100 & \\
\textit{show}/B-present & 0.2338 & 0.1462 & 1.8188 & 2.2169 & \\
\textit{is}/A-present & $-0.1087$ & $-0.1751$ & 1.9225 & 2.4750 & \\
\textit{is}/B-present & $-0.2006$ & $-0.2686$ & 2.0019 & 2.4944 & \\
\bottomrule
\end{tabular}}
\DIFaddendFL \end{table}

\FloatBarrier

\section{Questions, Scores, and Notation}
\label{app:protocol}

\paragraph{Templates and score.} Each concept is asked in six templates that cross two wordings, \emph{is there} (I) and \emph{does the image show} (W), with three answer interfaces, yes/no (Y), A-present (A), and B-present (B). A template is named by its wording and its interface, and the primary template is IY. At the first answer position, $\ell_{\mathrm{yes}}$ and $\ell_{\mathrm{no}}$ are the maximum log probabilities among the valid single-token case and leading-space variants of each answer, and the registered score is $P(\mathrm{yes})=\sigma(\ell_{\mathrm{yes}}-\ell_{\mathrm{no}})$, where $\sigma$ is the logistic function. Taking the best valid variant for each answer makes the score insensitive to which capitalisation or leading-space form receives the highest probability. The score measures the relative support for yes versus no at the first answer position, not the accuracy of a generated report. Questions are the rows of $W$ and directions its columns, so changing a direction never changes the question being scored.

\paragraph{Notation.} Table~\ref{tab:notation} collects the notation of the paper.

\begin{table}[H]
\centering
\caption{Notation used in the paper.}
\label{tab:notation}
\DIFdelbeginFL %DIFDELCMD < \fitwidth{%
%DIFDELCMD < \begin{tabular}{l>{\raggedright\arraybackslash}p{0.867\linewidth}}
%DIFDELCMD < \toprule
%DIFDELCMD < Symbol & Definition \\
%DIFDELCMD < \midrule
%DIFDELCMD < $c$, $q$, $d$ & A concept; the yes/no question about a concept; a write direction. Questions and directions are indexed by the concept they belong to. \\
%DIFDELCMD < cell, block & A checkpoint--concept pair on one dataset; a checkpoint--dataset pair. \\
%DIFDELCMD < $\hvec_{it}$, $\hvec_i$ & Output of the consumed block for image $i$ and consumed token $t$; its mean over the consumed tokens. \\
%DIFDELCMD < $D$ & Hidden width of the consumed block. \\
%DIFDELCMD < $R$, $\boldsymbol s$, $\boldsymbol\beta_c$ & Fixed $D\times512$ Gaussian projection; training-row feature scales; logistic coefficients of concept $c$'s probe. \\
%DIFDELCMD < $\hat\wvec_c$ & Label direction: the probe's logit gradient $R\,\mathrm{diag}(\boldsymbol s)^{-1}\boldsymbol\beta_c$ at unit length (Equation~\ref{eq:lift}). \\
%DIFDELCMD < $a_q$ & Answer direction: the ridge fit of the clean answer margin $\ell_{\mathrm{yes}}-\ell_{\mathrm{no}}$, lifted in the same way. \\
%DIFDELCMD < $\alpha$ & Signed dose; $|\alpha|$ is the write's norm as a fraction of each token's norm. \\
%DIFDELCMD < $S$ & Probe selectivity: concept AUROC minus the mean AUROC of 20 random-label control probes (Equation~\ref{eq:selectivity}). \\
%DIFDELCMD < $W_{q,d}$ & Mean change in question $q$'s $P(\mathrm{yes})$ when direction $d$ is written. \\
%DIFDELCMD < $O_q$ & Ownership, $W_{q,q}-\max_{d\neq q}W_{q,d}$ (Equation~\ref{eq:margin}). \\
%DIFDELCMD < $O^m_q$ & Ownership computed on changes in the logit margin $\ell_{\mathrm{yes}}-\ell_{\mathrm{no}}$. \\
%DIFDELCMD < $\kappa_d$ & Column selectivity, $W_{d,d}/\sum_q|W_{q,d}|$. \\
%DIFDELCMD < $F_{q,d}$ & Mean over checkpoints of $\max(W_{q,d},0)$, the band width in Figure~\ref{fig:write-flow}. \\
%DIFDELCMD < $\cos_\Sigma$ & Whitened cosine under the training-feature covariance $\Sigma$ (Appendix~\ref{app:cf-further}). \\
%DIFDELCMD < \bottomrule
%DIFDELCMD < \end{tabular}}
%DIFDELCMD < %%%
\DIFdelendFL \DIFaddbeginFL \resizebox{\linewidth}{!}{%
\begin{tabular}{l>{\raggedright\arraybackslash}p{0.867\linewidth}}
\toprule
Symbol & Definition \\
\midrule
$c$, $q$, $d$ & A concept; the yes/no question about a concept; a write direction. Questions and directions are indexed by the concept they belong to. \\
cell, block & A checkpoint--concept pair on one dataset; a checkpoint--dataset pair. \\
$\hvec_{it}$, $\hvec_i$ & Output of the consumed block for image $i$ and consumed token $t$; its mean over the consumed tokens. \\
$D$ & Hidden width of the consumed block. \\
$R$, $\boldsymbol s$, $\boldsymbol\beta_c$ & Fixed $D\times512$ Gaussian projection; training-row feature scales; logistic coefficients of concept $c$'s probe. \\
$\hat\wvec_c$ & Label direction: the probe's logit gradient $R\,\mathrm{diag}(\boldsymbol s)^{-1}\boldsymbol\beta_c$ at unit length (Equation~\ref{eq:lift}). \\
$a_q$ & Answer direction: the ridge fit of the clean answer margin $\ell_{\mathrm{yes}}-\ell_{\mathrm{no}}$, lifted in the same way. \\
$\alpha$ & Signed dose; $|\alpha|$ is the write's norm as a fraction of each token's norm. \\
$S$ & Probe selectivity: concept AUROC minus the mean AUROC of 20 random-label control probes (Equation~\ref{eq:selectivity}). \\
$W_{q,d}$ & Mean change in question $q$'s $P(\mathrm{yes})$ when direction $d$ is written. \\
$O_q$ & Ownership, $W_{q,q}-\max_{d\neq q}W_{q,d}$ (Equation~\ref{eq:margin}). \\
$O^m_q$ & Ownership computed on changes in the logit margin $\ell_{\mathrm{yes}}-\ell_{\mathrm{no}}$. \\
$\kappa_d$ & Column selectivity, $W_{d,d}/\sum_q|W_{q,d}|$. \\
$F_{q,d}$ & Mean over checkpoints of $\max(W_{q,d},0)$, the band width in Figure~\ref{fig:write-flow}. \\
$\cos_\Sigma$ & Whitened cosine under the training-feature covariance $\Sigma$ (Appendix~\ref{app:cf-ansdir}). \\
\bottomrule
\end{tabular}}
\DIFaddendFL \end{table}

\FloatBarrier

\DIFaddbegin \begin{DIFnomarkup}\cfnote{附录精简}{max-$T$ 公式从 pilot 协议移到这里；$\tau_j$ 的单独公式并入正文一句（分母 $B-1$）；判定规则（advantage / stronger competitor / unresolved）、百分位区间、各研究的抽样设置分成三段，每段先给结论。}\end{DIFnomarkup}

\DIFaddend \section{Joint bootstrap and simultaneous lower bounds}
\label{app:simultaneous-bounds}

\DIFdelbegin \DIFdel{Equation~\ref{eq:max-t} uses the standard deviation across bootstrap
estimates as a fixed scale for each statistic. With $\overline C^*_j$
denoting their mean, that scale is
}%DIFDELCMD < \begin{equation}
%DIFDELCMD < \tau_j=\left[\frac{1}{B-1}\sum_{b=1}^{B}
%DIFDELCMD <        (C^*_{bj}-\overline C^*_j)^2\right]^{1/2}.
%DIFDELCMD < \label{eq:bootstrap-scale}
%DIFDELCMD < \end{equation}
%DIFDELCMD < %%%
\DIFdel{The standardized deviations are centered on the observed }\DIFdelend \DIFaddbegin \DIFadd{For a family of $m$ statistics, let }\DIFaddend $\widehat C_j$ \DIFaddbegin \DIFadd{be statistic $j$'s estimate, $C^*_{bj}$ its value in joint bootstrap draw $b\in\{1,\dots,B\}$, and $\tau_j$ its bootstrap standard deviation with denominator $B-1$. The centred studentised draws, the common critical value}\DIFaddend , \DIFdelbegin \DIFdel{not on $\overline C^*_j$. For a component with $\tau_j\leq10^{-12}$, all its $Z_{bj}$ values are set to zero. If every component satisfies
that threshold, the critical value is zero. Otherwise, each draw's
maximum includes these zero components, and its }\DIFdelend \DIFaddbegin \DIFadd{and the simultaneous lower bounds are
}\begin{equation}
Z_{bj}=\frac{C^*_{bj}-\widehat C_j}{\tau_j},\qquad
c_{.95}=Q_{.95}\Bigl(\bigl\{\max_{1\leq j\leq m}Z_{bj}\bigr\}_{b=1}^{B}\Bigr),\qquad
L_j=\widehat C_j-c_{.95}\,\tau_j,
\label{eq:max-t}
\end{equation}
\DIFadd{where $Q_{.95}$ is the }\DIFaddend empirical 95th percentile \DIFdelbegin \DIFdel{uses }\DIFdelend \DIFaddbegin \DIFadd{with }\DIFaddend linear interpolation. \DIFdelbegin \DIFdel{All 5,000 draws are retained. The final
subtraction $L_j=\widehat C_j-c_{.95}\tau_j$ still uses the actual $\tau_j$, including for thresholded components. }\DIFdelend \DIFaddbegin \DIFadd{The draws are centred on the estimate, and each $\tau_j$ stays fixed across draws. A component with $\tau_j\leq10^{-12}$ gets $Z_{bj}=0$ but keeps its own $\tau_j$ in $L_j$; every draw is retained. Upper bounds use the same $c_{.95}$ with every contrast's sign reversed.
}\DIFaddend 

\DIFdelbegin \DIFdel{For the ownership matrix, we resample 400 patients with replacement using seed 20260904. Each draw uses the same patient indices for every question and intervention, preserving all paired baseline differences.
The clinical family comprises the }\DIFdelend \DIFaddbegin \paragraph{\DIFadd{Verdicts.}} \DIFadd{For the ownership family of }\DIFaddend 30 \DIFdelbegin \DIFdel{fixed }\DIFdelend contrasts $W_{q,q}-W_{q,d}$\DIFdelbegin \DIFdel{for the six questions and their five alternatives. Each question }\DIFdelend \DIFaddbegin \DIFadd{, question $q$}\DIFaddend 's simultaneous \DIFdelbegin \DIFdel{ownership }\DIFdelend lower bound is the minimum of its five contrast bounds. \DIFdelbegin \DIFdel{Simultaneous upper bounds use the same critical value with the sign of every contrast reversed. A question has }\DIFdelend \DIFaddbegin \DIFadd{The question has the advantage when that minimum is positive, }\DIFaddend a stronger competitor \DIFdelbegin \DIFdel{if any upper bound is below zero. A question is unresolved if it has neither a positive minimum lower bound nor a negativeupper bound. The shared-alias calculation uses a separate 12-statistic family: for each of six directions, its mean off-diagonal effect and its mean excess over the five corresponding matched-direction effects. The Mass confirmation applies the same construction to 200 independently selected patients, with seed 20260910 and a separate family of 20 Mass-minus-alternative contrasts across four coded prompt conditions.
}%DIFDELCMD < 

%DIFDELCMD < %%%
\DIFdel{Direct target-minus-maximum intervals summarise a different bootstrap quantity. For each draw , we recompute the largest competing mean and subtract it from the target mean. The displayed }\DIFdelend \DIFaddbegin \DIFadd{when any of its upper bounds is negative, and is unresolved otherwise. $O_q$ itself receives a percentile interval instead: every draw recomputes the largest competitor, the }\DIFaddend two-sided \DIFdelbegin \DIFdel{percentile intervals use }\DIFdelend \DIFaddbegin \DIFadd{interval takes }\DIFaddend the 2.5th and 97.5th percentiles\DIFdelbegin \DIFdel{of these margins. The earlier Effusion specificity supplement uses their }\DIFdelend \DIFaddbegin \DIFadd{, and the pilot's Effusion supplement takes the }\DIFaddend 5th \DIFdelbegin \DIFdel{percentile }\DIFdelend as its one-sided \DIFdelbegin \DIFdel{lower bound. These intervals are distinct from the simultaneous bounds on the fixed linear contrasts in Equation~\ref{eq:max-t}}\DIFdelend \DIFaddbegin \DIFadd{bound}\DIFaddend . \DIFdelbegin %DIFDELCMD < 

%DIFDELCMD < %%%
\DIFdelend The \DIFdelbegin \DIFdel{campaign applies the same construction to every block, using }\DIFdelend \DIFaddbegin \DIFadd{logit-scale grade (Appendix~\ref{app:cf-robustness}) applies both rules to margin changes.
}

\paragraph{\DIFadd{Draws.}} \DIFadd{The campaign resamples the }\DIFaddend 600 test rows \DIFdelbegin \DIFdel{and }\DIFdelend \DIFaddbegin \DIFadd{of each block in }\DIFaddend 2{,}000 shared \DIFdelbegin \DIFdel{bootstrap draws of rows }\DIFdelend \DIFaddbegin \DIFadd{draws }\DIFaddend (seed 2026090601). The \DIFdelbegin \DIFdel{30 contrasts $W_{q,q}-W_{q,d}$ receive simultaneous 95\% bounds from the centred studentised max-$T$ bootstrap. $O_q$ receives a direct percentile interval with the maximum recomputed in every draw. The logit-scale check in Table~\ref{tab:cf-scale} applies the same rules to margin-scale deltas.
}\DIFdelend \DIFaddbegin \DIFadd{pilot's ownership matrix resamples 400 patients in 5}{\DIFadd{,}}\DIFadd{000 draws (seed 20260904), with the same patient indices for every question and intervention, and bounds its shared-alias statistics as a separate family of 12. The Mass confirmation resamples 200 new patients (seed 20260910), with a family of 20 Mass-minus-alternative contrasts across four coded prompts.
}

\DIFaddend \section{Reproducibility specification}
\label{app:repro}

This appendix specifies each family's consumed block, the consumed-token mask, the projection and probe, the control constructions, the cohort roles, and the protocol record at the registered protocol's level of detail. The released code implements each component.

\subsection{Consumed loci per family}
\label{app:repro-loci}

The consumed block is the last visual block that feeds the consumer. The connector locus is the consumer's output that enters the language model. We identify the block from each checkpoint's own configuration when loading the checkpoint. Each block's preflight record stores the identified block, its hidden dimension, its consumer, and any branch that bypasses it.

\begin{itemize}[leftmargin=*,itemsep=1pt,topsep=2pt]
\item Qwen2.5-VL, Lingshu, and Qwen3-VL: the merger reads the last block of the vision transformer, which has 32 blocks in Qwen2.5-VL and Lingshu, 24 in Qwen3-VL 4B, and 27 in Qwen3-VL 8B and 32B. A Qwen2.5-VL or Lingshu image contributes 576 patch tokens, which the merger joins in $2\times2$ groups into 144; a Qwen3-VL image contributes at most 1{,}024. The Qwen3-VL DeepStack mergers, which read earlier blocks, are not written.
\item Gemma~3 and MedGemma: the projector reads the last of the 27 SigLIP encoder layers through the tower's final layer normalisation, and consumes all 4{,}096 patch tokens of the $896\times896$ image.
\item InternVL3.5: the projector reads the last InternViT encoder layer, of 24 at 8B and 14B and of 45 at 38B; the tower's final normalisation is an identity in the released configuration. An image contributes 1{,}024 patch tokens from a single $448\times448$ tile.
\item LLaVA-1.5 and LLaVA-Med: the projector reads the penultimate of the 24 CLIP encoder layers, and the final layer is computed and discarded. An image contributes 576 patch tokens.
\item Llama~3.2 Vision: the projector input concatenates the last of the eight global-encoder layers with five intermediate layers of the local encoder, which bypass the consumed block. An image contributes 1{,}601 tokens, the CLS token and 1{,}600 patches, from its one real $560\times560$ tile.
\end{itemize}
Table~\ref{tab:cf-models} lists the consumed-token count of every checkpoint.

\subsection{The consumed-token mask}
\label{app:repro-mask}

The write in Equation~\ref{eq:intervention} changes only the consumed tokens at the positions the consumer reads at the locus. We derive these positions for each batch from that batch's processor outputs:

\begin{itemize}[leftmargin=*,itemsep=1pt,topsep=2pt]
\item Qwen2.5-VL, Lingshu, and Qwen3-VL: the tokens of all images in a batch form one sequence in which every image occupies one contiguous span, whose bounds follow from the image's patch grid. At the connector locus the span is a quarter as long, one merged token per $2\times2$ patch group.
\item Gemma~3 and MedGemma: all 4{,}096 patch tokens; the tower produces no CLS token. The number of written tokens must equal the number of image placeholder tokens in the language-model input.
\item InternVL3.5: the patch tokens of the single $448\times448$ tile. The CLS token, which the consumer drops, is excluded. The number of written tokens must equal the number of image placeholder tokens.
\item LLaVA-1.5 and LLaVA-Med: the 576 patch tokens. The CLS token, which the checkpoint's feature selection drops, is excluded.
\item Llama~3.2 Vision: the 1{,}601 tokens, CLS and 1{,}600 patches, of the one real $560\times560$ tile. The seven padding tokens of each tile and the three padding tiles are excluded, in agreement with the model's own aspect-ratio and cross-attention masks.
\end{itemize}

Scoring stops with an error and no fallback when an image span does not exactly cover the locus, a mask length differs from the token count at the locus, an image has no consumed token, or one forward pass reaches the locus more than once.

Preflight check (C) verifies the mask. It captures the locus tensor on a clean pass and a pass steered at $\alpha=+0.25$ along the first concept direction. It reports the maximum absolute difference outside the mask and the maximum absolute changes in the consumer's output and answer logits. Across all \cfBlocks{} blocks, the change outside the consumed tokens is $0$ at both loci. The consumer's output and answer logits both change. The same records show a bitwise-identical repeated clean pass (check A) and a bitwise-identical $\alpha=0$ pass (check B) for every block.

\subsection{Projection, lift, scaler, and probe}
\label{app:repro-probe}

\DIFaddbegin \begin{DIFnomarkup}\cfnote{附录精简}{删去与正文 Eq.~2 重复的两句（单位化、正类系数定方向）。两种 lift 的特征值与条件数只保留在本节，A.7 引用这里。}\end{DIFnomarkup}

\DIFaddend We construct the projection $R\in\R^{D\times512}$ by drawing standard normal variates once from a PCG64 generator seeded with $0$, dividing by $\sqrt{512}$, and storing the result in float32. Its entries are independent $\mathcal{N}(0,1/512)$ variates. The seed is $0$ for every checkpoint, dataset, and locus. $D$ is the family's hidden dimension. We neither normalise nor orthogonalise the columns of $R$. \DIFdelbegin \DIFdel{We normalise the lifted direction of Equation~\ref{eq:lift} once to unit $\ell_2$ norm in model space. }\DIFdelend The Gram matrix $R^{\top}R/D$ has full rank 512 in all \cfAltdirdBlocks{} blocks that carry the displacement module, with eigenvalues $\cfAltdirdEigMin$--$\cfAltdirdEigMax$ and condition number \cfAltdirdCondMin{}--\cfAltdirdCondMax{}.

Projecting the token-mean read vector $\hvec_i$ gives features $R^{\top}\hvec_i$. A featurewise standardiser fitted only on training rows supplies the mean $\boldsymbol\mu$ and scale $\boldsymbol s$. Every row is standardised as $(R^{\top}\hvec-\boldsymbol\mu)\oslash\max(\boldsymbol s,10^{-8})$, where $\oslash$ and $\max$ act elementwise. Each concept probe is an $\ell_2$-regularised logistic regression with inverse regularisation strength $C=1$. We fit each probe by L-BFGS for at most 2{,}000 iterations with unweighted classes and seed $0$, using training rows with a known label for that concept. \DIFdelbegin \DIFdel{The positive-class coefficient fixes the direction's clinical pole. }\DIFdelend For each of the two refit seeds, we resample training units with replacement, retaining all rows of each sampled unit at its multiplicity. We then refit the scaler and six probes using the same projection. The random family, shams, and control probes remain at seed 0.

\subsection{Control constructions}
\label{app:repro-controls}

\paragraph{Random directions.} We draw one $119\times D$ standard-normal matrix using a PCG64 generator seeded with $0$ and normalise each row to unit $\ell_2$ norm. Every question in the block uses the same 119 directions.

\paragraph{Coordinate-permutation sham.} Immediately after the random draw, the same generator produces one permutation of $\{0,\dots,D{-}1\}$ per concept, in protocol concept order. We form concept $c$'s sham by rearranging the coordinates of $\hat\wvec_c$ with concept $c$'s own permutation. This preserves the direction's coordinate multiset and unit norm. We write each question with its own sham.

\paragraph{Matched random-label probes.} Each row has a type: view $\times$ sex $\times$ age decade on the chest sets, and aspect-ratio $\times$ area buckets on COCO. For control seed $k\in\{0,\dots,19\}$, we traverse the types in sorted order and use a PCG64 generator seeded with $k$ to draw one label per type uniformly from $\{0,1\}$. If all types receive the same label, we flip the first. Every row of a type receives that type's label. We give each control probe the concept probe's capacity and nuisance structure by fitting it on the same training rows and known-label subset with the same hyperparameters. A block fits control probes when its training rows span at least eight types. Selectivity is Equation~\ref{eq:selectivity}. Its bootstrap over calibration units retains a draw when the real AUROC is finite and at least ten of the twenty control AUROCs are finite. We grade a cell for readability when at least 1{,}900 of the 2{,}000 draws are retained.

\subsection{Cohort roles}
\label{app:repro-cohorts}

\DIFdelbegin \DIFdel{Table~\ref{tab:cf-roles} lists each role's size, its independent units, and its use. }\DIFdelend \DIFaddbegin \begin{DIFnomarkup}\cfnote{删表}{Cohort roles 表删去；各 cohort 的规模和互不重叠在 A.1 与本段已有文字说明。}\end{DIFnomarkup}

\DIFaddend Every chest image is a frontal study. One role holds more than one row per unit: the NIH training role, which takes every study of each training patient, \cfCohortNihTrainRows{} studies from \cfCohortNihTrainUnits{} patients, \cfCohortNihTrainMulti{} of whom contribute more than one and one of whom contributes \cfCohortNihTrainMax{}. Every other role of every dataset holds one row per unit, and no unit appears in two roles of a dataset. We use the calibration role only to determine readability and answerability. We neither fit probes nor score writes on it. Every write uses the test role.

\DIFdelbegin \begin{table}[htbp]
\DIFdeletedtable{%
{\small\textbf{Table (deleted).} \textbf{Cohort roles.} Rows per role, with independent units in parentheses (patients on the chest sets, images on COCO), and the role's use. The last two rows give the largest rows per unit and the largest overlap between two roles.\par}\medskip
\centering
\setlength{\tabcolsep}{3.5pt}

\fitwidth{%
\fitwidth{%
\begin{tabular}{>{\raggedright\arraybackslash}p{0.164\textwidth}lll>{\raggedright\arraybackslash}p{0.479\textwidth}}
\toprule
Role & NIH ChestX-ray14 & CheXpert Plus & COCO & Use \\
\midrule
train & \cfCohortNihTrainRows{} (\cfCohortNihTrainUnits{}) & \cfCohortChexTrainRows{} (\cfCohortChexTrainUnits{}) & \cfCohortCocoTrainRows{} (\cfCohortCocoTrainUnits{}) & Scaler, six concept probes, twenty control probes, the two nuisance probes of the chest sets, and the answer-direction ridge. No row of this role is scored. \\
preflight & \cfCohortNihPreflightRows{} (\cfCohortNihPreflightUnits{}) & \cfCohortChexPreflightRows{} (\cfCohortChexPreflightUnits{}) & \cfCohortCocoPreflightRows{} (\cfCohortCocoPreflightUnits{}) & The seven acceptance gates, at both loci. \\
calibration & \cfCohortNihCalibrationRows{} (\cfCohortNihCalibrationUnits{}) & \cfCohortChexCalibrationRows{} (\cfCohortChexCalibrationUnits{}) & \cfCohortCocoCalibrationRows{} (\cfCohortCocoCalibrationUnits{}) & Clean pass at $\alpha=0$: probe readability and answerability. \\
test & \cfCohortNihTestRows{} (\cfCohortNihTestUnits{}) & \cfCohortChexTestRows{} (\cfCohortChexTestUnits{}) & \cfCohortCocoTestRows{} (\cfCohortCocoTestUnits{}) & Every write matrix, dose, refit, template, locus, and addendum module. \\
valid & -- & \cfCohortChexValidRows{} (\cfCohortChexValidUnits{}) & -- & The regrade on radiologist consensus labels (Table~\ref{tab:cf-valid}). \\
\midrule
rows per unit & \cfCohortNihTrainMax{} (\cfCohortNihTrainMulti{} units) & \cfCohortChexTrainMax{} (\cfCohortChexTrainMulti{}) & \cfCohortCocoTrainMax{} (\cfCohortCocoTrainMulti{}) & Largest number of rows one unit contributes, with the units contributing more than one; the NIH training role is the campaign's only role above one. \\
shared units & \cfCohortNihMaxOverlap{} & \cfCohortChexMaxOverlap{} & \cfCohortCocoMaxOverlap{} & Largest number of units two roles of the dataset share, over all \cfCohortPairsAll{} role pairs of the three datasets. \\
\bottomrule
\end{tabular}}}
}
\end{table}
\begin{table}[htbp]
\DIFoldtable{%
{\small\textbf{Table (before revision).} \textbf{Small and fine-grained objects, and difficulty matching.} Top and middle: six fine-grained COCO categories graded like the six easy objects, against the easy objects of every COCO block and of the 6 blocks with fine-grained cells. Bottom: chest cells matched to natural-image cells on difficulty, and the same cells stratified on clean-answer AUROC.\par}\medskip
\centering
\setlength{\tabcolsep}{4pt}
\setlength{\abovecaptionskip}{2pt}
\setlength{\belowcaptionskip}{0pt}
\renewcommand{\arraystretch}{0.95}

\fitwidth{%
\fitwidth{%
\begin{tabular}{lrrrrrr}
\toprule
\multirow{2}{*}{Checkpoint} & \multicolumn{2}{c}{owned of six} & \multicolumn{2}{c}{median answer AUROC} & \multicolumn{2}{c}{median selectivity}\\
\cmidrule(lr){2-3}\cmidrule(lr){4-5}\cmidrule(lr){6-7}
& \multicolumn{1}{c}{fine} & \multicolumn{1}{c}{easy} & \multicolumn{1}{c}{fine} & \multicolumn{1}{c}{easy} & \multicolumn{1}{c}{fine} & \multicolumn{1}{c}{easy}\\
\midrule
Qwen2.5-VL-3B & 6 & 6 & 0.947 & 0.971 & 0.024 & 0.037 \\
Qwen3-VL-4B & 5 & 6 & 0.984 & 0.990 & 0.036 & 0.033 \\
InternVL3.5-8B & 6 & 6 & 0.981 & 0.978 & 0.104 & 0.117 \\
Gemma 3 4B & 4 & 5 & 0.948 & 0.924 & 0.076 & 0.070 \\
Lingshu 7B & 6 & 6 & 0.950 & 0.968 & 0.025 & 0.023 \\
LLaVA-1.5-7B & 3 & 5 & 0.956 & 0.979 & 0.266 & 0.285 \\
\midrule
\multicolumn{7}{l}{\textit{Paired within these blocks:} fine-grained minus easy is +0, -1, +0, -1, -2, +0 cells of six; fine higher in 0,} \\
\multicolumn{7}{l}{easy higher in 3, tied in 3; exact two-sided sign test $p=0.25$} \\
\midrule
\multicolumn{7}{l}{\textit{The four pools of cells}} \\
Cells & blocks & owned & rate & median answer AUROC & median selectivity & readable \\
\midrule
chest findings & 50 & 38/300 & 0.127 & 0.695 & 0.137 & 178/300 \\
easy objects, every COCO block & 25 & 113/150 & 0.753 & 0.977 & 0.064 & 90/150 \\
easy objects, these blocks & 6 & 34/36 & 0.944 & 0.980 & 0.065 & 24/36 \\
fine-grained objects, these blocks & 6 & 30/36 & 0.833 & 0.963 & 0.058 & 23/36 \\
\midrule
\multicolumn{7}{l}{\textit{Chest cells matched to natural-image cells on difficulty, under both estimators}} \\
Matching & chest cells & partners & chest rate & natural rate & difference & 95\% interval \\
\midrule
probe selectivity alone & 300/300 & 186/186 & 0.127 & 0.769 &  &  \\
\quad pooled &  &  &  &  & -0.642 & [-0.763, -0.512] \\
\quad matched &  &  &  &  & -0.539 & [-0.678, -0.359] \\
clean-answer AUROC alone & 300/300 & 186/186 & 0.127 & 0.769 &  &  \\
\quad pooled &  &  &  &  & -0.642 & [-0.749, -0.505] \\
\quad matched &  &  &  &  & -0.066 & [-0.304, -0.004] \\
both variables & 212/300 & 101/186 & 0.137 & 0.624 &  &  \\
\quad pooled &  &  &  &  & -0.487 & [-0.653, -0.201] \\
\quad matched &  &  &  &  & -0.024 & [-0.198, 0.076] \\
both variables, easy partners only & 211/300 & 79/150 & 0.137 & 0.570 &  &  \\
\quad pooled &  &  &  &  & -0.432 & [-0.625, -0.150] \\
\quad matched &  &  &  &  & -0.008 & [-0.198, 0.085] \\
\midrule
\multicolumn{7}{l}{\textit{Chest cells and natural-image cells stratified on clean-answer AUROC, owned of cells}} \\
Cells & $<$ 0.50 & 0.50--0.70 & 0.70--0.80 & 0.80--0.90 & $\geq$ 0.90 & all cells \\
\midrule
NIH ChestX-ray14 & 0/7 & 6/80 & 3/43 & 4/19 & 0/1 & 13/150 \\
CheXpert Plus & 0/13 & 6/52 & 2/20 & 6/29 & 11/36 & 25/150 \\
\textit{chest findings} & 0/20 & 12/132 & 5/63 & 10/48 & 11/37 & 38/300 \\
easy objects, every COCO block & 0/6 & 0/6 & 0/7 & 1/4 & 112/127 & 113/150 \\
fine-grained objects, these blocks & 0/0 & 0/0 & 0/0 & 1/1 & 29/35 & 30/36 \\
\textit{natural-image objects} & 0/6 & 0/6 & 0/7 & 2/5 & 141/162 & 143/186 \\
\multicolumn{7}{l}{\quad chest minus natural-image in the $\geq$ 0.90 bin: -0.573, 95\% interval [-0.763, -0.388]} \\
\bottomrule
\end{tabular}}}
}
\end{table}
\begin{table}[htbp]
\DIFoldtable{%
{\small\textbf{Table (before revision).} \textbf{Four endpoints no direction was fitted against.} Both direction families of Section~\ref{sec:ansdir}, written at the primary dose and read at the negated question, the forced choice in both orders, and a report continuation. Bottom: the ownership contrast added to a regression of the endpoint effect, cross-validated by concept and by block.\par}\medskip
\centering
\setlength{\tabcolsep}{3pt}

\fitwidth{%
\fitwidth{%
\begin{tabular}{lllcccc}
\toprule
\multirow{2}{*}{Checkpoint} & \multirow{2}{*}{Dataset} & \multirow{2}{*}{Endpoint} & \multicolumn{2}{c}{label direction} & \multicolumn{2}{c}{answer direction}\\
\cmidrule(lr){4-5}\cmidrule(lr){6-7}
&  &  & \multicolumn{1}{c}{owned} & \multicolumn{1}{c}{reference} & \multicolumn{1}{c}{owned} & \multicolumn{1}{c}{reference}\\
\midrule
Lingshu 32B & CheXpert Plus & negated question & 2/6 & 2/6 & 3/6 & 4/6 \\
Lingshu 32B & CheXpert Plus & forced choice, finding first & 3/6 & 3/6 & 4/6 & 4/6 \\
Lingshu 32B & CheXpert Plus & forced choice, competitor first & 4/6 & 4/6 & 5/6 & 5/6 \\
Lingshu 32B & CheXpert Plus & report continuation & 0/6 & 0/6 & 0/6 & 0/6 \\
Qwen2.5-VL-7B & NIH ChestX-ray14 & negated question & 0/6 & 1/6 & 3/6 & 6/6 \\
Qwen2.5-VL-7B & NIH ChestX-ray14 & forced choice, finding first & 0/6 & 0/6 & 3/6 & 3/6 \\
Qwen2.5-VL-7B & NIH ChestX-ray14 & forced choice, competitor first & 0/6 & 0/6 & 2/6 & 2/6 \\
Qwen2.5-VL-7B & NIH ChestX-ray14 & report continuation & 0/6 & 0/6 & 0/6 & 1/6 \\
Qwen2.5-VL-7B & CheXpert Plus & negated question & 1/6 & 1/6 & 5/6 & 6/6 \\
Qwen2.5-VL-7B & CheXpert Plus & forced choice, finding first & 0/6 & 0/6 & 3/6 & 4/6 \\
Qwen2.5-VL-7B & CheXpert Plus & forced choice, competitor first & 0/6 & 0/6 & 2/6 & 2/6 \\
Qwen2.5-VL-7B & CheXpert Plus & report continuation & 0/6 & 0/6 & 0/6 & 0/6 \\
\midrule
\textit{all} & \textit{3 blocks} & \textit{4 endpoints} & 10/72 & 11/72 & 30/72 & 37/72 \\
\midrule
\multicolumn{7}{l}{\textit{The three endpoint-specific tests}} \\
Checkpoint & Dataset & opposite on the negation & order gap & forced choice owned & report owned & core owned \\
\midrule
Lingshu 32B & CheXpert Plus & 4/6 & 0.026 & 3/6 & 0/6 & 5/6 \\
Qwen2.5-VL-7B & NIH ChestX-ray14 & 6/6 & 0.089 & 0/6 & 0/6 & 0/6 \\
Qwen2.5-VL-7B & CheXpert Plus & 6/6 & 0.111 & 0/6 & 0/6 & 0/6 \\
\midrule
\multicolumn{7}{l}{\textit{What ownership adds, held out: leaving out one concept at a time}} \\
Checkpoint & Dataset & Direction & held-out $R^2$ & with ownership & $\Delta R^2$ & permutation $p$ \\
\midrule
Lingshu 32B & CheXpert Plus & label direction & -0.034 & 0.147 & 0.181 & 0.39 \\
Lingshu 32B & CheXpert Plus & answer direction & 0.069 & -0.223 & -0.291 & 0.70 \\
Qwen2.5-VL-7B & CheXpert Plus & label direction & -0.139 & -0.149 & -0.010 & 0.55 \\
Qwen2.5-VL-7B & CheXpert Plus & answer direction & -3.168 & -0.112 & 3.057 & 0.18 \\
Qwen2.5-VL-7B & NIH ChestX-ray14 & label direction & -0.195 & -0.147 & 0.048 & 0.31 \\
Qwen2.5-VL-7B & NIH ChestX-ray14 & answer direction & -1.310 & -0.176 & 1.134 & 0.50 \\
\textit{pooled} & \textit{6 tests} & \textit{both} & -- & -- & 0.687 & 0.26 \\
\midrule
\multicolumn{7}{l}{\textit{The same, leaving out one block at a time}} \\
\midrule
\textit{all 3} & \textit{pooled} & label direction & -0.578 & -0.500 & 0.078 & 0.30 \\
\textit{all 3} & \textit{pooled} & answer direction & -0.461 & -0.398 & 0.063 & 0.26 \\
\textit{pooled} & \textit{2 tests} & \textit{both} & -- & -- & 0.070 & 0.26 \\
\midrule
\multicolumn{7}{l}{\textit{In-sample reference only, since adding a predictor cannot lower a training $R^2$}} \\
Checkpoint & Dataset & Direction & training $R^2$ & with ownership & $\Delta R^2$ & \\
\midrule
Lingshu 32B & CheXpert Plus & label direction & 0.769 & 0.788 & 0.019 & \\
Lingshu 32B & CheXpert Plus & answer direction & 0.368 & 0.423 & 0.055 & \\
Qwen2.5-VL-7B & NIH ChestX-ray14 & label direction & 0.241 & 0.246 & 0.004 & \\
Qwen2.5-VL-7B & NIH ChestX-ray14 & answer direction & 0.267 & 0.392 & 0.125 & \\
Qwen2.5-VL-7B & CheXpert Plus & label direction & 0.175 & 0.250 & 0.074 & \\
Qwen2.5-VL-7B & CheXpert Plus & answer direction & 0.473 & 0.509 & 0.036 & \\
\bottomrule
\end{tabular}}}
}
\end{table}
%DIFDELCMD < %%%
\DIFdelend \FloatBarrier

\DIFaddbegin \begin{DIFnomarkup}\cfnote{删去运行记录}{删去各 block 运行记录的整段（协议标识 cf-transfer-v1、run identifier、评估代码版本、库版本等）及同类表述：第 4 节的 at a pinned revision、复现声明中“标识与代码版本记录在 scorecard 中”、A.1 的 its specified code version 与 One analysis pass over the packaged outcomes 一句。“修订早于结果”原以 outcome shard 表述，改为 first result；小节标题由 Protocol record and amendments 改为 Protocol amendments。发布内容与随机种子的复现规格保留。}\end{DIFnomarkup}

\DIFaddend \subsection{Protocol \DIFdelbegin \DIFdel{record and }\DIFdelend amendments}
\label{app:repro-protocol}

\DIFdelbegin \DIFdel{Each block writes a run record containing the protocol identifier }\texttt{\DIFdel{cf-transfer-v1}}%DIFAUXCMD
\DIFdel{, a run identifier naming the checkpoint and dataset, the evaluation code version that produced the block, the pinned checkpoint and processor revisions, the library versions, the numeric precisions, and the resolved processing settings, including the rendered chat template. All }%DIFDELCMD < \cfBlocks{} %%%
\DIFdel{blocks carry the protocol identifier }\texttt{\DIFdel{cf-transfer-v1}}%DIFAUXCMD
\DIFdel{. The protocol file registers the code version it was frozen against.
}\DIFdelend \DIFaddbegin \begin{DIFnomarkup}\cfnote{删表}{20 行修订记录表删去，改为一句：正文报告的、计划模块之外的每项分析都作为带日期的修订写入协议文件，且每条修订早于对应模块的第一个结果。逐条日期留在发布的协议文件里。两处 amendment A15 改为 its amendment。}\end{DIFnomarkup}
\DIFaddend 

\DIFdelbegin \DIFdel{Table~\ref{tab:cf-amendments} lists every addition }\DIFdelend \DIFaddbegin \DIFadd{Every analysis the paper reports beyond the }\cfCoveragePlannedWord{} \DIFadd{planned modules was added }\DIFaddend to the protocol file \DIFaddbegin \DIFadd{as a dated amendment }\DIFaddend after the first wave, and \DIFdelbegin \DIFdel{the date the file records for it. Every analysis the paper reports that is not one of the planned modules appears here, with the date the file gives it . The }%DIFDELCMD < \cfCoveragePlannedWord{} %%%
\DIFdel{planned modules have no amendment date. }\DIFdelend \DIFaddbegin \DIFadd{each amendment precedes the first result of the module it registers. }\DIFaddend Each amendment entry names its own rows, references, and decision rule before it names a result\DIFdelbegin \DIFdel{, and the robustness tables in Appendix~\ref{app:cf-robustness} identify the blocks that carry each one.
A1 and A2 also carry the identifiers of the execution log.
}\DIFdelend \DIFaddbegin \DIFadd{.
}\DIFaddend 


\DIFdelbegin \begin{table}[htbp]
\DIFdeletedtable{%
{\small\textbf{Table (deleted).} \textbf{Protocol amendments.} Every addition to the registered protocol file after the first wave, with the date the file records for it. Each date precedes the first outcome shard of the module it registers.\par}\medskip
\centering
\setlength{\tabcolsep}{5pt}

\fitwidth{%
\fitwidth{%
\begin{tabular}{ll>{\raggedright\arraybackslash}p{0.886\textwidth}}
\toprule
Amendment & Date & What it adds \\
\midrule
A1 & 2026-09-14 & Per-block primary template: the single-template modules score the first eligible template in the order (IY, IB), so a checkpoint whose yes/no template fails the image-free mapping gate is graded on its B-positive template instead of losing its grade. Llama~3.2 Vision 11B re-ran on all three datasets. \\
A2 & 2026-09-14 & CheXpert gains the dose sweep, the two probe refits, the connector locus, the connector-locus calibration, and the five additional templates, with Effusion and Edema as its two template concepts, so the second chest set carries the same modules as the first. \\
A3 & 2026-09-14 & Four alternative direction constructions per concept: difference of class means, Haufe pattern, orthogonalised normal, and residualised normal. \\
A4 & 2026-09-14 & Extra clinical directions for every additional dataset label with at least 100 known positives and negatives, added to the competitor family. \\
A5 & 2026-09-14 & Softmax-weighted and top-quarter token weights for the six logistic directions. \\
A6 & 2026-09-14 & The write grid rescored with fp32 weights and forward pass, and in bf16 at batch size one. \\
A7 & 2026-09-14 & The answer direction: a ridge regression of the clean answer margin on the same projected, train-scaled features. \\
A8 & 2026-09-15 & The difference-of-means and pattern vectors lifted as displacements through the minimum-norm preimage. \\
A9 & 2026-09-15 & Three non-clinical attribute directions written with the six clinical directions as one nine-direction family. \\
A10 & 2026-09-15 & The answer directions fitted on the primary template written under the five held-out templates. \\
A11 & 2026-09-15 & The write grid regraded on the CheXpert validation split with radiologist consensus labels. \\
A12 & 2026-09-16 & The crossed tower and reader: one checkpoint's vision tower loaded behind the other's connector and language model, and all four combinations of tower and reader graded on the same rows. \\
A13 & 2026-09-16 & The identical-tensor replay: the consumed block's stored token tensors fed to every reader of a shared-tower group. \\
A14 & 2026-09-16 & The four semantic endpoints: the negated question, the counterbalanced two-finding forced choice, and the finding word's log-probability in a report continuation, each with its own clean baseline, sham, and random family, together with the incremental-validity test of the ownership contrast. \\
A15 & 2026-09-16 & The fine-grained object family: six further COCO categories chosen by the rule stated below, fitted, written, and graded as their own six-direction family in the same blocks and rows. \\
A16 & 2026-09-16 & The clean pass of the fine-grained categories on the calibration rows, so their clean-answer AUROC is measured on the same rows and by the same rule as every other cell and the difficulty matching compares like with like. \\
A17 & 2026-09-16 & The matched attribute reference: the protocol's own 119-direction random family written against the three attribute questions. Gives attribute and clinical questions the same steering reference. \\
A18 & 2026-09-16 & Three phrasings of each attribute question scored clean on the calibration rows, with the selection rule and its reason recorded per cell. \\
A19 & 2026-09-16 & The expert-label refits: the six concept probes refitted on the radiologist-labelled films, cross-fitted over patients, and written on the test rows. \\
A20 & 2026-09-16 & Projection seeds 1 and 2: the probes, the random family, and the sham redrawn in two further random projections and regraded against that projection's own references. \\
\bottomrule
\end{tabular}}}
}
\end{table}
%DIFDELCMD < 

%DIFDELCMD < %%%
\DIFdelend \FloatBarrier

\paragraph{The fine-grained category selection rule.} The rule that picks the six fine-grained COCO categories \DIFdelbegin \DIFdel{of amendment A15 }\DIFdelend is stated in the protocol file in full, and it reads on the labels alone. A category is eligible when its training positives are at least as many as the weakest of the six existing concepts, when it has at least ten positives among the 400 calibration rows and at least ten among the 600 test rows, and when its mean relative instance area is below the median of the six existing concepts. Eligible categories are ranked by calibration positives in descending order, ties by mean relative instance area in ascending order, and remaining ties by category identifier; the first six are taken. The rule uses the label counts and the instance areas of the fixed cohorts and nothing else: no probe, no answer, and no write enters it. The protocol file carries the rule under the date of \DIFdelbegin \DIFdel{amendment A15}\DIFdelend \DIFaddbegin \DIFadd{its amendment}\DIFaddend , and the module's first \DIFdelbegin \DIFdel{outcome shard }\DIFdelend \DIFaddbegin \DIFadd{result }\DIFaddend is later than that date, so the six categories and the two matching windows were fixed before any of them was scored.
\end{document}
